% Options for packages loaded elsewhere
\PassOptionsToPackage{unicode}{hyperref}
\PassOptionsToPackage{hyphens}{url}
\PassOptionsToPackage{dvipsnames,svgnames,x11names}{xcolor}
\documentclass[
  english,
  11pt,
  reqno]{amsart}
\usepackage{xcolor}
\usepackage{amsmath,amssymb}
\setcounter{secnumdepth}{-\maxdimen} % remove section numbering
\usepackage{iftex}
\ifPDFTeX
  \usepackage[T1]{fontenc}
  \usepackage[utf8]{inputenc}
  \usepackage{textcomp} % provide euro and other symbols
\else % if luatex or xetex
  \usepackage{unicode-math} % this also loads fontspec
  \defaultfontfeatures{Scale=MatchLowercase}
  \defaultfontfeatures[\rmfamily]{Ligatures=TeX,Scale=1}
\fi
\usepackage{lmodern}
\ifPDFTeX\else
  % xetex/luatex font selection
\fi
% Use upquote if available, for straight quotes in verbatim environments
\IfFileExists{upquote.sty}{\usepackage{upquote}}{}
\IfFileExists{microtype.sty}{% use microtype if available
  \usepackage[]{microtype}
  \UseMicrotypeSet[protrusion]{basicmath} % disable protrusion for tt fonts
}{}
\usepackage{longtable,booktabs,array}
\newcounter{none} % for unnumbered tables
\usepackage{calc} % for calculating minipage widths
% Correct order of tables after \paragraph or \subparagraph
\usepackage{etoolbox}
\makeatletter
\patchcmd\longtable{\par}{\if@noskipsec\mbox{}\fi\par}{}{}
\makeatother
% Allow footnotes in longtable head/foot
\IfFileExists{footnotehyper.sty}{\usepackage{footnotehyper}}{\usepackage{footnote}}
\makesavenoteenv{longtable}
% definitions for citeproc citations
\NewDocumentCommand\citeproctext{}{}
\NewDocumentCommand\citeproc{mm}{%
  \begingroup\def\citeproctext{#2}\cite{#1}\endgroup}
\makeatletter
 % allow citations to break across lines
 \let\@cite@ofmt\@firstofone
 % avoid brackets around text for \cite:
 \def\@biblabel#1{}
 \def\@cite#1#2{{#1\if@tempswa , #2\fi}}
\makeatother
\newlength{\cslhangindent}
\setlength{\cslhangindent}{1.5em}
\newlength{\csllabelwidth}
\setlength{\csllabelwidth}{3em}
\newenvironment{CSLReferences}[2] % #1 hanging-indent, #2 entry-spacing
 {\begin{list}{}{%
  \setlength{\itemindent}{0pt}
  \setlength{\leftmargin}{0pt}
  \setlength{\parsep}{0pt}
  % turn on hanging indent if param 1 is 1
  \ifodd #1
   \setlength{\leftmargin}{\cslhangindent}
   \setlength{\itemindent}{-1\cslhangindent}
  \fi
  % set entry spacing
  \setlength{\itemsep}{#2\baselineskip}}}
 {\end{list}}
\usepackage{calc}
\newcommand{\CSLBlock}[1]{\hfill\break\parbox[t]{\linewidth}{\strut\ignorespaces#1\strut}}
\newcommand{\CSLLeftMargin}[1]{\parbox[t]{\csllabelwidth}{\strut#1\strut}}
\newcommand{\CSLRightInline}[1]{\parbox[t]{\linewidth - \csllabelwidth}{\strut#1\strut}}
\newcommand{\CSLIndent}[1]{\hspace{\cslhangindent}#1}
\ifLuaTeX
\usepackage[bidi=basic,shorthands=off]{babel}
\else
\usepackage[bidi=default,shorthands=off]{babel}
\fi
\ifLuaTeX
  \usepackage{selnolig} % disable illegal ligatures
\fi
\setlength{\emergencystretch}{3em} % prevent overfull lines
\providecommand{\tightlist}{%
  \setlength{\itemsep}{0pt}\setlength{\parskip}{0pt}}
% Portable additions for the arXiv-oriented profile.
% Keep packages within the standard TeX Live distribution.
\usepackage{mathtools}
\usepackage{amssymb}
\usepackage{amsthm}
\usepackage{booktabs}
\usepackage{listings}
\usepackage{microtype}
\usepackage{centernot}
\usepackage{newunicodechar}

\ifPDFTeX
  % unicode-math supplies mathscr on Xe/LuaTeX; pdfTeX needs its own font.
  \usepackage{mathrsfs}
  % Explicit, meaning-preserving mappings only. Other unsupported input errors
  % are fatal; never discard characters or guess a different math expression.
  \newunicodechar{α}{\ensuremath{\alpha}}
  \newunicodechar{β}{\ensuremath{\beta}}
  \newunicodechar{γ}{\ensuremath{\gamma}}
  \newunicodechar{δ}{\ensuremath{\delta}}
  \newunicodechar{ε}{\ensuremath{\varepsilon}}
  \newunicodechar{θ}{\ensuremath{\theta}}
  \newunicodechar{λ}{\ensuremath{\lambda}}
  \newunicodechar{μ}{\ensuremath{\mu}}
  \newunicodechar{π}{\ensuremath{\pi}}
  \newunicodechar{σ}{\ensuremath{\sigma}}
  \newunicodechar{ω}{\ensuremath{\omega}}
  \newunicodechar{ℂ}{\ensuremath{\mathbb{C}}}
  \newunicodechar{ℕ}{\ensuremath{\mathbb{N}}}
  \newunicodechar{ℚ}{\ensuremath{\mathbb{Q}}}
  \newunicodechar{ℝ}{\ensuremath{\mathbb{R}}}
  \newunicodechar{ℤ}{\ensuremath{\mathbb{Z}}}
  \newunicodechar{∞}{\ensuremath{\infty}}
  \newunicodechar{∈}{\ensuremath{\in}}
  \newunicodechar{∅}{\ensuremath{\emptyset}}
  \newunicodechar{⊂}{\ensuremath{\subset}}
  \newunicodechar{⊆}{\ensuremath{\subseteq}}
  \newunicodechar{∩}{\ensuremath{\cap}}
  \newunicodechar{∪}{\ensuremath{\cup}}
  \newunicodechar{≤}{\ensuremath{\leq}}
  \newunicodechar{≥}{\ensuremath{\geq}}
  \newunicodechar{≠}{\ensuremath{\neq}}
  \newunicodechar{→}{\ensuremath{\rightarrow}}
  \newunicodechar{↦}{\ensuremath{\mapsto}}
  \newunicodechar{⇒}{\ensuremath{\Rightarrow}}
  \newunicodechar{⇔}{\ensuremath{\Leftrightarrow}}
  \newunicodechar{−}{\ensuremath{-}}
\fi
\newunicodechar{□}{\ensuremath{\square}}

% Portable negated hooked arrow used by several manuscripts.
\providecommand{\nhookrightarrow}{\centernot\hookrightarrow}

% U+220E END OF PROOF: render a small hollow square at the right margin.
\newcommand{\rightqedsquare}{%
  \ifhmode
    \unskip\nobreak\hfill\ensuremath{\square}%
  \else
    \noindent\hfill\ensuremath{\square}%
  \fi
}
\newunicodechar{∎}{\rightqedsquare}

% Reduce avoidable overfull lines without changing the source manuscript.
\setlength{\emergencystretch}{2em}
% Match amsart's theorem preskip while keeping Markdown statement bodies roman.
\makeatletter
\newcommand{\pandocstatementheading}[1]{%
  \par\addpenalty{\@beginparpenalty}%
  \addvspace{.5\baselineskip plus .2\baselineskip minus .2\baselineskip}%
  {\normalfont\bfseries #1\par}%
  \nobreak\@afterindenttrue\@afterheading
}
\makeatother
% House style for standard/arxiv-ams: match the References section heading.
% The native abstract box, margins and body font remain owned by amsart.
\usepackage{etoolbox}
\newcommand{\pandocabstractheading}{%
  \par\begingroup
    \normalsize\section*{\abstractname}%
  \endgroup
  \nobreak\noindent\ignorespaces
}
\makeatletter
\patchcmd{\abstract}
  {\item[\hskip\labelsep\scshape\abstractname.]}
  {\item[]\pandocabstractheading}
  {}%
  {\PackageError{studio-ams-abstract}
     {Unsupported amsart abstract definition}
     {Check the installed class before updating the abstract-heading patch.}}
\makeatother
% Common mathematics macros
% Keep this minimal and portable across engines.

% Math shortcuts
\newcommand{\N}{\mathbb{N}}
\newcommand{\Z}{\mathbb{Z}}
\newcommand{\Q}{\mathbb{Q}}
\newcommand{\R}{\mathbb{R}}
\newcommand{\C}{\mathbb{C}}
\newcommand{\eps}{\varepsilon}

% Theorem-like environments (lightweight)
% If you need full theoremstyles, consider amsthm or thmtools here.
% \usepackage{amsthm}   % already available via TeX Live if needed
% \newtheorem{theorem}{Theorem}
% \newtheorem{lemma}{Lemma}

% Code listings default (avoid minted to keep shell-escape off)
\lstset{basicstyle=\ttfamily\small,breaklines=true,columns=fullflexible}

\usepackage{bookmark}
\IfFileExists{xurl.sty}{\usepackage{xurl}}{} % add URL line breaks if available
\urlstyle{same}
\hypersetup{
  pdftitle={A projective twistor correspondence for local analytic Einstein germs},
  pdfauthor={Lior Isthmus},
  pdflang={en},
  colorlinks=true,
  linkcolor={blue},
  filecolor={Maroon},
  citecolor={blue},
  urlcolor={blue},
  pdfcreator={LaTeX via pandoc}}

\title{A projective twistor correspondence for local analytic Einstein
germs}
\author[Lior Isthmus]{Lior Isthmus\textsuperscript{1}}
\date{2026-11-13}

\begin{document}
\maketitle
\footnotetext[1]{Email: \href{mailto:lioristhmus@proton.me}{\nolinkurl{lioristhmus@proton.me}}}
\setcounter{footnote}{1}

\markboth{\MakeUppercase{Lior Isthmus}}{\MakeUppercase{A projective twistor correspondence for local analytic\ldots{}}}

\protect\phantomsection\label{sec:gg-title}

\protect\phantomsection\label{a-projective-twistor-correspondence-for-local-analytic-einstein-germs}

\protect\phantomsection\label{sec:gg-abstract}

\section*{Abstract}\label{abstract}

For each prescribed numerical cosmological constant, including zero, we
establish a convergent projective twistor correspondence recovering
every marked four-dimensional local analytic Einstein germ of any
algebraic Weyl type. The inputs are two independent analytic Čech
classes of weights \(2\) and \(-6\) on a fixed flat twistor germ \(T_0\)
along its whole marked zero line, encoding both complete chiral
curvature records. The fixed source correction and calibration select a
unit optical section, an intrinsic projective extension condition
recognizes Einstein geometry, and the prescribed normalized contact,
optical and affine maps and their inverses converge on common
neighborhoods of the whole physical sky for bounded holomorphic input
families. On independently admitted ordinary and operator domains we
prove both inverse comparisons, with the two source operations acting in
one projective analytic Heisenberg model. The same comparisons preserve
the nonlinear source, compact point objects, optical evaluation and
calibrated infinity forms, together with physical point and full
spin-frame changes and coupled Lorentz reality. For fixed positive
Schwarzschild mass, we obtain finite horizon-crossing presentations, a
calibrated local primitive comparison and optical approximation bounds
for the entire specified compact nonradial crossing family using finite
local readouts and the exact benchmark cuts, coordinate comparisons and
Sachs reference frames.

\tableofcontents

\protect\phantomsection\label{sec:gg-introduction}

\subsection{1. Introduction}\label{introduction}

Penrose described the googly obstruction as the problem of incorporating
the interacting helicity opposite to the one described by the ordinary
nonlinear-graviton construction, while retaining the same basic twistor
conventions (Penrose 2015, abstract). The classical correspondence
produces half-flat conformal geometry from holomorphic threefolds with
suitable rational curves (Penrose 1976). The problem concerns the full
nonlinear interaction of the two gravitational chiralities in local
analytic geometry.

Our inputs are two independent analytic Čech classes of weights \(2\)
and \(-6\) on the fixed germ \(T_0\) of
\(\operatorname{Tot}(\mathcal O(1)\oplus\mathcal O(1))\) along its whole
marked zero line, a numerical cosmological constant and the point and
full spin marking. The contact carrier, optical bundle and selected
Einstein scale are derived from these data; no independent dual-twistor
class, connection, scale or Einstein metric is supplied.

The three spaces have distinct roles: \(T_0\) is the fixed input germ,
\(Z\) the generated contact fivefold, and \(Z_+\) the selected half-flat
background of the independent ordinary presentation. The initial model
\(Z_0\subset\mathbb P(T^*T_0)\) supplies the starting contact charts; a
general mixed output need not retain its fixed cotangent projection.
Physical points are represented by a complete four-dimensional family of
compact Legendrian sky quadrics
\(Q_x\simeq\mathbb P^1\times\mathbb P^1\) in \(Z\). The rational lines
in \(Z_+\) describe the auxiliary half-flat metric, on which the
unrestricted opposite-chirality class is registered and read linearly.
Reconstruction of the full coupled Einstein geometry from the two
records is nonlinear. Section 7 constructs both ordinary counits, and
§8.3 gives covariant point and full spin-frame changes, which generally
change both the background and the registered class.

Theorem 2.4 recovers every marked local analytic Einstein germ, with
both complete chiral curvature records, every algebraic Weyl type and
numerical cosmological constant including zero, and proves both inverse
comparisons on the independently defined ordinary domain. To prove
convergence, the same prescribed normalizer is applied to finite source
execution and to an analytic Cartan witness: Theorem 6.3 identifies each
finite stored tuple with the amplitude Taylor restriction of one jointly
holomorphic tuple. Its prescribed contact, optical and affine
representatives, including their inverses, consequently converge on
common neighborhoods of the whole physical sky for bounded holomorphic
families and fixed finite conditioned diagrams.

The two bare contact-lift spaces fail to close under the mixed bracket,
and contact-cocycle compatibility alone does not ensure extension of the
selected scale (Propositions 9.1--9.2). The source law in §5 uses fixed
holomorphic cochain operations and the compatible inverse \(L_m\) to
remove the optical obstruction while retaining the prescribed entire
free record. Fresh incidence and optical computation follow each
correction; value, gradient and numerical scalar fix the six flat
homogeneous scale modes. The resulting unit whole-neighborhood optical
section determines the calibrated affine reduction of Proposition 5.7,
whose intrinsic projective extension condition recognizes Einstein
geometry by Proposition 5.4. Proposition 5.9 and §6.1 specify the
respective roles of finite source execution and the geometric comparison
used to prove convergence.

The free symmetric curvature coordinates are the classical exact-set
gravitational null data (Penrose 1980, secs. 2--3), whose description
already includes a cosmological term (p.233). The flat free-field
integral formulas of that paper's §4 are distinct from the nonlinear
existence and analytic estimates proved here in §§3--6 and Appendices
A--B.

The null-geodesic/contact correspondence, reconstruction from complete
normal sky quadrics and deformation framework are classical (LeBrun
1983, III.1--3 and III.6). The same paper's Appendix gives the
metric-cohomology correspondence underlying the affine torsor of §D.7
(LeBrun 1983, 230). Baston and Mason studied formal-neighborhood
criteria for conformal gravity and Einstein geometry in second-order
perturbation theory (Baston and Mason 1987, 815). LeBrun proved the
fully curved extension theorem: the canonical first-order Poisson
thickening extends uniquely through order four, its fifth-order
obstruction is the Bach tensor, and the additional sixth-order
obstruction is the Eastwood--Dighton tensor (LeBrun 1991, Main Theorem
and Theorems 9.1--9.3, p.8 and §9). The conformally Einstein
interpretation of their joint vanishing has the stated
algebraic-generality hypothesis on the Weyl curvature.

LeBrun's Einstein-bundle criterion recognizes conformally Einstein
metrics by nonvanishing holomorphic sections on the space of null
geodesics (LeBrun 1985). The ray bundle has rank two on the fivefold
\(Z\); the ordinary Einstein-bundle form takes values in the rank-three
bundle \(\Omega^1_{Z_+}\otimes L^2\) on \(Z_+\). Wolf explicitly
attributes the recognition to LeBrun and describes this intrinsic
ordinary sheaf (Wolf 2007, sec. 3.8). Here Proposition 5.3 constructs
the ray bundle from the tangent-sky curve, Proposition 5.4 identifies
its whole-neighborhood unit sections with scales, and Proposition J.4
compares its Cauchy and tractor realizations.

The unit-scale/parallel-tractor equivalence is standard (Gover 2004,
sec. 2, Theorem 2.1); Proposition 5.4 proves the form used here and
identifies it with the intrinsic optical data. Friedrich's conformal
characteristic construction gives formal null-data determination at past
time-like infinity (Friedrich 2013, arXiv v1, §1 and Lemma 6.1). It
concerns a conformal vacuum boundary problem and is not used as an
analytic convergence theorem for the present germs.

Penrose's bi-twistor framework combines twistor and dual-twistor
variables with canonical commutation relations, using their classical
limit to describe null-geodesic propagation in conformally curved
spacetimes (Penrose 2026, sec. 6, pp.8--9). Its concluding paragraph
anticipates a cohomological construction of general real spacetimes as a
future direction.

Theorems 9.10--9.12 construct the independent operator presentation in
one projective analytic Heisenberg model, with the negative source
acting through derivative variables. Both inverse comparisons hold on
the entire independently admitted operator domain. Theorem 9.12
identifies the nonlinear source, calibrated infinity forms, optical
evaluation and compact point objects through the same natural
comparisons, using raywise descent of the normal local-twistor bundle
and explicit point-plane and affine-slice comparisons. The point modules
represent the same physical skies \(Q_x\). These comparisons retain
physical point and full spin-frame changes and coupled Lorentz reality.

Adamo et al.~propose data on one twistor space, holomorphic constraints
and Einsteinity without explicitly solving field equations on twistor
space as three criteria for their Schwarzschild construction (Adamo et
al. 2026b, sec. 1, including footnote 2). Section 10 and Appendix K
apply the correspondence across the future Schwarzschild horizon, with a
calibrated local source-primitive comparison and enriched finite
overlaps. Proposition 10.5 gives optical approximation bounds on the
entire specified compact nonradial crossing family using finite local
readouts and the exact benchmark cuts, coordinate comparisons and Sachs
reference frames. Proposition 10.3 remains a conditional single-readout
estimate, and §§K.6--K.7 give a sufficiently short common-source
interval.

The proof separates finite holomorphic source execution from an
analytically bounded Cartan realization before comparing their
prescribed representatives. Ordinary and operator counits are proved on
their independent domains before the source, infinity and point diagrams
are collected in Theorem 9.12.

The main proof dependencies are:

{\def\LTcaptype{none} % do not increment counter
\begin{longtable}[]{@{}
  >{\raggedright\arraybackslash}p{(\linewidth - 2\tabcolsep) * \real{0.5000}}
  >{\raggedright\arraybackslash}p{(\linewidth - 2\tabcolsep) * \real{0.5000}}@{}}
\toprule\noalign{}
\begin{minipage}[b]{\linewidth}\raggedright
Assertion
\end{minipage} & \begin{minipage}[b]{\linewidth}\raggedright
Construction and proof
\end{minipage} \\
\midrule\noalign{}
\endhead
\bottomrule\noalign{}
\endlastfoot
Original input and curvature records & Theorem 3.1 and Appendix A \\
Finite source law and fixed compatible inverse & Theorem 5.5 and
Appendices C--D \\
Common convergence of its prescribed representatives & Theorems 3.3 and
6.3, Appendix B \\
Geometric, ordinary and physical inverse comparisons & Theorem 2.4,
§§7--8, Appendix E \\
Independent analytic operator and module inverse comparisons & Theorem
9.10, Appendices F--H \\
The same source, infinity structure and point diagram & Theorem 9.12,
Propositions I.3, J.1 and J.3--J.5, Theorems J.2 and J.6, and
Proposition J.7 \\
Finite Schwarzschild application & Section 10 and Appendix K, including
the full-family finite-readout approximation of Proposition 10.5, using
the preceding correspondence \\
\end{longtable}
}

The proofs of Theorem 2.4 and Theorems 9.10--9.11 use neither Theorem
9.12 nor Appendix J. The simultaneous source, infinity and point
comparison uses these preceding correspondences and explicit geometric
comparisons. The application row depends on these results and supplies
no premise for them.

All uniform assertions use the bounded original families and fixed
finite conditioned diagrams of Definition 2.1. The subject is local
analytic geometry: arbitrary global spacetime topology, smooth
nonanalytic data and uniform estimates for infinite unbounded
point/frame histories are outside the assertions. The notation index in
Appendix L records the source spaces, bundle conventions and comparison
maps.

\textbf{Computational supplement.} Reproducible SageMath computations
for selected identities and finite-order examples are supplied in the
computational supplement, version 0.2.0, revision 3 (Isthmus 2026). The
package includes code, inputs, recorded environments and results, and an
index linking each check to the manuscript formulas. It covers selected
curvature, optical and coordinate-change identities and low-order
curvature-record recursions; the general estimates and convergence
arguments are given in the text.

\protect\phantomsection\label{sec:gg-objects}

\subsection{2. Objects and the main
correspondence}\label{objects-and-the-main-correspondence}

Let \(S,S'\) be complex symplectic two-spaces and \(V=S\otimes S'\),
with its induced nondegenerate symmetric form \(\eta\). Fix the bases,
affine origin and full spin marking. In both ordered spin bases our
alternating matrices and index operations are
\begin{equation}\protect\phantomsection\label{eq:gg-spin-index-convention}{\begin{gathered}
\epsilon_{AB}=\epsilon_{A'B'}=
       \begin{pmatrix}0&1\\-1&0\end{pmatrix}, \\
\epsilon^{AB}=\epsilon^{A'B'}=
       \begin{pmatrix}0&-1\\1&0\end{pmatrix}, \\
v^A=\epsilon^{AB}v_B,\quad v_A=\epsilon_{AB}v^B .
\end{gathered}}\end{equation} The primed operations have the same order. Thus
\(\eta_{AA'BB'}=\epsilon_{AB}\epsilon_{A'B'}\),
\(\pi^{A'}=(-\pi_{1'},\pi_{0'})\), and, for the north lift
\(\pi_{A'}=(1,\zeta)\), \(\pi_{A'}\,\mathrm d\pi^{A'}=-\mathrm d\zeta\).

\protect\phantomsection\label{sec:gg-curvature-signature-conventions}

Our curvature convention is
\(R(a,b)=[\nabla_a,\nabla_b]-\nabla_{[a,b]}\), with
\(\operatorname{Ric}(b,c)=\operatorname{tr}(a\mapsto R(a,b)c)\). The
Lorentz restriction uses signature \((+,-,-,-)\). In a compatible
soldering frame the real vectors are Hermitian matrices and
\(\eta(v,v)=2\det(v^{AA'})\). Throughout, \(\Lambda\) means the
prescribed coefficient in \(\operatorname{Ric}(g)=\Lambda g\). For
example, the de Sitter metric of radius \(a\) in this signature has
\(\operatorname{Ric}=-3a^{-2}g\) with this curvature convention.
Appendix A.1 temporarily uses Gaussian witnesses of signature
\((-,+,+,+)\); its explicit metric negation and reversed auxiliary
scalar give the opposite-signature witnesses before the fixed spin
readout. Negating a metric fixes its connection and covariant Ricci
tensor and changes its Einstein coefficient's sign; no such change is
made silently in the source label.

In four dimensions our Schouten tensor is
\(P=\tfrac12(\operatorname{Ric}-\tfrac16\operatorname{Scal}\,g)\). For
the component comparison with Penrose's printed formulas (Penrose 2015,
secs. 4 and 6, especially p.12), denote their scalar by
\(\Lambda_{\mathrm P}\). We declare the curvature-component dictionary
\(R^{\mathrm P}_{abc}{}^d=[R(e_a,e_b)e_c]^d\) with the \(R\) just fixed.
The primary text contracts
\(\operatorname{Ric}^{\mathrm P}_{ac}=R^{\mathrm P}_{abc}{}^b\), whereas
our Ricci contraction is \(R^{\mathrm P}_{bac}{}^b\). Consequently
\begin{equation}\protect\phantomsection\label{eq:gg-original-infinity-curvature-dictionary}{
 \operatorname{Ric}^{\mathrm P}=-\operatorname{Ric},\qquad
 \Lambda_{\mathrm P}=-\Lambda,\qquad
 P^{\mathrm P}
   =\tfrac1{12}\operatorname{Scal}^{\mathrm P}g
                         -\tfrac12\operatorname{Ric}^{\mathrm P}
   =P .
}\end{equation} The primary article states that contraction but does not
specify its curvature commutator sign. The displayed component
identification explicitly fixes our dictionary; it is not a claim that
an unstated historical convention was independently verified. The
numerical \(\Lambda\) prescribed throughout this manuscript is
unchanged.

We lower the output of the curvature operator in the last slot:
\begin{equation}\protect\phantomsection\label{eq:gg-weyl-spinor-decomposition}{
 \begin{aligned}
 R_{abcd}&=g(R(e_a,e_b)e_c,e_d),\\
 C_{abcd}&=R_{abcd}
   -g_{ad}P_{bc}+g_{bd}P_{ac}+g_{ac}P_{bd}-g_{bc}P_{ad},\\
 C_{AA'BB'CC'DD'}
 &=-\Psi^+_{ABCD}\epsilon_{A'B'}\epsilon_{C'D'}
   -\Psi^-_{A'B'C'D'}\epsilon_{AB}\epsilon_{CD} .
 \end{aligned}
}\end{equation} In particular the component readout, inverse to the last
line, is
\begin{equation}\protect\phantomsection\label{eq:gg-weyl-spinor-readout}{
 \begin{aligned}
 \mathfrak p_+(C)_{ABCD}=\Psi^+_{ABCD}
   &=-\tfrac14\epsilon^{A'B'}\epsilon^{C'D'}
                       C_{AA'BB'CC'DD'},\\
 \mathfrak p_-(C)_{A'B'C'D'}=\Psi^-_{A'B'C'D'}
   &=-\tfrac14\epsilon^{AB}\epsilon^{CD}
                       C_{AA'BB'CC'DD'} .
 \end{aligned}
}\end{equation} Indeed \(\epsilon^{AB}\epsilon_{AB}=-2\), and
contraction of a symmetric spinor with \(\epsilon\) is zero. Thus, with
\(u=x^{00'}\), \(v=x^{01'}\) and \(q=x^{10'}\), the highest components
are \(\Psi^+_{0000}=-C_{uvuv}\) and \(\Psi^-_{0'0'0'0'}=-C_{uquq}\). The
chiral labels \(+\) and \(-\) mean these unprimed and primed spinors,
respectively; physical names for their helicities depend on the twistor
orientation. On the marked Lorentz slice this real component convention
gives \(\Psi^-=\overline{\Psi^+}\).

Put \(T_0=\operatorname{Tot}(\mathcal O(1)^{\oplus2})\to\mathbb P(S')\),
as a germ along its entire zero line. Its north and south coordinates
and incidence are
\begin{equation}\protect\phantomsection\label{eq:gg-flat-incidence}{
 \zeta=\pi_{1'}/\pi_{0'},\quad
 w^A=\omega^A/\pi_{0'},\qquad
 \xi=\zeta^{-1},\quad v^A=w^A/\zeta,\qquad
 \omega^A=x^{AA'}\pi_{A'}.
}\end{equation}

\pandocstatementheading{\protect\phantomsection\label{def:gg-original-input}Definition 2.1.}

For an integer \(m\), let \(\mathscr H_m(T_0)\) be the two-chart
analytic Čech group in the direct limit over neighborhoods of the whole
zero line, with coboundaries
\(a_0(w,\zeta)-\zeta^m a_1(w/\zeta,1/\zeta)\). In the north frame, fix
the ordered overlap and Čech differential by
\begin{equation}\protect\phantomsection\label{eq:gg-original-cech-orientation}{
 F_{NS}=F,\qquad
 (\delta a)_{NS}(w,\zeta)
   =\zeta^m a_S(w/\zeta,1/\zeta)-a_N(w,\zeta).
}\end{equation} This differential has the displayed coboundary image,
since reversing its sign does not change that subspace. When both
cochains are already expressed in one common frame, we write the same
differential as \(a_S-a_N\). The original free input at a numerical
\(\Lambda\in\mathbb C\) is
\begin{equation}\protect\phantomsection\label{eq:gg-original-input}{
 \mathcal A_\Lambda=\mathscr H_2(T_0)\oplus\mathscr H_{-6}(T_0).
}\end{equation} A bounded holomorphic family has representatives bounded
on one common annular fiber neighborhood, bounded \(|\Lambda|\), and
fixed positive parent-chart and bounded spin/inverse-spin margins. The
parameter base is a finite-dimensional complex manifold germ; the
representatives, numerical function \(\Lambda\) and markings depend
holomorphically on it. The finite-jet assertions are restrictions of
these holomorphic families to finite Taylor neighborhoods. The nilpotent
amplitude stages used in §5 are the separate flat-residue source
filtration specified there. Uniform output assertions always mean a
common shrink for each such bounded family and each fixed finite diagram
of permitted point/frame operations. A conditioned finite diagram has
positive coordinate and conic Cauchy buffers; bounded maps and inverses;
a positive lower bound for each determinant, contact multiplier, root
coefficient and frame coefficient that must be a unit; and images of all
its specified finite composites contained in their assigned parents with
positive buffer. These conditions are imposed on the supplied
representatives. A common restriction reserves the minima of these
finitely many buffers and unit margins, while retaining the entire
compact parameter set under consideration, including the amplitude path
through one when it is used.

Write \(\mathfrak E_\Lambda\) for holomorphic Einstein four-germs with
\(\operatorname{Ric}=\Lambda g\), the fixed tangent form at the marked
point, and the full spin marking. Objects are considered up to
fixed-mark isometry. Moving-point and frame instructions are retained as
physical arrows, with their specified spin lifts; they are not erased
when their tensor endpoints agree. Their short representatives and
finite compositions are specified in §8.

The geometric output has an intrinsic description. A sky incidence
consists of a contact fivefold \(Z\) with a complete regular
four-dimensional family of embedded compact Legendrian quadrics
\(Q_x\simeq\mathbb P^1\times\mathbb P^1\), their ordered rulings, and
correspondence \(p:F\to M\), \(q:F\to Z\). The contact line restricts to
\(\mathcal O(1,1)\) on each sky, and the one-dimensional \(q\)-fibers
project nontrivially to \(M\).

\pandocstatementheading{\protect\phantomsection\label{def:gg-ray-object}Definition 2.2.}

A regular-normal affine ray object consists of this incidence, a
rank-six holomorphic nondegenerate quadratic bundle \((\mathcal V,H)\)
on \(Z\), a section \(I\) with \(H(I,I)=-\Lambda/3\), and an affine line
subbundle \(\mathscr A\subset\mathcal V\). The bundle is trivial on the
central sky, and the translation line of \(\mathscr A\) restricts to
\(\mathcal O(-1,-1)\). Every \(Y\in\mathscr A\) satisfies \(H(Y,Y)=0\)
and \(H(I,Y)=1\). The initial full spin lift is retained.

Complete-sky evaluation and inverse null transport reconstruct the
rank-six bundle \(\mathcal T=p_*q^*\mathcal V\), its connection \(D\),
and the unique affine-torsor section \(X\), as proved in §4. Put
\(W_{\mathrm{scr}}=(\mathbb CI+\mathbb CX)^\perp\), the rank-four middle
bundle of tangent Spin type \(S\otimes S'\), and \(\theta(v)=D_vX\).
Regular-normal means that \(\theta:TM\to W_{\mathrm{scr}}\) is
invertible, \(F^DX=0\), and the Ricci contraction of the projected
ambient curvature is zero. Explicitly, with
\(\operatorname{pr}_{W_{\mathrm{scr}}}\) the \(H\)-orthogonal
projection, put
\begin{equation}\protect\phantomsection\label{eq:gg-ray-middle-curvature-definition}{\begin{gathered}
\mathcal R_{\mathrm{mid}}(v,w)u
  =\theta^{-1}\operatorname{pr}_{W_{\mathrm{scr}}}
                   F^D(v,w)\theta(u), \\
\operatorname{Ric}_{\mathrm{mid}}(w,u)
  =\operatorname{tr}\bigl(v\mapsto\mathcal R_{\mathrm{mid}}(v,w)u\bigr)=0.
\end{gathered}}\end{equation} This is the middle block of the ambient curvature,
rather than the curvature of the connection obtained by projecting \(D\)
to \(W_{\mathrm{scr}}\). Denote these objects and their enriched,
spin-lifted geometric comparisons by \(\mathfrak S_\Lambda\).

The connection in this definition is reconstructed, not an additional
arbitrary input. These recognition conditions are proved for every
generated output. There is also a presentation without a supplied
rank-six bundle or normality condition: take a regular contact seed and
a unit whole-neighborhood section of its intrinsic projective optical
bundle, with the same central calibration and derived numerical scalar.
Proposition 5.3 and Appendix D construct that bundle independently;
Proposition 5.4 and the canonical reconstruction of §4 give this ray
presentation and its inverse.

\pandocstatementheading{\protect\phantomsection\label{def:gg-ordinary-object}Definition 2.3.}

An ordinary object consists of a holomorphic threefold germ \(Z_+\)
along a complete regular rational line \(\ell\), its complete nearby
line family with \(N_{\ell/Z_+}=\mathcal O(1)^{\oplus2}\), a degree-one
line bundle \(L\) with an actual root identification
\(\alpha:L^{-4}\xrightarrow{\sim}K_{Z_+}\), an actual form
\(\tau\in H^0(Z_+,\Omega^1\otimes L^2)\), and an arbitrary analytic
class \(\beta\in H^1(Z_+,L^{-6})\), with the retained point, tangent and
full spin marking.

The marking identifies \(\ell\) with \(\mathbb P(S'_o)\) and induces the
adjunction identification
\(d_o:\mathcal O(-4)\xrightarrow{\sim}K_{Z_+}|_\ell\). The degree-one
calibration is an actual isomorphism
\begin{equation}\protect\phantomsection\label{eq:gg-ordinary-degree-one-calibration}{
 \kappa:L|_\ell\xrightarrow{\sim}\mathcal O_{\mathbb P(S'_o)}(1),
 \qquad \alpha|_\ell=d_o\circ\kappa^{-4}.
}\end{equation} The central restriction of \(\tau\), read using
\(\kappa^2\), is exactly
\begin{equation}\protect\phantomsection\label{eq:gg-central-ordinary-calibration}{
 \tau|_{T\ell}=\pi_{A'}\,\mathrm d\pi^{A'}.
}\end{equation} The classical complete-line construction reads a density
\(\sigma\) from this same form, with \(\sigma(o)=1\). Its canonical
prolongation \(I_\tau\) is given in §7. The numerical label is
\(\Lambda=-3H(I_\tau,I_\tau)\), without rescaling the form or class.
Objects are taken modulo fixed-mark ordinary comparisons preserving all
these data, including \(\alpha\) and \(\kappa\). Thus a line comparison
\(f_L\) at a fixed mark satisfies \(\kappa_t\circ f_L|_\ell=\kappa_s\).
Denote them by \(\mathfrak W_\Lambda\). Their physical arrow admission
is defined independently in §8. For a moving comparison with specified
primed spin map \(B:S'_s\to S'_t\), let
\(j_B:\mathcal O_s(1)\to\mathcal O_t(1)\) be the induced degree-one map
over \(\mathbb P(B)\). The calibration condition is
\(\kappa_t\circ f_L|_{\ell_s}=j_B\circ\kappa_s\). In particular \(B=-1\)
acts by \(-1\) on \(\mathcal O(1)\), and its specified full-spin arrow
is retained.

The fixed central value removes a real sign ambiguity: opposite unit
sections can determine the same selected metric. The definition does not
identify them by a scalar gauge. It is a condition in ordinary
holomorphic geometry, not membership in a known generated image.

The ordinary self-dual Einstein correspondence is classical (Ward 1980).
Its Einstein-bundle formulation is used here in the bosonic local form
of (Wolf 2007, Propositions 3.4--3.7). The author's degree-minus-one
universal bundle is the inverse of our \(L\). The actual form has the
same bundle type at zero cosmological constant. The selected density and
the numerical constant are read with our calibration, as in §7.

\pandocstatementheading{\protect\phantomsection\label{thm:gg-main}Theorem 2.4.}

For each fixed numerical \(\Lambda\in\mathbb C\) and every original
analytic pair in \(\mathcal A_\Lambda\), the holomorphic generating law
of §5, with the specified normalizer of §6, produces convergent contact
and optical transition maps and inverse maps on a neighborhood of the
whole compact physical sky. Their exact cocycles, complete
four-dimensional incidence, nondegenerate solder and canonical normal
connection are generated by this law. Its distinguished whole-sky
optical section is the same canonical Einstein scale, and its two
complete free curvature records are precisely the original ones.

The resulting object maps
\(\mathcal G:\mathcal A_\Lambda\to\mathfrak S_\Lambda\) and
\(\mathcal E:\mathfrak S_\Lambda\to\mathcal A_\Lambda\) recover every
marked holomorphic Einstein germ and satisfy
\begin{equation}\protect\phantomsection\label{eq:gg-main-generator-inverses}{
 \mathcal E\circ\mathcal G=1_{\mathcal A_\Lambda},\qquad
 \mathcal G\circ\mathcal E\simeq1_{\mathfrak S_\Lambda}.
}\end{equation} There are independently defined ordinary objects and
actual representative maps
\(\Gamma:\mathcal A_\Lambda\to\mathfrak W_\Lambda\) and
\(d_{\mathrm{rec}}:\mathfrak W_\Lambda\to\mathcal A_\Lambda\) with
\begin{equation}\protect\phantomsection\label{eq:gg-main-ordinary-inverses}{
 d_{\mathrm{rec}}\circ\Gamma=1_{\mathcal A_\Lambda},\qquad \Gamma\circ d_{\mathrm{rec}}\simeq1_{\mathfrak W_\Lambda}.
}\end{equation} Consequently
\(\mathcal G\circ d_{\mathrm{rec}}:\mathfrak W_\Lambda\to\mathfrak S_\Lambda\)
and \(\Gamma\circ\mathcal E:\mathfrak S_\Lambda\to\mathfrak W_\Lambda\)
have both inverse object comparisons. At the physical-arrow level,
Theorem 8.3 equips the original records with the short
point/full-spin-frame instructions, their source-dependent composition
and their inverses. With precisely those arrows, the correspondences are
bijective on physical spin-isometry Hom sets.

All assertions hold for both independent chiralities, all algebraic Weyl
types and numerical \(\Lambda=0\), with common domains for the bounded
holomorphic families and fixed finite diagrams of Definition 2.1. The
original Lorentz-real coefficient condition and real \(\Lambda\) produce
exactly the Lorentz-real Einstein germs, with coupled reality and
point/frame operations on the ordinary objects. No Weyl determinant or
cosmological constant is inverted.

The proof is completed in §8. Sections 3--6 establish the actual
generator and its analytic estimates; §7 proves the ordinary counit on
every independently admitted object.

\protect\phantomsection\label{sec:gg-records}

\subsection{3. Original classes and convergent record
transport}\label{original-classes-and-convergent-record-transport}

For \(r\ge0\) put
\(W_r^+=\operatorname{Sym}^{r+4}S^*\otimes\operatorname{Sym}^rS'^*\) and
\(W_r^-=\operatorname{Sym}^rS^*\otimes\operatorname{Sym}^{r+4}S'^*\).
The free records of an Einstein metric are formed from the spinors in
eq.~\ref{eq:gg-weyl-spinor-readout}:
\begin{equation}\protect\phantomsection\label{eq:gg-free-record-projection}{
 U_r^+=\operatorname{Sym}_{A_1\cdots A_{r+4}}
       \operatorname{Sym}_{A'_1\cdots A'_r}
       \nabla_{A_1A'_1}\cdots\nabla_{A_rA'_r}
                     \Psi^+_{A_{r+1}\cdots A_{r+4}}\big|_o .
}\end{equation} The negative formula exchanges the spin factors and uses
\(\Psi^-\). Each symmetrization is an average. A subscript \(a,j\)
counts slots equal to one in the fixed component bases; it is not a
coefficient of a generating polynomial.

\pandocstatementheading{\protect\phantomsection\label{thm:gg-residue}Theorem 3.1.}

The original analytic classes are equivalent to independent complete
lists \(U_r^\pm\in W_r^\pm\) satisfying, for some \(C,R_0>0\),
\begin{equation}\protect\phantomsection\label{eq:gg-original-factorial-bound}{
 \|U_r^\pm\|_{\mathrm{ev}}\le Cr!R_0^{-r},\qquad
 \|U_r^+\|_{\mathrm{ev}}
  =\sup_{|\xi|,|\pi|\le1}|U_r^+(\xi^{r+4},\pi^r)|.
}\end{equation} The analogous definition exchanges the spin factors for
\(U^-\). The exact encoders are
\begin{equation}\protect\phantomsection\label{eq:gg-original-encoders}{
 \begin{aligned}
 F_2[U^+]&=\sum_{r,a,j}
 \frac{U^+_{r;a,j}}{(r+4-a)!a!}
 (w^0)^{r+4-a}(w^1)^a\zeta^{-j-1},\\
 F_{-6}[U^-]&=\sum_{r,a,j}
 \frac{U^-_{r;a,j}}{(r-a)!a!}
 (w^0)^{r-a}(w^1)^a\zeta^{-j-1}.
 \end{aligned}
}\end{equation} In the first sum \(0\le a\le r+4\), \(0\le j\le r\); in
the second \(0\le a\le r\), \(0\le j\le r+4\). The correspondence is on
actual analytic quotients and has common smaller representatives and
coboundaries for bounded holomorphic families.

\emph{Proof.} At fiber degree \(n\), a weight-\(m\) overlap has a
Laurent expansion. The north coboundary removes powers \(h\ge0\) and the
south coboundary removes \(h\le m-n\). Exactly
\begin{equation}\protect\phantomsection\label{eq:gg-surviving-laurent}{
 m-n+1\le h\le-1
}\end{equation} remain. Thus the normal representative is
\[\sum_{n\ge\max(0,m+2)}\sum_{a=0}^n\sum_{j=0}^{n-m-2}
c^{(m)}_{n,a,j}(w^0)^{n-a}(w^1)^a\zeta^{-j-1}\]. The two tails have no
common surviving monomial, proving uniqueness.

Choose fixed nested circles on either side of \(|\zeta|=1\) and a
smaller fiber polydisk contained in both chart domains. Fiber Taylor
estimates and the Laurent Cauchy formula bound all surviving
coefficients by \(AB^n\) for finite \(A,B\). Conversely such a bound
makes the normal series converge on a smaller annular fiber
neighborhood, since \(a,j\) have only polynomially many choices and
\(|\zeta|^{-j-1}\) grows at most exponentially in \(n\). The removed
tails are Cauchy integrals on the two fixed circles. Their geometric
denominators are bounded by the buffers; after the south substitution
\(w=\zeta v\), its additional degree factors are at most exponential. A
single fiber restriction therefore makes both tails actual holomorphic
cochains. This proves the analytic quotient assertion.

The flat contour fields use the counterclockwise Laurent contour in the
north frame, with the ordered overlap \(F_{NS}=F\) and Čech differential
fixed by eq.~\ref{eq:gg-original-cech-orientation}. In this
trivialization its differential is
\(\mathrm d\zeta=-\pi_{A'}\,\mathrm d\pi^{A'}\). They are
\begin{equation}\protect\phantomsection\label{eq:gg-flat-contours}{
 \begin{aligned}
 C^+_{ABCD}(x)
 &=\frac1{2\pi i}\oint
   \partial_{w^A}\partial_{w^B}
   \partial_{w^C}\partial_{w^D}F_2(x\pi,\zeta)\,
   \mathrm d\zeta,\\
 C^-_{A'B'C'D'}(x)
 &=\frac1{2\pi i}\oint
   \pi_{A'}\pi_{B'}\pi_{C'}\pi_{D'}
   F_{-6}(x\pi,\zeta)\,\mathrm d\zeta .
 \end{aligned}
}\end{equation} Symmetry and an alternating spin contraction give their
flat massless equations. Differentiating at \(x=0\) shows that the
component multipliers are exactly \((r+4-a)!a!\) and \((r-a)!a!\),
respectively: the primed slots contribute \(\zeta^j\) and the residue
picks \(\zeta^{-j-1}\). This proves the formulas and their inverses.

For completeness, put the auxiliary Hermitian norms inherited from the
full spin tensor powers on symmetric tensors. Coherent evaluation gives
\begin{equation}\protect\phantomsection\label{eq:gg-coherent-norm}{
 \|T\|_{\mathrm{ev}}\le\|T\|
 \le\sqrt{(p+1)(q+1)}\,\|T\|_{\mathrm{ev}},
 \qquad
 T\in\operatorname{Sym}^pS^*\otimes\operatorname{Sym}^qS'^* .
}\end{equation} Indeed the squared coherent integral on the two unit
spheres is \(\|T\|^2/((p+1)(q+1))\), by phase orthogonality and
\(\int_0^1 t^j(1-t)^{p-j}\mathrm dt=j!(p-j)!/(p+1)!\). The reverse
inequality follows by the supremum and the forward one by
Cauchy--Schwarz. The factor for \(W_r^\pm\) is at most \(r+3\).

Consequently the factorial bound implies exponential bounds for both
normal coefficient arrays: use \(r!/((r-a)!a!)=\binom ra\) and
\(r!/((r+4-a)!a!)=\binom{r+4}a/((r+1)\cdots(r+4))\). Polynomial factors
are absorbed by one further exponential. Conversely an exponential
coefficient bound and the exact extraction multipliers bound every
component by a constant times \((r+4)!B^r\). The full tensor norm is at
most \(2^{r+2}\) times the largest component. Since \((r+4)!/r!\) is
polynomial, another fixed radius loss gives the stated \(r!\) bound. All
bounds are uniform on the fixed input buffers. Every splitting is a
bounded complex-linear operation there, so finite-dimensional
holomorphic dependence is preserved. ∎

The contour fields encode the entire free list at the mark. They are not
identified away from that mark with the nonlinear Weyl fields of the
final metric.

\pandocstatementheading{\protect\phantomsection\label{lem:gg-completion}Lemma 3.2.}

There is a unique polynomial ordered completion
\(\mathcal C_r(U,\Lambda)\) of the two free lists into Einstein
curvature jets. The original infinitesimal displacement and spin-frame
operations are polynomial derivations \(\mathsf H_a,\mathsf V_B\), with
\begin{equation}\protect\phantomsection\label{eq:gg-cartan-derivations}{\begin{aligned}
[\mathsf H_a,\mathsf H_b]&=-\mathsf V_{R_U(a,b)}, \\
[\mathsf V_B,\mathsf H_a]&=\mathsf H_{Ba}, \\
[\mathsf V_B,\mathsf V_C]&=\mathsf V_{[B,C]}, \\
\mathsf H_aU_r&=i_aU_{r+1}+Q_{r,a}(U,\Lambda), \\
\mathsf H_a\Lambda&=\mathsf V_B\Lambda=0, \\
\mathsf H_a\mathcal C_r&=\mathcal C_{r+1}(a;\,\cdot).
\end{aligned}}\end{equation} Here \(Q_0=0\), \(Q_r\) uses only
\(U_0,\ldots,U_{r-1},\Lambda\), and
\begin{equation}\protect\phantomsection\label{eq:gg-algebraic-curvature}{
 R_U(a,b)c=W_U(a,b)c+
       \frac{\Lambda}{3}\{a\,\eta(b,c)-b\,\eta(a,c)\}.
}\end{equation} Here \(\eta(W_U(a,b)c,d)\) is the tensor in
eq.~\ref{eq:gg-weyl-spinor-decomposition} with \(\Psi^\pm=U_0^\pm\).
This specifies the inverse component conversion in every curvature
action. The Bianchi identities and frame-equivariance identities hold
polynomially on the independent complex lists, with the curvature action
on every derivative slot retained. Complete free records and \(\Lambda\)
determine the formal normal metric uniquely.

\emph{Proof.} The complete proof is given in Appendix A. The highest
difference of two ordered completions has commuting derivative slots and
zero Bianchi contractions. The two-spinor identity makes it separately
symmetric, hence precisely the two free tensors of that order. Adding
their projection to the linear compatibility system makes its constant
matrix injective. Finite analytic Einstein witnesses for every real
finite prefix prove consistency. A fixed constant left inverse therefore
gives polynomial coefficients; maximal-real-form polynomial continuation
supplies independent complex inputs.

Every asserted derivation identity uses finitely many components and
holds on those witnessed prefixes as actual covariant differentiation
and frame change. It is consequently a polynomial identity. The
necessary finite prefix includes every derivative appearing in that
identity; no infinite analytic witness is used in this argument. The
radial Jacobi recurrence fixes the formal normal metric from the
completed ordered curvature. Appendix A supplies the witness
construction, all normal-form hypotheses and the formal uniqueness
argument. ∎

The infinitesimal frame convention is the readout at the changed frame:
\(\mathsf V_B\) inserts \(B\) in each tensor argument. Its finite
version is the pullback of the tensor by the specified spin element. The
first Bianchi identity and its differential form imply the Lie--Rinehart
Jacobi identity, including the curvature-dependent coefficients in the
first bracket.

Let \(X_0=(U_0^+,U_0^-,\Lambda)\) and \(X_r=(U_r^+,U_r^-)\) for \(r>0\),
and define
\begin{equation}\protect\phantomsection\label{eq:gg-record-banach-norm}{
 \|X\|_R=\sum_{r\ge0}\frac{R^r}{r!}\|X_r\|.
}\end{equation} The original bounds give
\(\|X\|_{R_0/2}\le16C+|\Lambda|\) by the preceding coherent estimate.
This is an initial radius conversion, not a norm change repeated at each
nonlinear derivative.

\pandocstatementheading{\protect\phantomsection\label{thm:gg-record-flows}Theorem 3.3.}

On every bounded set of this analytic-record space, the actual
\(\mathsf H_a\) have normally convergent flows into spaces of smaller
positive radius, for bounded \(a\) and small common complex time.
Inverses, derivatives and every specified finite word have common
positive domains. Finite-dimensional holomorphic parameter dependence
and bounded full-spin actions are preserved. For the coefficient
involution \(\mathfrak c\) of §8.4, the flows satisfy
\(\mathfrak c\Phi_a^t(X)=\Phi_{\mathfrak ca}^{\bar t}(\mathfrak cX)\).
Its fixed set is preserved for \(\mathfrak ca=a\) and real \(t\), and by
full Lorentz-real Spin actions.

\emph{Proof.} We give the coefficient estimate and scalar analytic
argument; Appendix B proves the estimate for the original derivation,
including its linear-transfer and insertion constants. With
\(X_0=(U_0^+,U_0^-,\Lambda)\) and \(X_j=(U_j^+,U_j^-)\) for \(j>0\),
write the recursively ordered expansion
\begin{equation}\protect\phantomsection\label{eq:gg-ordered-polynomial-expansion}{
 Q_{r,a}(X)=\sum_{d\ge2}\ 
 \sum_{\substack{\mathbf j=(j_1,\ldots,j_d)\in\mathbb N_0^d\\
                 |\mathbf j|=r+3-2d}}
 B^{[d]}_{r,\mathbf j}(a;X_{j_1},\ldots,X_{j_d}).
}\end{equation} The input positions retain their derivative orders and
chiralities; a scalar \(\Lambda\) position is a component of \(X_0\).
There is no \(1/d!\) symmetrization. The coefficients are the ordered
multilinear maps of Appendix B.2, grouped using the fixed auxiliary
direct-sum Hermitian norms. There are at most \(3^d\) input-type choices
and two output chiralities. Hence the passage from a typed coefficient
bound to a grouped bound costs at most \(2\cdot3^d\); replacing its
constants \(a,b\) by \(18a,3b\) absorbs this once, since
\(2\cdot3^d a b^{d-2}=18a(3b)^{d-2}\). For \(\|a\|\le1\) the grouped
coefficients therefore satisfy, for structural \(a_0,b_0>0\),
\begin{equation}\protect\phantomsection\label{eq:gg-actual-coefficient-budget}{
 \|B^{[d]}_{r,\mathbf j}\|
 \le a_0b_0^{d-2}r\,\frac{r!}{\prod_i j_i!},
 \qquad
 |\mathbf j|=r+3-2d,\qquad d\ge2.
}\end{equation} Appendix B computes the normalized one-derivative
recursion for \(\mathsf H_a\), uses one contraction mass for all
operations, and retains the inverse-order gain of each nonlinear
insertion. Its transfer costs within the recurrence for \(Q\) are
\begin{equation}\protect\phantomsection\label{eq:gg-linear-transfer-product}{
 c_r=\frac{r(r+4)}{r(r+4)-1},
 \qquad
 \prod_{r\ge1}c_r\le e^{11/18}.
}\end{equation} This accumulated cost is already included in the
displayed coefficient budget; the separate leading term is a tensor
contraction, so \(\|i_aU_{r+1}\|\le\|a\|\,\|U_{r+1}\|\).

For \(\|a\|\le A\) and initial bound \(M\), divide the record components
by \(r!\). Their formal flow series are dominated by the nonnegative
scalar recursion
\begin{equation}\protect\phantomsection\label{eq:gg-flow-majorant}{\begin{gathered}
\partial_\tau F
  =A\left(F_z+z\frac{\mathrm d}{\mathrm dz}V(z,F)\right), \\
V(z,w)=\frac{a_0zw^2}{1-b_0z^2w}, \\
F(0,z)=\frac M{1-z/R}.
\end{gathered}}\end{equation} The derivative of \(V(z,F)\) is total. To verify the
majorization, write \(X_r(\tau)=\sum_{k\ge0}X_{r,k}\tau^k\) and
\(F(\tau,z)=\sum_{k,r\ge0}f_{k,r}\tau^kz^r\). The initial norm bound
gives \(\|X_{r,0}\|/r!\le f_{0,r}=MR^{-r}\). Apply norms to each time
coefficient. If the coefficients through order \(k\) are dominated, the
linear shift is bounded by the corresponding coefficient of \(AF_z\),
and the ordered degree-\(d\) part is bounded by
\begin{equation}\protect\phantomsection\label{eq:gg-flow-coefficient-majorization}{
 A a_0b_0^{d-2}r
 [\tau^k]\!
 \sum_{j_1+\cdots+j_d=r+3-2d}
       \prod_{i=1}^d[ z^{j_i}]F(\tau,z)
 = [\tau^kz^r]\,A z\frac{\mathrm d}{\mathrm dz}
       \bigl(a_0b_0^{d-2}z^{2d-3}F^d\bigr).
}\end{equation} For fixed \(r\), this is a finite sum over \(d\) and the
ordered input list. The coefficient equation
\((k+1)X_{r,k+1}=[\tau^k](\mathsf H_aX)_r\) and the same scalar equation
for \(f_{k+1,r}\) therefore prove \(\|X_{r,k+1}\|/r!\le f_{k+1,r}\). All
coefficients of the scalar recursion are nonnegative, so induction in
\(k\), simultaneously for every \(r\), proves formal coefficientwise
domination. This argument precedes any assertion of convergence.

Let \(W_*=4M+1\) and choose \(0<\rho\le R/8\) with
\(16b_0\rho^2W_*\le1\). The finite characteristic equations are
\begin{equation}\protect\phantomsection\label{eq:gg-majorant-characteristics}{
 \dot z=-A(1+zV_w),\qquad
 \dot w=AzV_z,\qquad
 z(0)=\zeta,\quad w(0)=M/(1-\zeta/R).
}\end{equation} On \(|z|\le2\rho\), \(|w|\le W_*\), the denominator is
bounded away from zero and \(|zV_w|\le16a_0\rho^2W_*\),
\(|zV_z|\le5a_0\rho W_*^2\). In the scaled coordinates
\((z/\rho,w/W_*)\), the vector field and its first derivatives are
bounded. The ordinary integral equation is a contraction on a small
common time disk, with displacement less than \(1/8\). For
\(|\zeta|\le\rho\) its initial scaled second coordinate is less than
\(2/7\), so trajectories remain in the displayed domain.

On \(|\zeta|=3\rho/4\), the displacement bound
\(|z(\tau,\zeta)-\zeta|\le\rho/8\) and Rouché's theorem give a unique
simple inverse for \(|z_*|<\rho/2\). Thus
\(F(\tau,z_*)=w(\tau,\zeta(\tau,z_*))\) is jointly holomorphic and
bounded by \(W_*\). Its Taylor series is the nonnegative formal
recursion just used. The proved coefficientwise domination now makes the
double record series converge normally on strict interior time disks and
record radii, with termwise derivatives.

Each coefficient is a polynomial in finitely many initial records. The
same domination gives holomorphic maps into the smaller-radius sequence
space; alternatively any supplied finite-dimensional holomorphic family
is jointly holomorphic by uniform convergence. Formal reversal and
composition identities become analytic identities. A bounded spin flow
acts on \(2r+4\) slots with bound \(\exp(c(r+2)|tB|)\), absorbed by one
strict radius decrease. For a fixed finite word choose its successive
output bounds and radii as the next inputs. A finite number of positive
restrictions, also for the inverse word, has a positive minimum. The
polynomial tensor identities give
\(\mathfrak c(\mathsf H_aX)=\mathsf H_{\mathfrak ca}(\mathfrak cX)\).
Conjugating the convergent time series and using uniqueness of the
initialized flow proves the stated equivariance. Real directions, real
times and Lorentz-real Spin elements therefore preserve the fixed set. ∎

\protect\phantomsection\label{sec:gg-cartan}

\subsection{4. Cartan geometry and its intrinsic ray
description}\label{cartan-geometry-and-its-intrinsic-ray-description}

The record-flow theorem gives a finite construction of the geometry
needed to prove convergence of the generating law. It also supplies
physical point and spin operations before any ordinary counit is used.

\pandocstatementheading{\protect\phantomsection\label{prop:gg-finite-cartan}Proposition 4.1.}

The original analytic record and \(\Lambda\) determine a holomorphic
Einstein four-germ by record flows and finite matrix ODEs. Its entire
free record is the original one. For bounded holomorphic input families,
its normal coordinates and radial spin frames have common positive
domains. The construction is unique up to the prescribed marked
isometry.

\emph{Proof.} Choose a constant basis \(e_1,\ldots,e_{10}\) of
\(V\oplus\mathfrak{so}(V,\eta)\), write \(D_i=\rho(e_i)\) for the
derivations of Lemma 3.2 and \([e_i,e_j]=c_{ij}^{\,k}(U)e_k\) for their
Lie--Rinehart coefficients. These coefficients depend only on
\(U_0,\Lambda\) and fixed representation matrices.

Successively set \(U_n(x_1,\ldots,x_n)=\Phi_n^{x_n}(U_{n-1})\). Theorem
3.3 gives these maps on common smaller neighborhoods, including
derivatives and parameters. Construct columns \(\alpha_i\) by the finite
ODEs
\begin{equation}\protect\phantomsection\label{eq:gg-finite-matrix-transport}{\begin{gathered}
\partial_n\alpha_i=-C_n(U_n)\alpha_i\quad(i<n), \\
\alpha_i|_{x_n=0}=(\alpha_{n-1})_i, \\
\alpha_n=e_n, \\
(C_n)^k{}_l=c_{nl}^{\,k}.
\end{gathered}}\end{equation} We verify the anchor equation
\(\partial_i f(U_n)=\alpha_i^k(D_kf)(U_n)\) on every cylinder polynomial
and the zero-defect equations
\begin{equation}\protect\phantomsection\label{eq:gg-coframe-defect}{
 F_{ij}:=\partial_i\alpha_j-\partial_j\alpha_i
              +c(U_n)(\alpha_i,\alpha_j)=0.
}\end{equation} For the anchor residual \(Z_i(f)\), the defining flow
and matrix equation give \(\partial_nZ_i(f)=Z_i(D_nf)\), with zero
initial value. Induction makes every \(x_n\) Taylor coefficient zero.
Holomorphy gives \(Z_i(f)=0\). The mixed defects \(F_{ni}\) vanish
directly. For \(i,j<n\), differentiate the defect and use the anchor
equation. The terms outside \(-C_nF_{ij}\) are the coefficient Jacobi
expression
\begin{equation}\protect\phantomsection\label{eq:gg-coefficient-jacobi}{
 D_nc_{lm}^{\,k}+D_lc_{mn}^{\,k}+D_mc_{nl}^{\,k}
 +c_{lm}^{\,q}c_{nq}^{\,k}
 +c_{mn}^{\,q}c_{lq}^{\,k}
 +c_{nl}^{\,q}c_{mq}^{\,k}=0.
}\end{equation} Thus \(\partial_nF_{ij}=-C_nF_{ij}\) with zero initial
data. Ordinary finite linear uniqueness completes the induction.

At \(n=10\) the coframe matrix \(\alpha=(\alpha_i)\) is the identity at
the origin, regardless of the rank of the record orbit. It remains
invertible after a common shrink. Split it as \((\theta,\omega)\). The
zero-defect equations say
\begin{equation}\protect\phantomsection\label{eq:gg-realized-cartan}{
 \mathrm d\theta+\omega\wedge\theta=0,\qquad
 \mathrm d\omega+\omega\wedge\omega
       =\tfrac12R_U(\theta,\theta),\qquad
 \mathrm df(U)=\rho(\theta,\omega)f(U).
}\end{equation} The involutive rank-six vertical distribution
\(\ker\theta\) has a local four-dimensional holomorphic quotient.
Equivalently, take a transverse four-slice and extend it by the six
finite vertical flows; their derivative is a unit at the mark. The
identity \(\mathcal L_{Z_B}\theta=-B\theta\) makes
\(\eta(\theta,\theta)\) basic, hence a nondegenerate quotient metric
\(g\). The first equation identifies \(\omega\) with its Levi--Civita
connection. The second identifies its curvature with \(R_U\), whose
Ricci trace is \(\Lambda\eta\). The anchor equation gives every ordered
covariant curvature derivative, including its derivative-slot action,
and hence the entire original free record.

All inverse matrices are units at the fixed mark. On common smaller
polydisks they have norm at most twice their marked inverse bound. The
record-flow bounds, finite linear ODE estimates and finite inverse
theorem therefore give uniform coordinate and metric/inverse-metric
bounds. Apply the finite geodesic endpoint ODE, its inverse, and radial
spin parallel transport to obtain common normal coordinates and spin
frames, with \(g(0)=\eta\) and \(\partial g(0)=0\).

For uniqueness, equal free lists give equal completed ordered curvature
jets. In marked exponential coordinates and parallel frames, the
curvature along each radial geodesic has identical Taylor series. The
radial Jacobi equation \(J_v''+R(J_v,x)x=0\), \(J_v(0)=0\),
\(J_v'(0)=v\) then gives identical endpoint Jacobi fields, and
\(g_x(v,w)=\eta(J_v(1),J_w(1))\). An isometry with the prescribed point
and tangent map must intertwine these exponential maps, proving
uniqueness. This argument works formally over parameter coefficient
rings as well. ∎

For any given analytic Einstein metric its free records satisfy the
original factorial bound. Evaluate its Weyl spinors along the geodesics
with initial velocities \(\xi\otimes\pi\), transporting the frame. The
unit spinor pairs are compact, so finite holomorphic ODE estimates give
one affine disk and curvature bound. Cauchy's formula gives exactly the
bound of Theorem 3.1. Thus Proposition 4.1 covers all marked analytic
Einstein germs.

In the selected Einstein scale the canonical rank-six metric and
connection are
\begin{equation}\protect\phantomsection\label{eq:gg-canonical-tractor}{\begin{aligned}
H((\sigma,\mu,\rho),(\sigma',\mu',\rho'))
   &=\sigma\rho'+\rho\sigma'+g(\mu,\mu'), \\
D_v(\sigma,\mu,\rho)
   &=\bigl(v\sigma-g(v,\mu),\\
 &\quad \
      \nabla_v\mu+((\Lambda/6)\sigma+\rho)v,\\
 &\quad \
      v\rho-(\Lambda/6)g(v,\mu)\bigr), \\
X&=(0,0,1),\qquad I=(1,0,-\Lambda/6).
\end{aligned}}\end{equation} In particular \(DI=0\), \(H(I,I)=-\Lambda/3\),
\(H(X,X)=0\), \(H(I,X)=1\), and the Gram determinant of \((I,X)\) is
\(-1\). These formulas identify the normal standard-tractor connection:
for a general auxiliary scale its middle coefficient is the Schouten
tensor, which here is \((\Lambda/6)g\), and its curvature has Weyl
middle block and zero Cotton block. It is not an arbitrary connection
preserving a vector.

The null-geodesic contact space and its complete family of normal sky
quadrics are the classical correspondence of (LeBrun 1983, III.1--3).
Propositions 4.2--4.3 give its local representatives together with the
calibrated rank-six bundle data used here.

\pandocstatementheading{\protect\phantomsection\label{prop:gg-ray-construction}Proposition 4.2.}

The metric of Proposition 4.1 has a holomorphic contact ray neighborhood
of its whole compact sky, its canonical descended rank-six bundle and
the affine line of \(X\) along every ray. They form an object of
\(\mathfrak S_\Lambda\). No chirality or algebraic-type restriction is
required.

\emph{Proof.} On its frame bundle fix a nonzero null \(k\). For
\(\alpha_k=\eta(k,\theta)\), the torsion equation gives
\begin{equation}\protect\phantomsection\label{eq:gg-ray-levi-pairing}{
 \begin{aligned}
 \alpha_k(X_v+Z_B)&=\eta(k,v),\\
 \mathrm d\alpha_k(X_v+Z_B,X_w+Z_C)
       &=\eta(Bk,w)-\eta(Ck,v).
 \end{aligned}
}\end{equation} The \(\mathfrak{so}(V)\) orbit tangent of \(k\) is
\(k^\perp\): for \(w\perp k\) choose \(\ell\) with \(\eta(\ell,k)=1\)
and use \(Bv=w\eta(\ell,v)-\ell\eta(w,v)\). Thus the characteristic
kernel of this pairing on \(\ker\alpha_k\) is
\begin{equation}\protect\phantomsection\label{eq:gg-ray-characteristic}{
 \mathcal K_k=\mathbb CX_k
       \oplus\{Z_B:Bk\in\mathbb Ck\}.
}\end{equation} It has dimension five. The vertical stabilizer is a Lie
algebra, and its bracket with \(X_k\) is a multiple of \(X_k\); hence it
is involutive. The quotient hyperplane has nondegenerate Levi form,
proving contactness. Curvature has not been restricted in this
calculation.

For a whole-sky neighborhood use the projective null-direction bundle
over a small four-neighborhood. Choose finitely many compact direction
cores with larger angular parents. On each parent choose a fixed
coordinate covector \(b\) nonzero on its directions and normalize the
initial affine tangent by \(b(k)=1\). Holomorphic geodesic flows and
their direction lifts exist on one buffered time disc \(|t|<T\). After
shrinking these parents, their spacetime curves satisfy
\(|(b\circ\gamma)'(t)-1|\le1/2\) there. The hit on \(b(x)=0\) is
therefore unique and holomorphic, with holomorphic inverse flow
coordinates. Denote the transverse hit map, including the direction at
that hit, by \(\chi\).

We also reserve a connected-fiber margin before forming the quotient.
Fix the derivative bounds on larger spatial and angular parents first.
With \(u=b(x)\), write the spacetime projection of each relevant leaf as
\(G_z(u)\), where \(z\) is its transverse hit. On a whole convex parent
\(u\)-disc the finite ODE bounds give \(|G_z'|\ge a>0\),
\(|G_z''|\le C\) and \(|G_z|\le R_p\). Choose that parent small enough
that \(R_pC<a^2/2\). The real Hessian of \(|G_z|^2\) is then bounded
below by \(2(a^2-R_pC)\) times the Euclidean form. Choose an inner
spatial ball \(B_r\) with \(\|b\|r\) below half the free-coordinate
parent radius and \(4\|b\|r<T/2\). Every point over \(B_r\) in a
direction core has a hit in the retained transverse parent, by the same
inverse-flow bound. Thus each graph sublevel \(|G_z(u)|<r\) is convex
and connected on its whole parent. In a normal spin frame the connection
is \(O(|x|)\), so the direction change along its \(O(r)\) coordinate
segments is \(O(r^2)\); choose \(r\) so this is smaller than every
angular buffer. A leaf starting in a core cannot leave this graph while
remaining over \(B_r\): its hit is constant, its free coordinate stays
strictly inside its parent, and its direction stays inside the larger
angular parent. Uniqueness of the flow then keeps every continuation on
that graph. Conversely the convex free-coordinate segment joins any two
graph points over \(B_r\).

Let \(F_{\rm in}\) be the projective null-direction bundle over \(B_r\),
and let \(R_{\rm short}\) mean membership in the same connected
characteristic leaf in \(F_{\rm in}\). The preceding argument identifies
each such leaf with one whole graph sublevel. For two related points
whose initial direction is in a fixed core, let \(t\) be their
connecting affine time in its parent. Integration along the complex time
segment gives
\begin{equation}\protect\phantomsection\label{eq:gg-short-ray-time-buffer}{
 |t|\le2|b(\gamma(t)-\gamma(0))|
       \le4\|b\|r<T/2.
}\end{equation} Consequently \(R_{\rm short}\) is relatively closed in
\(F_{\rm in}\times F_{\rm in}\). Indeed, for a convergent sequence of
related pairs choose a subsequence whose first directions lie in one
fixed core. Their connecting times have a convergent subsequence in the
compact disc \(|t|\le T/2\), and the ODE depends continuously on the
initial point and time there. The buffered hit maps extend at both
limits; from \(\chi(p_n)=\chi(q_n)\) we obtain \(\chi(p)=\chi(q)\).
Their common graph sublevel is connected, so the limiting pair is again
related. This asserts relative closedness of the relation, not
closedness of an open inner plaque in a larger parent.

The quotient projection is open, since saturation of an open set is the
union of its images under local holomorphic leaf flows. For an unrelated
pair, relative closedness supplies a product neighborhood disjoint from
\(R_{\rm short}\); its two open quotient images are disjoint. The
quotient is Hausdorff. The hit maps are its holomorphic submersion
charts near the central sky, and their transition maps are the actual
hits on the next transversal along the same short leaf. Uniqueness gives
inverse maps and exact cocycles, without a further long-ray
identification. All choices use finitely many ODE bounds, unit hit
derivatives and angular buffers; they are uniform on the conditioned
families used in §6.1. This is the one-dimensional characteristic-flow
version of the graph and no-exit construction in §E.6.

Directions at one point form \(Q=\mathbb P(S)\times\mathbb P(S')\) and
embed in the quotient. Its tangent is killed by \(\alpha_k\), so it is
Legendrian; the contact line is the dual tautological null line
\(\mathcal O(1,1)\). The contact normal sequence has kernel
\(\Omega_Q^1\otimes\mathcal O(1,1)=\mathcal O(-1,1)\oplus\mathcal O(1,-1)\)
and quotient \(\mathcal O(1,1)\). The kernel has zero cohomology in
every degree. Hence \(H^0(N_{Q/Z})=H^0(\mathcal O(1,1))\) has dimension
four and \(H^1(N_{Q/Z})=0\). The compact-support completeness of Lemma
D.1 gives a universal smooth four-dimensional embedded family. The
rulings persist: each ruling line has trivial normal line in \(Q\) and
vanishing first obstruction, and the two resulting fibrations identify
nearby surfaces with \(\mathbb P^1\times\mathbb P^1\). Their contact
line retains bidegree \((1,1)\). On each surface the contact restriction
is a section of \(\Omega_Q^1\otimes\mathcal O(1,1)\), whose \(H^0\) is
zero, so all these surfaces are Legendrian. The point-family derivative
is \(v\mapsto[k\mapsto g(k,v)]\), an isomorphism onto that
four-dimensional section space. The finite inverse theorem therefore
identifies it with the complete family.

Parallel transport of the canonical bundle along the same short ray
segments gives holomorphic rank-six transitions and exact cocycles.
Evaluation at the point trivializes it on each sky. In its ray-parallel
trivialization, \(D_k^2X=0\) for an affinely parametrized null geodesic,
directly from the displayed canonical connection. Thus \(X\) traces an
affine line whose translation is the tautological null line, restricting
to \(\mathcal O(-1,-1)\). The selected \(I\) is its same descended
parallel section. The canonical curvature gives both normality
conditions. ∎

\pandocstatementheading{\protect\phantomsection\label{prop:gg-intrinsic-recognition}Proposition 4.3.}

Every independently supplied object of \(\mathfrak S_\Lambda\) is the
canonical ray geometry of its recovered analytic Einstein metric. The
identification retains the actual selected section and the full spin
mark, and is fully faithful.

\emph{Proof.} First reconstruct the bundle and connection. Proper base
change in the form for a holomorphic bundle on a proper compact-fiber
family gives locally free \(p_*q^*\mathcal V\) and the evaluation
isomorphism near the central sky (Grauert 1960, sec. 7, Sätze 3, 5--6;
Grauert 1963). The base is a smooth germ, the incidence is proper and
smooth, and the pulled-back bundle is coherent and flat over that base.
Here the central bundle is \(\mathcal O_Q^{\,6}\) and
\(H^j(Q,\mathcal O_Q)=0\) for \(j>0\). Upper semicontinuity makes these
higher groups vanish nearby; the Euler characteristic is six, so
\(h^0=6\). Evaluation is an isomorphism on the compact central fiber and
therefore on a common nearby family. This verifies all required
hypotheses. For a finite local Artin extension, factor it into finitely
many small extensions \(0\to J\to A'\to A\to0\) with
\(\mathfrak m_{A'}J=0\). At each step the kernel is
\(J\otimes_{\mathbb C}\mathcal O_Q^{\,6}\); its \(H^1\) vanishes, so the
six sections lift and their basis evaluation remains a unit. Finite
induction gives the assertion for the original extension.

The affine torsor on a sky has translation \(\mathcal O(-1,-1)\), whose
\(H^0\) and \(H^1\) both vanish. It has exactly one section. Its image
in the sky-trivial bundle is constant on the compact sky, defining
\(X\in\mathcal T\). Analytic dependence can also be seen without a
torsor base-change assertion: in an evaluation frame the
affine-membership equations are linear, their homogeneous kernels have
zero intersection, and finitely many rows have rank six at the central
sky. Solve these unit rows; fiberwise uniqueness verifies all the
remaining equations. This works over finite parameter algebras.

For the inverse transport, the point-family tangent is
\(H^0(Q,\mathcal O(1,1))\). At a sky point, the projected ray line is
the kernel of its first-jet evaluation; a section with zero value and
first derivative is the product of the two linear factors vanishing at
that point. These kernels form \(\mathcal O(-1,-1)\) and give the
complete Segre null quadric. In an evaluation frame the given
\(q\)-constant partial connection, divided by this tautological ray
direction, has coefficients in
\(H^0(Q,\operatorname{End}\mathcal T\otimes\mathcal O(1,1))\). Each is a
unique linear function of the tangent direction. It consequently defines
one connection \(D\) on \(\mathcal T\), with the usual Leibniz rule. Its
\(q\)-parallel transport is the supplied one. The descended \(H,I\)
satisfy \(DH=DI=0\) since their derivatives vanish in every null
direction.

Put \(g(v,w)=H(\theta(v),\theta(w))\) and project \(D\) to
\(W_{\mathrm{scr}}=(I,X)^\perp\). The nondegenerate Gram matrix of
\(I,X\) makes this projection a holomorphic algebraic operation, even at
\(\Lambda=0\). Pairing derivatives with \(I,X\) shows
\begin{equation}\protect\phantomsection\label{eq:gg-recognition-connection}{
 D_v\theta(w)
  =\theta(\nabla_vw)-g(v,w)I
                         -(\Lambda/3)g(v,w)X.
}\end{equation} The projected connection, defined by
\(\theta(\nabla_vw)=\operatorname{pr}_{W_{\mathrm{scr}}}D_v\theta(w)\),
is metric. Since \(H\) and the scale plane are nondegenerate, so are
\(H|_{W_{\mathrm{scr}}}\) and \(g\). Antisymmetrizing
eq.~\ref{eq:gg-recognition-connection} gives the correctly typed torsion
identity
\begin{equation}\protect\phantomsection\label{eq:gg-recognition-torsion-type}{
 \theta\bigl(T^\nabla(v,w)\bigr)
 =D_v\theta(w)-D_w\theta(v)-\theta([v,w])
 =F^D(v,w)X.
}\end{equation} It vanishes by regular-normality, so \(\nabla\) is
Levi--Civita. Alternating two derivatives of
eq.~\ref{eq:gg-recognition-connection} and projecting to
\(W_{\mathrm{scr}}\) gives the ambient middle curvature of Definition
2.2:
\begin{equation}\protect\phantomsection\label{eq:gg-recognition-curvature}{
 R^g(v,w)u-\frac{\Lambda}{3}
                 \{g(w,u)v-g(v,u)w\}.
}\end{equation} Taking the trace of
\(v\mapsto\mathcal R_{\mathrm{mid}}(v,w)u\) in dimension four gives
\(\operatorname{Ric}_{\mathrm{mid}}=\operatorname{Ric}(g)-\Lambda g\).
Its vanishing is exactly \(\operatorname{Ric}(g)=\Lambda g\). In the
basis \(X\), \(I+(\Lambda/6)X\) and \(W_{\mathrm{scr}}\) the connection
is the canonical one above.

Under the given \(q\)-parallel identification, \(X\) stays in its
prescribed affine line. Its tangent is the corresponding null direction.
Differentiating once more in an affine parameter and using the
connection formula proves that the projected ray is the null geodesic of
\(g\). The contact structure is the canonical Hilbert contact structure
as well. Both contain the tangent planes of all skies through a ray. The
planes from two distinct sufficiently close points span the
four-dimensional contact plane: their Jacobi-field pairing has leading
term \((t_2-t_1)\eta_{\mathrm{screen}}\), which is invertible for small
nonzero separation. Hence the two contact hyperplanes agree. The
supplied ray bundle and affine line are identified by their same null
transport, not by replacement with another connection.

An enriched comparison induces the recovered isometry. Conversely an
isometry transports its null rays, canonical bundle, \(I\) and affine
line. There is no further bundle ambiguity: \(I,X\) and the four solder
vectors span rank six with nonzero Gram determinant. The initial spin
lift is retained separately, including its central kernel. This proves
full faithfulness. ∎

Section 5 defines the source gluing law; §6 uses the Cartan construction
above to prove its common convergence.

\protect\phantomsection\label{sec:gg-generation}

\subsection{5. Holomorphic generation and the compatible optical
correction}\label{holomorphic-generation-and-the-compatible-optical-correction}

We now define the source law. All its finite stages use the original
classes and fixed holomorphic operations. Geometry recovered from those
stages is derived data.

We distinguish coordinate-algebra comparisons from point maps. The
overlap operator is \(E_{ij}:\mathcal O_j\to\mathcal O_i\), whereas the
stored point map is \(\Phi_{ij}:Z_j\to Z_i\), with
\begin{equation}\protect\phantomsection\label{eq:gg-overlap-operator-point-convention}{
 E_{ij}=(\Phi_{ij}^{-1})^*,\qquad
 E_{ij}=\exp(\varepsilon X_{h_{ij}})
       \ \Longleftrightarrow\
 \Phi_{ij}=\operatorname{Fl}_{-\varepsilon X_{h_{ij}}}.
}\end{equation} All logarithms, Baker--Campbell--Hausdorff products and
Hamiltonian correction cochains below refer to \(E\). All graph actions,
contact multipliers and optical transitions refer to the actual point
maps \(\Phi\). Proposition 5.8 verifies the sign against the incidence
metric.

Start with homogeneous \(Z=(\omega^A,\pi_{A'})\) on \(T_0\) and its
marked linear dual \(W=(\kappa_A,\eta^{A'})\). The open incidence
\(W\cdot Z=0\), \(\pi\ne0\), \(\kappa\ne0\) is an open part of
\(\mathbb PT^*T_0\), with contact form \(W\,\mathrm dZ\). Its zero sky
is \(Q=\mathbb P(S^*)\times\mathbb P(S')\). The dual structure is fixed
by the original linear marking.

The original negative class determines, rather than accompanies, a dual
positive class. If its coefficients in Theorem 3.1 are \(c^-_{r,a,j}\),
set
\begin{equation}\protect\phantomsection\label{eq:gg-derived-dual-class}{
 G(\eta,\varrho)=
 \sum_{r\ge0}\sum_{a=0}^r\sum_{j=0}^{r+4}
 (-1)^r\frac{(r-a)!a!}{(r+4-j)!j!}\,
 c^-_{r,a,j}(\eta^0)^{r+4-j}(\eta^1)^j\varrho^{-a-1}.
}\end{equation} The sign is fixed by the dual incidence
\(\eta=-\kappa x\). Both extraction multipliers are those in Theorem
3.1, so this preserves every \(U_r^-\). This is a bijection of analytic
germs: for \(|c^-_{r,a,j}|\le M\delta^{-r}\) and
\(1/2\le|\varrho|\le2\), with \(b=|\eta^0|+|\eta^1|\),
\begin{equation}\protect\phantomsection\label{eq:gg-dual-convergence}{
 |G|\le\frac{Mb^4}{12(1-2b/\delta)^2},
 \qquad 2b<\delta.
}\end{equation} Indeed
\(\sum_a(r-a)!a!|\varrho|^{-a-1}\le(r+1)r!2^{r+1}\), the remaining
monomial sum is \(b^{r+4}/(r+4)!\), and \(r!/(r+4)!\le1/24\). Conversely
the inverse coefficient multiplier is bounded by
\((r+1)(r+2)(r+3)(r+4)2^r\). Fixed-radius Cauchy splitting gives actual
cochains in both directions. The same inequalities prove the family
statement.

In the chart \(\pi_{0'}=\kappa_0=1\), put \(\zeta=\pi_{1'}\),
\(\varrho=\kappa_1\), \(s=\eta^{1'}\). Then
\(\eta^{0'}=-w^0-\varrho w^1-\zeta s\) and
\begin{equation}\protect\phantomsection\label{eq:gg-cotangent-contact-form}{
 \alpha=\mathrm dw^0+\varrho\,\mathrm dw^1+s\,\mathrm d\zeta.
}\end{equation} For a Hamiltonian \(h\), the exact contact field is
\begin{equation}\protect\phantomsection\label{eq:gg-contact-hamiltonians}{
 \begin{aligned}
 X_h={}&(h-\varrho h_\varrho-sh_s)\partial_{w^0}
 +h_\varrho\partial_{w^1}+h_s\partial_\zeta\\
 &+(\varrho h_{w^0}-h_{w^1})\partial_\varrho
 +(s h_{w^0}-h_\zeta)\partial_s,\\
 \{h,k\}&=X_hk-h_{w^0}k,\qquad [X_h,X_k]=X_{\{h,k\}}.
 \end{aligned}
}\end{equation} The two first-order lifts are
\begin{equation}\protect\phantomsection\label{eq:gg-two-contact-lifts}{
 h_+=\partial_{w^1}F_2-\varrho\partial_{w^0}F_2,\qquad
 h_-=-\partial_{\eta^{1'}}G+\zeta\partial_{\eta^{0'}}G.
}\end{equation} They start in normal degree three and are contact-line
cochains on the product cover. The positive lift is the cotangent lift
of the Hamiltonian deformation on \(T_0\); the negative lift is its
determined dual lift. Their cohomology classes have exactly the two
original free lists, with the normalization verified in Appendix C. They
have a mixed interaction: for \(F=(w^0)^4/\zeta\) and
\(G=(\eta^{0'})^4/\varrho\) the bracket is
\(48(w^0)^2(\varrho w^1+\zeta s)(w^0+\varrho w^1+\zeta s)^2\), which is
not zero.

Use an amplitude \(\varepsilon\), with prescribed complete record
\(\varepsilon U\) and numerical coefficient
\(\lambda(\varepsilon)=\varepsilon\Lambda\). At finite order work over
\(A_N=\mathbb C[\varepsilon]/(\varepsilon^{N+1})\); finite parameter
algebras with the same flat residue can be used for the same induction.
Products, logarithms and exponentials relative to the reference atlas
are finite nilpotent operations. No logarithm of an arbitrary analytic
contact map is asserted to converge.

\pandocstatementheading{\protect\phantomsection\label{lem:gg-contact-extension}

\protect\phantomsection\label{lem:gg-contact-extension-cubic}Lemma 5.1.}

Every finite source contact stage in the cubic Hamiltonian gauge
\(H_{ij}\in\mathcal I_Q^3\mathcal L\) admits a contact-cocycle lift in
that same gauge to the next parameter order, with holomorphic spatial
coefficients. Here \(\mathcal I_Q\) is the ideal of the reference zero
sky, and \(H_{ij}\) is the nilpotent logarithm relative to the reference
atlas. Closed contact corrections in each normal degree have the
cohomology described in Appendix C. The prescribed source execution
remains in this gauge at every stage.

\emph{Proof.} The flat normal bundle has
\(0\to K\to V\otimes\mathcal O_Q\to N\to0\), where
\(K=\mathcal O(-1,-1)\) and the contact line is \(\mathcal L=K^*\).
Degree-\(n\) contact Hamiltonians form
\(\mathcal F_n=\mathcal L\otimes\operatorname{Sym}^nN^*\). For \(n\ge1\)
the normalized derivative \((1/n)\partial_k\) gives
\begin{equation}\protect\phantomsection\label{eq:gg-contact-normal-sequence}{
 0\longrightarrow\mathcal F_n
 \longrightarrow\mathcal L\otimes\operatorname{Sym}^nV^*
 \longrightarrow\mathcal L^2\otimes\operatorname{Sym}^{n-1}V^*
 \longrightarrow0 .
}\end{equation} Both right-hand bundles have no higher cohomology, hence
\(H^2(\mathcal F_n)=0\). Appendix C gives a fixed explicit two-cocycle
repair with norm at most \(C_0(n+1)C_1^n\) on buffered charts. It
therefore repairs an entire analytic normal series after one
restriction.

Lift the nilpotent contact logarithms by coefficient lifts within
\(\mathcal I_Q^3\mathcal L\). This is an invariant source gauge. In
adapted coordinates \(\alpha=dz+p\,dq\), with \(\mathcal I_Q=(z,p)\),
the Jacobi formula gives
\begin{equation}\protect\phantomsection\label{eq:gg-cubic-source-filtration}{
 \{\mathcal I_Q^a\mathcal L,\mathcal I_Q^b\mathcal L\}
       \subset\mathcal I_Q^{a+b-1}\mathcal L,
       \qquad a,b\ge1.
}\end{equation} Tangential derivatives preserve normal order, normal
derivatives lower it by at most one, and a term with two normal
derivatives has a normal-coordinate factor. Thus cubic Hamiltonians are
closed under the contact brackets and finite Baker--Campbell--Hausdorff
operations. Reference substitutions preserve the ideal. The triple
defect is consequently a cubic contact Hamiltonian two-cocycle with
coefficients in the highest kernel \(J\); associativity makes it closed
because \(\mathfrak mJ=0\). Its degree-preserving repair uses only
\(n\ge3\). Subtracting that repair gives the required exact cocycle,
with literal inverse orientations and holomorphic coefficients.

The initial source lifts displayed above are cubic. The
zero-new-coefficient trial rule preserves this property. The optical
correction has degree \(n=m+1\ge3\), and its fixed representatives
preserve that degree. Induction therefore proves the gauge assertion for
the actual source execution. This statement concerns the
source-correction representatives; the later stored normalizer is
compared by its own deterministic metric-normalization argument. Only
finitely many derivatives and the convergent fixed spatial splittings
occur at each parameter order. ∎

The cubic condition is stronger than fixing the sky and the point-map
first jet. For example, \(h=z^2\) has flow
\((z,p,q)\mapsto(z/(1-tz),p/(1-tz)^2,q)\) for \(\alpha=dz+p\,dq\). It
has identity first point jets along the sky while its Hamiltonian is
quadratic. The source gauge in the lemma is established by the preceding
induction, rather than deduced from those point jets.

The compatible optical correction is a further condition; it does not
follow from this contact lemma alone. To state it we first construct its
incidence and bundle from the maps.

For contact graph coordinates put \(q=(\varrho,\zeta)\),
\(u=w^0+\varrho w^1\), \(p=(w^1,-s)\), so
\(\alpha=\mathrm du-p\,\mathrm dq\). A scalar \(f(q)\) represents its
Legendrian graph \(j^1f\). For a contact map \(\Phi\), let
\(\psi_{\Phi,f}(q)=\operatorname{pr}_q\Phi(j^1f(q))\) and define its
scalar graph action by
\(\mathscr A_\Phi(f)(\psi_{\Phi,f}(q))=\operatorname{pr}_u\Phi(j^1f(q))\).
The direction projection is a unit at the reference and therefore
invertible at each finite stage. These are exact actions under
composition.

Fix four direction points \((0,0),(1,0),(0,1),(1,1)\) and their line
frames. Let \(\mathfrak e\) be their four evaluations followed by the
inverse evaluation matrix of \(1,\varrho,\zeta,\varrho\zeta\). The flat
scalar sky is
\(s_x=x^{00'}+\varrho x^{10'}+\zeta x^{01'}+\varrho\zeta x^{11'}\).

\pandocstatementheading{\protect\phantomsection\label{prop:gg-finite-incidence}Proposition 5.2.}

Every finite contact-cocycle stage has a unique scalar family satisfying
\(\mathscr A_{\Phi_{ij}}f_j=f_i\) and \(\mathfrak e(f)=x\). It is a
complete family of embedded polarized Legendrian quadrics over its
finite algebra. Its conformal metric and actual ray derivatives are
determined by this same family.

\emph{Proof.} Write \(f_s=f+sv\) and
\((Q_s,U_s,P_s)=\Phi(q,f_s(q),\mathrm df_s(q))\), with
\(F_s=\mathscr A_\Phi(f_s)\). At fixed old direction \(q\), the contact
identity \(\Phi^*\alpha=c_\Phi\alpha\) gives
\(\delta U-P\,\delta Q=(c_\Phi\circ j^1f)v\), where
\(\delta=\partial_s|_{s=0}\). The image is Legendrian, so
\(U_s=F_s(Q_s)\) and \(P_s=\mathrm dF_s(Q_s)\). Differentiating the
first identity and subtracting \(P\,\delta Q\) therefore gives
\((\mathrm d\mathscr A_\Phi|_f(v))\circ\psi_{\Phi,f}
=(c_\Phi\circ j^1f)v\). Consequently, at fixed new direction,
\begin{equation}\protect\phantomsection\label{eq:gg-graph-linearization}{
 \mathrm d\mathscr A_\Phi|_f(v)
       =\bigl((c_\Phi\circ j^1f)v\bigr)
                         \circ\psi_{\Phi,f}^{-1}.
}\end{equation} The terms involving \(\mathrm dv\) in \(\delta Q\) and
\(\delta U\) cancel in this contact variation.

Thus at the highest parameter degree the new graph residual is an
\(\mathcal O(1,1)\) one-cocycle, without a direction derivative of the
highest unknown. The tensor product of the fixed two-circle Laurent
contractions solves this cocycle. Subtract its global section with the
same four evaluations. This produces the unique primitive of zero
four-evaluation; \(H^1(\mathcal O(1,1))=0\) gives existence and the
four-dimensional \(H^0\) gives uniqueness. Iterate over the finite
parameter algebra. The source terms and graph inverses are finite
analytic operations, so the coefficients are holomorphic on a possibly
order-dependent spatial neighborhood.

At each finite extension the abstract direction-chart change is a
cocycle in \(T_Q=\mathcal O(2,0)\oplus\mathcal O(0,2)\), whose \(H^1\)
vanishes. The contact-line deformation is trivial because
\(H^1(\mathcal O_Q)=0\). Hence the abstract polarized type remains
\((Q,\mathcal O(1,1))\); this does not trivialize its embedded contact
deformation. The four variations \(\partial_{x^a}f_i\) are its complete
contact-line section system.

Its inverse conformal metric can be read by one finite matrix inverse.
Evaluate \(\ell_\nu=(\partial_{x^a}f_i)\) at the nine directions
\((\varrho,\zeta)\in\{-1,0,1\}^2\), in fixed frames, and solve
\begin{equation}\protect\phantomsection\label{eq:gg-nine-direction-metric}{
 G^{ab}(\ell_\nu)_a(\ell_\nu)_b=0\quad(\nu=1,\ldots,9),
 \qquad \eta_{ab}G^{ab}=4 .
}\end{equation} At the residue the restrictions of quadratics to
\((1,\varrho,\zeta,\varrho\zeta)\) form all nine bidegree-\((2,2)\)
polynomials. The tensor product of the two three-point Vandermonde
matrices is invertible, and its quadratic kernel is the single null-form
line. The last normalization is nonzero on that line. The ten-by-ten
matrix is consequently a unit. The actual moving quadric has this same
one-dimensional quadratic ideal, so \(G\) is its inverse conformal form,
not a fit to only nine directions. Put \(g_{\mathrm{aux}}=G^{-1}\).

We also identify the actual ray curves before imposing an Einstein
condition. Write \(g=g_{\mathrm{aux}}\) in this calculation. Let \(F\)
be this complete incidence, with projections \(p:F\to M\) and
\(q:F\to Z\). A local contact frame \(\alpha\) pulls back to a one-form
vanishing on \(\ker dp\), because the skies are Legendrian. It therefore
defines a nonzero covector \(\ell_f\) at \(p(f)\) by
\(q^*\alpha(Y)=\ell_f(dp\,Y)\). On every whole sky its projective class
is the evaluation embedding of the complete \(H^0(\mathcal O(1,1))\)
system. Thus
\begin{equation}\protect\phantomsection\label{eq:gg-trial-characteristic-null-incidence}{
 \iota:F\xrightarrow{\ \sim\ }
 \mathcal C_g=\{(x,[p])\in\mathbb P(T^*M):
                         g^{ab}p_ap_b=0\},\qquad
 \iota(f)=(p(f),[\ell_f]).
}\end{equation} This is a fiberwise Segre isomorphism over \(M\), with
unit coordinate minors over the finite algebra. It carries the
pulled-back contact value to the tautological covector value. On
\(\ker q^*\alpha\), the characteristic line of
\(d(q^*\alpha)=q^*d\alpha\) is precisely \(\ker dq\), since \(q\) is a
submersion and the contact Levi form on \(Z\) is nondegenerate.

For the affine normalization, lift \(\iota\) by the covector \(\ell_f\)
to the nonprojective null cone in \(T^*M\). Write
\(\vartheta=p_a\,dx^a\) and \(H_g=g^{ab}p_ap_b/2\) there. The Hamilton
field with \(\iota_{X_{H_g}}d\vartheta=-dH_g\) is
\begin{equation}\protect\phantomsection\label{eq:gg-trial-affine-hamiltonian}{
 X_{H_g}=g^{ab}p_b\,\partial_{x^a}
          -\tfrac12(\partial_a g^{bc})p_bp_c\,\partial_{p_a}.
}\end{equation} On the null cone \(\vartheta(X_{H_g})=2H_g=0\), and
\(-dH_g\) vanishes on its tangent space. Its nonzero projective image
therefore spans the characteristic line already identified above; its
spacetime projection is \(g^{-1}p\). Let \(D\) be the generator of
\(\ker dq\) with \(dp(D)=g^{-1}\ell_f\); its scaling exists and is
unique by the preceding characteristic identification. In the covector
lift, \(\iota_Dd\vartheta=q^*(\iota_{dqD}d\alpha)=0\) on every incidence
tangent, while \(\iota_{X_{H_g}}d\vartheta=-dH_g\) also pulls back to
zero. Their difference has zero spacetime projection, so its contraction
with \(d\vartheta\) is a vertical covector variation
\(\delta p_a\,dx^a\). Since \(p:F\to M\) is a submersion, this can
vanish on all incidence tangents only if \(\delta p=0\). The lift of
\(D\) is therefore exactly \(X_{H_g}\), rather than that field plus an
Euler rescaling. Its spacetime equation is the affinely parametrized
null-geodesic equation of \(g_{\mathrm{aux}}\).

This argument uses regular contact incidence and the unit conformal
metric only. All its equalities, kernels and inverse minors commute with
finite parameter reduction, so it applies to the trial stage as well as
a corrected Einstein stage. ∎

\pandocstatementheading{\protect\phantomsection\label{prop:gg-intrinsic-optical}Proposition 5.3.}

The regular complete contact incidence determines a rank-two holomorphic
bundle \(\mathcal E_J\) independently of a chosen scalar metric. Its
solutions are defined by a projective ray operator on the intrinsic
density line. With the same density calibration, its Cauchy transitions
and inverse transitions are exactly the matrices \(B_{ij}\) below. Their
Wronskian supplies a specified isomorphism
\(\bigwedge^2\mathcal E_J\xrightarrow{\sim}\mathcal L\), which is the
identification used in subsequent formulas. The construction is natural
under marked contact/incidence comparisons and finite parameter
restriction.

\emph{Proof.} Put \(\mathcal K=\ker\mathrm dq\). First-jet evaluation
gives \(\mathcal K|_{Q_x}=\mathcal O(-1,-1)\), so
\(\mathcal K\otimes q^*\mathcal L\) descends by compact-fiber evaluation
to a line \(\mathcal B\) on \(M\). Take its root pair
\(\mathcal E[1]^2=\mathcal B\) on the small marked smooth moduli germ,
with the specified central sign. This uses no square root of
\(\mathcal O(1,1)\).

For any ray coordinate \(u\) and contact-line frame \(l\), let
\(m_u^2=\partial_u\otimes l\). The tangent-sky graph defines the
potential by the matrix Schwarzian in the next displayed formula.
Appendix D proves that, under \(v=f(u)\) and \(l_v=h\,l\),
\begin{equation}\protect\phantomsection\label{eq:gg-projective-optical-definition}{\begin{gathered}
\kappa_v=f'^{-2}\big(\kappa_u-\tfrac12\{f,u\}\big), \\
y_v=\sqrt{f'/h}\,y_u, \\
(\partial_v^2+\kappa_v)y_v
  =h^{-1/2}f'^{-3/2}(\partial_u^2+\kappa_u)y_u .
\end{gathered}}\end{equation} Hence these are one intrinsic line-valued ray operator.
Its two-dimensional solution spaces descend by short-ray ODE uniqueness.
The Wronskian takes values in \(\mathcal L\). Appendix D proves
graph-frame invariance directly and compares the identical calibrated
Cauchy frames, including inverse matrices and scale changes. No Schouten
tensor or Einstein germ is an input of this definition. ∎

Choose local contact forms \(\alpha_i\). For comparison with the
intrinsic construction, take the density frame \(e\) specified by
\(\boldsymbol g=g_{\mathrm{aux}}e^2\); this agrees with the central
mark. On the incidence there is a unique generator \(D_i\) of
\(\ker\mathrm dq\) whose projection \(k_i\) satisfies
\(g_{\mathrm{aux}}(k_i,\mathrm dp\,Y)=q^*\alpha_i(Y)\). The
characteristic and covector-lift computation of Proposition 5.2
identifies this generator with the Hamilton field of the null covector,
so it is affine for the trial \(g_{\mathrm{aux}}\) without an Einstein
assumption. If \(\Phi_{ij}^*\alpha_i=c_{ij}\alpha_j\), then
\(D_i=c_{ij}D_j\), with \(c_{ij}\) constant along a ray.

The tangent-sky curve in the fixed contact plane is a regular symmetric
graph \(S_i(t)\). Its optical potential is defined from that curve:
\begin{equation}\protect\phantomsection\label{eq:gg-optical-equation}{
 \kappa_i=\frac12\operatorname{tr}
   \left\{\frac12(S_i')^{-1}S_i'''
           -\frac34\big((S_i')^{-1}S_i''\big)^2\right\},
 \qquad D_i^2y+\kappa_i y=0 .
}\end{equation} To identify this potential geometrically, write
\(g=g_{\mathrm{aux}}\), fix one short affine null geodesic \(\gamma\)
with tangent \(k\), and write \(u\) for its affine parameter and
\(\mathcal D_z=\ker\alpha_z\) for the contact plane. Proposition 5.2, or
the independent reconstruction in §D.7, identifies \(z\in Z\) with this
unparametrized geodesic. A variation through null geodesics is a Jacobi
field. Conversely the initial-value theorem integrates every Jacobi
initial datum satisfying the linearized null constraint to such a
variation. The contact condition is \(g(k,J)=0\). Changes of affine
origin and speed are the fields \((a+bu)k\), so the contact plane is the
four-dimensional solution space
\begin{equation}\protect\phantomsection\label{eq:gg-contact-jacobi-solution-space}{
 \mathcal D_z\simeq
 \frac{\{J:\nabla_k^2J+R(J,k)k=0,\quad
             g(k,J)=g(k,\nabla_kJ)=0\}}
      {\{(a+bu)k:a,b\in\mathbb C\}}.
}\end{equation} Here the orthogonality conditions can equivalently be
imposed at one initial point. They persist by the Jacobi equation. This
is a space of solutions along \(\gamma\), rather than the
two-dimensional screen \(k^\perp/\langle k\rangle\) at one point. In the
covector normalization of Proposition 5.2 its conserved symplectic
pairing is
\(W(J_1,J_2)=g(J_1,\nabla_kJ_2)-g(\nabla_kJ_1,J_2)=-d\alpha(J_1,J_2)\).
Indeed differentiating the null Hamilton flow gives the displayed Jacobi
equation, and differentiating \(\alpha(J)=g(k,J)\) gives the last
equality. Both pairings pass to the quotient by \((a+bu)k\). Varying a
ray while keeping \(x=\gamma(u)\) fixed is precisely a variation
vanishing there; modulo the tangential fields this says
\(T_zQ_x=\ker(\operatorname{ev}_u:\mathcal D_z\to k^\perp/\langle k\rangle)\).
These kernels are Lagrangian. This is the contact realization of the sky
incidence (LeBrun 1983, III.2--3).

Choose a parallel orthonormal screen frame. Its Jacobi equation is
\(j''+\mathcal Rj=0\), where \(\mathcal R[j]=[R(J,k)k]\) is symmetric.
Complete \(k\) to a null frame to take the Ricci trace: the two
null-direction contributions vanish by \(R(k,k)=0\) and
\(g(R(n,k)k,k)=0\). Thus
\(\operatorname{tr}\mathcal R=\operatorname{Ric}(k,k)\) with exactly the
curvature and Ricci conventions of §2. Choose a symplectically
normalized fundamental pair of screen solution matrices \(Y,Z\), so
\begin{equation}\protect\phantomsection\label{eq:gg-geometric-jacobi-wronskians}{
 Y^{\mathsf T}Z'-(Y')^{\mathsf T}Z=1_2,\qquad
 Y^{\mathsf T}Y'=(Y')^{\mathsf T}Y,\qquad
 Z^{\mathsf T}Z'=(Z')^{\mathsf T}Z.
}\end{equation} A solution is \(j=Ya+Zb\). On a chart where \(Y\) is a
unit, its vanishing at \(u\) is \(a=-Y^{-1}Zb\). The row identity
\(YZ^{\mathsf T}=ZY^{\mathsf T}\), a consequence of the displayed
symplectic normalization, makes \(S=Y^{-1}Z\) symmetric. In the constant
symplectic coordinates \((b,-a)\) the actual tangent-sky curve is
precisely this graph, with the stated sign convention. In these
coordinates \(d\alpha=d(-a)\wedge db=-W\), consistently with the contact
normalization above. The Wronskian identities give
\(S'=Y^{-1}Y^{-{\mathsf T}}\) and
\((S')^{-1}S''=-2(Y^{-1}Y')^{\mathsf T}\). Writing \(A=(S')^{-1}S''\)
now gives \(\frac12A'-\frac14A^2=(Y^{-1}\mathcal RY)^{\mathsf T}\). The
trace in the optical definition is consequently
\(\kappa=\frac12\operatorname{Ric}(k,k)=P(k,k)\), since \(g(k,k)=0\) in
dimension four. Any other constant symplectic graph frame gives the same
trace by §D.2. The proof uses only null-geodesic incidence, metric
compatibility and finite-dimensional ODE uniqueness, so it applies to
the trial metric before Einstein admission and commutes with finite
parameter reduction.

For actual transverse slices let \(\tau_{ij}\) be the time in \(D_j\)
needed to move from slice \(j\) to slice \(i\) over \(\Phi_{ij}z\). The
optical transitions are
\begin{equation}\protect\phantomsection\label{eq:gg-actual-optical-transitions}{
 \begin{aligned}
 \partial_tF_j(t,z)
  &=\begin{pmatrix}0&1\\-\kappa_j(t,z)&0\end{pmatrix}F_j(t,z),
       &F_j(0,z)&=1,\\
 B_{ij}(z)
  &=\begin{pmatrix}1&0\\0&c_{ij}(z)\end{pmatrix}
                         F_j(\tau_{ij}(z),z).
 \end{aligned}
}\end{equation} They satisfy \(B_{ik}=(B_{ij}\circ\Phi_{jk})B_{jk}\),
because both routes continue the same solution to the same final Cauchy
state. Their determinant is \(c_{ij}\), not a globally imposed unit.
Thus the Wronskian identifies
\(\bigwedge^2\mathcal E_J\simeq\mathcal L\) by the specified Wronskian
isomorphism.

Under \(\widehat g=\Omega^2g\) the density convention is
\(\widehat y=\Omega y\), \(\widehat D=\Omega^{-2}D\). Direct
differentiation gives
\begin{equation}\protect\phantomsection\label{eq:gg-optical-density-covariance}{
 \begin{aligned}
 \widehat\kappa
   &=\Omega^{-4}\{\kappa-D^2\log\Omega+(D\log\Omega)^2\},\\
 C_{\Omega,i}
   &=\begin{pmatrix}
       \Omega&0\\\Omega^{-1}D_i\log\Omega&\Omega^{-1}
      \end{pmatrix},&
 \widehat B_{ij}
   &=(C_{\Omega,i}\circ\Phi_{ij})B_{ij}C_{\Omega,j}^{-1}.
 \end{aligned}
}\end{equation} These coherent formulas retain the selected section
under a scale comparison.

\pandocstatementheading{\protect\phantomsection\label{prop:gg-optical-einstein}Proposition 5.4.}

A holomorphic section of this optical bundle on a neighborhood
containing a whole compact sky is exactly an analytic density satisfying
\((\nabla^2\sigma+P\sigma)_0=0\). If its value at the mark is a unit,
its selected metric \(\sigma^{-2}g_{\mathrm{aux}}\) is Einstein and its
prolongation is the same canonical parallel section \(I\).

\emph{Proof.} Evaluate the section on the incidence. On each compact
connected sky it is a holomorphic function with values in the fixed
one-dimensional density fiber, hence constant. These constants descend
holomorphically to \(\sigma(x)\). The ray equation implies
\(k^ak^b(\nabla_a\nabla_b\sigma+P_{ab}\sigma)=0\) in every null
direction. A symmetric quadratic form zero on the complete nondegenerate
null quadric is a multiple of \(g\). Conversely that trace-free equation
supplies raywise solutions whose Cauchy states agree under the
\(B_{ij}\), proving both inverses.

Set \(J_g=\operatorname{Scal}(g)/6\) and
\begin{equation}\protect\phantomsection\label{eq:gg-selected-scale-prolongation}{
 \mu=\nabla\sigma,\qquad
 \rho=-\tfrac14(\Delta\sigma+J_g\sigma),\qquad
 \lambda_\sigma=-3(2\sigma\rho+|\mu|^2).
}\end{equation} The trace-free equation gives
\(\nabla_a\mu_b+P_{ab}\sigma+g_{ab}\rho=0\). Its divergence, using
\(\nabla^aP_{ab}=\partial_bJ_g\) and \(\operatorname{Ric}=2P+J_gg\),
gives \(\nabla_b\rho=P_b{}^a\mu_a\). These are exactly the normal
tractor parallel equations. Their norm is constant. The Schouten
transformation is
\begin{equation}\protect\phantomsection\label{eq:gg-optical-einstein-schouten}{
 P^{\sigma^{-2}g}
  =P^g+\sigma^{-1}\nabla^2\sigma
               -\tfrac12\sigma^{-2}|\mathrm d\sigma|^2g
  =\left(-\sigma\rho-\tfrac12|\mu|^2\right)\sigma^{-2}g .
}\end{equation} It proves
\(\operatorname{Ric}(\sigma^{-2}g)=\lambda_\sigma\sigma^{-2}g\) and
identifies the same \(I=(\sigma,\mu,\rho)\), without dividing by
\(\lambda_\sigma\). Normalizing the evaluation vector by \(\sigma\)
makes its position affine: if \(X''+\kappa X=0\) and
\(\sigma''+\kappa\sigma=0\), then \(Y=X/\sigma\),
\(\mathrm d\widehat t/\mathrm dt=\sigma^{-2}\) satisfy
\(\mathrm dY/\mathrm d\widehat t=\sigma X'-\sigma'X\), whose derivative
is zero. The translation line is \(\mathcal L^{-1}\). The dual optical
plane embeds in the normal bundle as \(\langle X,D_kX\rangle\), by its
Cauchy evaluation, so this is precisely the affine line of Definition
2.2. ∎

The requirement is on a whole neighborhood, not only the central sky.
Indeed \(\mathcal E_J|_Q\simeq\mathcal O\oplus\mathcal O(1,1)\) has five
central sections. The sixth flat scale freedom is invisible in these
central Cauchy values and is detected transversely.

We can now correct the actual bundle defect and the entire record
simultaneously. At a small extension with highest ideal \(J\), first
express the lower selected density as the Einstein scale and lift that
scale to the trial stage. For arbitrary lifted local states define the
target-minus-source defect
\begin{equation}\protect\phantomsection\label{eq:gg-optical-defect-sign}{
 r_{ij}=\widetilde s_i\circ\Phi_{ij}-B_{ij}\widetilde s_j .
}\end{equation} Its highest weighted piece of weight \(m\) belongs to
the flat sheaf \(\mathcal S_m=\ker(D^2:P_m\otimes\mathcal O_Q
\to P_{m-2}\otimes\mathcal L^2)\). A value has degree \(m\) and a slope
degree \(m-1\). This grading is required for the actual transport, and
is not ordinary normal degree on the ray space. Fixed scalar Čech
splitting \(\delta u=r\) gives \(D^2u=A_{ab}(x)k^ak^b\) and the class
\begin{equation}\protect\phantomsection\label{eq:gg-optical-obstruction-class}{
 a_m=[A]\in
 \mathcal Q_m=
 \frac{P_{m-2}\otimes\operatorname{Sym}^2_0V^*}
                   {\operatorname{Hess}_0P_m}.
}\end{equation} Contact degree \(m+1\) has class space
\(\mathcal C_{m+1}=H^1(Q,\mathcal F_{m+1})\) and free projection
\(\Pi_{m+1}:\mathcal C_{m+1}\to W_{m-2}^+\oplus W_{m-2}^-\). The fixed
linear vacuum inclusion is \(\iota_{m+1}\). Appendix D proves the exact
sequence of the specified maps
\begin{equation}\protect\phantomsection\label{eq:gg-optical-symbol-exact-sequence}{\begin{gathered}
0\longrightarrow W_{m-2}^+\oplus W_{m-2}^-
 \xrightarrow{\iota}\mathcal C_{m+1}
 \xrightarrow{o_m}\mathcal Q_m
 \xrightarrow{d_m}Z^2_{m-4}\longrightarrow0, \\
o_mJ[h]=[P(h)^\circ], \\
d_m[A]=\mathrm d\,\operatorname{div}A.
\end{gathered}}\end{equation} Negative polynomial degrees mean zero. The inverse
\(L_m:\ker d_m\to\ker\Pi_{m+1}\) of \(o_m\) is a fixed finite
constant-matrix operation.

\pandocstatementheading{\protect\phantomsection\label{thm:gg-compatible-correction}Theorem 5.5.}

For every trial contact lift of an already corrected lower stage the
actual optical class satisfies \(d_ma_m=0\). If \(Z_{m-2}\) is the
difference between the prescribed and current entire free record at that
weight, the closed contact correction class is
\begin{equation}\protect\phantomsection\label{eq:gg-simultaneous-correction}{
 \Delta c_{m+1}=-L_ma_m+\iota_{m+1}Z_{m-2}.
}\end{equation} A fixed analytic cochain representative, fresh incidence
and optical recomputation, and the uniquely normalized remaining scalar
correction give the next stage. It has a whole optical section, its same
canonical Einstein scale, numerical coefficient \(\varepsilon\Lambda\)
and every prescribed free coefficient. It is independent, up to the
marked selected-section comparison, of the trial lift.

\emph{Proof.} Appendix D derives the optical-class map using the actual
bundle above. In the lifted lower Einstein scale the comparison scalar
\(1\) has ray residual \(P^{\mathrm{trial}}(k,k)\) in \(J\). Local
double primitives make its highest class exactly
\([P^{\mathrm{trial}\,\circ}]\). The contracted Bianchi identity
\(\nabla^aP^\circ_{ab}=\tfrac34\partial_bJ_g\) has, at this highest
coefficient, flat derivative: all lower \(P^\circ\) terms vanish and the
lower scalar is constant. Hence \(\mathrm d\,\operatorname{div}a_m=0\).
This uses the moving incidence and Cauchy frames; an arbitrary rank-two
defect would not have this compatibility.

A closed contact class changes the obstruction and free record by
\((a_m,\operatorname{Rec})\mapsto
(a_m+o_m\Delta c,\operatorname{Rec}+\Pi_{m+1}\Delta c)\). Every product
of a \(J\)-valued variation with a coefficient of positive parameter
degree vanishes, since \(\mathfrak mJ=0\) for the maximal ideal of this
small extension. The exact sequence and \(\Pi\iota=1\) prove the two
simultaneous corrections, including every derivative order. Use the
fixed representative \(\ell_{m+1}\Delta c\) constructed in Appendix C
and apply the positive highest exponential to the coordinate-algebra
overlaps \(E\). The corresponding stored point correction is the inverse
flow \(\operatorname{Fl}_{-X_{\ell_{m+1}\Delta c}}\). Recompute the
graphs from those point maps, as fixed by Proposition 5.8.

Recompute the graph incidence, tangent-sky potential and optical
matrices after this change. Its remaining defect is exact in
\(\mathcal S_m\). If \(\delta u=r'\) and \(D^2u=A'\), the vanishing
class gives \(A'=\operatorname{Hess}_0f_m\). Then \(v_i=u_i-f_m\) has
\(D^2v_i=0\) and \(\delta v=r'\). Adding these local homogeneous ray
solutions cancels the actual target-minus-source defect. Appendix D
supplies the homogeneous Hessian primitive explicitly.

The remaining flat global kernel is \(a+b\cdot x+\gamma\eta(x,x)\). Fix
the value and gradient in the lifted lower selected scale, then compute
the numerical coefficient \(\lambda_*\) of this same section from its
prolongation. Since \(\Delta\eta(x,x)=8\), changing by
\(\gamma\eta(x,x)\) changes the highest numerical coefficient by
\(12\gamma\). Set
\begin{equation}\protect\phantomsection\label{eq:gg-numerical-normalization}{
 \gamma=\frac{\lambda_{\mathrm{target}}-\lambda_*}{12}.
}\end{equation} Thus the original numerical label, including zero, is
retained. No desired scalar curvature has been assigned to the off-shell
optical potential. These conformal highest corrections do not change the
highest Weyl record.

A closed difference \(v\) of trial lifts changes \((a,Z)\) into
\((a+o_mv,Z-\Pi v)\). Since \(L_mo_m=1-\iota\Pi\), the new correction is
the old one minus \(v\). The resulting marked geometric comparison
preserves the selected density: its remaining difference is a scalar
coboundary plus the six global modes just uniquely fixed. This proves
trial independence at the needed geometric level, without asserting
equality of arbitrary raw corner gauges.

All spatial coefficients at this finite stage are analytic. The
nontrivial inverse estimate is proved in Appendix C: a compatible
analytic \(A\) gives a conserved Ricci source, an explicit bounded
potential with that Schouten tensor, and removal of its whole flat
vacuum record. Applying the same fixed cochain lift \(\ell\) represents
precisely \(L_ma_m\) with a bound after one fixed spatial restriction.
The scalar splittings, double primitives and Hessian primitives have
polynomial-times-exponential degree bounds. They can therefore be summed
in normal degree. Finitely many operations complete each finite
parameter stage, possibly on a smaller stage-dependent domain. No
infinite-stage convergence is inferred here. ∎

For the source execution sequence fix the trial rule as well. In each
independent overlap orientation, expand the Hamiltonian logarithm
relative to the fixed reference substitutions as
\(H_{ij}^{\le N}=\sum_{k=1}^N\varepsilon^kh_{ij,k}\). Its lift to
\(A_{N+1}\) is this same coefficient polynomial, with zero new
coefficient. Exponentiate it and use the fixed two-cocycle repair of
Lemma 5.1; store opposite orientations by exact inversion. Lift the
already selected local scalar/Cauchy representatives by the same
coefficient rule, and recompute all unit inverses, scale changes,
incidence and optical data from those representatives before applying
Theorem 5.5. This fixes a raw sequence without additional free
coefficients. Other trial lifts remain allowed for the geometric
independence assertion of that theorem; the prescribed normalized output
is the same. For a finite algebra without a chosen coefficient grading,
any fixed linear lift can be used, with normalized comparison and
restriction supplied by §6.

Starting with the original first-order lifts and the flat scale, this
trial rule, Lemma 5.1 and Theorem 5.5 define all stages \(\mathcal C_N\)
compatibly. The next section specifies their stored normalization and
proves its convergence on one common domain.

\pandocstatementheading{\protect\phantomsection\label{def:gg-independent-affine}Definition 5.6.}

A calibrated affine object consists of a regular contact seed
\((Z,Q,\mathcal L)\) with \(\mathcal L|_Q=\mathcal O(1,1)\), its ordered
rulings and original marking, and an affine-line torsor \(A\to Z\) with
a specified translation identification \(T(A)=\mathcal L^{-1}\). Its
unique lifted compact sky must have normal bundle \(\mathcal O_Q^4\).
Let \(j_A:F\to A\) be the unique family of lifted skies. Require that
\(j_A\iota^{-1}\) extend from the physical short incidence to a
holomorphic fiberwise projective isomorphism
\begin{equation}\protect\phantomsection\label{eq:gg-independent-projective-extension}{
 \overline J:\mathbb P(\mathcal E_J^*)\longrightarrow\overline A ,
}\end{equation} where \(\overline A\) is the fiberwise projective
completion of the affine torsor and \(\iota\) is projective evaluation
for Proposition 5.3. Both completions are defined from the supplied
holomorphic objects.

Retain the root pair \(\beta:\mathcal E[1]^2\to\mathcal B\) of that
proposition and its nonzero central density value \(e_o\) fixed by the
marking. With \(b=(\mathrm d_qj_A)^{-1}\) descended to \(\mathcal B\),
its marked value must satisfy \(b(o)=\beta(e_o^2)\); the selected root
has \(\sigma(o)=e_o\). The finite-jet invariant of Proposition 5.4 at
the mark is the requested numerical \(\Lambda\). The complete family,
its lift \(j_A\), and \(\sigma\) are derived; no metric or scale is an
independent input.

\pandocstatementheading{\protect\phantomsection\label{prop:gg-affine-equivalence}Proposition 5.7.}

These independently admitted affine objects are equivalent to the
calibrated contact/optical objects with a whole unit section. Their
overlap maps have the local holomorphic law
\begin{equation}\protect\phantomsection\label{eq:gg-affine-local-law}{
 (z,\tau_j)\longmapsto
       \bigl(\Phi_{ij}(z),\,c_{ij}^{-1}(z)\tau_j+a_{ij}(z)\bigr),
 \qquad \Phi_{ij}^*\alpha_i=c_{ij}\alpha_j.
}\end{equation} For the generated objects the translations are
determined by the original Cauchy matrices, without new free input. The
same finite correction of Theorem 5.5 corrects precisely the obstruction
to this calibrated affine reduction. Both representative maps and their
inverses converge on the same conditioned parent domains as the stored
optical data.

\emph{Proof.} Appendix D.7 proves uniqueness and holomorphic dependence
of the lifted compact skies, and identifies their normal-bundle
condition with invertibility of \(\mathrm d_qj_A\). Appendix D.8
identifies the projective extension with the intrinsic Schwarzian
equality, obtains the compatible marked root and the same whole optical
section, and proves both inverse comparisons. Appendix D.9 gives the
explicit matrix calculation and its cochain obstruction. The convergence
statement follows from the finite formulas there and the common analytic
witness in §6. Thus the Einstein recognition in Proposition 5.4 is a
consequence of this holomorphic admission. ∎

\pandocstatementheading{\protect\phantomsection\label{prop:gg-overlap-orientation}Proposition 5.8.}

With equation (\ref{eq:gg-overlap-operator-point-convention}), the
contact class in Appendix C has the same sign as the actual conformal
metric variation. A highest closed operator correction \(d\) therefore
changes the optical class by \(+\mathfrak o d\) and the free record by
\(+\Pi d\), as used in Theorem 5.5.

\emph{Proof.} At first order write the scalar sky graphs as
\(f_i=f_0+\varepsilon v_i\). The point generator of \(\Phi_{ij}\) is
\(-X_{h_{ij}}\), so \(\mathscr A_{\Phi_{ij}}f_j=f_i\) gives
\(v_i-v_j=-h_{ij}|_{f_0}\). If \(h_{ij}=a_j-a_i\) and
\(Da_i=h_{\mathrm{met}}(k,k)/2\) as in Appendix C, one can take
\(v_i=a_i\) up to a global contact-line section, which is the
infinitesimal coordinate freedom. The null covector changes by
\(\mathrm d_xv_i\). For a covariant metric variation
\(h_{\mathrm{met}}\), the first variation of its null equation is
\begin{equation}\protect\phantomsection\label{eq:gg-overlap-metric-sign-check}{
 2D v_i-h_{\mathrm{met}}(k,k)=0 .
}\end{equation} This is exactly the stated primitive equation. The
auxiliary scalar normalization changes only the trace; neither that
trace nor the coordinate freedom changes the conformal class, optical
Hessian quotient or linear free Weyl record. At a highest ideal all
positive-parameter products vanish, so the same calculation gives the
signs of \(\mathfrak o\) and \(\Pi\) at every finite stage.

The convention also respects the exact cocycle order:
\(E_{ij}E_{jk}=(\Phi_{jk}^{-1}\circ\Phi_{ij}^{-1})^*
=(\Phi_{ij}\circ\Phi_{jk})^{-1*}\). Assigning the positive point flow to
the same positive operator generator would instead produce
\(-h_{\mathrm{met}}\); it is not an equivalent notation. For
\(F=(w^0)^4/\zeta\), Appendix C gives \(h_{vv}=-24u^2\). The
linearization of our last-slot curvature is
\begin{equation}\protect\phantomsection\label{eq:gg-last-slot-linear-curvature}{
 R^{(1)}_{abcd}(h)=\tfrac12
 (\partial_a\partial_ch_{bd}+\partial_b\partial_dh_{ac}
  -\partial_a\partial_dh_{bc}-\partial_b\partial_ch_{ad}).
}\end{equation} Hence \(C^{(1)}_{uvuv}=R^{(1)}_{uvuv}=-24\), and
eq.~\ref{eq:gg-weyl-spinor-readout} gives \(\Psi^+_{0000}=+24\), exactly
the positive contour of Theorem 3.1. For the negative source
\(F_{-6}=1/\zeta\), Appendix C gives \(h_{qq}=-u^2\), hence
\(C^{(1)}_{uquq}=-1\) and \(\Psi^-_{0'0'0'0'}=+1\), the negative
contour. Reversing the stored point velocity changes both metric
variations and therefore both of these records by a minus sign. ∎

\pandocstatementheading{\protect\phantomsection\label{prop:gg-finite-source-execution}Proposition 5.9.}

For the fixed trial rule above, each finite parameter order \(N\) and
finite transverse jet order \(d\) of the corrected stage and its
prescribed stored normalization is determined by finitely many
coefficients of the two original residue representatives, \(\Lambda\)
and the fixed marking/presentation choices. Transverse jets mean the
normal ideal of the reference compact sky on the contact carrier, and
the corresponding marked central incidence on the derived components.
They are computed by the finite cochain, matrix, derivative,
unit-inverse and initialized ODE operations specified here. An
independently solved nonlinear Einstein metric is not an input of this
rule.

\emph{Proof.} Let \(\mathcal I\) be the normal ideal of the reference
compact sky. In its adapted contact coordinates, a degree-\(n\)
Hamiltonian has normal vector components of degree \(n\) and angular
components of degree \(n-1\). In particular
\begin{equation}\protect\phantomsection\label{eq:gg-source-execution-filtration}{
 \{\mathcal I^a,\mathcal I^b\}
             \subset\mathcal I^{a+b-1},\qquad a,b\ge1.
}\end{equation} The source Hamiltonians start in degree three. A fixed
normal jet of a finite nilpotent product or logarithm consequently sees
only finitely many homogeneous input degrees. The fixed repairs \(S_n\)
and representatives \(\ell_n\) preserve their displayed homogeneous
degrees. Each matrix \(L_m\) and \(\iota_{m+1}\) acts within one
displayed degree, with constants fixed before any input is supplied.

The scalar graph inverse and four-evaluation normalization are unit
implicit operations. A specified finite jet of their output uses a
finite jet of their source maps. The nine-direction metric inverse is
another unit matrix inverse; covariant curvature readout and the
tangent-sky projective operator use only finitely many further
derivatives. The weighted optical splitting, double primitive and
Hessian primitive preserve the indicated value/slope grading. Thus, at
each of the finitely many parameter steps, a fixed output normal jet
needs only finitely many trial graph jets and free record components.
Terms above that finite homogeneous cutoff cannot enter through an
unperformed infinite sum. Induction gives a finite original prefix for
every \((N,d)\).

The normalization of §6 has the same property. Its physical metric is
the metric already read from the corrected incidence and its own
selected density. Formal geodesic development, radial spin transport,
the identity-initialized Moser flow and their unit inverses have unique
finite Taylor recursions. Their fixed finite jets therefore require only
finite jets of that derived metric. The fresh optical and affine
components are finite derivatives and initialized linear ODEs of the
same incidence. No new geometry is chosen to execute this normalization.

The matrix \(L_m\) is the substantive inverse of the optical obstruction
on the zero-record complement. Appendix C uses a conserved linear
potential to prove that this same fixed matrix has its analytic cochain
bound; the execution rule evaluates its matrix on the derived defect.
The nonlinear CK witnesses of Appendix A and the direct Cartan witness
of §6 establish consistency and convergence, respectively, and are not
arguments supplied to these finite cochain operations. This proves the
asserted source dependence. ∎

\protect\phantomsection\label{sec:gg-common-convergence}

\subsection{6. The selected normalization and common
convergence}\label{the-selected-normalization-and-common-convergence}

\protect\phantomsection\label{sec:gg-bounded-witness}

\subsubsection{6.1. A direct bounded geometric
witness}\label{a-direct-bounded-geometric-witness}

The finite source procedure is the cochain construction of §5 and
Proposition 5.9. For its convergence proof, Theorem 3.3 and Proposition
4.1 construct a bounded geometric comparison from the same original
records. Section 6.2 defines one prescribed normalizer for both
constructions, and Theorem 6.3 identifies every finite stored tuple with
the corresponding Taylor coefficients of the comparison. The Cartan
geometry is a proof witness, not an additional input to the finite
source procedure.

Fix \(1<T_*<T_c<T_b<T_{\mathrm{in}}\) and consider
\((\varepsilon U^+,\varepsilon U^-;\varepsilon\Lambda)\) for
\(|\varepsilon|<T_{\mathrm{in}}\). Scaling the numerical constant as
well as the records makes the zero-amplitude residue flat, as required
by the highest-ideal calculation in §5. The final value is
\(\varepsilon=1\).

\pandocstatementheading{\protect\phantomsection\label{prop:gg-common-cartan-witness}Proposition 6.1.}

For every bounded original family this amplitude family has a direct
Cartan realization, its marked normal coordinates and canonical ray
geometry on common spatial domains over \(|\varepsilon|\le T_b\). Their
fixed finite derivatives and inverse maps are bounded. The construction
uses Theorem 3.3 and the finite equations of §4.

\emph{Proof.} The initial record bound is at most
\(T_{\mathrm{in}}(16C+|\Lambda|)\) on the fixed smaller record radius.
Apply Theorem 3.3 successively to the four horizontal and six vertical
basis directions used in Proposition 4.1. Use each intermediate output
bound and positive radius as the next input, also for the reversed
finite word. Their finite minimum gives a common ten-variable
neighborhood. The coframe is constructed by the finite matrix equations
\(\partial_n\alpha_i=-C_n(\mathcal U_n)\alpha_i\), \(\alpha_n=e_n\),
with \(\alpha(0,\varepsilon)=1\). The matrices depend on finitely many
bounded record components. The ODE bound and interior Cauchy estimates
therefore control \(\alpha\) and its derivatives uniformly on the
amplitude disk. Shrinking space gives \(\|\alpha-1\|\le1/2\),
\(\|\alpha^{-1}\|\le2\).

The anchor and structure equations proved in §4 apply to these flows.
With horizontal directions first, the remaining six directions span
\(\ker\theta\), and its quotient has the derived metric. Its connection
is torsion free and its curvature is \(R_{\mathcal U}\), so its Ricci
trace is exactly \(\varepsilon\Lambda g_\varepsilon\).

The finite geodesic and parallel-spin ODEs now give its marked
exponential coordinates and radial spin frame. Their initial derivative
is the identity for every amplitude. A finite implicit inverse and a
spatial restriction give common normal-coordinate and inverse domains,
with \(g_\varepsilon(0)=\eta\) and \(\partial_xg_\varepsilon(0)=0\). The
metric, inverse and any specified finite derivatives have common bounds.

In the physical scale the canonical sections are
\begin{equation}\protect\phantomsection\label{eq:gg-common-canonical-sections}{\begin{gathered}
I_\varepsilon=(1,0,-\varepsilon\Lambda/6), \\
X=(0,0,1), \\
H(I_\varepsilon,I_\varepsilon)=-\varepsilon\Lambda/3,\quad
 H(I_\varepsilon,X)=1 .
\end{gathered}}\end{equation} The scale plane has determinant \(-1\) for every
amplitude and constant. For a null affine tangent \(k\), \(D_k^2X=0\),
so its actual transported positions form an affine line. The direct ray
construction of §4 has a finite direction cover, common transverse-hit
margins, short ODE and inverse-transport bounds. To specify the
connected-fiber margin, fix the derivative bound \(C\) on larger flow
parents first. Where \(k^\nu\) is a unit, use \(u=x^\nu\) and write each
geodesic as a graph \(G_z(u)\) on a convex parent disc, with
\(|G_z'|\ge a>0\) and \(|G_z''|\le C\). Restrict the entire graph family
so that \(|G_z|\le R_p\) on that parent and \(R_pC<a^2/2\). The real
Hessian of \(|G_z|^2\) there is at least \((2a^2-2R_pC)\) times the
Euclidean form. Its inner-ball sublevel is therefore convex and
connected. Choose the inner spatial radius below the parent
\(u\)-radius; the free coordinate cannot exit that parent while the
graph remains in the inner ball. In a normal spin frame the connection
is \(O(|x|)\), so the \(O(r)\) short hit changes the projective
direction by \(O(r^2)\); a fixed angular buffer keeps the entire
retained plaque in one such parent. This establishes the assertion on
the whole convex parent, rather than inferring connectedness from a
Hessian bound only inside the sublevel. The hit comparisons and quotient
are those of this short relation; their equality sets are closed by the
same finite transverse coordinates, giving the Hausdorff germ. The
ordered complete skies and actual initial spin lift are retained. The
vector \(I_\varepsilon\) varies with amplitude; it is not orthogonally
transported to a fixed vector of another norm. ∎

\protect\phantomsection\label{sec:gg-prescribed-normalizer}

\subsubsection{6.2. The prescribed stored
representative}\label{the-prescribed-stored-representative}

Fix once the finite direction cover and transverse slices, their flat
identification with the contact charts of §5, the contact-line frames,
the four evaluation directions of Proposition 5.2 and the normal central
spin gauge. These are presentation choices, denoted by \(P\). The same
choices are used over finite algebras and analytic parameter families.

The normalizer \(\mathcal N_P\) is defined as follows. Reconstruct the
selected normal geometry and its actual \(I,X\) from the corrected
contact incidence and whole section, using §§4--5. Develop its marked
physical metric into geodesic normal coordinates and the radial spin
frame. Form the prescribed short-ray charts and Hilbert contact forms
there. Identify them with the fixed flat charts by the contact Moser
path \(\alpha_t=\alpha_0+t(\alpha_g-\alpha_0)\). Write
\(\gamma=\alpha_g-\alpha_0\) and solve
\begin{equation}\protect\phantomsection\label{eq:gg-defined-contact-moser}{
 \alpha_t(Y_t)=0,\qquad
 (\iota_{Y_t}\mathrm d\alpha_t)|_{\ker\alpha_t}
                          =-\gamma|_{\ker\alpha_t}.
}\end{equation} Let \(M_i\) be this identity-initialized point flow from
the fixed model chart into the physical chart, so \(M_i^*\alpha_g\) is a
unit multiple of \(\alpha_0\). Use its inverse and its actual contact
multiplier as well. The normalized point overlaps are
\(M_i^{-1}\Phi_{ij}^{\mathrm{geom}}M_j\), and their algebra operators
are their inverse pullbacks. The normal-coordinate expansion gives
\(\gamma\) of second normal order along the compact sky, hence \(Y_t\)
of second normal order. The construction fixes that sky and its first
jets.

Next send the actual complete skies through these same inverse maps
\(M_i^{-1}\) and take their scalar graphs. Use their same four
evaluations to choose point coordinates, by the inverse of the map
\(E_4(x,\varepsilon)=\mathfrak e(f(x,\varepsilon))\). Finally recompute
the nine-direction conformal form, the density representing the same
selected \(I\), the actual ray generators, tangent-sky potentials, hit
times and the Cauchy matrices of Proposition 5.3. The affine data of
Proposition 5.7 are then obtained by its fixed finite formulas. No old
optical matrix is copied in place of this fresh computation.

The augmented stored tuple consists of these contact maps and inverses,
contact multipliers and inverses, scalar graph and evaluation inverses,
actual density and Cauchy states, moving hit maps, optical matrices and
inverses, and the derived affine transitions and inverse transitions.
This definition fixes the representative, not just its geometric
isomorphism class. It also defines its finite-algebra version by the
same ODEs and unit inverses.

\pandocstatementheading{\protect\phantomsection\label{prop:gg-bounded-prescribed-normalizer}Proposition 6.2.}

For the witness of Proposition 6.1, the prescribed normalizer produces
one jointly holomorphic augmented tuple \(Y_C(\varepsilon)\) with a
common finite bound on retained parent domains over
\(|\varepsilon|\le T_b\). It retains the same actual canonical section.

\emph{Proof.} Screen projection is bounded without a cosmological
denominator. With \(W_{\mathrm{scr}}=(\mathbb CI+\mathbb CX)^\perp\), it
is
\begin{equation}\protect\phantomsection\label{eq:gg-bounded-screen-projection}{
 \Pi_{\mathrm{scr}}z=z-H(I,z)X-H(X,z)I+H(I,I)H(X,z)X .
}\end{equation} For middle vectors with \(H(e_a,X)=0\), the projected
vectors \(e_a-H(I,e_a)X\) have exactly their original Gram matrix. The
finite normal-development, spin-transport and transverse-hit operations
therefore have the fixed unit matrices and bounds supplied by
Proposition 6.1.

The Hilbert forms and their derivatives are uniformly controlled on the
finite parents. Since \(\gamma\) has second normal order, a spatial
restriction keeps the entire fixed Moser path contact and its time-one
flow and inverse inside their buffers. The contact multiplier is an
exponential of a bounded scalar integral, so it and its inverse are
bounded. This restriction changes space, not the amplitude disk.

Direction projection for the pulled-back skies and the four-point map
\(E_4\) have their fixed invertible derivatives at the central sky for
all amplitudes. The finite implicit theorem supplies common inverses.
Consequently the scalar family is literally
\(f(E_4^{-1}(x,\varepsilon),\varepsilon)\) with the chosen
four-evaluation normalization. The ten-by-ten system of Proposition 5.2
has its same uniform inverse, so it reads a bounded \(g_{\mathrm{aux}}\)
and inverse. Write \(g_{\mathrm{aux}}=\Omega^2g_{\mathrm{phys}}\) in
these point coordinates. The fixed mark gives
\(\Omega(0,\varepsilon)=1\). A common spatial restriction gives
\(|\Omega-1|\le1/2\), making the branch holomorphic over the retained
amplitude disk. The selected auxiliary density is this same \(\Omega\),
its inverse is bounded by two, and its fixed finite derivatives have
interior Cauchy bounds.

The tangent-sky graph and its required three derivatives are bounded on
a fixed invertible graph minor. Thus \(\kappa\) is bounded by some \(K\)
on parent ray discs of length \(\ell\). The fundamental matrices and
inverses have bound \(\exp(\ell(1+K))\) in a fixed induced matrix norm.
The contact-unit factors give
\(\|B_{ij}^{\pm1}\|\le C_c\exp(\ell(1+K))\), including their moving
arguments. The short projective evaluation margin and the matrices
\(G_i\) of §D.9 give the affine maps and their inverse domains with the
bounds proved there.

Every operation used a finite ODE, a specified finite inverse, or a
holomorphic ratio with a common unit margin. It preserves the actual
canonical \(I\) and position \(X\) under the derived geometry
comparison. In a changed auxiliary scale the Cauchy states transform by
\(C_\Omega\), so the value and gradient are recomputed in that scale
rather than treated as unchanged components. The same deterministic
formulas make all operations commute with finite truncation. ∎

\protect\phantomsection\label{sec:gg-literal-convergence}

\subsubsection{6.3. Coefficient comparison with the finite
generator}\label{literal-comparison-with-the-finite-generator}

At order \(N\) let \(\mathcal C_N\) be the corrected stage of §5 over
\(A_N=\mathbb C[\varepsilon]/(\varepsilon^{N+1})\), and store
\(X_N=\mathcal N_P(\mathcal C_N)\). The normalization always retains the
selected whole section, including its value, gradient and numerical
scalar.

\pandocstatementheading{\protect\phantomsection\label{thm:gg-literal-stored-convergence}Theorem 6.3.}

The defined stored tuples satisfy the equality
\(X_N=j_\varepsilon^NY_C(\varepsilon)\). In particular their coefficient
series and every included inverse converge at \(\varepsilon=1\) on
common bounded-family domains. If \(\|Y_C\|\le M_C\) on
\(|\varepsilon|=T_c>1\), then for every fixed \(1<T_*<T_c\) their tails
obey
\begin{equation}\protect\phantomsection\label{eq:gg-literal-convergence-tail}{
 \|X(1)-X_N(1)\|
 \le \frac{M_CT_*^{-N-1}}{1-T_*^{-1}} .
}\end{equation} Here \(X=Y_C\) denotes the joint stored tuple in this
estimate, so \(X(1)=Y_C(1)\). The expression \(X_N(1)\) means evaluation
of its degree-\(N\) coefficient polynomial, not a homomorphism from the
nilpotent algebra.

\emph{Proof.} The finite corrected geometry has exactly the whole lists
\(\varepsilon U_r^\pm\) and constant \(\varepsilon\Lambda\) in \(A_N\),
that is, modulo \(\varepsilon^{N+1}\), including the highest retained
coefficient. Its selected normal geometry and the truncated Cartan
witness consequently have identical completed ordered jets by
Proposition A.3. Proposition A.4, whose Jacobi recursion works
coefficientwise over \(A_N\), makes their metrics identical in marked
normal coordinates. Spatial analyticity of each finite coefficient
promotes that formal spatial identity to equality of actual coefficient
germs.

Their canonical bundles, \(I\) and \(X\) then agree, not just their
scalar norm. Apply the same \(P\) on both sides. Normal development,
prescribed Moser path, initialized flow, graph projection and the same
four-evaluation inverse have unique identical solutions. Their fresh
tangent-sky graphs, potentials, hit maps, selected Cauchy states and
optical fundamental matrices agree by the same finite ODE uniqueness.
The affine matrices are the same finite functions of these states. This
proves equality of every augmented component and inverse. It also proves
compatibility under truncation.

Put \(a_n=[\varepsilon^n]Y_C\), so
\(Y_C(\varepsilon)=\sum_{n\ge0}a_n\varepsilon^n\) and
\(X_N=\sum_{n=0}^Na_n\varepsilon^n\). Coefficient Cauchy estimates on
\(|\varepsilon|=T_c\) give \(\|a_n\|\le M_CT_c^{-n}\le M_CT_*^{-n}\).
Hence \(\|Y_C(1)-X_N(1)\|\le\sum_{n>N}\|a_n\|\), and summing the
geometric series proves the displayed tail. A finite-stage coefficient
initially defined on a smaller spatial neighborhood has a common
representative because, as a germ, it equals the corresponding
coefficient of the single holomorphic tuple \(Y_C\). The exact cocycles
and inverse identities hold for all truncations and hence for the common
analytic sum.

All finite parameter restrictions and derivatives are taken on this one
joint analytic tuple. Interior Cauchy estimates give any specified
finite transverse jet after its usual radius gap. Fixed finite
conditioned point/spin diagrams add only finitely many ODE and hit-map
restrictions. No bound for arbitrary unnormalized formal gauges,
unbounded histories or norms of growing finite algebras is needed. ∎

The finite-stage source dependence is Proposition 5.9, and the analytic
bound for its fixed linear correction is proved in Appendix C. Theorem
6.3 supplies the common limit of the normalized stored tuples.

\protect\phantomsection\label{sec:gg-ordinary-comparison}

\subsection{7. The independent ordinary
presentation}\label{the-independent-ordinary-presentation}

The decoder uses the independently supplied data of Definition 2.3: the
threefold, its complete line family, calibrated Einstein-bundle form and
unrestricted analytic opposite-chirality class. We first recover the
selected background and then register the negative class; Theorem 7.4
proves both counits on the entire ordinary domain.

\protect\phantomsection\label{sec:gg-ordinary-background}

\subsubsection{7.1. Reading the background and the selected
scale}\label{reading-the-background-and-the-selected-scale}

\pandocstatementheading{\protect\phantomsection\label{prop:gg-independent-ordinary-background}Proposition 7.1.}

An ordinary object determines an actual marked analytic Einstein germ
\(g_+\) with vanishing negative Weyl spinor, together with its numerical
constant \(-3H(I_\tau,I_\tau)\). Its complete-line correspondence,
calibrated line and actual form are the ones recovered from \(g_+\). The
construction is holomorphic on common smaller representatives of
conditioned bounded families and commutes with their finite Taylor
reductions.

\emph{Proof.} The completeness theorem (Kodaira 1962), with
\(H^1(\ell,\mathcal O(1)^{\oplus2})=0\), gives a smooth four-dimensional
curve-moduli germ. Its Kodaira map is
\(T_xM=H^0(\ell_x,N_{\ell_x/Z_+})\). A section of
\(\mathcal O(1)^{\oplus2}\) vanishes somewhere precisely when its \(2\)
by \(2\) coefficient matrix has zero determinant. These null vectors
define a nondegenerate conformal quadric, with the original tangent/spin
mark. Vectors vanishing at a fixed point of \(\ell_x\) form one of its
null two-planes. The incidence fibers over \(Z_+\) integrate these
planes.

For completeness, these null surfaces are totally geodesic. With
\(X,Y,Z\) tangent to such a surface, \(g(\nabla_XY,Z)\) is symmetric in
\(X,Y\) by integrability and torsion-freeness, and alternating in
\(Y,Z\) by nullity and metric compatibility. The two symmetries force it
to vanish. Thus its projective spinor is parallel along the surface. The
resulting alpha distribution is exactly the projective spin-connection
distribution; its Frobenius curvature is the negative Weyl spinor. That
spinor is zero. This is the local complete-line correspondence of
(Penrose 1976), with the chosen chirality.

The incidence is the same actual germ: the map sending an incident pair
to its curve point identifies its alpha leaves with the supplied
quotient fibers, and the initialized transversal hits give the inverse.
The whole-parent and angular argument in §E.2 uses only integrability
and unit frame blocks, so applies already to this conformal geometry,
before an Einstein scale has been selected. It supplies short connected
simply connected quotient fibers and the required common
representatives.

The calibrated root is retained in this identification. Let
\(L_{\rm sp}^{-1}\) be the line obtained by parallel descent of the
tautological spinor line. The incidence determinant sequences give
\(q_+^*K_{Z_+}\simeq\mathcal O(-4)\) from \(T_{F_+/M}=\mathcal O(2)\)
and \(T_{F_+/Z_+}=p_+^*S\otimes\mathcal O(-1)\). One must also check
that this determinant identification is parallel in the leaf directions.
In a chosen homogeneous spin frame put
\(\Omega_F=\operatorname{vol}_{\mathrm{aux}}\wedge(\pi_{A'}D\pi^{A'})\),
read as a local ordinary top form in that frame. Its direction factor is
\(-d\zeta\) in the north frame and \(d\xi\) in the south frame; its
horizontal terms vanish against the spacetime volume. These are the
coherent spin factors, rather than separately chosen coordinate
differentials. Put \(s=\iota_{L_1}\iota_{L_0}\Omega_F\). Its spin
transition weight is four times that of the tautological spinor line,
and it is a local frame of the pullback canonical bundle. At any chosen
incidence point take a Levi-Civita-normal spin frame. There
\(D_\alpha\pi=0\), the coefficient divergence of each \(L_A\) against
\(\Omega_F\) is zero, and the horizontal projection of \([L_0,L_1]\) is
zero. Involutivity and injectivity of the horizontal projection on the
alpha distribution make the whole bracket zero there. Consequently
\begin{equation}\protect\phantomsection\label{eq:gg-spin-determinant-bott-descent}{
 \mathcal L_{L_C}s
 =\iota_{[L_C,L_1]}\iota_{L_0}\Omega_F
  +\iota_{L_1}\iota_{[L_C,L_0]}\Omega_F
  +\iota_{L_1}\iota_{L_0}\mathcal L_{L_C}\Omega_F=0
}\end{equation} at that point, in the same tautological parallel frame.
The difference of these two partial line connections is tensorial.
Choosing a normal frame at each point proves it is zero everywhere. Thus
the determinant identification is Bott-parallel and descends along the
connected short alpha leaves, giving the actual isomorphism
\(K_{Z_+}\simeq L_{\rm sp}^{-4}\).

Retain the descended determinant isomorphism
\(\alpha_{\rm sp}:L_{\rm sp}^{-4}\xrightarrow{\sim}K_{Z_+}\) and its
degree-one calibration \(\kappa_{\rm sp}\) induced by the same spin
mark. The quotient neighborhood retracts onto \(\ell\). The supplied
pair \((L,\alpha)\) and \((L_{\rm sp},\alpha_{\rm sp})\) differ by a
\(\mu_4\) torsor; \(H^1(\mathbb P^1,\mu_4)=0\). Hence there is an actual
root-preserving comparison \(f:L\to L_{\rm sp}\). On \(\ell\), the two
calibrations have the same fourth power by Definition 2.3. There is a
unique adjustment of \(f\) by a constant in \(\mu_4\) such that
\(\kappa_{\rm sp}\circ f|_\ell=\kappa\). Two such comparisons have ratio
\(c\) with \(c^4=1\) and \(c|_\ell=1\), hence \(c=1\) on a smaller
connected neighborhood. In particular the residual \(c=-1\), invisible
to \(\tau\) and \(\beta\) alone, is excluded at a fixed mark. Over a
finite local coefficient algebra the same uniqueness follows from
\((c-1)(c^3+c^2+c+1)=0\), since the second factor has unit residue four.
The construction is holomorphic in the retained parameters. For moving
full-spin marks the calibration is transported by \(j_B\) as in
Definition 2.3, so this uniqueness does not identify opposite specified
spin lifts. No rescaling of \(\tau\) or \(\beta\) is used.

Apply the bosonic Einstein-bundle correspondence (Wolf 2007,
Propositions 3.4--3.7). In that source the degree-minus-one universal
line is our \(L^{-1}\), so its bundle is precisely
\(\Omega^1_{Z_+}\otimes L^2\). It identifies the supplied section with
an almost-Einstein density in the conformal germ. Only the branch in
which that density is a unit is used here: its central value is one by
the exact restriction in Definition 2.3. The argument uses the bundle
section at zero cosmological constant as well; integrability of
\(\ker\tau\) is not used to infer a scalar value for a nilpotent family.

The pullback identification can also be checked directly. On the
homogeneous spinor incidence a weight-two one-form annihilating the
Euler field and the alpha distribution has the form
\(\sigma\vartheta-\theta^A\pi^{A'}\mu_{AA'}\). Holomorphy on each
complete projective line makes its vertical coefficient \(\sigma\)
independent of the direction and its remaining coefficients degree-one
polynomials in \(\pi\). In a Levi-Civita-normal spin frame, the
coefficient of \(\vartheta\) in its basicness equation is
\(\pi^{A'}(\nabla_{AA'}\sigma-\mu_{AA'})\). It vanishes in every
direction exactly when \(\mu=\nabla\sigma\). Thus the identification
reads
\begin{equation}\protect\phantomsection\label{eq:gg-ordinary-section-density}{
 q_+^*\tau=\sigma\,\vartheta
       -\theta^A\pi^{A'}\nabla_{AA'}\sigma,\qquad
 \vartheta=\pi_{A'}D\pi^{A'},\qquad
 \theta^A=e^{AA'}\pi_{A'} .
}\end{equation} With the component conventions of §2, the flat north
expressions are \(\theta^A=du^A+\zeta\,dv^A\) and \(\vartheta=-d\zeta\).
For a change \(\widehat g=\Omega^2g\), use the compatible spin frames
\(\widehat e^{AA'}=\Omega e^{AA'}\),
\(\widehat\pi_{A'}=\Omega^{-1/2}\pi_{A'}\) and
\(\widehat\sigma=\Omega\sigma\); the local square root has the existing
central spin calibration. With \(\Upsilon=\mathrm d\log\Omega\), the
torsion-free spin-connection change gives
\begin{equation}\protect\phantomsection\label{eq:gg-ordinary-section-scale-invariance}{
 \widehat\theta^A=\Omega^{1/2}\theta^A,\qquad
 \widehat\vartheta=\Omega^{-1}
   (\vartheta+\theta^A\pi^{A'}\Upsilon_{AA'}),\qquad
 \widehat\nabla_{AA'}\widehat\sigma
          =\nabla_{AA'}\sigma+\sigma\Upsilon_{AA'}.
}\end{equation} Substitution cancels the two \(\Upsilon\) terms, proving
invariance of the pullback form. Under \(\pi\mapsto a\pi\) it is
homogeneous of degree two and annihilates the Euler field, as required
for an \(L^2\)-valued one-form. On the alpha-integrable locus, the
remaining component of the basicness equation, obtained by
differentiating in a normal spin frame, is
\begin{equation}\protect\phantomsection\label{eq:gg-ordinary-section-basicness}{
 \iota_{L_A}\mathrm d(q_+^*\tau)
 =-\theta^B\pi^{A'}\pi^{B'}
       (\nabla_{AA'}\nabla_{BB'}\sigma+P_{AA'BB'}\sigma).
}\end{equation} The terms containing \(D\pi\) cancel after
\(\mu=\nabla\sigma\). To fix the curvature coefficient and its sign,
choose at the point an auxiliary scale with \(\Omega=1\), \(\Upsilon=0\)
and \(\nabla\Upsilon=P\). Then \(\widehat P=0\) and the primed spin
curvature vanishes there, since the negative Weyl spinor is zero.
Differentiating the scale-change formula at that point gives
\(\iota_{L_A}d\vartheta=-\theta^B\pi^{A'}\pi^{B'}P_{AA'BB'}\), which
gives the displayed identity after differentiating the density term. A
metric trace is killed by \(\epsilon_{A'B'}\pi^{A'}\pi^{B'}=0\). The
remaining polynomial vanishes in every direction exactly when
\((\nabla^2\sigma+P\sigma)_0=0\). The vertical restriction reads
\(\sigma\) itself, retaining the supplied scalar normalization.

Proposition 5.4 now gives \(I_\tau=(\sigma,\nabla\sigma,\rho)\) and
\(g_+=\sigma^{-2}g_{\mathrm{aux}}\), with
\(\operatorname{Ric}(g_+)=-3H(I_\tau,I_\tau)g_+\). This fixes our
numerical normalization rather than importing a cosmological-constant
convention from the cited source. In the physical scale \(\sigma=1\),
the same actual form is \(\pi_{A'}D\pi^{A'}\). All moduli inverses,
normal coordinates, unit roots, transversal hits and prolongations used
here are finite holomorphic operations with retained unit margins. Their
restrictions to parameter jets are therefore the same operations. ∎

\pandocstatementheading{\protect\phantomsection\label{lem:gg-pure-chiral-background}Lemma 7.2.}

If every negative free record is zero, the generated metric has
identically zero negative Weyl spinor. Its alpha quotient is a smooth
threefold germ along a complete line with normal bundle
\(\mathcal O(1)^{\oplus2}\), and the physical form \(\pi_{A'}D\pi^{A'}\)
descends with exactly the central normalization of Definition 2.3. The
statement with the spin factors interchanged also holds.

\emph{Proof.} Differentiate the Einstein Bianchi identity for the
negative Weyl spinor. Every commutator term contains a lower negative
Weyl derivative, multiplied by known curvature derivatives. If all lower
negative derivatives are zero, these terms vanish. The homogeneous
highest-derivative kernel is the negative free module of Appendix A.2.
Its zero free component therefore forces the next whole derivative to
vanish. Starting with order zero proves that all negative Weyl
derivatives at the mark vanish. Analyticity makes the field zero in the
germ. Positive curvature and the scalar do not supply an independent
term in this homogeneous identity.

On \(F_+=\mathbb P(S')\) the horizontal fields are
\(L_A=\pi^{A'}\nabla_{AA'}\). Their bracket modulo their span is the
primed spin curvature contracted with the two tangent factors
\(\pi^{A'}\pi^{B'}\) and the tautological spinor. The negative Weyl
contribution is zero. The trace-free Ricci contribution has the factor
\(\epsilon_{A'B'}\pi^{A'}\pi^{B'}=0\), independently of Einsteinity;
scalar contributions likewise have an alternating contraction of two
copies of \(\pi\). Hence the alpha distribution is integrable at every
numerical \(\Lambda\), using only the vanished negative Weyl spinor.
Section E.6 first constructs the separated short-leaf quotient from
these fields, without assuming an existing \(q_+\). Its initialized
finite transversal flows, closed short-leaf relation and open quotient
give the Hausdorff germ and embedded whole line. The subsequent
whole-fiber constructions of §E.2 apply to that quotient.

At the central line the normal infinitesimal quotient is
\((T_oM\otimes\mathcal O)/(\,S\otimes\mathcal O(-1)\,)
=S\otimes\mathcal O(1)\). Its four global sections are the four
spacetime variations, so this is the complete local line family. On the
homogeneous spinor lift \(D\pi=0\) along \(L_A\), and the form
\(\vartheta=\pi_{A'}D\pi^{A'}\) annihilates that field. Its exterior
derivative now has only one alpha tangent inserted, so its trace-free
Ricci contribution need not vanish by the preceding alternating
identity. Equivalently, setting \(\sigma=1\) in the basicness formula
above gives
\(\iota_{L_A}d\vartheta=-\theta^B\pi^{A'}\pi^{B'}P_{AA'BB'}\). Here
Einsteinity is used: \(P=(\Lambda/6)g\), so this expression is zero. The
physical form is therefore basic, including when \(\Lambda=0\). Its
value on the central line is literally \(\pi_{A'}d\pi^{A'}\) in the
marked frame. The spin-incidence determinant identity gives the same
calibrated root. ∎

\protect\phantomsection\label{sec:gg-negative-registration}

\subsubsection{7.2. Registering the negative
class}\label{registering-the-negative-class}

Let \(d_{\mathrm{rec}}\) decode an independently supplied ordinary
object as follows. Proposition 7.1 first reads \(g_+\) and its complete
positive free record \(U^+\). The contour \(P_Z\beta\) is then an
analytic negative massless spinor \(\psi\) on that same curved
background. Read all its symmetrized covariant free derivatives in the
original spin frame, calling them \(U^-\). Apply the original inverse
residue formulas of Theorem 3.1 to these two whole sequences, and keep
the numerical value \(-3H(I_\tau,I_\tau)\). These operations define
\(d_{\mathrm{rec}}(Z_+,L,\tau,\beta)\) without a supplied ray
construction or nonlinear Einstein metric with both chiralities.

Conversely define \(\Gamma\) on an original pair. Generate the
pure-positive background from \((U^+,0;\Lambda)\) by §§5--6 and take the
quotient and exact form of Lemma 7.2. To define the negative ordinary
class at finite order, differentiate the \emph{linear} massless Bianchi
equation and commute its derivatives on this known background. At order
\(n\) its inhomogeneous terms use only background jets and lower field
jets. The homogeneous highest kernel is \(W_n^-\) by the same
permutation argument as Appendix A.2. Choose the fixed
representation-theoretic splitting on the image, and prescribe its free
component to be \(U_n^-\). Apply the finite-jet formula of the one
operator \(K_Z\) from Appendix E to these field jets. This is the
registration rule for \(\beta\). It uses neither a nonlinear Einstein
inverse nor a newly chosen primitive at each order.

\pandocstatementheading{\protect\phantomsection\label{prop:gg-negative-registration-convergence}Proposition 7.3.}

This registration rule is consistent at every finite order and is the
Taylor rule of an actual ordinary analytic class. Its representative has
a common domain for bounded original input families, is holomorphic in
their parameters, and has exactly the original negative free records and
factorials.

\emph{Proof.} Use the joint analytic family already proved in Theorem
6.3,
\begin{equation}\protect\phantomsection\label{eq:gg-negative-registration-witness}{
 g(\varepsilon,t)=g(\varepsilon U^+,\varepsilon tU^-;\varepsilon\Lambda),
 \qquad
 \psi_\varepsilon=
 \left.\partial_t W^-\!\left(g(\varepsilon,t)\right)\right|_{t=0}.
}\end{equation} To apply Theorem 6.3, regard \((U^+,tU^-;\Lambda)\) on a
fixed disk \(|t|\le t_0\) as one bounded holomorphic input family of
Definition 2.1. Its common output has an amplitude disk of radius
greater than one and that same fixed positive \(t\) radius. Here spinors
are identified by the initialized physical normal spin frames of this
family. By Lemma 7.2, \(W^-(g(\varepsilon,0))=0\) identically.
Differentiating the Einstein Bianchi identity consequently gives
\(\nabla^{AA'}\psi_{\varepsilon,A'B'C'D'}=0\) on \(g(\varepsilon,0)\):
every connection or frame variation term multiplies a zero background
negative Weyl derivative. The same observation shows that its full free
record is \(\varepsilon U^-\).

The ordered jets of this witness obey exactly the finite linear
recurrence just specified. At each order, equality of the inhomogeneous
equations and free component, with the highest-kernel description, gives
equality of the whole jet. This proves consistency of the finite rule
without putting the witness into that rule as an input. Lemma E.2 and
Theorem E.1 then make its registered coefficients literally the
coefficients of the one representative \(K_Z\psi_\varepsilon\). Its norm
is bounded on common domains by Theorem E.1 and the interior parameter
Cauchy estimate for \(\partial_tg\). Evaluation at \(\varepsilon=1\)
gives the asserted ordinary class and convergence.

The normalization can also be checked on the line-normal graded pieces:
\begin{equation}\protect\phantomsection\label{eq:gg-negative-ordinary-graded-piece}{\begin{gathered}
\operatorname{gr}^n\!\left(L^{-6}\right)
   =\operatorname{Sym}^nS^*\otimes\mathcal O(-6-n), \\
\dim H^1=(n+1)(n+5), \\
H^0=H^2=0.
\end{gathered}}\end{equation} In the original normal monomial and Laurent bases,
changing only degree \(n\) gives
\(\Delta U^-_{n,ab}=(n-a)!\,a!\,\Delta\beta_{n,ab}\). This is the same
contour coefficient as Theorem 3.1. Curved lower-order terms have
already been fixed and do not alter this associated graded map. The
actual representative and its bound come from \(K_Z\), not from an
assumption that formal cohomology classes converge. ∎

The negative class is a massless field on a curved selected half-flat
background. It is not identified with the negative Weyl tensor of the
fully coupled metric away from the mark. The latter is reconstructed by
the full generator; registration reads its complete free input through
the derivative in the displayed witness.

\protect\phantomsection\label{sec:gg-ordinary-both-counits}

\subsubsection{7.3. Both inverse comparisons on the independent
domain}\label{both-inverse-comparisons-on-the-independent-domain}

\pandocstatementheading{\protect\phantomsection\label{thm:gg-independent-ordinary-inverses}Theorem 7.4.}

The independently defined decoder and registration satisfy
\begin{equation}\protect\phantomsection\label{eq:gg-independent-ordinary-counits}{
 d_{\mathrm{rec}}\Gamma=1,\qquad
 \Gamma d_{\mathrm{rec}}\simeq1
       \quad\text{on all of }\mathfrak W_\Lambda .
}\end{equation} The second comparison preserves the actual selected
form, the calibrated root and the ordinary class by an actual
holomorphic coboundary.

\emph{Proof.} For an original pair, Lemma 7.2 and Proposition 7.3
recover its entire positive and negative records, with the original
factors in Theorem 3.1. The scalar is unchanged. This proves the first
inverse.

Now start with any independently admitted ordinary object \(W\).
Proposition 7.1 constructs its selected background directly from its
curves and \(\tau\). That background and the pure-positive background
generated from \(d_{\mathrm{rec}}W\) have the same whole positive free
record, zero whole negative record, the same scalar and the same
point/spin mark. Appendix A.2 and the normal-metric uniqueness of
Proposition A.4 give the actual marked normal isometry. The
complete-line correspondence carries it to an actual ordinary threefold
comparison, with the calibrated root identification described in
Proposition 7.1.

Under this comparison the two selected metrics coincide. For an
auxiliary coframe their two unit densities therefore have ratio \(r\)
with \(r^2=1\) and \(r(o)=1\). On a smaller connected germ \(r+1\) is a
unit, so \(r=1\). The Einstein-bundle section map is injective, hence
the forms coincide, including their derivatives. This argument also
works over a finite local coefficient algebra because \(r+1\) has unit
residue. It does not replace preservation of the form by preservation of
its kernel or its squared norm.

The two negative massless fields are now on this same selected
background. They have the same whole free covariant jet. The linear
recurrence used in §7.2 shows inductively that their ordered jets agree:
its lower commutator terms are known, and the highest-kernel difference
is zero. Analyticity makes the fields equal. Let \(\beta'\) be the
generated representative transported by this actual ordinary comparison.
Naturality of the contour under the spin/line identification gives
\(P_Z(\beta'-\beta)=0\). Apply the zero-contour construction of §E.4
directly to that difference. It gives \(\beta'-\beta=\delta v\) with an
actual bounded ordinary zero-cochain on common smaller representatives.
This proves the second inverse for the transported representatives
themselves, with the original class rather than only its associated
graded or formal completion.

Every comparison in this proof is an initialized normal map,
complete-line hit map, unit root or the fixed operation \(K_Z\). The
proofs of §§6 and E give their common domains, holomorphic parameter
dependence and literal finite-jet restrictions. ∎

The independent ray correspondence is equally onto. Its decoder reads
the metric and full free records by §4. Normal-metric uniqueness gives
\(\mathcal G\mathcal E\simeq1\) on the independently admitted objects of
Definition 2.2, while exact record generation gives
\(\mathcal E\mathcal G=1\). Thus \(\Gamma\mathcal E\) and
\(\mathcal Gd_{\mathrm{rec}}\) are inverse object comparisons between
the ray/affine presentations and the independent ordinary presentation.
Source changes, physical arrows and reality are established next.

\protect\phantomsection\label{sec:gg-source-arrows-reality}

\subsection{8. Source comparisons, classical limits and physical
arrows}\label{source-comparisons-classical-limits-and-physical-arrows}

\protect\phantomsection\label{sec:gg-actual-source-comparisons}

\subsubsection{8.1. Actual representative
comparisons}\label{actual-representative-comparisons}

An original source is a cohomology class on the fixed marked seed. The
generating rule first applies the fixed residue and middle-Laurent
projection of Theorem 3.1. Adding an ordinary source coboundary
therefore gives exactly the same canonical pair and the same prescribed
normal output. This comparison is the identity. A raw analytic
presentation can also be retained when it satisfies the regular
contact-seed conditions and the whole unit-section condition of
Proposition 5.4, equivalently the calibrated affine admission of
Definition 5.6 and Proposition 5.7. The comparison below applies to such
retained presentations and to the independently admitted ray objects.

\pandocstatementheading{\protect\phantomsection\label{prop:gg-actual-affine-source-comparisons}Proposition 8.1.}

Two admitted analytic presentations with the same whole original
records, numerical scalar and mark have actual holomorphic
contact-affine comparisons. They intertwine the full transition maps,
compact sky lifts, projective extensions and selected root. They are
coherent under composition and holomorphic on conditioned common
parents.

\emph{Proof.} Let their selected metrics be \(g_1,g_2\). The
normal-metric uniqueness of Proposition A.4 gives the marked isometry
\(f=\operatorname{Exp}_{2}\operatorname{Exp}_{1}^{-1}\). The tangent
identification in this formula is the retained marked one. Its inverse
and spin lift are fixed by the same initialized normal maps. Let \(C_f\)
be its induced null-ray map. In ray charts \(\psi_i^a\) set
\(K_i=\psi_i^2C_f(\psi_i^1)^{-1}\) and define the unit \(d_i\) by
\(K_i^*\alpha_i^2=d_i\alpha_i^1\).

Let \(\tau_i^a\) denote the affine coordinate along the canonical lifted
incidence of §D.7. The comparison must have the form
\begin{equation}\protect\phantomsection\label{eq:gg-actual-source-affine-comparison}{
 \widehat K_i(z,\tau)
   =\bigl(K_i(z),d_i(z)^{-1}\tau+u_i(z)\bigr).
}\end{equation} Choose any one of the fixed local transversal hits on
the short ray. Define \(u_i\) there as
\(\tau_i^2\circ f^F-d_i^{-1}\tau_i^1\). This value is independent of the
hit. Indeed the contact unit and the kinematic normalization of the ray
fields give
\begin{equation}\protect\phantomsection\label{eq:gg-source-affine-slice-independence}{
 f^F_*(\sigma_1^2D_1)
       =d_i^{-1}\sigma_2^2D_2,\qquad
 D_a\tau_i^a=\sigma_a^{-2},\qquad
 D_1(\tau_i^2\circ f^F-d_i^{-1}\tau_i^1)=0 .
}\end{equation} The second density and field in the first formula are
evaluated at the image point. The unit \(d_i\) is a ray function and
hence constant under \(D_1\). The last equality follows directly from
the first two. Whole short-ray connectedness therefore makes \(u_i\) a
holomorphic ray function with the stated transversal value.

These definitions give \(\widehat K_i j_1=j_2 f^F\) on the incidence,
and hence intertwine all compact lifted skies. In point-map notation,
\begin{equation}\protect\phantomsection\label{eq:gg-actual-source-overlap-conjugacy}{
 \Phi^2_{ij}=K_i\Phi^1_{ij}K_j^{-1},\qquad
 \widehat\Phi^2_{ij}
       =\widehat K_i\widehat\Phi^1_{ij}\widehat K_j^{-1}.
}\end{equation} All functions in these compositions are evaluated at
their moving arguments. The projective extensions and their central
compatible roots are unique by §D.8, so these are comparisons of the
same whole section, not just its squared norm. In particular the
numerical scalar is retained.

The comparisons compose as follows. For successive comparisons write
\((K_{12},d_{12},u_{12})\) and \((K_{23},d_{23},u_{23})\). Their
composite has
\begin{equation}\protect\phantomsection\label{eq:gg-source-comparison-coherence}{
 K_{13}=K_{23}K_{12},\quad
 d_{13}=(d_{23}\circ K_{12})d_{12},\quad
 u_{13}=(d_{23}^{-1}\circ K_{12})u_{12}+u_{23}\circ K_{12}.
}\end{equation} These are the direct composition formulas for the
displayed affine maps. The normal isometries compose, and the two
composite affine lifts agree on the canonical incidence and therefore
agree everywhere in their chart germs. Identity and inverse comparisons
follow in the same way. Every map here is a finite exponential, parallel
transport, transversal hit, unit ratio or initialized linear ODE. The
bounds of §§6 and D give the asserted common domains. ∎

The source logarithm convention remains that of Proposition 5.8:
\(E=(\Phi^{-1})^*\) has generator \(+X_h\), whereas the stored point map
has infinitesimal velocity \(-X_h\). In an identity local model the
affine \emph{operator} generator and its bracket are
\begin{equation}\protect\phantomsection\label{eq:gg-affine-operator-bracket}{
 \widehat X_{h,\chi}=X_h+(\chi-r_h\tau)\partial_\tau,\qquad
 [\widehat X_{h,\chi},\widehat X_{k,\psi}]
  =\widehat X_{\{h,k\},
       X_h\psi-X_k\chi+r_h\psi-r_k\chi}.
}\end{equation} Here \(\mathcal L_{X_h}\alpha=r_h\alpha\). If \(a_s\) is
the actual point-translation with \(a_0=0\), then
\(\chi=-\partial_sa_s|_0\). In the section-adapted matrix of §D.9,
\(a=-b/c\), so this is \(\chi=\partial_sb_s|_0\) when \(c_0=1\). Thus
the affine coefficient is computed from the same Cauchy matrices and
section; it is not new free source data. The stored affine point
velocity is the negative of this operator vector field. Direct
differentiation gives the displayed bracket. Its mixed contact term
remains the one in Proposition 9.1, including the negative-source
self-interaction.

\protect\phantomsection\label{sec:gg-classical-cotangent-comparison}

\subsubsection{8.2. The two classical chiral
limits}\label{the-two-classical-chiral-limits}

\pandocstatementheading{\protect\phantomsection\label{thm:gg-actual-chiral-cotangent-atlas}Theorem 8.2.}

On either pure chiral axis the generated ray contact germ is actually
isomorphic, along the whole physical sky, to the projectivized cotangent
contact germ of the corresponding classical ordinary twistor threefold.
The compact sky goes to the projectivized conormal of its ordinary line.
The contact, line, selected-section and affine comparisons commute.

\emph{Proof.} Suppose \(U^-=0\). Lemma 7.2 constructs the ordinary
quotient \(Z_+\) from \(P(S')\) and identifies its whole complete line
family and exact selected form. Each null geodesic lies in a unique
alpha surface: choose its primed projective spinor at one point and use
the unique integral surface there. Total geodesy keeps the whole short
geodesic in that surface. This gives an actual holomorphic submersion
\(\rho:Z_{\mathrm{ray}}\to Z_+\). In the null incidence its kernel has
the two alpha directions and the additional unprimed projective
direction. Quotienting by the null-geodesic direction leaves
\(\ker d\rho\) of rank two.

Write \(\mathcal D\) for the ray contact distribution and
\(\mathcal V=\ker d\rho\). A ray variation within one alpha surface has
spacetime variation tangent to that null surface. Its contact value is
\(g(k,\delta x)=0\). Hence \(\mathcal V\subset\mathcal D\). The fibers
are integral, so \(\mathcal V\) is Legendrian. The quotient
\(d\rho(\mathcal D)\) is therefore a two-dimensional hyperplane in
\(T Z_+\). Its annihilator defines
\begin{equation}\protect\phantomsection\label{eq:gg-chiral-contact-annihilator-map}{
 C:Z_{\mathrm{ray}}\longrightarrow\mathbb P(T^*Z_+),\qquad
 C(z)=\operatorname{Ann}\bigl(d\rho(\mathcal D_z)\bigr).
}\end{equation} By definition its tautological contact value pulls back
to the ray contact value, with the induced value-line isomorphism. The
contact Levi pairing \(\mathcal V\times\mathcal D/\mathcal V\) is
perfect. It is exactly the derivative of this annihilator line along a
\(\rho\)-fiber, so \(dC\) is injective on \(\mathcal V\). Since \(\rho\)
is a submersion, \(C\) is a local biholomorphism.

For a fixed spacetime point \(x\), \(\rho\) sends its whole sky \(Q_x\)
to \(\ell_x\). The contact form is zero on \(TQ_x\), so \(C(Q_x)\) lies
in \(\mathbb P(N^*_{\ell_x/Z_+})\). On each unprimed ruling its source
contact line has degree one, and is the pullback of the tautological
contact line. Thus its map
\(\mathbb P^1\to\mathbb P(N^*_{\ell_x/Z_+}|_z)\simeq\mathbb P^1\) has
degree one and is an isomorphism. The base map is the line
identification, so \(C\) restricts to an isomorphism of whole compact
skies with the conormal surface.

A local biholomorphism injective on this compact surface is injective on
some neighborhood of it. Otherwise two sequences of different points
with the same image could approach the surface. Compactness would give
limiting points with the same surface image, hence the same point; its
local inverse excludes those sequences. The open image contains the
whole target surface. This proves the asserted germ isomorphism.

The form on \(Z_+\) is the descended \(\vartheta\) of Lemma 7.2. In the
physical scale of the source metric, the pullback formula of Proposition
7.1 therefore reads the density \(\sigma=1\); Proposition 5.4 evaluates
the selected ray section as this same density. The contact comparison
\(C\) preserves the tautological contact value and the sky incidence, so
the construction of §D.1 identifies their density lines with the same
central root. For each fixed ray, its Hamilton affine parameter and
these density frames then identify the two optical equations and their
Cauchy data. ODE uniqueness identifies their full solutions, transition
matrices and Wronskian isomorphism, including inverse transitions.
Finally the affine coordinate is the ratio of the same companion
solution to \(\sigma\), with derivative \(\sigma^{-2}\). The unique lift
of each compact sky in §D.7 fixes its additive constant, and §D.8 gives
the unique projective extension. This proves the contact,
selected-section and affine commutative comparison with the same
numerical scalar.

If \(U^+=0\), interchange the spin factors and the two rulings. The same
curvature contraction, compact-line and annihilator arguments give the
opposite ordinary quotient and its cotangent comparison. In the original
convention its apparent positive dual Hamiltonian is obtained from
\(F_{-6}\) by the fixed residue dictionary, not by adding another input.
The finite flow and inverse estimates used above are the ones of §§6 and
E.2, so the diagrams have common representatives in bounded families at
every numerical scalar. ∎

In the mixed case the construction retains its derived contact-affine
carrier. The theorem asserts a three-dimensional ordinary quotient on
each classical chiral axis; the independent ordinary presentation of §7
encodes the mixed input on a selected curved background with its
registered class.

\protect\phantomsection\label{sec:gg-physical-spin-arrows}

\subsubsection{8.3. Point changes and full spin
arrows}\label{point-changes-and-full-spin-arrows}

Let \(H_a=a^bH_b\) be the horizontal constant-direction Cartan record
flow of §3. For a small marked displacement \(a\) and a full complex
spin element \(A\), define, directly on whole records,
\begin{equation}\protect\phantomsection\label{eq:gg-independent-source-physical-action}{
 R_{a,A}U=A^*\!\left(\operatorname{Fl}_{H_a}^{\,1}U\right).
}\end{equation} Its admissibility means that the indicated short flow
and marked maps lie in the retained parents. This definition precedes
the ray or ordinary decoder. It retains the full spin element even when
its tensor action has a kernel.

\pandocstatementheading{\protect\phantomsection\label{thm:gg-full-spin-physical-homs}Theorem 8.3.}

The generated metric, ray and contact-affine presentations have the same
local physical arrow groupoid as the independently defined source
action. Each actual orientation-preserving metric isometry with a
specified full spin lift is represented by exactly one admissible
\((a,A)\) in these marked representatives. The independent ordinary
presentation has this same physical action through \(d_{\mathrm{rec}}\),
with actual representative comparisons and coherent finite diagrams.

\emph{Proof.} Set \(x_a=\operatorname{Exp}_o(a)\) and let \(\nu_a\) be
the initialized spin parallel transport along its radial geodesic. If
\(V=R_{a,A}U\), the record-flow equations say that \(V\) is precisely
the full free record at \(x_a\) in the frame \(\nu_a A\). Normal-metric
uniqueness therefore gives the actual isometry
\begin{equation}\protect\phantomsection\label{eq:gg-physical-arrow-normal-isometry}{
 \Xi_{a,A}:g[V]\longrightarrow g[U],\qquad
 v\longmapsto
 \operatorname{Exp}_{x_a}^{\,g[U]}(\nu_a A v).
}\end{equation} Its spin differential is the specified initialized lift.
The covariant source instruction \(U\to V\) uses \(\Xi_{a,A}^{-1}\); the
displayed metric map has the opposite direction. Their ray and affine
maps are the maps of Proposition 8.1 with that direction.

For the group law put \(h_U(a,c)=\operatorname{Exp}_o^{-1}
 \operatorname{Exp}_{x_a}(\nu_ac)\). Let \(K_U(a,c)\) be the spin
rotation comparing direct radial transport to transport first to \(x_a\)
and then to this endpoint along the short second geodesic. Explicitly it
is direct-radial transport inverted, composed with broken-path
transport; it is initialized to the identity. Actual exponential
composition gives
\begin{equation}\protect\phantomsection\label{eq:gg-source-physical-composition}{\begin{gathered}
(b,B)\circ(a,A)=
   \bigl(h_U(a,Ab),K_U(a,Ab)AB\bigr), \\
(a,A)^{-1}=\bigl(-A^{-1}a,A^{-1}\bigr)\quad\text{at }V.
\end{gathered}}\end{equation} Associativity follows from map composition, including
its initialized spin transports. On the reversed radial segment the
direct and broken transports cancel, which proves the inverse formula.
For any metric isometry with spin lift, its image of the marked point
determines \(a\) and its differential relative to \(\nu_a\) determines
the full \(A\). An isometry is fixed by these data since it intertwines
the initialized exponential maps. This proves existence, uniqueness and
the bijection of physical Hom sets. An element in the kernel of the
tensor representation still has its distinct spin lift and is not
removed.

For independently supplied ordinary objects \(W,W'\), define physical
admission by the independently computed condition
\begin{equation}\protect\phantomsection\label{eq:gg-independent-ordinary-physical-admission}{
 d_{\mathrm{rec}}W'=R_{a,A}\,d_{\mathrm{rec}}W .
}\end{equation} This uses complete-line reconstruction, the ordinary
section and the contour on the supplied objects. It does not invoke an
image test. The actual action on representatives is computed by
\(\Gamma R_{a,A}d_{\mathrm{rec}}\) together with the two counits of
Theorem 7.4. Its composition and inverse diagrams are the ones just
proved, followed by their ordinary comparisons. In a mixed field this
operation rebuilds both the selected background and its registered
class; there need not be a point map of one fixed old background
threefold.

The enriched ray maps preserve \(I,X\) and \(DX\). These generate the
canonical six-dimensional fiber because the scale-plane Gram determinant
is \(-1\) and the solder is nondegenerate. Thus no independent vertical
comparison remains once the physical map is fixed. The full spin lift is
retained separately; it cannot be inferred from its action on tensor
cochains. All constructions in this proof are finite exponential,
transport, hit and unit operations. On any fixed finite conditioned
point/spin diagram the estimates of §6 give common representatives and
finite-jet compatibility. ∎

\protect\phantomsection\label{sec:gg-lorentz-reality}

\subsubsection{8.4. Lorentzian reality}\label{lorentzian-reality}

The fixed Lorentz spin identification, with the signature and scalar
conventions of §2, defines an antilinear involution \(\mathfrak c\) on
the original whole records: conjugate coefficients, exchange the two
spin factors and the two chiral records, and conjugate \(\Lambda\).
Transfer it to the original source classes by the exact two residue
dictionaries of Theorem 3.1. Their different line-bundle weights \(2\)
and \(-6\) are retained; both source classes have cohomological degree
one.

\pandocstatementheading{\protect\phantomsection\label{prop:gg-coupled-lorentz-reality}Proposition 8.4.}

For real numerical \(\Lambda\), the \(\mathfrak c\)-fixed original data
generate all marked real analytic Lorentzian Einstein germs, with their
full real spin arrows. The ray and affine constructions carry their
actual induced reality. On the independent ordinary presentation the
same reality is the coupled operation
\(\Gamma\mathfrak c\,d_{\mathrm{rec}}\).

\emph{Proof.} The Cartan record equations and the canonical normal
connection are real under the declared spin identification and real
tangent form. Their initialized analytic ODEs therefore give conjugate
coframes from conjugate data. Uniqueness makes the coframe and metric
real on the Lorentz tangent slice for fixed data. Their nondegenerate
central Lorentz form fixes the signature after shrinking. Conversely the
curvature records of every real analytic Lorentz Einstein germ have
exactly this reality, and the normal-metric inverse reproduces that
germ.

Actual null transport, contact forms, optical solutions and compact
affine lifts are natural under this conjugation. Unit roots have the
fixed central branch. The contact-affine comparisons of Proposition 8.1
provide the same induced reality on all admitted representative charts.
Real point changes and full real spin lifts are the restrictions of
Theorem 8.3.

The ordinary decoder is independent, so its reality condition is that
its whole decoded records are fixed by \(\mathfrak c\). Applying
\(\Gamma\) gives its actual real representative; Theorem 7.4 supplies
the comparison to the original object. The resulting operation is
involutive through that actual counit and obeys the same coherent
diagrams. It generally changes the selected chiral background and
registered class together. It is not a requirement for a pointwise
antiholomorphic involution on a single old background. ∎

\protect\phantomsection\label{sec:gg-main-proof}

\subsubsection{8.5. Proof of the main
correspondence}\label{proof-of-the-main-correspondence}

\emph{Proof of Theorem 2.4.} The two original analytic residue
dictionaries are Theorem 3.1, with both factorial normalizations fixed.
Lemma 3.2 and Appendix A supply the polynomial curvature completion.
Theorem 3.3 proves convergence of its whole-record Cartan flows with all
derivative slots and the ordered reserve of Appendix B.

The source Hamiltonians and their interaction are the original ones of
§5 and Proposition 9.1. The explicitly fixed cochain splittings and
zero-record symbol inverse are proved in Appendix C; they correct the
optical defect while retaining the complete free record. Proposition 5.8
fixes the operator/point orientation. Theorem 5.5 is the finite
generating law. Proposition 5.9 fixes its trial execution and proves
finite-prefix source dependence, including the normalization of its own
derived geometry. The prescribed normalization and coefficient
comparison of Theorem 6.3 prove common convergence of its maps, inverse
maps, section, incidence and canonical connection.

Propositions 4.1--4.3 and 5.4 recognize the canonical normal Einstein
geometry and selected scalar. Proposition A.4 proves uniqueness from the
whole free record and scalar. Every marked analytic Einstein germ
supplies analytic original classes by Theorem 3.1; applying the
generator and this uniqueness gives its marked inverse. The independent
ray recognition and the same argument give both ray inverses. Definition
5.6 and Proposition 5.7 and Appendix D give the intrinsic optical and
affine presentations and their two inverse comparisons.

Proposition 7.1 decodes all independently admitted ordinary objects.
Proposition 7.3 and Appendix E give the registered class and its common
representative inverse. Theorem 7.4 proves both ordinary counits on
their entire independent domain. Proposition 8.1 proves source
independence and coherent chart comparisons. Theorem 8.2 proves both
classical chiral cotangent limits by atlas maps along the whole compact
sky. Theorem 8.3 gives the independent point action, its full spin Hom
sets and the ordinary physical action. Proposition 8.4 gives the Lorentz
real restriction.

Every scalar occurrence in these constructions is polynomial, an
initialized finite ODE coefficient or a unit expression with a
scalar-independent central residue. The scale-plane Gram determinant is
\(-1\); no division by \(\Lambda\) occurs. All assertions therefore
include numerical zero and the finite Taylor restrictions of holomorphic
scalar families. The common-domain bounds concern bounded original
inputs and conditioned finite diagrams, as stated; they make no
assertion about arbitrary unnormalized gauges or unbounded histories.
These observations establish all the claimed comparisons. ∎

\protect\phantomsection\label{sec:gg-active-realization}

\subsection{9. Active twistor
realization}\label{active-twistor-realization}

Penrose's active-patching proposal uses local Heisenberg algebras,
structure-preserving comparisons, nonlocal point recognition and an
infinity structure (Penrose 2015, secs. 3--6). We first analyze the
source interactions and calibrated affine preservation, then pass to the
analytic operator presentation in Propositions 9.6--9.9. Theorems
9.10--9.11 give its independent-domain correspondence and the
simultaneous structural conclusions.

The following calculations explain the nonlinear correction and its
intrinsic realization. The calibrated affine structure of Proposition
5.7 supplies explicit local models, overlap and inverse laws, and an
independently defined holomorphic admission. Its source comparisons and
classical limits concern these actual representatives.

\pandocstatementheading{\protect\phantomsection\label{prop:gg-mixed-lift-nonclosure}Proposition 9.1.}

The sum of the two bare contact-lift subspaces in equation
(\ref{eq:gg-two-contact-lifts}) is not closed under the actual mixed
contact bracket. Negative self-brackets do close in their dual
Hamiltonian subalgebra.

\emph{Proof.} Rename the contact coordinates \((a,b,r,s,\zeta)\), with
\(\alpha=\mathrm da+r\,\mathrm db+s\,\mathrm d\zeta\),
\(e=-a-rb-\zeta s\) and \(c=rb+\zeta s\). Positive lifts \(F_b-rF_a\)
are independent of \(s\). Negative lifts \(-G_s+\zeta G_e\), with
\(G=G(e,s,r)\), are annihilated by \(T=\partial_b-r\partial_a\). Since
\([T,\partial_s]=0\), every sum satisfies \(T\partial_s h=0\).

For \(F=a^4/\zeta\), \(G=e^4/r\), the exact Hamiltonians are
\(h=-4ra^3/\zeta\) and \(k=-4\zeta(a+c)^3/r\). Substitution in the
contact field gives \(X_hk=48a^2(a+c)^2(a+2c)\) and
\(-h_a k=-48a^2(a+c)^3\). Hence
\begin{equation}\protect\phantomsection\label{eq:gg-mixed-nonclosure}{
 M=\{h,k\}=48a^2c(a+c)^2,\qquad
 T\partial_sM=-288r\zeta\,ac(a+c)\not\equiv0 .
}\end{equation} The second formula uses \(T=r(\partial_c-\partial_a)\)
and \(\partial_s=\zeta\partial_c\) on these polynomials. It rules out
membership in the two bare lift spaces even locally on the indicated
overlap.

For negative self-interaction the dual contact form is
\(\mathrm de+\zeta\,\mathrm ds+b\,\mathrm dr=-\alpha\). The base field
\(v_G=G_s\partial_e-G_e\partial_s\) has the stated negative cotangent
lift. Functoriality of cotangent lift gives
\begin{equation}\protect\phantomsection\label{eq:gg-negative-self-bracket}{
 \{h_-(G),h_-(K)\}=h_-(B(G,K)),\qquad
 B(G,K)=G_sK_e-G_eK_s .
}\end{equation} For \(G=e^4/r\), \(K=s^4/r\), this is
\(B=-16e^3s^3/r^2\ne0\). ∎

This is a statement about the specified local lift spaces. It does not
exclude additional operators, nonlinear corner comparisons, a larger
ordinary structure or a noncommutative realization. The mixed operation
presently lives in the derived contact algebra; the linear ordinary
registration alone is not its geometric realization.

\pandocstatementheading{\protect\phantomsection\label{prop:gg-mixed-einstein-defect}Proposition 9.2.}

There are two exact opposite-chirality vacuum waves whose
contact-compatible mixed trial has a nonzero obstruction to extending
their selected Einstein scale. The obstruction is removable at the mixed
order.

\emph{Proof.} In complex coordinates take
\begin{equation}\protect\phantomsection\label{eq:gg-two-wave-trial}{
 \eta=2\,\mathrm du\,\mathrm dv-2\,\mathrm dx\,\mathrm dy,\qquad
 g=\eta+t f(u)\,\mathrm dx^2+s q(x)\,\mathrm dv^2 .
}\end{equation} The inverse has \(g^{uv}=1\), \(g^{xy}=-1\),
\(g^{uu}=-sq\), \(g^{yy}=-tf\), with their symmetric entries and no
others. Its determinant is one. The entire nonzero Christoffel list,
apart from lower-index symmetry, is
\begin{equation}\protect\phantomsection\label{eq:gg-two-wave-christoffels}{
 \begin{aligned}
 \Gamma^v_{xx}&=-tf'/2,&
 \Gamma^u_{xx}&=tsqf'/2,&
 \Gamma^y_{ux}&=-tf'/2,\\
 \Gamma^y_{vv}&=sq'/2,&
 \Gamma^u_{vx}&=sq'/2.
 \end{aligned}
}\end{equation} In the order \((x,v,u,y)\) all connection matrices are
strictly triangular. Thus their traces and product traces vanish. The
Ricci formula reduces to the divergence of this list, giving exactly
\begin{equation}\protect\phantomsection\label{eq:gg-two-wave-ricci}{
 \operatorname{Ric}(g)=\tfrac12tsqf''\,\mathrm dx^2,
 \qquad \operatorname{Scal}(g)=0 .
}\end{equation} Each single wave is an exact vacuum metric. Their
curvature two-forms are proportional to \(\mathrm du\wedge\mathrm dx\)
and \(\mathrm dx\wedge\mathrm dv\) respectively. Exact null coframes for
the two metrics are
\(\left(\begin{smallmatrix}\mathrm du&\mathrm dx\\\mathrm dy-(tf/2)\mathrm dx&\mathrm dv\end{smallmatrix}\right)\)
and
\(\left(\begin{smallmatrix}\mathrm du+(sq/2)\mathrm dv&\mathrm dx\\\mathrm dy&\mathrm dv\end{smallmatrix}\right)\),
with metric twice their determinant. The two curvature forms occupy a
common row and a common column, respectively, and therefore have
opposite chiralities.

Put \(f=u^3\), \(q=x^2\) and work over \(A=\mathbb C[t,s]/(t^2,s^2)\)
with highest ideal \((ts)\). Then
\begin{equation}\protect\phantomsection\label{eq:gg-mixed-cotton}{
 [ts]P^\circ=\tfrac32ux^2\,\mathrm dx^2=:A_5,\qquad
 [ts](\nabla_uP_{xx}-\nabla_xP_{ux})=\tfrac32x^2 .
}\end{equation} The lower Schouten tensors vanish, so the highest
derivative is flat. The tensor \(A_5\) has zero flat divergence, but its
class in \(\mathcal Q_5\) is nonzero. Indeed if
\(A_5=\operatorname{Hess}_0\phi\), then \(\eta_{xx}=\eta_{ux}=0\) gives
\(\phi_{xx}=3ux^2/2\) and \(\phi_{ux}=0\). Commuting their next
derivatives contradicts \(\phi_{uxx}=3x^2/2\ne0\).

This is an admitted contact trial, rather than an assumed Einstein
output. The short null-geodesic and transverse-hit construction applied
to the displayed nondegenerate analytic metric gives actual contact
cocycles and compact skies, without an Einstein hypothesis. Its first
variations have classes \(J[h_{\mathrm{wave}}]\),
\(J[k_{\mathrm{wave}}]\), where \(h_{\mathrm{wave}}=f(u)\mathrm dx^2\)
and \(k_{\mathrm{wave}}=q(x)\mathrm dv^2\). The linear vacuum-record
isomorphism and analytic corner splittings of Appendix C identify these
with the original two encoded lift classes and make their
representatives agree. Pulling the lower selected states through the
same comparison retains the highest nonzero optical class. This
establishes an allowed contact-cocycle trial; it does not assert that a
particular fixed two-cocycle repair selects this exact intermediate
coefficient.

For an explicit mixed correction put \(H=3u^2vx^2/2\). The tensor
\(H\,\mathrm dx^2\) has zero flat trace and divergence, and
\begin{equation}\protect\phantomsection\label{eq:gg-mixed-explicit-repair}{
 \Box_\eta H=6ux^2,\qquad
 \operatorname{Ric}^{(1)}(H\,\mathrm dx^2)
                           =-3ux^2\,\mathrm dx^2 .
}\end{equation} Thus \(g+tsH\,\mathrm dx^2\) is exactly Ricci-flat over
\(A\). Any highest free-record change can subsequently be removed by the
flat vacuum inclusion of Appendix C. This last operation uses that
proved inclusion; it is not part of the displayed finite Ricci
calculation. ∎

The inverse in Theorem 5.5 has a precise content:
\begin{equation}\protect\phantomsection\label{eq:gg-meaning-of-inverse}{
 L_m=\left(o_m|_{\ker\Pi_{m+1}}\right)^{-1}
       :\ker d_m\longrightarrow\ker\Pi_{m+1},\qquad
 o_mJ[h]=[P(h)^\circ].
}\end{equation} It removes the linearized Einstein-scale obstruction
modulo homogeneous Hessians, with the entire free record fixed. It is
not an inverse on arbitrary sources or on all metric gauges. Once the
actual whole optical section glues, Einsteinity follows by Proposition
5.4; establishing that gluing and deriving its Einstein consequence are
distinct steps.

A tempting simplification would attach the optical bundle to the local
contact structure without the complete family. The next obstruction
rules out a specific such attempt.

\pandocstatementheading{\protect\phantomsection\label{prop:gg-contact-local-optical-obstruction}Proposition 9.3.}

No rank-two bundle associated to a holomorphic representation of a
finite-order full contact-jet frame group of bare complex contact
five-manifolds can be the optical bundle on the flat ray germ, with the
identification intertwining every local conformal comparison between
neighboring flat sky germs, including small translations and conformal
dilations with their density-line action.

\emph{Proof.} In symmetric Darboux coordinates \((z,v)\),
\(v\in H=\mathbb C^4\) symplectic, write
\(\alpha=\mathrm dz+\omega(v,\mathrm dv)/2\). The isotropy contains
\(\operatorname{Sp}_4(\mathbb C)\). Its action on a two-dimensional
fiber is trivial: a nonzero Lie algebra map from the simple
ten-dimensional \(\mathfrak{sp}_4\) to \(\mathfrak{gl}_2\) would be
injective, which is impossible. Connectedness gives the same assertion
for the group.

Every allowed horizontal Hamiltonian module
\(z^a\operatorname{Sym}^dH^*\) with \(d>0\) consequently acts by zero.
Its map to the trivially acted-on fiber endomorphisms would be an
invariant functional; none exists, since the symplectic group is
transitive on \(H\setminus0\). For horizontal linear forms \(\ell,m\)
with nonzero symplectic pairing, the contact bracket is
\begin{equation}\protect\phantomsection\label{eq:gg-contact-isotropy-bracket}{
 \{z^a\ell,z^bm\}
  =c\,z^{a+b}
       +\tfrac{b-a}{2}z^{a+b-1}\ell m,\qquad c\ne0 .
}\end{equation} Taking \(a=1\), \(b=k-1\) kills the pure powers \(z^k\)
for every \(k\ge2\), because the inputs and the second output already
act by zero. The identity holds in every finite jet quotient where that
output survives. The remaining action factors through contact dilation.
The symplectic and unipotent isotropy parts are connected, and the
quotient is \(\mathbb C^*\), so the representation is a sum of integer
contact-line characters: \(E=L^a\oplus L^b\).

On a flat sky \(L|_Q=\mathcal O(1,1)\) and the optical restriction is
\(\mathcal O\oplus\mathcal O(1,1)\). With generators \(h_1,h_2\) its
Chern classes imply \(a+b=1\) and \(2ab\,h_1h_2=0\). Hence
\(\{a,b\}=\{0,1\}\). The trivial character gives a nonzero natural
global section invariant under the contact symmetries.

The actual whole-neighborhood optical sections are
\(\sigma(x)=A+B_ax^a+C\eta(x,x)\). This follows by integrating
\(\operatorname{Hess}_0\sigma=0\). Write
\(\partial_a\partial_b\sigma=\tau\eta_{ab}\) and commute \(\partial_c\)
with \(\partial_a\) to get
\((\partial_c\tau)\eta_{ab}=(\partial_a\tau)\eta_{cb}\). Contracting
\(a,b\) gives \(4\partial_c\tau=\partial_c\tau\), so \(\tau\) is
constant. Translation invariance forces such a section to be constant.
Under the conformal dilation \(x\mapsto rx\), expressed back in the
fixed flat representative, its weight-one action is
\(\sigma(x)\mapsto r^{-1}\sigma(rx)\). No nonzero constant is invariant.
The natural identification would carry the preceding nonzero invariant
section to an impossible one. Small translations compare neighboring sky
germs; no large global quotient is needed. ∎

The hypotheses specify the finite-contact-jet setting and the conformal
comparisons covered by Proposition 9.3. The present optical construction
uses the complete compact family and falls outside them. This
proposition says nothing against constructions using that family, a
Legendrian seed, infinite-order operations or noncommutative algebras.

The intrinsic bundle of Proposition 5.3 and the extension condition of
Definition 5.6 use this complete compact family. The following corollary
expresses their calibrated affine preservation law.

\pandocstatementheading{\protect\phantomsection\label{prob:gg-intrinsic-active-realization}

\protect\phantomsection\label{cor:gg-active-affine-realization}Corollary 9.4.}

The coupled source generator has the explicit calibrated affine
realization of Proposition 5.7 on its derived contact carrier. Its
preservation problem is the independently defined holomorphic admission
of Definition 5.6. Its highest obstruction is the actual cochain in
equation (\ref{eq:gg-affine-obstruction-cochain}), and the stated
contact/optical correction solves that same problem. These
representatives and their inverses have the common analytic bounds of
§6.

\emph{Proof.} The section equation makes \(T_{ij}\) triangular with
fixed first column and determinant \(c_{ij}\). Equation
(\ref{eq:gg-affine-translation-from-source}) gives its affine coordinate
and the exact overlap law, including all moving arguments. Conversely
that affine law recovers the original Cauchy matrix. Sections D.7--D.8
derive the compact lift, root and whole section from the independent
admission, so the preservation condition is defined before Einstein
recognition. Section D.9 identifies its failure with the original
target-minus-source defect; Theorem 5.5 therefore supplies the same
correction at every finite order. Theorem 6.3 gives its analytic sum.

Projection to the contact component preserves composition and brackets.
Thus the negative self-operation and mixed operation already present in
the source law are retained. They are not two separately reconstructed
metrics: Proposition 9.2 exhibits their nonzero common obstruction
before correction. The normal-bundle and projective-extension
requirements are essential geometric restrictions on the coupled gluing.
∎

The ordinary counits and physical comparisons of §§7--8 apply to this
same contact-affine construction and its complete families.

\pandocstatementheading{\protect\phantomsection\label{prop:gg-fixed-source-projection-obstruction}Proposition 9.5.}

The displayed original contact lifts need not descend through the fixed
cotangent projection to point transformations of \(T_0\).

\emph{Proof.} In the marked chart let
\(\operatorname{pr}(w^0,w^1,\zeta,\varrho,s)=(w^0,w^1,\zeta)\). The
positive lift has base components
\((\partial_{w^1}F_2,-\partial_{w^0}F_2,0)\), independent of the
cotangent directions, and is projectible. For the allowed original
negative class \(F_{-6}=24/\zeta\), equation
(\ref{eq:gg-derived-dual-class}) gives \(G=(\eta^{0'})^4/\varrho\). Put
\(A=w^0+\varrho w^1+\zeta s\). Then
\begin{equation}\protect\phantomsection\label{eq:gg-fixed-source-projection-obstruction}{
 h_-=-4\zeta A^3/\varrho,\qquad
 X_{h_-}\zeta=-12\zeta^2A^2/\varrho,\qquad
 \partial_s(X_{h_-}\zeta)=-24\zeta^3A/\varrho\ne0 .
}\end{equation} Thus the base component varies along a cotangent fiber
on the overlap where the displayed denominators are units. No vector
field on \(T_0\) can be \(\operatorname{pr}\)-related to this contact
field. Negating it for the stored point velocity does not remove that
obstruction. ∎

The obstruction concerns the fixed cotangent projection. On the opposite
pure axis, Theorem 8.2 uses the dual polarization for the classical
comparison. In the mixed case the analytic bijection
\(F_{-6}\leftrightarrow G\) determines \(G\) from the original negative
class, independently of \(F_2\). Section F.1 represents the resulting
nonprojectible operation by derivative variables in the same original
twistor operator model; the independent admission and point comparisons
are proved in Appendix H.

\pandocstatementheading{\protect\phantomsection\label{prop:gg-finite-source-operator-realization}Proposition 9.6.}

Both original contact source operations have order-one microdifferential
representatives in the same original twistor coordinates. At every
finite amplitude order, the corrected active overlaps have compatible
filtered, adjoint-preserving operator comparisons and exact coherent
inner triple comparisons. Their first negative-order unit sheaf
transports precisely the calibrated affine torsor. Each fixed admitted
analytic contact-affine object also has its canonical structured
quantization and framed simple point objects on its complete skies.

\emph{Proof.} Appendix F.1 constructs the source operators and proves
their analytic symbol bounds. Theorem F.1 proves the finite comparisons
and all triple and quadruple identities. Proposition F.2 identifies the
derived-sheaf torsor transport with the original affine law, retaining
the moving arguments. Section F.4 applies the fixed-manifold
quantization and structured uniqueness theorems with the filtration,
full symbols and adjoint coherences retained. Proposition F.3 proves the
framed point assertion using the symplectization of each sky. ∎

This proposition gives finite-amplitude operator descent and a
fixed-object enrichment. Its torsor transport is the derived-sheaf
statement of Proposition F.2; the quantum commutator correction and the
distinct unit-neutralization problem are described in §F.3. Analytic
relative overlaps and moving point modules are supplied by Propositions
9.7--9.8.

\pandocstatementheading{\protect\phantomsection\label{prop:gg-analytic-source-operator-realization}Proposition 9.7.}

The actual corrected contact overlaps have analytic relative filtered
operator comparisons and exact normalized inner coherence on common
bounded-family parents, including amplitude one. Their classical symbols
and calibrated affine cochains are the existing actual source tuples.
Their amplitude jets agree with the auxiliary finite operator
construction up to compatible normalized filtered inner equivalence.

\emph{Proof.} Theorem F.4 proves analytic transport on the classical
tubes already furnished by Theorem 6.3, with coefficientwise locality
and factorial symbol bounds. Lemma F.5 constructs the analytic
negative-order units needed for coherence. Theorem F.6 applies these
results to the actual overlap family, retaining exact inverse
orientation and the same classical symbols. Proposition F.7 gives the
finite comparison and the literal preservation of the affine cochains. ∎

The moving supported modules for these analytic overlaps are constructed
next.

\pandocstatementheading{\protect\phantomsection\label{prop:gg-relative-sky-operator-modules}Proposition 9.8.}

The analytic operator realization has actual holomorphic families of
supported simple point modules on its complete skies, with the
calibrated relative half-form and exact associator-compatible descent.
They have common bounded-family estimates and natural finite parameter
jets. With the actual whole-parameter framing, the constructed locally
standard modules are unique up to unique framed comparison.

\emph{Proof.} The sky embeddings are the actual complete families of
Theorem 6.3. Their fixed ruled parameterization and marking supply the
relative root \(\mathcal O(-1,-1)\) and its calibrated frames, including
any retained base factor. Lemma G.1 constructs the analytic local module
and proves division by its finite defining ideal. Lemma G.2 constructs
tangent transport with uniform factorial bounds. Theorem G.3 constructs
joint Hom generators and proves the exact equation
\(\tau_{ij}\tau_{jk}=C_{ijk}\tau_{ik}\), keeping the full operator
associator. Proposition G.4 proves the framed relative uniqueness. ∎

The framing is an actual nonzero holomorphic section along the
calibrated cone lift for all parameters. A single value at a central
parameter would be insufficient. These modules are constructed from the
same complete source sky family; they do not assert neutralization of
the algebroid by a global free object.

\pandocstatementheading{\protect\phantomsection\label{prop:gg-op-physical-source-comparisons}Proposition 9.9.}

The operator realization and its moving supported modules carry the
actual physical point/full-spin comparisons of §8, their coherent finite
diagrams and the coupled Lorentz operation. Their symbols, affine data,
selected section and full spin labels are the original ones.

\emph{Proof.} Lemma F.8 gives actual relative normalized inner
comparisons for maps with the same full symbols. Theorem G.5 applies it
to the actual source physical maps, with the nonlinear physical
amplitude and independent internal construction amplitude distinguished.
It retains both associators in descent, constructs the module
comparisons using the actual flattened generators, and gives the
antilinear square and its framed-module coherence with the compatible
Lorentz root calibration. The physical Hom assertion is over exactly the
enriched admission of Theorem 8.3; full spin labels are retained
separately from their possible scalar kernels. ∎

Appendix H defines an independent operator domain and proves the two
actual counits needed for the following correspondence. Its data retain
the analytic source operations and complete calibrated point families of
Propositions 9.6--9.9.

\pandocstatementheading{\protect\phantomsection\label{thm:gg-independent-source-operator-correspondence}Theorem 9.10.}

The actual source operator construction gives a correspondence between
the whole original analytic domain \(\mathcal A_\Lambda\) and the
independently admitted operator data \(\mathfrak O_\Lambda\) of
Definition H.1. Its source decoder and constructor have both actual
inverse comparisons on their entire stated domains, and are bijective on
physical point/full-spin Hom classes. The inverse preserves the same
coupled higher preservation law, selected optical section and derived
numerical scalar. It includes both complete independent curvature
records, all algebraic Weyl types, every complex numerical \(\Lambda\)
including zero, bounded holomorphic families, their finite jets,
conditioned finite comparison diagrams and the coupled Lorentz
restriction.

\emph{Proof.} Theorem H.4 gives the two actual inverses between the
independently specified operator domain and Definition 5.6's calibrated
affine domain. Definition 5.6 and Proposition 5.7, D.7--D.8 and Theorem
2.4 give the latter's inverse comparisons with all marked Einstein germs
and their original analytic records. Composing these actual comparisons
supplies the decoder on every independently admitted operator object.

For the other direction use the actual source operators of F.1, their
corrected analytic overlaps of Theorem F.6 and the calibrated module
family of Appendix G. Their full symbols and first-unit torsor are
literally the original actual contact-affine tuple. The counit of H.4
compares this supplied source operator object with the constructed
quantization of that tuple. Thus composing the classical counit and
H.4's counit compares the decoded-and-regenerated source object to the
supplied independent object, with both associators and all module maps.
The opposite source composite is the initialized comparison with the
same whole records. Theorem H.4's uniqueness proves the unit/counit and
finite naturality diagrams.

Proposition H.5 identifies the actual source operation and the full
holomorphic preservation obstruction, with the same \(o_m\), \(L_m\) and
fresh correction. Einsteinity follows from the independent admission by
Proposition 5.4. The physical/full-spin and coupled Lorentz assertions
are those of H.4 and §8, with the labels retained. The source estimates
of §§3--6 and the conditioned operator/module estimates of F--H give the
common families and their actual finite jets. No Weyl-type determinant
or \(\Lambda\) is inverted. ∎

The next theorem collects the operator, infinity-structure and
point-recognition properties of this correspondence. Its proof uses the
explicit comparisons in Appendix I.

\pandocstatementheading{\protect\phantomsection\label{thm:gg-original-positive-realization}

\protect\phantomsection\label{thm:gg-structured-googly-realization}Theorem 9.11.}

For every \(\Lambda\in\mathbb C\), let \(\mathcal S_{\mathrm{op}}\) and
\(\mathcal D_{\mathrm{op}}\) be the actual source constructor and
decoder of Theorem 9.10. They have the following simultaneous
properties.

\begin{enumerate}
\tightlist
\item
  The two unrestricted original classes in \(\mathcal A_\Lambda\) act by
  skew order-one analytic operators in the same half-density model of
  F.1. Their coupled gluing retains the nonzero mixed principal brackets
  and lower-order microdifferential commutator terms; no independent
  dual-twistor input is added.
\item
  Einstein admission is the independently holomorphic
  projective-extension predicate of Definitions 5.6, H.1 and I.2. Under
  Proposition I.3 it is equivalent to the calibrated parallel normal
  tractor and local-twistor infinity forms. Its pre-admission defect and
  correction are the same actual \(o_m\), \(L_m\) and six normalizations
  of Theorem 5.5 and Proposition H.5.
\item
  A physical point is recognized from a compact supported simple object.
  Proposition H.3 constructs its complete four-dimensional family and
  full null Segre quadric, with actual associator-compatible module
  comparisons.
\item
  On the entire original and independently admitted domains,
\end{enumerate}

\begin{equation}\protect\phantomsection\label{eq:gg-structured-source-operator-inverses}{
 \mathcal D_{\mathrm{op}}\mathcal S_{\mathrm{op}}=1_{\mathcal A_\Lambda},
 \qquad
 \mathcal S_{\mathrm{op}}\mathcal D_{\mathrm{op}}
                \simeq1_{\mathfrak O_\Lambda}.
}\end{equation} The first equality is on the original marked input
classes and their numerical scalar; the comparison on operator objects
and all arrows is the actual one of Theorem 9.10.

The comparisons are bijective on physical point/full-spin Hom classes
and retain both complete curvature records, every Weyl type, numerical
\(\Lambda=0\), the bounded families and their actual finite jets,
conditioned finite diagrams and the coupled Lorentz restriction of
Theorem 9.10. The convergent stored representatives are those selected
in Theorem 6.3.

\emph{Proof.} Theorem I.4 proves the operator, infinity-structure and
point-recognition assertions, using the actual source formulas and
comparisons of F--H and Proposition I.3. Lemma I.1 and Proposition H.5
identify the pre-admission reduction defect with the same optical
cochain and fixed inverse, followed by fresh recomputation. Theorem 9.10
supplies both inverse comparisons on the whole independent domains, with
the stated physical Hom sets. Theorem 6.3 and the relative estimates of
F--H supply the analytic family and jet assertions. All hypotheses are
the original input and independent admission conditions already stated.
∎

\pandocstatementheading{\protect\phantomsection\label{thm:gg-unified-projective-gravitational-correspondence}Theorem 9.12.}

For the original input of Definition 2.1 and the independent operator
domain of Definition H.1, the correspondence of Theorems 2.4 and
9.10--9.11 preserves, through one natural family of comparisons, all of
the following:

\begin{enumerate}
\tightlist
\item
  The same two original classes, their nonlinear source operations and
  six scale normalizations, with the convergent prescribed
  representatives of Theorem 6.3.
\item
  The complete compact point-module family, the rank-four raywise
  local-twistor bundle \(\mathcal T_{\mathrm{LT,ray}}\) with its volume
  and flag, and the optical point planes
  \(K_{x,\gamma}=\ker\mathcal F_{\mathrm{spin}}(X(x))\).
\item
  The same torsor/root-selected infinity pair, optical evaluation and
  Wronskian, with its affine slice \(H(I,v)=1\) and compact evaluation
  \(X/\sigma\).
\item
  Both actual inverse comparisons on the whole independently defined
  domains, including the associator-compatible module maps and physical
  point/full-spin Hom classes.
\end{enumerate}

Einstein admission is the holomorphic projective-extension predicate;
under the comparison it is equivalent to descent of that same selected
infinity form. Every marked local analytic Einstein germ, with arbitrary
Weyl type and numerical \(\Lambda\in\mathbb C\) including zero, is
recovered. On the Lorentz restriction the rank-four comparison is the
antilinear duality
\((\mathcal T_{\mathrm{LT,ray}})_y\to (\mathcal T_{\mathrm{LT,ray}}^*)_{\iota(y)}\)
of Proposition J.7, with the coupled square of G.5. These assertions
hold for the bounded holomorphic families, actual finite jets and
conditioned finite diagrams of Definition 2.1.

\emph{Proof.} Theorems 2.4 and 9.10 provide the source, geometric and
independent operator correspondences. Proposition J.1, Theorem J.2 and
Proposition J.3 construct the pre-admission normal spin transport and
raywise local-twistor descent, with the same selected density and actual
comparison data. Proposition J.4 identifies its optical evaluation with
the point-plane pencil and supplied affine slice. Proposition J.5
identifies the pre-admission source defect and its correction. The
natural comparisons of Theorem J.6 combine these identifications with
the actual counits of H.4 and 9.10 on every independent object.
Proposition J.7 proves the typed Lorentz compatibility. The common
domains and restrictions are those of Theorem 6.3 and Appendices F--H
after the finite buffered restrictions in J.5. ∎

Theorem 9.12 gives a local analytic gravitational realization of
Penrose's active-patching proposal in the projective operator model of
Definition I.2 (Penrose 2015, secs. 5--6). Its two independent classes
lie on one \(T_0\); their coupled source generates the contact fivefold
and acts in one operator primitive, with the negative class represented
by derivative variables. The holomorphic admission retains the affine
torsor, root and calibration, and its selected optical section yields
the Einstein scale and preserved infinity form by Propositions 5.4 and
I.3. The fixed optical correction of Theorem 5.5 and Proposition H.5
constructs that admission from the source; compact point modules and
raywise local-twistor descent supply the compared nonlocal incidence. A
literal identification with every unreduced Palatial presentation
remains a separate presentation question. The assertion is the stated
local correspondence and preservation diagram, with no
historical-priority claim.

\protect\phantomsection\label{sec:gg-schwarzschild-future-horizon}

\subsection{10. Schwarzschild near the future
horizon}\label{schwarzschild-near-the-future-horizon}

We apply the correspondence to a fixed positive mass \(m>0\) and a
finite neighborhood crossing the future Schwarzschild horizon. In
ingoing Eddington--Finkelstein coordinates, with the curvature and
signature conventions of §2, put
\begin{equation}\protect\phantomsection\label{eq:gg-schwarzschild-main-metric}{
 g_m=f(r)\,dv^2-2\,dv\,dr-r^2g_{S^2},\qquad
 f(r)=1-\frac{2m}{r},\qquad g_c=r^{-2}g_m.
}\end{equation} Fix a compact core
\(K=[v_-,v_+]\times[2m-d,2m+d]\times S^2\), where \(v_\pm\) are finite
and \(0<d<2m\). All representatives below have real-analytic parents
with \(r>0\) and smaller complex parents when the holomorphic
correspondence is used. Take this finite core large enough to contain
the ray family below. A point and a full spin frame, including its lift,
are retained at each mark.

\pandocstatementheading{\protect\phantomsection\label{prop:gg-schwarzschild-main-comparison}Proposition 10.1.}

The complete curvature records of \(g_m\) determine analytic inputs in
\(\mathcal A_0\). Their generated objects have actual marked comparisons
with the canonical contact, calibrated affine and projective operator
objects of \(g_m\). The comparisons preserve the compact point modules,
optical evaluation and Wronskian, density root and selected scale,
normal local-twistor transport, infinity pair and point planes.

The core \(K\) admits a finite system of these short representatives
covering its exterior, horizon and interior portions. The system has the
common parents, positive buffers and unit margins of Definition 2.1,
with actual inverse comparisons and coherent finite compositions,
including their full-spin and coupled Lorentz data.

In the representative \(g_c\), the Einstein density and its normal
tractor are
\begin{equation}\protect\phantomsection\label{eq:gg-schwarzschild-main-scale}{
 \sigma=\frac1r,\qquad
 I=L\sigma=\left(\frac1r,-\frac{dr}{r^2},\frac{f}{2r}\right),
 \qquad DI=0,\qquad H(I,I)=0.
}\end{equation} The middle slot in this display is written as a
covector; raising it with \(g_c\) gives the vector
\(\partial_v+f\partial_r\) used by the canonical tractor convention of
§4.

For \(r_0=r(o)\), the normalization \(\sigma_o(o)=1\) preserving the
selected physical metric is
\begin{equation}\protect\phantomsection\label{eq:gg-schwarzschild-main-normalization}{
 g_{\rm aux}=r_0^2g_c,\qquad \sigma_o=\frac{r_0}{r},\qquad
 \sigma_o^{-2}g_{\rm aux}=g_m.
}\end{equation} \emph{Proof.} The \((v,r)\) metric block has determinant
\(-1\) at every \(r>0\). Direct computation gives
\(\operatorname{Ric}(g_m)=0\). Regular orthonormal coframes and local
sections of one fixed spin bundle extend across \(r=2m\). Extract every
free record by the averaged spin projection of \(\nabla^j\Psi^\pm\) in
§3. On the finite buffered point/frame parents, parallel transport along
complex null directions and Cauchy's formula give a common bound
\(Cr!R^{-r}\) for the order-\(r\) records. Theorem 3.1 supplies the
exact two analytic classes. Theorem 2.4 and the normal-metric uniqueness
of Appendix A identify the generated metric with \(g_m\) through the
initialized exponential maps. Their cotangent and calibrated ray maps,
together with Theorems 9.10--9.12, give the stated enriched comparisons.

The finite cover is constructed on compact point cores together with
their whole skies. Finite null flow boxes, nonzero transversal
derivatives, and bounded normal inverses supply the ray parents and
buffers. Changes of sections of the fixed spin bundle supply compatible
lifts; the operator and module comparisons retain the actual
associators. The overlap identities follow from the physical metric maps
and the normalized comparison uniqueness of Appendix H. These are finite
physical comparisons and do not require a single global ray quotient or
operator ring.

For the scale calculation, \(|dr|_{g_m}^2=-f\). The conformal Schouten
formula gives
\begin{equation}\protect\phantomsection\label{eq:gg-schwarzschild-main-scale-check}{
 \nabla_{g_c}^2\sigma+P^{g_c}\sigma=-\frac{f}{2r}g_c,
 \qquad
 d\left(\frac f{2r}\right)=P^{g_c}\bigl((d\sigma)^{\sharp_{g_c}},\cdot\bigr).
}\end{equation} These are all components of \(DI=0\), and
\(2\sigma\rho+|d\sigma|_{g_c}^2=0\). In the normalized auxiliary
splitting the same tractor is \((r_0/r,-r_0dr/r^2,f/(2r_0r))\); its
selected Einstein splitting is \((1,0,0)\). The conformal frame change
has dilation one at \(o\), so the specified full spin mark is preserved
there. The marked positive root fixes the sign of \(\sigma_o\).
Propositions K.1--K.4 give the scale and frame calculations; Theorems
K.7--K.8 give the actual enriched maps and finite cover construction. ∎

The normalization in eq.~\ref{eq:gg-schwarzschild-main-normalization}
changes the auxiliary representative together with the density
component. Changing only \(\sigma\) to \(r_0\sigma\) would select
\(r_0^{-2}g_m\). The scalar curvature of the defined representative
\(g_c\) is \(-12m/r\) in our convention; hence
\(\sigma=-\operatorname{Scal}(g_c)/(12m)\).

\pandocstatementheading{\protect\phantomsection\label{prop:gg-schwarzschild-main-optics}Proposition 10.2.}

There is a compact nonradial family of affinely parametrized null
segments crossing the future horizon from \(r=5m/2\) to \(r=3m/2\), with
\begin{equation}\protect\phantomsection\label{eq:gg-schwarzschild-main-ray-family}{
 r(0)=\frac52m,\quad v(0)=\phi(0)=0,\quad \theta=\frac\pi2,
 \qquad E\in[1,2],\quad L/m\in[1/4,1/2].
}\end{equation} Choose the ingoing branch and physical affine parameter
\(u\). A parallel orthonormal screen basis for the positive screen
metric \(-g_m\) has tidal matrix
\begin{equation}\protect\phantomsection\label{eq:gg-schwarzschild-main-tide}{
 \mathcal R_u=\frac{3mL^2}{r^5}
 \begin{pmatrix}1&0\\0&-1\end{pmatrix}.
}\end{equation} Thus the transverse Jacobi phase space has dimension
four and nonzero Weyl focusing. In the dimensionless time \(t=u/m\),
with phase \((x,p)=(x,m\,dx/du)\), its transfer matrix \(\Phi(t,s)\) is
symplectic and composes exactly on every finite subdivision.

Writing
\(\Phi(\tau,0)=\left(\begin{smallmatrix}A&B\\C&D\end{smallmatrix}\right)\)
at the final time, one has
\begin{equation}\protect\phantomsection\label{eq:gg-schwarzschild-main-transfer-margin}{
 \frac{10}{21}<\tau<\frac{100}{99},\qquad
 \|\Phi(\tau,0)\|<3,\qquad
 \sigma_{\min}(B)>\frac37,\qquad \|B^{-1}\|<\frac73.
}\end{equation} Finite stationary quadratic generating families retain
the complete phase relation when an endpoint-position projection is
singular. In particular, on these same segments, put \(t_c=\tau/2\) and
let \(A_c,B_c\) be the position blocks of \(\Phi(t_c,0)\). The initial
Lagrangian \(p=P_0x\), where \(P_0=-B_c^{-1}A_c=P_0^T\), propagates to
the vertical Lagrangian at \(t_c\).

\emph{Proof.} Put \(q=(E^2-fL^2/r^2)^{1/2}\). The regular first
integrals are
\begin{equation}\protect\phantomsection\label{eq:gg-schwarzschild-main-first-integrals}{
 \frac{dr}{du}=-q,\qquad
 \frac{dv}{du}=\frac{L^2}{r^2(E+q)},\qquad
 \frac{d\phi}{du}=\frac L{r^2}.
}\end{equation} On the stated radial interval,
\(124/125\le q^2\le109/27\), so \(99/100<q<21/10\). This gives the time
bounds, unique transverse crossing and a common regular finite
coordinate tube. The direct parallel-screen calculation in Appendix K
gives eq.~\ref{eq:gg-schwarzschild-main-tide}; the rotational Jacobi
fields \(r\sin\phi\) and \(r\cos\phi\) independently check its focusing
sign.

With \(\rho=r/m\), \(\ell=L/m\) and \(\omega=3\ell^2/\rho^5\), the phase
generator is
\(\left(\begin{smallmatrix}0&I\\\operatorname{diag}(-\omega,\omega)&0\end{smallmatrix}\right)\),
where \(6/3125\le\omega\le8/81\). Its norm is one and it is Hamiltonian.
The iterated-integral estimate gives \(\|\Phi(t,s)\|\le e^{t-s}\). The
Volterra equation gives \(\|B(h)-hI\|<h/10\) on every subsegment of
length \(0<h\le100/99\), proving the endpoint margin. ODE uniqueness
proves finite composition. Short segments have invertible position
blocks and quadratic endpoint actions. Keeping intermediate positions
makes stationarity exactly momentum matching, even when the total
projection is singular. Finally \(A_c+B_cP_0=0\) and
\(C_c+D_cP_0=-B_c^{-T}\) prove the asserted vertical image. Appendix K
supplies the generating functions and their finite composition proof. ∎

The scalar optical bundle of §5 has rank two and affine equation
\(y''=0\) here. It differs from the four-dimensional transverse phase
space in Proposition 10.2. The chosen point-source block \(B\) is
nonsingular on this short family. The displayed caustic concerns the
specified initial linear Lagrangian condition; it asserts neither a
further point-source conjugate point nor a nonlinear congruence. The
full phase relation survives that projection failure, whereas an inverse
endpoint block does not.

\pandocstatementheading{\protect\phantomsection\label{prop:gg-schwarzschild-main-errors}Proposition 10.3.}

Fix the finite coordinate and frame parents of Propositions 10.1--10.2.
Suppose a single prescribed representative of Theorem 6.3 admits the
finite metric readout of Proposition 6.2 on one fixed complex coordinate
parent \(\Omega^+\), jointly holomorphically in amplitude for
\(|a|<T_b\), where \(T_b>1\), with the retained inverse-unit margins.
Choose an inner polydisc \(\Omega\) with coordinatewise Cauchy gap
\(\rho>0\) from \(\Omega^+\). Assume the closure of an inflated position
tube around the chosen ray family lies inside \(\Omega\) with a further
positive position buffer. Separately retain a bounded state tube for
\(Z=(x,k,n,e_1,e_2)\), with compact closure and positive first-exit
buffer around the exact geodesic and parallel-frame states; its position
projection lies in that position tube. Let \(g_N\) be the degree-\(N\)
amplitude Taylor polynomial of the physical metric in these fixed
coordinates, evaluated at \(a=1\), using the coupled Lorentz-real
amplitude family. For \(1<T<T_b\), there is a finite \(H_g\) such that
\begin{equation}\protect\phantomsection\label{eq:gg-schwarzschild-main-taylor-error}{
 \|D^j(g_N-g_m)\|\le e_{j,N}:=
 4^j j!\rho^{-j}\frac{H_gT^{-N-1}}{1-T^{-1}},
 \qquad 0\le j\le3.
}\end{equation} All norms are fixed Euclidean coordinate or multilinear
operator norms. If \(G\ge\|g_m^{-1}\|\) and \(e_{0,N}\le(2G)^{-1}\),
then \(g_N\) is Lorentzian with \(\|g_N^{-1}\|\le2G\). For sufficiently
large \(N\), the retained initial null-frame and domain margins also
persist. The normalized approximate rays and parallel screens then exist
on the same individual finite affine intervals. There is a finite
constant \(C_{\rm opt}\), depending only on the fixed mass, parents,
metric derivative and inverse bounds, frame bounds and buffers, such
that their dimensionless phase coefficients differ by at most
\begin{equation}\protect\phantomsection\label{eq:gg-schwarzschild-main-coefficient-error}{
 \epsilon_N=C_{\rm opt}\sum_{j=0}^3e_{j,N}.
}\end{equation} For \(\epsilon_N\le1/10\), the exact transfers for those
metric systems obey
\begin{equation}\protect\phantomsection\label{eq:gg-schwarzschild-main-optical-error}{
 \|\Phi_N-\Phi\|\le\delta_N:=4\epsilon_N,\qquad
 |m^2\det B_N-m^2\det B|\le9m^2\delta_N.
}\end{equation} If also \(\delta_N\le3/14\), then
\begin{equation}\protect\phantomsection\label{eq:gg-schwarzschild-main-inverse-error}{
 \|B_N^{-1}\|\le\frac{14}{3},\qquad
 \|B_N^{-1}-B^{-1}\|\le\frac{98}{9}\delta_N.
}\end{equation} \emph{Proof.} Apply the finite readout and initialized
normal development of Lemma K.9 on the prescribed coordinate parent
\(\Omega^+\) with its retained inverse-unit margins. The parent and
inner domain \(\Omega\) in the hypothesis are retained when applying the
amplitude and spatial Cauchy estimates, which give
eq.~\ref{eq:gg-schwarzschild-main-taylor-error}. Lorentz reality makes
\(g_N\) real symmetric on the real tube. The inverse series and constant
inertia along \(g_m+s(g_N-g_m)\) give the signature and inverse bound.

The Christoffel formula and one further derivative bound the connection
and curvature errors by finite linear combinations of
\(e_{0,N},e_{1,N},e_{2,N}\). Analytic normalization supplies nearby
initial null frames. Gronwall's inequality for the joint
geodesic/parallel-frame equation keeps the solutions inside the retained
state buffer. Curvature evaluation in those frames then gives
eq.~\ref{eq:gg-schwarzschild-main-coefficient-error}, including the
physical-to-dimensionless factor \(m^2\). Lemmas K.9--K.11 and Theorem
K.12 specify these constants and the precise smallness inequalities
excluding a first exit.

Duhamel's formula on \(\tau<100/99\) gives the transfer bound
\(\epsilon_N\tau e^{(1+\epsilon_N)\tau}<4\epsilon_N\). The
two-dimensional determinant estimate
\(|\det(B+E)-\det B|\le2\|B\|\|E\|+\|E\|^2\) gives the forward area
bound. The singular-value margin in
eq.~\ref{eq:gg-schwarzschild-main-transfer-margin} and the inverse
difference identity give
eq.~\ref{eq:gg-schwarzschild-main-inverse-error}. ∎

Here \(H_g\) and \(C_{\rm opt}\) are finite constants specified by
analytic parents and compact-frame margins; the rational constants in
eq.~\ref{eq:gg-schwarzschild-main-transfer-margin}--eq.~\ref{eq:gg-schwarzschild-main-inverse-error}
are explicitly evaluated bounds. The single-readout position-tube
inclusion assumed in Proposition 10.3 has not been established for the
whole \(r:5m/2\to3m/2\) family of Proposition 10.2; this outstanding
spatial-domain condition is separate from the numerical evaluation of
\(H_g\), \(C_{\rm opt}\) and a sufficient \(N\). Proposition 10.5
supplies a finite-readout alternative on that entire family. Sections
K.6--K.7 instead choose the common parents, inner \(C^3\) core and
passive-displacement margin first and then obtain a positive short
interval inside them. Comparisons in Proposition 10.3 use fixed affine
endpoints. Reading data instead on \(r=3m/2\) requires the transverse
hitting-time margin and finite extension stated in Appendix K. The
metric Taylor polynomials are taken after the exact chart comparisons
are fixed. They supply local approximations and do not replace the exact
gluing data. The forward phase and determinant estimates require no
inverse projection; inverse estimates require their stated positive
margin.

\pandocstatementheading{\protect\phantomsection\label{prop:gg-schwarzschild-primitive-comparison}Proposition 10.4.}

Fix the primitive restricted Dolbeault projector and positively oriented
contour in §K.5, with its printed prefactor \(1/(2\pi)\), the source
parameter \(\mu_{\rm src}=m/2\), and the continued signed radial sheet
\(\varrho=-r\). On a sufficiently small simply connected neighborhood of
\((v,r,\theta,\phi)=(0,2m,\pi/2,0)\), retain the local angular inverse
branches and \(x+\mathrm iy\ne0\). The actual map
\begin{equation}\protect\phantomsection\label{eq:gg-schwarzschild-primitive-map}{
 \Phi_{\rm pr}(v,r,\theta,\phi)
 =\bigl(\mathrm i(r-v),r\sin\theta\cos\phi,
                   r\sin\theta\sin\phi,r\cos\theta\bigr)
}\end{equation} into the complexification of the primitive metric domain
has an actual local inverse and satisfies
\begin{equation}\protect\phantomsection\label{eq:gg-schwarzschild-primitive-identity}{
 \Phi_{\rm pr}^*G_{\rm int}=h_*:=2r^{-2}g_m,\qquad
 \sigma_*=\sqrt2/r,\qquad \sigma_*^{-2}h_*=g_m.
}\end{equation} The selected density, its square identification and
central root value, and the normal tractor agree with those of
Proposition 10.1 after the stated constant auxiliary-frame change.
Transporting the specified EF full-spin frame gives actual marked
comparisons. The selected primitive metric and the record-generated
objects therefore have the actual enriched comparisons of Theorem K.7,
with two-sided inverses, on their common local representatives.

\emph{Proof.} Proposition K.14 evaluates the projector and contour
throughout the stated chart, including the factor
\(\omega_{\rm int}=2\omega_0\). Theorem K.15 gives the metric pullback,
density, tractor, root and full-spin comparison, using the source
parameter and signed sheet exactly as stated. The selected physical
source tensor is \(g_{\rm src}=(-\sqrt2/\varrho)^{-2}G_{\rm int}\). If
\(F:g[A_g]\to g_m\) is the initialized isometry in Theorem K.7, then
\(F_{\rm src}=\Phi_{\rm pr}\circ F\) and
\(F_{\rm src}^{-1}=F^{-1}\circ\Phi_{\rm pr}^{-1}\). The isometry
identifies Levi--Civita connections, their iterated curvature
derivatives and hence every averaged free spin record, in the
transported full-spin frame. The canonical ray, affine, root,
local-twistor, infinity and structured operator/module comparisons of
that theorem apply to this same geometry, with the stipulated
associators. Section K.6 gives compatible physical amplitude families
for finite comparisons. ∎

This primitive comparison is local in the spacetime angular and root
variables. The full complex sky at each such base point is retained, but
it is distinct from the spacetime angular \(S^2\) in \(K\). Proposition
10.1 continues to cover that standard-metric core with finitely many
charts; Proposition 10.4 asserts no native gluing over all of it. With
the conventions fixed in §K.5, the contour integral in (Adamo et al.
2026a, equation (2.32)) equals twice its displayed evaluated two-form;
Proposition K.14 gives the residue calculation and the resulting metric
normalization. The result uses neither an identity with its fixed SDTN
ordinary background nor a direct class construction from its quadric.

The nonlinear compatible family is
\(V_i(\varepsilon)=R_{a_i,A_i}(\varepsilon U;0)\), with a separate
internal construction amplitude at each vertex, as in §K.6. A single
such family admits the additionally initialized short optical example of
§K.7 after a positive source-domain restriction. Its rational
full-interval enclosure is a separate calculation from the earlier
family of Proposition 10.2. The finite-readout proof in §K.4 performs
inverses and initialized ODEs in the truncated amplitude algebra before
evaluating the retained coefficient polynomial at one. The metric and
optical error constants remain theoretical finite constants; no
numerical sufficient amplitude order is asserted.

\pandocstatementheading{\protect\phantomsection\label{prop:gg-schwarzschild-finite-readout}Proposition 10.5 (Finite-readout optical approximation).}

Fix \(m>0\) and the full nonradial family of Proposition 10.2, with its
specified entrance, \(E\in[1,2]\), \(L/m\in[1/4,1/2]\), physical affine
clock \(u\). The construction uses the exact benchmark cuts, initialized
coordinate comparisons and polar-first parallel Sachs reference frame
specified in §K.8. There are a finite parameter cover, a finite
local-readout schedule and positive spatial, inverse-unit and frame
margins, all chosen independently of \(N\), with the following property.
There are \(T_*>1\), finite constants \(C_Z,K_{\rm tid},H_{\rm rot}\)
and an integer \(N_0\) such that for every \(N\ge N_0\) the scheduled
degree-\(N\) metric readouts and their prescribed state and covariant
Cauchy resets exist along every original individual interval
\(0\le t=u/m\le\tau_\lambda\). Put
\begin{equation}\protect\phantomsection\label{eq:fro-main-effective-error}{
 \beta_N=\frac{T_*^{-N-1}}{1-T_*^{-1}},\qquad
 \varepsilon_N=\left(K_{\rm tid}+\frac{H_{\rm rot}}{\tau_-}\right)\beta_N,
 \qquad \tau_-=10/21.
}\end{equation} Use the maximum block norm of
\(Z=(X,K,\mathcal N,E_1,E_2)\) in dimensionless EF coordinates, with
\(g_0=m^{-2}g_m\), \(t=u/m\) and physical frame
\((K,\mathcal N,E_A)/m\). The complete piecewise state has error at most
\(C_Z\beta_N\), on both sides of every switch and inside every slab. Its
full Cauchy transport \(U_N\) is symplectic. For \(N\) sufficiently
large that \(\varepsilon_N\le1/10\), at the original exact terminal
affine time it obeys
\begin{equation}\protect\phantomsection\label{eq:fro-main-full-family-estimates}{\begin{gathered}
\|U_N-\Phi_\lambda(\tau_\lambda,0)\|
 \le4\varepsilon_N=:\delta_N, \\
\|U_N\|<4, \\
|m^2\det B_N-m^2\det B|\le9m^2\delta_N.
\end{gathered}}\end{equation} Here \(B_N\) and \(B\) are the endpoint
position-response blocks in the same dimensionless-clock phases. The
physical-slope response blocks are \(mB_N\) and \(mB\). If also
\(\delta_N\le3/14\), then
\begin{equation}\protect\phantomsection\label{eq:fro-main-full-family-inverses}{
 \sigma_{\min}(B_N)\ge3/14,\qquad
 \|B_N^{-1}\|\le14/3,\qquad
 \|B_N^{-1}-B^{-1}\|\le(98/9)\delta_N.
}\end{equation} All conclusions are uniform over the original entire
parameter rectangle. The forward phase and conditional expansion-trace
estimates of Proposition K.25, and the Lagrangian-projection estimates
eq.~\ref{eq:fr-opt-c-lagrangian-projection-error}, hold as well. The
retained exact twistor, full-spin, operator and compact point-module
comparisons are those of the original physical atlas.

\emph{Proof.} Theorem K.18 constructs the finite cover, enlarged slabs
and buffered node overlaps for the whole parameter rectangle, using the
original exact family. Proposition K.19 produces the actual
finite-algebra metric readouts and a common tail rate \(T_*>1\); take
\(T_{\min}=T_*\) in §K.8 and enlarge the finite constants \(D_{i,s}\) if
needed. Proposition K.20 gives their common-coordinate overlap defects
and retains the original exact enriched atlas. Lemma K.21 supplies the
actual null/frame and covariant Cauchy resets in the same affine clock,
including their computed screen rotations. Proposition K.22 proves
all-slab and all-switch existence by first-exit induction and gives the
state bound. The tidal estimate in §K.8 and Proposition K.23 sum every
slab and jump error. Corollary K.24 then gives the transfer estimate
with the displayed \(\varepsilon_N\). Proposition K.25 gives the
physical area and conditional inverse bounds from the original exact
endpoint margin. The finite minima and maxima include all parameter
cells. Since \(\beta_N\to0\), one sufficiently large \(N\) meets all
these finitely many strict domain and non-strict observation margins. ∎

This proposition provides a finite-readout alternative on the whole
original crossing family. Proposition 10.3 remains a conditional
single-readout result; the present construction does not establish its
single-parent spatial inclusion. The short common-source construction of
§§K.6--K.7 remains a separate one-source result. The constants, radii
and a sufficient truncation degree of this finite-readout construction
are theoretical finite data that have not been numerically evaluated.
The constructed object is the piecewise metric and Cauchy model of Lemma
K.21, rather than a single approximate metric or the differential of the
nonlinear state-reset algorithm.

\protect\phantomsection\label{app:gg-ordered-jets}

\subsection{Appendix A. Finite witnesses and polynomial curvature
completion}\label{appendix-a.-finite-witnesses-and-polynomial-curvature-completion}

\protect\phantomsection\label{sec:gg-finite-witnesses}

\subsubsection{A.1. Every finite free list has an analytic
witness}\label{a.1.-every-finite-free-list-has-an-analytic-witness}

Only this appendix uses the analytic constraint and evolution equations
to establish finite algebraic consistency. The generator of §5 does not
take their solutions as input. No common domain as the prefix length
increases is asserted here.

We use the analytic Cauchy theorem in normal form: a finite system
solved for its highest derivatives transverse to an analytic
hypersurface, with analytic right-hand side and analytic
noncharacteristic Cauchy data, has a unique local analytic solution
(Kowalevsky 1875). The complex version with additional undifferentiated
parameters gives joint holomorphic dependence after a common
restriction. The normal principal matrices are verified below.

\pandocstatementheading{\protect\phantomsection\label{lem:gg-finite-witness}Lemma A.1.}

Every finite Lorentz-real list \(U_r^\pm\), \(0\le r\le N\), and real
numerical \(\Lambda\) has a real analytic Einstein witness with its
prescribed tangent metric and spin marking. Its neighborhood may depend
on \(N\) and the list.

\emph{Proof.} First construct a finite family near the flat metric in
signature \((-,+,+,+)\). A temporal transverse trace-free linear wave
has spatial initial data \(A,B\) with
\(\operatorname{tr}A=\operatorname{tr}B=0\) and
\(\partial^jA_{ij}=\partial^jB_{ij}=0\). Here
\(B=\partial_0h^{(1)}|_0/2\). For a finite set of real amplitude
parameters \(c\), put \(\bar\gamma=\delta+\sum c_\nu A_\nu\) and let
\(\Sigma\) be the \(\bar\gamma\)-trace-free part of
\(\sum c_\nu B_\nu\). Use bilinear tensor contractions in the
complexified equations. With
\((\bar LZ)_{ij}=\bar\nabla_iZ_j+\bar\nabla_jZ_i
-\frac23\bar\gamma_{ij}\bar\nabla_kZ^k\), solve
\begin{equation}\protect\phantomsection\label{eq:gg-finite-constraint-solve}{
 \overline{\operatorname{div}}\,\bar LZ
       =-\overline{\operatorname{div}}\Sigma,\qquad
 \bar\Delta\phi-\frac18R(\bar\gamma)\phi
       +\frac18|\Sigma+\bar LZ|_{\bar\gamma}^2\phi^{-7}
       +\frac{\Lambda}{4}\phi^5=0
}\end{equation} with \(Z=\partial_1Z=0\), \(\phi=1\),
\(\partial_1\phi=0\) on \(y^1=0\). The first equation is independent of
\(\phi\). Its normal principal matrix at the flat seed is
\(I+\frac13(\mathrm dy^1)^\sharp\otimes\mathrm dy^1\), with eigenvalues
\(1,1,4/3\). The second has normal coefficient \(\bar\gamma^{11}\) and
is analytic near the unit \(\phi=1\). Both systems therefore have the
stated local holomorphic solutions.

Set \(\gamma=\phi^4\bar\gamma\), \(K=\phi^{-2}(\Sigma+\bar LZ)\). The
conformal connection difference \(2\phi^{-1}(\delta_i^k\partial_j\phi+
\delta_j^k\partial_i\phi-\bar\gamma_{ij}\bar\gamma^{k\ell}\partial_\ell\phi)\)
gives
\begin{equation}\protect\phantomsection\label{eq:gg-finite-constraint-transform}{
 \begin{aligned}
 R(\gamma)&=\phi^{-5}(-8\bar\Delta\phi+R(\bar\gamma)\phi),\\
 \operatorname{div}_\gamma K
       &=\phi^{-6}\overline{\operatorname{div}}(\Sigma+\bar LZ),\\
 |K|_\gamma^2&=\phi^{-12}|\Sigma+\bar LZ|_{\bar\gamma}^2,\qquad
 \operatorname{tr}_\gamma K=0 .
 \end{aligned}
}\end{equation} Consequently
\(R(\gamma)+(\operatorname{tr}K)^2-|K|^2=2\Lambda\) and
\(\operatorname{div}K-\mathrm d\operatorname{tr}K=0\). At \(\Lambda=0\),
the first amplitude variations of \(Z\) and \(\phi-1\) vanish: their
linear sources are respectively \(\partial^jB_{ij}=0\) and
\(\partial_i\partial_jA_{ij}-\Delta\operatorname{tr}A=0\), with zero
Cauchy data. Thus the initial metric and second fundamental form have
the prescribed first variations.

Evolve, with \(\tau=\operatorname{tr}_\gamma K\), by
\begin{equation}\protect\phantomsection\label{eq:gg-finite-gaussian-evolution}{
 \partial_0\gamma_{ij}=2K_{ij},\qquad
 \partial_0K_{ij}=-\operatorname{Ric}(\gamma)_{ij}
                  -\tau K_{ij}+2K_i{}^kK_{kj}
                  +\Lambda\gamma_{ij}.
}\end{equation} Eliminating \(K\) solves for \(\partial_0^2\gamma\) with
analytic right-hand side in the invertible metric and tangential
derivatives of order at most two. The analytic Cauchy theorem applies
again. For \(g=-(\mathrm dx^0)^2+\gamma_{ij}\mathrm dy^i\mathrm dy^j\),
the time-containing Christoffel symbols are \(\Gamma^0_{ij}=K_{ij}\) and
\(\Gamma^i_{0j}=K^i{}_j\). They give
\begin{equation}\protect\phantomsection\label{eq:gg-finite-gaussian-ricci}{
 \begin{aligned}
 \operatorname{Ric}(g)_{ij}
    &=\operatorname{Ric}(\gamma)_{ij}+\partial_0K_{ij}
                             +\tau K_{ij}-2K_i{}^kK_{kj},\\
 \operatorname{Ric}(g)_{00}&=-\partial_0\tau-|K|^2,\qquad
 \operatorname{Ric}(g)_{0i}=\nabla^jK_{ij}-\partial_i\tau .
 \end{aligned}
}\end{equation} The spatial equations hold. For their remaining
residuals \(\mathcal H=\operatorname{Ric}_{00}+\Lambda\) and
\(\mathcal M_i=\operatorname{Ric}_{0i}\), contracted Bianchi gives
\begin{equation}\protect\phantomsection\label{eq:gg-finite-constraint-propagation}{
 \partial_0\mathcal H
       =2\nabla^i\mathcal M_i-2\tau\mathcal H,\qquad
 \partial_0\mathcal M_i=\frac12\partial_i\mathcal H-\tau\mathcal M_i .
}\end{equation} Their initial data vanish by the constraints, so
analytic uniqueness gives the full Einstein equation. Linearization at
the flat seed is the specified wave, by uniqueness of its linear
analytic Cauchy problem. Reality follows from conjugation and
uniqueness. The spatial metric is positive on a smaller real
neighborhood.

We next choose a finite wave family whose curvature-jet map is
invertible. Put \(\kappa(z)=(1,z)\) and take the real null covector
\(k(z)_{AA'}=\kappa_A(z)\bar\kappa_{A'}(z)\). For each \(r\) the tensors
\(\kappa(z)^{r+4}\bar\kappa(z)^r\) span \(W_r^+\) over \(\mathbb C\).
Indeed an annihilating functional is a polynomial in \(z,\bar z\);
substitution \(z=x+iy\) makes it a polynomial vanishing for all real
\(x,y\), so it is zero. A finite spanning set valid for all \(r\le N\)
is the square grid \(x,y\in\{0,\ldots,2N+4\}\), by interpolation in each
real variable. Choose a basis separately for each \(r\).

The moment matrix \(((i\tau_\ell)^j)_{0\le j,\ell\le N}\) for distinct
nonzero real \(\tau_\ell\) is invertible. Its inverse columns give
finite profiles \(f_r(t)=\sum_\ell b_\ell^{(r)}e^{i\tau_\ell t}\) with
\(f_r^{(j)}(0)=\delta_{jr}\) for \(j\le N\). For each chosen direction
put \(k_{AA'}=\kappa_A\lambda_{A'}\), \(\lambda=\bar\kappa\), and choose
a complex null transverse polarization \(p_{AA'}=\kappa_A\mu_{A'}\) with
\(d=\lambda_{0'}\mu_{1'}-\lambda_{1'}\mu_{0'}\ne0\). In the manuscript
tangent convention, the linear wave
\(h_{ab}=a p_ap_be^{i\tau_\ell k\cdot x}\) has, by
eq.~\ref{eq:gg-last-slot-linear-curvature},
\begin{equation}\protect\phantomsection\label{eq:gg-finite-wave-component-readout}{
 R^{(1)}_{abcd}=-\tfrac12a\tau_\ell^2
 (k_ap_b-k_bp_a)(k_cp_d-k_dp_c)e^{i\tau_\ell k\cdot x},\qquad
 \Psi^+=\tfrac12a\tau_\ell^2d^2\kappa^4
                                      e^{i\tau_\ell k\cdot x} .
}\end{equation} The second equality uses
eq.~\ref{eq:gg-weyl-spinor-readout}. For the auxiliary opposite
signature, the metric negation below sends \(h\) to \(-h\) in this
component formula and changes the signed factor by \(-1\). In either
signature divide the amplitude by its actual nonzero factor. Replacing
\(p\) by \(p-(p_0/k_0)k\) removes its time component without changing
curvature, nullness or transversality; \(k_0\ne0\) for a nonzero real
null covector. Pair the wave with its Lorentz conjugate. Thus its
positive curvature spinor can be normalized to
\(c\kappa^4f_r(k\cdot x)\), with the conjugate negative spinor, and its
order-\(j\) free list is \(\delta_{jr}c\kappa^{r+4}\bar\kappa^r\).
Complex \(c\) supplies two real parameters.

The resulting finite family has \(2\sum_{r=0}^N(r+5)(r+1)\) real wave
parameters. Read its free lists in an analytic orthonormal spin frame at
the origin. At flat space all background Weyl derivatives vanish, so
variations of that frame or connection do not alter the displayed linear
curvature map. The analytic map from wave parameters and \(\Lambda\) to
their free lists and the unchanged \(\Lambda\) therefore has invertible
real derivative. The finite analytic inverse theorem realizes all
sufficiently small lists. The construction also gives their complexified
families. For an opposite Lorentz signature, run the construction with
constant \(-\Lambda\), negate the metric and use that signature's fixed
spin readout before taking this inverse. The linear wave-jet map is
still invertible and the resulting Einstein constant is exactly the
requested \(\Lambda\).

For an arbitrary finite target choose \(s>0\) large enough that
\((s^{-r-2}U_r,\allowbreak s^{-2}\Lambda)\) lies in the small neighborhood. If \(h\)
is that witness, set \(g=s^{-2}\delta_s^*h\), \(\delta_s(x)=sx\). Then
\(g(x)=h(sx)\), its tangent metric is unchanged, its orthonormal
curvature derivatives scale by \(s^{r+2}\), and its Einstein constant
scales by \(s^2\). Hence it has the exact target list and constant. This
final scaling only proves finite existence; it supplies no
infinite-prefix bound. ∎

\protect\phantomsection\label{sec:gg-polynomial-completion}

\subsubsection{A.2. The highest linear compatibility
system}\label{a.2.-the-highest-linear-compatibility-system}

\pandocstatementheading{\protect\phantomsection\label{lem:gg-highest-kernel}Lemma A.2.}

For fixed lower curvature derivatives, the homogeneous difference of two
order-\(r\) Einstein curvature completions, subject to their
differentiated algebraic and differential Bianchi identities and
commutator identities, is exactly \(W_r^+\oplus W_r^-\). Its two
separately symmetric projections are the identity on that space.

\emph{Proof.} For \(r=0\), use eq.~\ref{eq:gg-weyl-spinor-decomposition}
and its inverse eq.~\ref{eq:gg-weyl-spinor-readout}, with the Ricci part
fixed. At every higher order the free projection is
eq.~\ref{eq:gg-free-record-projection}. For \(r>0\), the highest
difference has commuting derivative slots, since their commutator is
curvature acting on a lower derivative. Its Ricci derivatives vanish.
For the positive Weyl spinor the differentiated Einstein Bianchi
equation contracts any derivative unprimed index against a Weyl unprimed
index to zero. In a two-space, a two-index antisymmetrization is its
alternating trace. These zero traces therefore imply invariance under
every interchange of a derivative unprimed index with a Weyl index, with
the remaining slots fixed. Together with Weyl symmetry, those
interchanges generate all unprimed permutations. Symmetry under exchange
of derivative pairs consequently makes their primed indices symmetric
too. This is precisely
\(\operatorname{Sym}^{r+4}S^*\otimes\operatorname{Sym}^rS'^*\). The
negative spinor exchanges the factors. Conversely these completely
symmetric spinors obey all homogeneous contractions and commuting-slot
identities. Their projection is unchanged. ∎

\pandocstatementheading{\protect\phantomsection\label{prop:gg-polynomial-completion}Proposition A.3.}

The polynomial completion and the original Cartan derivations of Lemma
3.2 exist uniquely and obey all the identities stated there, on
independent complex lists.

\emph{Proof.} At order \(r\), collect the tensor symmetries,
differentiated Bianchi equations, Ricci equations, commutator identities
and prescribed free projection into a constant finite linear system for
the ordered curvature derivative. Its right-hand side is polynomial in
lower completed derivatives, \(U_r^\pm\) and \(\Lambda\). Lemma A.2
makes its coefficient matrix injective. Choose a constant left inverse.
Induction therefore produces a unique candidate polynomial
\(\mathcal C_r\).

For every finite Lorentz-real list, Lemma A.1 supplies an actual
witness, so this candidate satisfies the whole compatibility system, not
just its left-inverse projection. The residual is polynomial and
vanishes on the maximal real form \(U^-=\overline{U^+}\),
\(\Lambda\in\mathbb R\). Writing \(U^+=x+iy\), \(U^-=x-iy\) proves that
it vanishes identically on the independent complex lists. This also
proves independence of the left inverse and frame equivariance.

Define \(\mathsf H_aU_r\) by projecting the next completed covariant
derivative, and define \(\mathsf V_B\) by inserting the specified
infinitesimal spin-frame action in every slot. On finite witnesses these
are exactly horizontal displacement readout and vertical spin readout.
Their commutators and their action on each \(\mathcal C_r\) consequently
hold as polynomial identities: each uses only a finite prefix, and the
same maximal-real-form argument applies. This proves all Cartan
identities, including the action on derivative slots.

Dilation gives \(\operatorname{wt}U_j=j+2\),
\(\operatorname{wt}\Lambda=2\), and
\(\operatorname{wt}(\mathsf H_aU_r)=r+3\). Its highest linear term is
\(i_aU_{r+1}\) with coefficient one, by averaged separate symmetry.
Every nonlinear term therefore uses only \(U_0,\ldots,U_{r-1},\Lambda\)
and has index sum \(r+3-2d\) at curvature degree \(d\). In particular
\(Q_0=0\). The original Lorentz involution and all fixed spin actions
commute with the construction by the same uniqueness. No quotient of the
specified spin lift by its orthogonal image is made. ∎

\protect\phantomsection\label{sec:gg-formal-normal-uniqueness}

\subsubsection{A.3. Formal marked normal
uniqueness}\label{a.3.-formal-marked-normal-uniqueness}

\pandocstatementheading{\protect\phantomsection\label{prop:gg-formal-normal-uniqueness}Proposition A.4.}

The completed curvature list determines a unique formal metric in marked
geodesic normal coordinates. This formal metric has its prescribed
Einstein constant and ordered curvature derivatives. Two analytic marked
normal metrics with the same complete free lists and numerical constant
are equal as germs.

\emph{Proof.} Along the formal radial geodesic of initial vector \(x\),
express curvature in the parallel initial frame as
\(R_x(t)=\sum_{r\ge0}t^r\mathcal C_r(x,\ldots,x;\,\cdot)/r!\). For each
\(v\) solve the formal finite Jacobi equation
\begin{equation}\protect\phantomsection\label{eq:gg-formal-radial-jacobi}{
 J_x''(t;v)+R_x(t)(J_x(t;v),x)x=0,\qquad
 J_x(0;v)=0,\quad J_x'(0;v)=v .
}\end{equation} Its coefficient recursion is unique. Define
\(g_x(v,w)=\eta(J_x(1;v),J_x(1;w))\). Every fixed coefficient uses only
finitely many completed derivatives. At a fixed spatial degree the
Jacobi recursion has only finitely many terms and divides only by fixed
positive integers, so its metric coefficients are polynomials in a
finite free prefix and \(\Lambda\). Formal inversion at the fixed
tangent metric has the same property. Hence every coefficient of its
curvature, Einstein and normal-coordinate residuals is also such a
polynomial. On a finite Lorentz-real witness with the required longer
prefix, the normal-coordinate Jacobi formula gives this coefficient and
every residual vanishes. The maximal-real-form argument of Proposition
A.3 makes these polynomial identities valid on independent complex lists
and therefore over every commutative complex coefficient algebra. This
explicitly supplies the formal parameter-ring use made in §6.

It follows, coefficient by coefficient, that the constructed formal
metric has the prescribed curvature, Einstein equation,
normal-coordinate conditions, tangent metric and mark. Conversely every
such metric satisfies the same Jacobi recursion. Analytic metrics with
these jets have the same convergent Taylor series and hence the same
germ. This argument neither takes a limit of the finite witnesses nor
assumes their neighborhoods are uniform. ∎

\protect\phantomsection\label{app:gg-record-budget}

\subsection{Appendix B. The coefficient estimate for the original record
derivations}\label{appendix-b.-the-coefficient-estimate-for-the-original-record-derivations}

\protect\phantomsection\label{sec:gg-budget-recursion}

\subsubsection{B.1. The normalized one-derivative
equation}\label{b.1.-the-normalized-one-derivative-equation}

Use the full-tensor Hermitian norms of §3, with the fixed alternating
spin bases orthonormal. Symmetrizations are averaging projections. All
adjoints below are adjoints of constant finite matrices; they do not
introduce conjugated record variables.

For either chirality and \(r\ge0\) put \(\mathcal D_r=V^*\otimes W_r\).
For \(r\ge1\), regrouping one spinor of each kind gives an isometric
inclusion \(j_r:W_r\to V^*\otimes W_{r-1}\). Let \(J_r=1\otimes j_r\)
and let \(\tau\) exchange the two outer derivative pairs after that
inclusion. Set \(\delta_r=(1-\tau)J_r\), without dividing by two. For
the positive chirality write \(n=r+4,m=r\); exchange them for the
negative one. The upper and lower projections in the two spin factors
are denoted by \(P_{++},P_{-+},P_{+-},P_{--}\).

\pandocstatementheading{\protect\phantomsection\label{lem:gg-budget-spencer}Lemma B.1.}

For the original derivation \(\mathsf H_aU_r=i_aU_{r+1}+Q_{r,a}\),
\(r\ge0\), the initial case is \(Q_0=0\). For \(r\ge1\) one has the
following identities, with the last sum interpreted as zero when
\(r=1\):
\begin{equation}\protect\phantomsection\label{eq:gg-budget-actual-recursion}{
 \begin{aligned}
 \ker\delta_r&=\operatorname{im}P_{++},\\
 M_r&=\frac n{n+1}P_{-+}+\frac m{m+1}P_{+-}
                         +\frac{nm}{nm-1}P_{--},\\
 Q_r&=M_rJ_r^*(\tau\,\mathsf H Q_{r-1}-K_r),\\
 K_r(a,b)&=\mathsf V_{R_U(a,b)}U_{r-1},\\
 (\mathsf H Q_{r-1})(a,b)
  &=\sum_{j=0}^{r-2}D_{U_j}Q_{r-1,b}
                         [i_aU_{j+1}+Q_{j,a}].
 \end{aligned}
}\end{equation} \emph{Proof.} On \(A\otimes\operatorname{Sym}^nA\),
compression of an exchange of the outer slot with an inner slot is
\(C_n=((n+1)P_n-1)/n\). Indeed full symmetrization is
\((1+\sum_{i=1}^n\tau_{0i})/(n+1)\), and the compressed inner exchanges
agree. Its two eigenvalues are \(1,-1/n\). Thus
\begin{equation}\protect\phantomsection\label{eq:gg-budget-spencer-spectrum}{
 \delta_r^*\delta_r=2(1-C_n\otimes C_m)
}\end{equation} has nonzero eigenvalues \(2(1+1/n)\), \(2(1+1/m)\) and
\(2(1-1/(nm))\). Their minimum is at least \(8/5\). On antisymmetric
sources the complementary inverse of \(\delta_r\) is \(M_rJ_r^*\).

Apply the commutator in Lemma 3.2 to \(U_{r-1}\). The completely upper
\(U_{r+1}\) term cancels and gives
\(\delta_rQ_r=-K_r-(1-\tau)\mathsf H Q_{r-1}\). Differentiating the
constant zero-upper-projection identity for \(Q_{r-1}\) gives
\(J_r^*\mathsf H Q_{r-1}=0\). The inverse proves the formula. The last
line is the ordinary polynomial product rule; \(\mathsf H_a\Lambda=0\).
The curvature in \(K_r\) acts on all \(2r+2\) spinor slots, including
the derivative slots. Consistency is supplied by the already proved
polynomial Cartan identity, not surjectivity onto arbitrary
antisymmetric tensors. ∎

For a degree-\(q\) free-index polynomial \(T_a(u)\) in one spin factor,
with \(q\ge0\), define the upper inclusion below. For \(q\ge1\) also
define the lower inclusion and \(f\):
\begin{equation}\protect\phantomsection\label{eq:gg-budget-strip}{\begin{gathered}
(U_{q+1}g)_a=\frac{\partial_ag}{q+1}, \\
(\iota_qf)_0=-u^1f,\quad(\iota_qf)_1=u^0f, \\
g=u^aT_a,\quad
 f=\frac{\partial_0T_1-\partial_1T_0}{q+1}.
\end{gathered}}\end{equation} For \(q\ge1\), \(T=U_{q+1}g+\iota_qf\); for \(q=0\),
\(T=U_1g\). The upper inclusion is isometric for every \(q\ge0\). For
\(q\ge1\), the lower inclusion on degree \(q-1\) has squared norm
\((q+1)/q\le2\): on \((u^0)^{q-1}\) its two components have squared
norms \(1/q\) and \(1\), and equivariance proves the general formula.

A normalized contraction diagram has ordered symmetric input tensors,
specified alternating contractions between input slots, and averaged
symmetrization of all free outputs of each spin type. Selecting \(h\)
polynomial derivatives on a degree-\(p\) input means dividing them by
\((p)_h=p(p-1)\cdots(p-h+1)\). Thus each selected derivative is a
tensor-slot operation. A diagram with \(c\) alternating contractions has
norm at most \(2^{c/2}\). A contraction joining two slots of the same
symmetric input is zero.

\pandocstatementheading{\protect\phantomsection\label{lem:gg-budget-opening}Lemma B.2.}

Replace a degree-\(p\) input, incident to \(h\) contractions, by a lower
inclusion \((\iota_pf)_c\). Let \(D_c(u)\) be the resulting degree-\(N\)
free-index output. Opening its \(j\)th incident edge gives a normalized
diagram \(E_j\) with orientation sign \(\epsilon_j\). Then
\begin{equation}\protect\phantomsection\label{eq:gg-budget-opened-edge}{\begin{gathered}
u^cD_c=\frac1p\sum_{j=1}^h\epsilon_jE_j, \\
\frac1N u^c\partial_aD_c
 =\frac{N+1}{N}(U_{N+1})_a
                 \left(\frac1p\sum_j\epsilon_jE_j\right)
                      -\frac1ND_a\quad(N\ge1).
\end{gathered}}\end{equation} For a linear-shift insertion \(U_{p+1}g\), contraction
of its new outer index with \(u\) gives exactly the old diagram with
that input degree raised by one. An old lower output has transfer mass
at most one. For an old upper output, its lower leakage is a sum of
diagrams adding one edge to another input \(j\), with coefficient of
absolute value \((p_j-h_j)/(N(N+1))\). The total leakage mass is at most
\(1/(N+1)\).

\emph{Proof.} The normalized product rule is
\begin{equation}\protect\phantomsection\label{eq:gg-budget-normalized-product}{
 \frac{\partial_{a_1}\cdots\partial_{a_h}
                    (\epsilon_{bc}u^bf)}{(p)_h}
 =\frac{p-h}{p}\epsilon_{bc}u^b
          \frac{\partial_{a_1}\cdots\partial_{a_h}f}{(p-1)_h}
 +\frac1p\sum_{j=1}^h\epsilon_{a_jc}
       \frac{\prod_{\ell\ne j}\partial_{a_\ell}f}{(p-1)_{h-1}} .
}\end{equation} When \(h=p\), omit the first term rather than using its
zero denominator; when \(h=0\), the sum is empty. Multiplication by
\(u^c\) kills the first term. On an opened edge the paired alternating
tensors give a signed identity. Multiplying its released partner
derivative by \(u\) cancels the falling-factorial ratio, leaving exactly
\(E_j/p\), with no partner-degree multiplier. Differentiating that
identity proves the second formula.

For a linear shift the selected derivatives are
\(\partial_{a_1}\cdots\partial_{a_h}\partial_ag/(p+1)_{h+1}\).
Contracting \(a\) with \(u\) gives unit total coefficient, since
\((p+1-h)(p+1)_h/((p+1)(p)_h)=1\). For an old upper output the lower
stripping cancels the second derivatives and leaves
\((\partial_1D_0-\partial_0D_1)/(N(N+1))\). The derivative on the
distinguished input vanishes by antisymmetry. On another input \(j\) it
selects one of its \(p_j-h_j\) remaining slots and adds the stated
normalized edge. Their sum is at most \(N\). These arguments keep the
raw lower inclusions and their orientation signs. ∎

\protect\phantomsection\label{sec:gg-budget-mass}

\subsubsection{B.2. A single mass for all
transfers}\label{b.2.-a-single-mass-for-all-transfers}

The recursion generates finite normalized diagram representations by
induction: its curvature formulas, Spencer projections, shifts,
insertions and strippings use only alternating contractions and averaged
permutations. Proposition A.3 gives weight \(j+2\) to every order-\(j\)
record and weight two to \(\Lambda\). After stripping the raw lower
output factors, the positive output blocks \(A=-+\), \(B=+-\), \(C=--\)
have respective valences \((r+3,r+1)\), \((r+5,r-1)\), \((r+3,r-1)\). A
curvature-degree-\(d\) coefficient containing \(\ell\) scalar factors
has total input valence \(2r+6-4\ell\). Consequently its bare edge
counts are \(1-2\ell\) in \(A,B\) and \(2-2\ell\) in \(C\); negative
counts give zero. There are at most two bare edges independently of
arity.

Fix an ordered input list, including each record's order and chirality
and each scalar \(\Lambda\) factor. For a block \(\nu\in\{A,B,C\}\) let
\(\widetilde B_\nu\) denote the coefficient after its raw lower output
inclusions have been stripped. A bare normalized diagram specifies which
slots of distinct symmetric inputs are paired by alternating tensors;
all remaining slots of each spin type are averaged onto the indicated
symmetric output. On a degree-\(p\) input, selecting \(h\) slots means
\(\partial^h/(p)_h\), as in Lemma B.2. Slot order and contraction
orientations are part of the diagram. Raw lower inclusions are then
restored in their original spin factors. Only diagrams with the
specified input/output valences are used.

For each block define its atomic mass by
\begin{equation}\protect\phantomsection\label{eq:gg-budget-atomic-mass}{
 m_{r,\nu}(B)=\inf\left\{\begin{gathered}
\sum_{D\in\mathcal F}|a_D|:
        \widetilde B_\nu=\sum_{D\in\mathcal F}a_DD, \\
\mathcal F\text{ a finite set of such diagrams}
\end{gathered}\right\}.
}\end{equation} The infimum is taken on the finite-dimensional span
generated by the actual recursion, not on a different tensor norm. The
recursion supplies a finite representation at every coefficient; the
wire reductions below stay in this span. Scalar inputs have no spin
slots. For negative chirality the two spin degrees and lower-inclusion
roles are exchanged. In the tables below \(A=-+\) and \(B=+-\) always
denote the physical unprimed/primed signs; after chirality exchange
relabel the inherited \(A,B\) names accordingly. Their equal weights
make this relabelling immaterial to the mass. Put
\begin{equation}\protect\phantomsection\label{eq:gg-budget-common-mass}{
 m_r=2m_{r,A}+2m_{r,B}+m_{r,C}.
}\end{equation} Missing blocks have zero mass. Concatenation gives the
triangle inequality. The preceding diagram and raw-inclusion estimates
imply \(\|B\|_{\mathrm{op}}\le4m_r(B)\). In particular a nonzero
coefficient has nonzero mass. The recursive ordered expansion is used
throughout: inserting a polynomial replaces one specified position by
its ordered inner input list.

\pandocstatementheading{\protect\phantomsection\label{lem:gg-budget-transfers}Lemma B.3.}

For an insertion of a coefficient of \(Q_t\), \(t\ge1\), into one
specified position of a coefficient of \(Q_{r-1}\), the actual recursion
satisfies
\begin{equation}\protect\phantomsection\label{eq:gg-budget-transfer-bounds}{
 \begin{aligned}
 m_r(M_rJ_r^*\tau I_i(B,C))
       &\le\frac{K_0}{t+1}m_{r-1}(B)m_t(C),&K_0&=2^{24},\\
 m_r(M_rJ_r^*\tau L_iB)
       &\le c_rm_{r-1}(B),&
 c_r&=\frac{r(r+4)}{r(r+4)-1}.
 \end{aligned}
}\end{equation} The transferred curvature source has mass at most
\(C_*(r+1)\) for a structural \(C_*\). Moreover
\(1\le P_r=\prod_{s=1}^rc_s\le P_\infty\le e^{11/18}\).

\emph{Proof.} Work first with one finite diagram representation in each
block. All estimates are homogeneous in its absolute coefficients, so
passing to infima at the end is legitimate. A degree-\(d\) coefficient
with \(\ell\) scalar inputs has record orders summing to \(r+3-2d\). Its
total spin valence is therefore \(2(r+3-2d)+4(d-\ell)=2r+6-4\ell\).
Subtracting the stripped output valences in B.2 and dividing by two
gives \(1-2\ell\) bare edges in either single-lower block and
\(2-2\ell\) in the double-lower block. Thus each inner or outer
coefficient has at most two bare edges in both spin factors together.
Negative counts mean the coefficient is zero; no edge can join two slots
of one symmetric input.

For \(h\) selected slots distributed among symmetric factors of degrees
\(p_i\), the normalized assignment weights are
\begin{equation}\protect\phantomsection\label{eq:gg-budget-assignment-weights}{
 \binom h{\alpha_1,\ldots,\alpha_d}
                  \frac{\prod_i(p_i)_{\alpha_i}}{(\sum_i p_i)_h}.
}\end{equation} Their sum is one by the falling-factorial Vandermonde
identity: both sides count ordered selections of \(h\) distinct slots. A
zero denominator corresponds to an unavailable selection, which is
omitted. Thus splitting derivatives or averaging permutations incurs no
factor depending on the number of inputs.

The following list accounts for all alternating tensors in a composition
before its final wire reduction.

{\def\LTcaptype{none} % do not increment counter
\begin{longtable}[]{@{}
  >{\raggedright\arraybackslash}p{(\linewidth - 8\tabcolsep) * \real{0.1579}}
  >{\raggedleft\arraybackslash}p{(\linewidth - 8\tabcolsep) * \real{0.2105}}
  >{\raggedleft\arraybackslash}p{(\linewidth - 8\tabcolsep) * \real{0.2105}}
  >{\raggedleft\arraybackslash}p{(\linewidth - 8\tabcolsep) * \real{0.2105}}
  >{\raggedleft\arraybackslash}p{(\linewidth - 8\tabcolsep) * \real{0.2105}}@{}}
\toprule\noalign{}
\begin{minipage}[b]{\linewidth}\raggedright
Operation
\end{minipage} & \begin{minipage}[b]{\linewidth}\raggedleft
Outer coefficient
\end{minipage} & \begin{minipage}[b]{\linewidth}\raggedleft
Inner coefficient
\end{minipage} & \begin{minipage}[b]{\linewidth}\raggedleft
Final projection
\end{minipage} & \begin{minipage}[b]{\linewidth}\raggedleft
Total bound
\end{minipage} \\
\midrule\noalign{}
\endhead
\bottomrule\noalign{}
\endlastfoot
Bare alternating contractions & \(2\) & \(2\) & \(0\) & \(4\) \\
Raw lower output inclusions & \(2\) & \(2\) & \(0\) & \(4\) \\
Normal\-ized lower stripping contractions & \(0\) & \(0\) & one per spin
factor & \(2\) \\
\end{longtable}
}

Indeed the inserted coefficient replaces one specified input box of the
outer coefficient. Each box has the bare-edge bound just proved and at
most one raw lower inclusion in each spin factor. The maps \(\tau\) and
\(J_r^*\) only permute or average slots. Upper projection is averaging,
while lower projection has stripping factor \(q/(q+1)\le1\) times one
normalized slot contraction before restoring the declared raw lower
output. Those restored output factors are not extra internal wires. No
other operation in \(M_rJ_r^*\tau I_i\) introduces an alternating
tensor.

Reduce a fixed assignment by following each wire. Paired alternating
tensors give signed identities; a closed wire gives a scalar of absolute
value two; a path ending twice in one symmetric input or in the
symmetric output vanishes. The remaining paths give bare diagrams of the
final valence. There are at most ten alternating tensors, hence the
coefficient bound \(2^{10}\) is valid. The assignment weights still sum
to one. Including the three target block weights, whose sum is five, and
the Spencer factor, at most \(5/4\), gives at most
\(5(5/4)2^{10}<2^{13}\). We retain the coarser bound \(2^{20}\) for
these compositions and the opened-wire representations below. This
counts both spin factors together; it is not a separate norm-equivalence
loss at each arity or derivative order.

To obtain the inverse-order gain, choose a lower spin factor of the
inner coefficient before taking absolute values. All choices are
exhausted by the following table, with the block signs referring to the
unprimed and primed factors in that order.

{\def\LTcaptype{none} % do not increment counter
\begin{longtable}[]{@{}
  >{\raggedright\arraybackslash}p{(\linewidth - 8\tabcolsep) * \real{0.2000}}
  >{\raggedright\arraybackslash}p{(\linewidth - 8\tabcolsep) * \real{0.2000}}
  >{\raggedright\arraybackslash}p{(\linewidth - 8\tabcolsep) * \real{0.2000}}
  >{\raggedright\arraybackslash}p{(\linewidth - 8\tabcolsep) * \real{0.2000}}
  >{\raggedright\arraybackslash}p{(\linewidth - 8\tabcolsep) * \real{0.2000}}@{}}
\toprule\noalign{}
\begin{minipage}[b]{\linewidth}\raggedright
Inner chirality
\end{minipage} & \begin{minipage}[b]{\linewidth}\raggedright
Degrees in \(W_t\)
\end{minipage} & \begin{minipage}[b]{\linewidth}\raggedright
Inner block
\end{minipage} & \begin{minipage}[b]{\linewidth}\raggedright
A lower factor available for opening
\end{minipage} & \begin{minipage}[b]{\linewidth}\raggedright
Its degree \(p\)
\end{minipage} \\
\midrule\noalign{}
\endhead
\bottomrule\noalign{}
\endlastfoot
\(+\) & \((t+4,t)\) & \(A=-+\) & unprimed & \(t+4\) \\
\(+\) & \((t+4,t)\) & \(B=+-\) & primed & \(t\) \\
\(+\) & \((t+4,t)\) & \(C=--\) & unprimed or primed & \(t+4\) or
\(t\) \\
\(-\) & \((t,t+4)\) & \(A=-+\) & unprimed & \(t\) \\
\(-\) & \((t,t+4)\) & \(B=+-\) & primed & \(t+4\) \\
\(-\) & \((t,t+4)\) & \(C=--\) & unprimed or primed & \(t\) or
\(t+4\) \\
\end{longtable}
}

For the double-lower block choose the unprimed factor, for definiteness;
do not sum over the two choices. There is no inner \(++\) block, and
\(Q_0=0\), so \(t\ge1\) and \(p\ge t\ge1\). In the chosen factor write
the raw inclusion as \((\iota_pf)_c=\epsilon_{bc}u^bf\). The normalized
product formula of Lemma B.2 has unhit coefficient \((p-h)/p\) and hit
coefficients exactly \(1/p\). After multiplication by \(u^c\), the unhit
term vanishes by \(\epsilon_{bc}u^bu^c=0\). Each hit opens one of the at
most two outer bare edges incident to this input. Multiplication of its
released partner derivative by \(u\) cancels the partner's
falling-factorial ratio. Thus no partner-degree factor survives.

The outer block may have either sign in the chosen factor. The complete
outer-block correspondence is

{\def\LTcaptype{none} % do not increment counter
\begin{longtable}[]{@{}llll@{}}
\toprule\noalign{}
Chosen factor & Outer \(A=-+\) & Outer \(B=+-\) & Outer \(C=--\) \\
\midrule\noalign{}
\endhead
\bottomrule\noalign{}
\endlastfoot
unprimed & lower & upper & lower \\
primed & upper & lower & lower \\
\end{longtable}
}

These rows hold for either outer chirality. To evaluate the exact slot
map, write an upper outer block as \(B_a=\partial_ag/(q+1)\), where its
old \(W_{r-1}\) degree is \(q\). Substitute the chosen coefficient into
\(g\) and leave its new derivative index \(c\) free, calling the
resulting polynomial \(D_c\). Substitution retains its prescribed output
degree. After \(J_r^*\tau\), the output in this factor is
\begin{equation}\protect\phantomsection\label{eq:gg-budget-exact-slot-degree}{
 T'_a=\frac{u^c\partial_aD_c}{N},\qquad
 \deg D_c=N=q+1\in\{r,r+4\},\qquad
 1\le t\le r-2.
}\end{equation} The last range follows because \(Q_{r-1}\) uses only
\(U_0,\ldots,U_{r-2}\) and \(Q_0=0\). The averaging in \(J_r^*\) is
multiplication by \(u^c\) upon evaluation on equal symmetric arguments;
it introduces no degree factor. Thus \(1/N\le1/(t+1)\).

Put \(E=u^cD_c\). Euler's identity gives the exact upper projection
\begin{equation}\protect\phantomsection\label{eq:gg-budget-exact-upper-projection}{
 u^aT'_a=\frac{u^cu^a\partial_aD_c}{N}=E,
 \qquad P_+T'=U_{N+1}E.
}\end{equation} Its coefficient mass is at most \(2/p\) by the
opened-edge computation above. In particular, the two upper
contributions obtained by applying Lemma B.2 separately cancel to this
expression. For the lower final projection, the mixed second derivatives
cancel:
\begin{equation}\protect\phantomsection\label{eq:gg-budget-upper-lower-cancellation}{
 \begin{aligned}
 \partial_0(u^c\partial_1D_c)
       -\partial_1(u^c\partial_0D_c)
      &=\partial_1D_0-\partial_0D_1,\\
 \operatorname{strip}_N T'
      &=\frac{\partial_1D_0-\partial_0D_1}{N(N+1)}
        =-\frac1N\operatorname{strip}_N D .
 \end{aligned}
}\end{equation} The remaining stripping is a normalized one-slot
operation with scalar factor \(N/(N+1)\le1\). Its graph cost is included
in the common \(2^{20}\) bound, and it leaves the explicit gain
\(1/N\le1/(t+1)\). For a lower outer block, its factor
\(\epsilon_{ba}u^b\) is independent of the record input. The output is
\(\epsilon_{ba}u^b E\), so its final upper projection is zero and its
lower coefficient has the same opened mass \(2/p\le4/(t+1)\). These
identities cover both final projections. The other spin factor uses its
ordinary normalized projections and the same common graph bound.
Together the two tables cover every inner and outer chirality and block,
with no multiplicity for the specified position.

The nonlinear transfer mass is consequently at most
\(2^{22}m_{r-1}(B)m_t(C)/(t+1)\le K_0m_{r-1}(B)m_t(C)/(t+1)\). All
coefficients are measured in the stipulated atomic mass: the assignment
weights, ten-tensor reduction, target block weights and Spencer factors
were counted above. Passing from finite representations to their
defining infima proves the estimate with no dependence on arity or
derivative order. The number of eligible positions is counted separately
in Theorem B.4.

For the linear transfer use the exact leakage coefficients, not the
coarse graph bound. On raw block masses the componentwise transfer
matrix is
\begin{equation}\protect\phantomsection\label{eq:gg-budget-linear-matrix}{
 \begin{pmatrix}
 n/(n+1)&0&0\\
 0&m/(m+1)&0\\
 c_r/(m+1)&c_r/(n+1)&c_r
 \end{pmatrix},\qquad n=r+4,\quad m=r .
}\end{equation} The first old block stays lower in its first factor and
leaks only from its upper second factor; the second exchanges these
roles; the third stays lower in both. This gives all entries from Lemma
B.2. Its column sums with weights \((2,2,1)\) are at most \(c_r\) times
the input weights. For example the only potentially stronger comparison
reduces to \(m(n+1)\le2(m+1)^2\), whose difference is \(r^2-r+2>0\). The
other comparison is \(n(m+1)\le2(n+1)^2\), with difference
\(r^2+15r+46>0\). Exchanging the two spin factors exchanges these
comparisons and the first two block weights, so the same bound holds for
negative chirality. This proves the linear bound without repeated fixed
graph losses. Finally
\begin{equation}\protect\phantomsection\label{eq:gg-budget-linear-product}{
 \log c_r\le\frac1{r(r+4)-1}\le\frac1{r(r+3)},\qquad
 \sum_{r\ge1}\frac1{r(r+3)}=\frac{11}{18}.
}\end{equation} The curvature source has one order-zero curvature factor
and acts on the \(2r+2\) slots of \(U_{r-1}\). Each slot action uses a
fixed finite alternating-tensor formula. Normalized stripping has the
same assignment bound, and final valence still permits at most two bare
edges. The fixed curvature expressions and bounded Spencer factors
therefore give \(C_*(r+1)\). This includes the scalar-curvature factor
and the action on derivative slots. ∎

Rescale the same mass by \(\gamma_r=4(P_\infty/P_r)m_r\). It dominates
the operator norm, makes each specified linear transfer nonexpanding,
and changes a nonlinear insertion by factor
\(P_{r-1}P_t/(4P_\infty P_r)\le1/4\). Lemma B.3 bounds each insertion
with fixed input and output chiralities uniformly, so its rescaled
constant is at most \(K_0/4\). The product rule in Lemma B.1 sums over
the two possible chiralities of the intermediate record; a scalar
\(\Lambda\) position contributes zero. Let \(K\) include this finite
sum. Since \(2(K_0/4)=K_0/2\le K_0\), the gain \(K/(t+1)\) holds with a
structural \(K\le K_0\). Eligible ordered insertion positions are still
summed, not averaged away.

\protect\phantomsection\label{sec:gg-budget-coefficients}

\subsubsection{B.3. Ordered coefficient
bounds}\label{b.3.-ordered-coefficient-bounds}

\pandocstatementheading{\protect\phantomsection\label{thm:gg-budget-coefficients}Theorem B.4.}

There are structural \(a,b>0\) such that the original ordered
degree-\(d\) coefficients satisfy
\begin{equation}\protect\phantomsection\label{eq:gg-budget-coefficient-conclusion}{
 \gamma_r(B^{[d]}_{r,\mathbf j}),\ 
 \|B^{[d]}_{r,\mathbf j}\|_{\mathrm{op}}
 \le ab^{d-2}r\,\frac{r!}{\prod_i j_i!},
 \qquad |\mathbf j|=r+3-2d,\quad d\ge2 .
}\end{equation} \emph{Proof.} At degree two, a source introduced at
order \(t\) has indices \((0,t-1)\) and mass at most \(C(t+1)\). The
unit linear transfers give at final indices \(j+k=r-1\)
\begin{equation}\protect\phantomsection\label{eq:gg-budget-quadratic-sum}{
 C\sum_{t=1}^{r-j}(t+1)\binom{r-t}{j}
   =C\left\{\binom{r+1}{j+2}+\binom r{j+1}\right\}.
}\end{equation} Multiplying by \(j!k!/r!\) bounds this by
\(C\{(r+1)/((j+1)(j+2))+1/(j+1)\}\le2Cr\). Choose \(a\) to include the
fixed curvature orderings and labels.

Normalize a degree-\(d\) coefficient by \(\prod j_i!/r!\) and bound it
by \(A_dr\). Put \(k=|\mathbf j|=r+3-2d\). The preceding-order linear
terms contribute at most \(A_d((r-1)/r)\sum_i j_i\le A_dk\). For an
inner degree \(e\) inserted into outer degree \(f\), \(e+f=d+1\), the
old position index is \(t=\sum_{\mathrm{inner}}j_i+2e-3\ge1\). Its
factorial denominator cancels the inner numerator. After summing the
intermediate chirality, the remaining normalized term is at most
\(KA_fA_e(r-1)t/(r(t+1))\le KA_fA_e\). There are \(f\) positions for the
fixed ordered lists. The reserve is \(r-k=2d-3\), so the induction
closes with
\begin{equation}\protect\phantomsection\label{eq:gg-budget-degree-recursion}{
 A_2=a,\qquad
 A_d=\frac K{2d-3}\sum_{\substack{e+f=d+1\\e,f\ge2}}fA_eA_f
 \le\frac{2K}{3}\sum_{\substack{e+f=d+1\\e,f\ge2}}A_eA_f .
}\end{equation} The last inequality symmetrizes \(e,f\) and uses
\((d+1)/(2(2d-3))\le2/3\) for \(d\ge3\). The binary-tree recurrence
gives \(A_d\le a(2Ka/3)^{d-2}\operatorname{Cat}_{d-2}
\le ab^{d-2}\) with \(b=8Ka/3\). All steps refer to the recursively
ordered expansion, so no unique ordered expression of a commutative
polynomial is assumed. ∎

For completeness these estimates define the actual nonlinear field on
smaller analytic-record radii. For the grouped records \(X_j\) of
Theorem 3.3 set \(a_0=18a\) and \(b_0=3b\): the at most \(3^d\) input
types and two output chiralities give
\(2\cdot3^d a b^{d-2}=a_0b_0^{d-2}\). This is the same single conversion
used there. If \(M=\|X\|_R\), \(0<s<R\) and \(\vartheta=b_0s^2M<1\),
then
\begin{equation}\protect\phantomsection\label{eq:gg-budget-nonlinear-sum}{
 \|Q(X)\|_s
 \le a_0sM^2\left(2+\frac{s}{R-s}\right)
                     \frac{2-\vartheta}{(1-\vartheta)^2}.
}\end{equation} Indeed put \(c_j=R^j\|X_j\|/j!\) and \(q=s/R\).
Unrestricted product sums are at most \(M^d\), and their \(k=\sum j_i\)
weighted sum is at most \(dM^dq/(1-q)\). Since \(r=k+2d-3\), summing the
bound gives the display using
\(\sum_{d\ge2}d\vartheta^{d-2}=(2-\vartheta)/(1-\vartheta)^2\). Strict
interior inequalities give normal convergence and parameter holomorphy.
Lemma B.1 identifies these coefficients with the original polynomial
derivation. The scalar characteristic argument in Theorem 3.3 therefore
proves convergence of that derivation's flows and their finite words,
rather than flows of a replacement field.

\protect\phantomsection\label{app:gg-contact-linear-operators}

\subsection{Appendix C. Contact cochains and the analytic linear
correction}\label{appendix-c.-contact-cochains-and-the-analytic-linear-correction}

\protect\phantomsection\label{sec:gg-contact-cochain-repair}

\subsubsection{C.1. The fixed cochain
operations}\label{c.1.-the-fixed-cochain-operations}

The contact normal-jet resolution and its conformal-metric deformation
interpretation are classical (LeBrun 1991, secs. 8, pp.29--31 and
Proposition 8.1). We fix the following cochain representatives and
normalizations.

At the flat reference let \(K=\mathcal O_Q(-1,-1)\) and
\(0\to K\to V\otimes\mathcal O_Q\to N\to0\). Put
\(L=K^*=\mathcal O(1,1)\) and \(F_n=L\otimes\operatorname{Sym}^nN^*\).
The symmetric derivative gives
\begin{equation}\protect\phantomsection\label{eq:gg-contact-tautological-resolution}{
 0\longrightarrow F_n\longrightarrow
 A_n=L\otimes\mathcal P_n
 \xrightarrow{d_n}B_n=L^2\otimes\mathcal P_{n-1}
 \longrightarrow0,\qquad d_n=\frac1n k^a\partial_{x^a}.
}\end{equation} For \(n=0\) omit the last term. The kernel consists of
polynomials on \(V/K\); local integration in \(k\) proves surjectivity.
We use the usual alternating Čech differential, so
\((\delta a)_{ij}=a_j-a_i\), with reference transitions understood.

Both fixed positive line bundles \(L,L^2\) have vanishing \(H^1,H^2\).
Their product two-circle Laurent contractions are bounded on fixed
nested annuli independently of the constant polynomial coefficient
space. A local holomorphic right inverse is
\begin{equation}\protect\phantomsection\label{eq:gg-contact-holomorphic-right-inverse}{
 (R_n^{\mathrm{hol}}b)(x)
   =n\lambda(x)\int_0^1
        b(x-(1-t)\lambda(x)k)\,\mathrm dt,\qquad \lambda(k)=1 .
}\end{equation} Its bound is \(nC^n\) in symmetric coefficient norms on
the fixed finite direction charts. Given a two-cocycle \(c\) in \(F_n\),
split it in \(A_n\) as \(\delta a=c\); split \(d_na\) in \(B_n\) as
\(\delta b=d_na\); then
\begin{equation}\protect\phantomsection\label{eq:gg-contact-fixed-repair-bound}{\begin{gathered}
S_nc=a-\delta(R_n^{\mathrm{hol}}b), \\
\delta S_nc=c,\quad d_nS_nc=0,\quad
 \|S_nc\|_{\mathrm{inner}}
       \le C_0(n+1)C_1^n\|c\|_{\mathrm{outer}}.
\end{gathered}}\end{equation} This proves the repair used in Lemma 5.1, including its
analytic estimate. Polynomial factors and fixed exponentials are
summable after one strict normal-radius restriction.

We also specify the closed-cochain lift \(\ell_n\) rather than choose
uncontrolled degreewise representatives. Give the constant tensor spaces
their full Hilbert norms and the two fixed line bundles their product
Fubini--Study metrics. Identify \(\mathscr C_n=H^1(F_n)\) with the
quotient \(H^0(B_n)/d_nH^0(A_n)\) and give it the quotient norm. For a
class \(u\) choose its unique minimal global representative \(b_u\). The
pointwise smooth right inverse \(R_n=d_n^\dagger(d_nd_n^\dagger)^{-1}\)
has \(\|R_n\|\le\sqrt n\): in a unitary frame with \(k=e_0\), normalized
contraction on a monomial with \(j\) copies of that slot has singular
value \(\sqrt{j/n}\). In particular \(d_nd_n^\dagger\ge n^{-1}\). Fixed
angular derivatives of \(d_n\) are bounded independently of \(n\);
repeated differentiation of this inverse gives polynomial bounds in
\(n\) for every fixed derivative of \(R_n\).

The closed form \(\mu_u=\bar\partial R_nb_u\) takes values in \(F_n\).
On each fixed buffered chart solve \(\bar\partial w_i=\mu_u\) in \(A_n\)
by the componentwise two-variable Cauchy--Green contraction. Its
operator bounds concern only the fixed line bundle and chart. Since
\(d_nw_i\) is holomorphic, put \(\phi_i=w_i-R_n^{\mathrm{hol}}d_nw_i\).
Then \(\phi_i\) takes values in \(F_n\) and
\(\bar\partial\phi_i=\mu_u\). The sections \(a_i=R_nb_u-\phi_i\) are
holomorphic in \(A_n\), with \(d_na_i=b_u\). Define
\(\ell_nu=(a_j-a_i)_{ij}\). This fixes the sign and the actual
holomorphic representative. The same formulas prove
\begin{equation}\protect\phantomsection\label{eq:gg-contact-fixed-class-lift-bound}{
 \delta\ell_nu=0,\qquad [\ell_nu]=u,\qquad
 \|\ell_nu\|_{\mathrm{inner}}
          \le C(n+1)^bC_2^n\|u\|_{\mathscr C_n}
}\end{equation} for a fixed sufficiently high angular Sobolev norm in
the intermediate estimate, followed by Sobolev embedding on the compact
real four-dimensional \(Q\). All degree dependence has been exhibited in
\(R_n\) and the local symmetric powers; no uniform Green estimate for
all the \(F_n\) is assumed. Minimal representatives are complex-linear
in their quotient classes, since their orthogonal complements are
complex subspaces. The fixed choices therefore preserve holomorphic
input parameters and commute with finite truncation. These definitions
fix the cochains used in §5.

\protect\phantomsection\label{sec:gg-contact-conformal-identification}

\subsubsection{C.2. The contact class and its two original
records}\label{c.2.-the-contact-class-and-its-two-original-records}

Global sections of \(L\) are linear functions of the null vector \(k\),
and those of \(L^2\) are its quadratic functions modulo \(\eta(k,k)\).
Consequently, for \(n\ge3\),
\begin{equation}\protect\phantomsection\label{eq:gg-contact-conformal-quotient}{
 \mathscr C_n\cong
 \frac{\mathcal P_{n-1}\otimes\operatorname{Sym}^2_0V^*}
      {\mathsf K(\mathcal P_n\otimes V)},\qquad
 (\mathsf K\xi)_{ab}
   =\partial_a\xi_b+\partial_b\xi_a
                      -\frac12\partial^c\xi_c\,\eta_{ab}.
}\end{equation} The same sequence identifies \(H^0(F_n)\) with
homogeneous conformal Killing fields of degree \(n\). Their higher
vanishing, used in §D.6, is explicit. If
\(\partial_a\xi_b+\partial_b\xi_a=2\vartheta\eta_{ab}\) and
\(\vartheta=\operatorname{div}\xi/4\), differentiation and permutation
give
\begin{equation}\protect\phantomsection\label{eq:gg-contact-killing-finite-type}{
 \partial_a\partial_b\xi_c
  =\eta_{ac}\partial_b\vartheta+\eta_{bc}\partial_a\vartheta
                                  -\eta_{ab}\partial_c\vartheta .
}\end{equation} Diverging in \(c\) gives
\(2\partial_a\partial_b\vartheta=-\eta_{ab}\Box\vartheta\). Taking the
trace gives \(\Box\vartheta=0\) and then its Hessian vanishes. Thus
every third derivative of \(\xi\) is zero. The degree-zero fields are
translations; degree one consists of skew linear fields and a dilation;
degree two is \(2(b\cdot x)x-\eta(x,x)b\). In particular \(H^0(F_n)=0\)
for \(n\ge3\).

The normalization is the connecting map \(J_n\) from
\(b_h=h_{ab}k^ak^b/(2n)\). A global middle section \(k^a\xi_a\) maps to
\(k^ak^b(\mathsf K\xi)_{ab}/(2n)\). This proves the quotient and its
kernel. Local holomorphic lifts are
\begin{equation}\protect\phantomsection\label{eq:gg-contact-metric-ray-primitive}{
 a_i(x,k)=\frac12\int_0^{\lambda_i(x)}
      h_{ab}(x-\lambda_i(x)k+sk)k^ak^b\,\mathrm ds,\qquad
 J_n[h]=[\delta a].
}\end{equation} They give exactly the first variation of the null
hitting maps: changing \(\frac12\eta^{ab}p_ap_b\) to that Hamiltonian
minus \(h^{ab}p_ap_b/2\) is removed locally by these Hamiltonian
primitives. Their overlap difference is therefore the contact
deformation class. A trace has zero null contraction. The finite
direction cover bounds the integrals on a smaller tube. Conversely embed
a contact cocycle into \(A_n\) and split it by the fixed \(L\) Laurent
contraction; its \(d_n\) image is global in \(B_n\) and reads
\(h/(2n)\). Normalized differentiation is bounded on coefficient tensors
and fixed coordinate substitutions cost at most \(C^n\). Both directions
therefore act on actual analytic germs after one fixed restriction.

The first-order lifts of §5 have the original, rather than freely
rescaled, curvature records. In
\(\eta=2\,\mathrm dx^{00'}\mathrm dx^{11'}-
2\,\mathrm dx^{01'}\mathrm dx^{10'}\), take \(F=(w^0)^{r+4}/\zeta\).
Splitting its lifted Hamiltonian \(-(r+4)\varrho(w^0)^{r+3}/\zeta\)
gives the metric representative
\begin{equation}\protect\phantomsection\label{eq:gg-contact-positive-normalization}{
 h=-2(r+4)(r+3)(x^{00'})^{r+2}(\mathrm dx^{01'})^2 .
}\end{equation} For an explicit check of its contact class, put
\(p=r+4\), \(A=x^{00'}\), \(B=x^{01'}\) and
\(k=(\zeta\varrho,-\varrho,-\zeta,1)\). In a common north frame the two
middle-bundle primitives are
\begin{equation}\protect\phantomsection\label{eq:gg-contact-highest-vector-primitives}{
 u_N=\frac{p\varrho}{\zeta}
             \bigl((A+\zeta B)^{p-1}-A^{p-1}\bigr),\qquad
 u_S=-\frac{p\varrho}{\zeta}A^{p-1}.
}\end{equation} Both have derivative
\(-\!p(p-1)\varrho^2A^{p-2}=h(k,k)/2\) along \(k\). The first is
holomorphic at \(\zeta=0\) by divisibility; the second, after the \(L\)
south-frame factor \(\zeta^{-1}\), is holomorphic at infinity. The
ordered difference \(u_S-u_N\) is exactly
\(-\!p\varrho(A+\zeta B)^{p-1}/\zeta\), the prescribed lifted
Hamiltonian. The dependence on \(\varrho\) is a global section of its
degree-one line. This verifies the class.

Its trace, divergence and wave operator vanish, hence its linearized
Ricci tensor vanishes. By eq.~\ref{eq:gg-last-slot-linear-curvature} its
last-slot component is
\(C^{(1)}_{00'\,01'\,00'\,01'}=\partial_{00'}^2h_{01'01'}/2=-(r+4)!(x^{00'})^r/r!\).
The minus sign in the component readout in
eq.~\ref{eq:gg-weyl-spinor-readout} therefore gives
\(\Psi^+_{0000}=-\partial_{00'}^2h_{01'01'}/2=(r+4)!(x^{00'})^r/r!\).
The order-\(r\) free component is exactly \((r+4)!\), the extraction
multiplier in Theorem 3.1. The maps are equivariant under the two spin
factors, and the orbit of this highest pure tensor spans each separately
symmetric space, so the identity holds for every component.

For the negative highest source \(F_{-6}=(w^0)^r/\zeta\), the flat
contour is \(\Psi^-_{0'0'0'0'}=(x^{00'})^r\) and its record is \(r!\).
Its derived class (\ref{eq:gg-derived-dual-class}) is
\(G=c_r(\eta^{0'})^{r+4}/\varrho\), where \(c_r=(-1)^rr!/(r+4)!\). Put
\(u=x^{00'}\), \(q=x^{10'}\) and \(a_r=r!/(r+3)!\). The factor exchange
\(\widetilde A=-u\), \(\widetilde B=-q\), \(\widetilde\zeta=\varrho\),
\(\widetilde\varrho=-\zeta\) gives the primitives
\begin{equation}\protect\phantomsection\label{eq:gg-contact-negative-highest-primitives}{
 \widetilde u_N=\frac{a_r\zeta}{\varrho}
                  \bigl((u+\varrho q)^{r+3}-u^{r+3}\bigr),\qquad
 \widetilde u_S=-\frac{a_r\zeta}{\varrho}u^{r+3} .
}\end{equation} The ordered difference \(\widetilde u_S-\widetilde u_N\)
is the prescribed \(h_-=-a_r\zeta(u+\varrho q)^{r+3}/\varrho\). For the
same \(k=(\zeta\varrho,-\varrho,-\zeta,1)\) their derivatives are both
\(-r!\zeta^2u^{r+2}/(r+2)!\). Consequently the metric representative and
its component readout are
\begin{equation}\protect\phantomsection\label{eq:gg-contact-negative-component-readout}{
 h^-=-\frac{2r!}{(r+2)!}u^{r+2}(\mathrm dq)^2,\qquad
 C^{(1)}_{uquq}=-u^r,\qquad
 \Psi^-_{0'0'0'0'}=-C^{(1)}_{uquq}=u^r .
}\end{equation} The negative representative also has zero trace,
divergence and wave operator. Thus the \(r\)th derivative is exactly
\(r!\), including the \((-1)^r\) from \(\eta=-\kappa x\) and the stated
dual coefficient. The same orbit argument gives all negative components.
Thus the two actual cohomology classes are the two linear vacuum
inclusions with exactly the original component and factorial
conventions.

\protect\phantomsection\label{sec:gg-linear-vacuum-inclusion}

\subsubsection{C.3. The vacuum inclusion and its
kernel}\label{c.3.-the-vacuum-inclusion-and-its-kernel}

Write \(\mathsf R(h)\) for the linearized Ricci tensor and
\(\mathsf P(h)=\frac12(\mathsf R(h)-\operatorname{tr}\mathsf R(h)\eta/6)\).
The linear Cotton tensor is \(\mathrm d\mathsf P(h)\). For \(r=n-3\),
the free Weyl projection is the same geometric readout as in §3:
\begin{equation}\protect\phantomsection\label{eq:gg-contact-free-projection-readout}{
 \Pi_n[J_nh]=
 \left.\operatorname{Sym}\partial^r
       \bigl(\mathfrak p_+(C^{(1)}(h)),
             \mathfrak p_-(C^{(1)}(h))\bigr)\right|_{x=0} .
}\end{equation} Here the symmetrizations and both component contractions
are exactly eq.~\ref{eq:gg-free-record-projection} and
eq.~\ref{eq:gg-weyl-spinor-readout}. It and the linear Cotton tensor
descend to \(\mathscr C_n\): symmetric gradients have zero curvature, a
trace has zero Weyl tensor, and \(\mathsf P(2f\eta)=-\partial^2f\).

If \(m=n-1\) and \(\mathrm d\mathsf P(h)=0\), two Euler homotopies give
\(f=x^ax^b\mathsf P(h)_{ab}/(m(m-1))\) with
\(\partial^2f=\mathsf P(h)\). Hence \(h+2f\eta\) is Ricci-flat and
represents the same contact class.

Conversely, for \(r=n-3\) insert \(r\) copies of \(x\) into
\(U_r^\pm/r!\) to form the flat Weyl spinor polynomials. Apply the
inverse component formula eq.~\ref{eq:gg-weyl-spinor-decomposition} to
form \(C\), and put
\(\mathcal W^a{}_b=\tfrac12\eta^{af}C_{cdbf}\,\mathrm dx^c\wedge\mathrm dx^d\).
Their separate symmetry gives zero flat divergence. This curvature
two-form obeys \(\mathrm d\mathcal W=0\) and
\(\mathcal W^a{}_b\wedge\mathrm dx^b=0\). For \(E=x^a\partial_a\),
define
\begin{equation}\protect\phantomsection\label{eq:gg-linear-vacuum-radial-homotopy}{
 \omega^a{}_b=\frac{\iota_E\mathcal W^a{}_b}{r+2},\qquad
 e^a=\frac{x^b\omega^a{}_b}{r+3},\qquad
 v(U)_{ab}=\eta_{ac}e^c{}_b+\eta_{bc}e^c{}_a .
}\end{equation} The Euler identity gives \(\mathrm d\omega=\mathcal W\)
and \(\iota_E\omega=0\). Applying it to the closed form
\(\omega^a{}_b\wedge\mathrm dx^b\) gives
\(\mathrm de^a+\omega^a{}_b\wedge\mathrm dx^b=0\). Since \(\omega\) is
skew, it is the linearized Levi--Civita connection of \(v(U)\). Thus
\(v(U)\) is Ricci-flat, has precisely \(U\), and is in marked radial
normal gauge. Set \(\iota_nU=J_n[v(U)_0]\); taking the trace-free
representative only passes to the conformal quotient.

For a homogeneous Ricci-flat variation, flat Bianchi and commuting
derivatives make its Weyl derivatives separately symmetric by Lemma A.2.
Its one homogeneous free list therefore determines its full curvature.
If that list is zero, the Christoffel one-forms are closed. Two
polynomial Poincaré integrations, using torsion symmetry, give
\(\Gamma^a{}_{bc}=\partial_b\partial_c\xi^a\). The remaining
\(h-\mathcal L_\xi\eta\) is constant and vanishes by positive degree.
This proves
\begin{equation}\protect\phantomsection\label{eq:gg-contact-cotton-kernel}{
 \ker(\mathrm d\mathsf P:\mathscr C_n\to\operatorname{im}\mathrm d\mathsf P)
       =\operatorname{im}\iota_n,\qquad \Pi_n\iota_n=1 .
}\end{equation} The radial construction is also analytic for a whole
factorially bounded free list. Tensor-norm conversion costs a polynomial
in \(r\), the fixed spin-to-coordinate change costs at most \(C^r\), and
division by \(r!\) and the two positive Euler degrees gives a normally
convergent metric series after a fixed restriction. The same reasoning
and Cauchy estimates apply to the entire free projection of any analytic
metric variation, including a nonvacuum one.

\protect\phantomsection\label{sec:gg-conserved-linear-potential}

\subsubsection{C.4. A bounded potential for the conserved linear
source}\label{c.4.-a-bounded-potential-for-the-conserved-linear-source}

\pandocstatementheading{\protect\phantomsection\label{lem:gg-conserved-linear-ricci-potential}Lemma C.1.}

For \(\eta=\operatorname{diag}(\epsilon,1,1,1)\), \(\epsilon=\pm1\), a
symmetric holomorphic tensor \(S\) with
\(\partial^aS_{ab}=\frac12\partial_b\operatorname{tr}S\) has a fixed
complex-linear potential \(G S\) satisfying \(\mathsf R(GS)=S\). There
are structural \(c,C>0\) with
\begin{equation}\protect\phantomsection\label{eq:gg-conserved-linear-potential-bound}{
 \|GS\|_{\mathbb D_{cR}^4}\le CR^2\|S\|_{\mathbb D_R^4}.
}\end{equation} Homogeneous degree \(d\) maps to degree \(d+2\). The
operation preserves bounded holomorphic parameter families and real
coefficients in Lorentz coordinates.

\emph{Proof.} Put \(T=S-\frac12(\operatorname{tr}S)\eta\) and
\(\bar h=h-\frac12(\operatorname{tr}h)\eta\). Then
\(\operatorname{div}T=0\). Solve \(\Box\bar h=-2T\) with
\(\beta_b=\partial^a\bar h_{ab}=0\). We specify compatible initial data
at \(t=0\). Let \(\Delta_yf=-2\epsilon T_{00}(0,y)\), with
\(f=\partial_1f=0\) on \(y^1=0\). Set \(A_{ij}=\delta_{ij}f\),
\(A_{00}=A_{0i}=0\) and
\begin{equation}\protect\phantomsection\label{eq:gg-linear-source-compatible-data}{\begin{aligned}
B_{12}&=2\int_0^{y^1}T_{02}(0,s,y^2,y^3)\,\mathrm ds, \\
B_{13}&=2\int_0^{y^1}T_{03}(0,s,y^2,y^3)\,\mathrm ds, \\
B_{11}&=\int_0^{y^1}
       (2T_{01}-\partial_2B_{12}-\partial_3B_{13})
                        (0,s,y^2,y^3)\,\mathrm ds, \\
B_{22}&=B_{23}=B_{33}=0, \\
B_{0i}&=-\epsilon\partial_if, \\
B_{00}&=0 .
\end{aligned}}\end{equation} Then \(\partial^jB_{ji}=2T_{0i}(0,y)\). For
\(\epsilon\partial_t^2\bar h+\Delta_y\bar h=-2T\) with initial data
\(A,B\),
\begin{equation}\protect\phantomsection\label{eq:gg-linear-source-gauge-data}{
 \beta_b|_0=\epsilon B_{0b}+\partial^iA_{ib}=0,\qquad
 \partial_t\beta_b|_0
       =-\Delta_yA_{0b}-2T_{0b}+\partial^iB_{ib}=0 .
}\end{equation} The scalar equation for \(f\) proves the second equality
for \(b=0\); the spatial divergence of \(B\) proves it for \(b=i\). Also
\(\Box\beta=-2\operatorname{div}T=0\), so analytic uniqueness gives
\(\beta=0\). The linearized Einstein tensor is now \(-\Box\bar h/2=T\),
hence trace reversal yields \(\mathsf R(h)=S\).

The estimates can be proved directly. If \(I_1\) denotes integration in
\(y^1\) with zero initial value, then
\begin{equation}\protect\phantomsection\label{eq:gg-linear-source-poisson-series}{
 f=-2\epsilon\sum_{j\ge0}(-1)^j
       I_1^{2j+2}(\partial_2^2+\partial_3^2)^jT_{00}(0,y).
}\end{equation} On a fixed tangential buffer \(\delta R\) the
differentiated numerator is bounded by
\(2^j(2j)!(\delta R)^{-2j}\|T\|\); the integrations give
\(|y^1|^{2j+2}/(2j+2)!\). A fixed restriction makes the sum geometric
and bounds \(f\) by \(CR^2\|T\|\). Interior derivatives and the finite
integrals give \(\|A\|\le CR^2\|T\|\), \(\|B\|\le CR\|T\|\) on a common
three-polydisk.

The wave solution has homogeneous-data sums
\(t^{2j}(-\epsilon\Delta_y)^jA/(2j)!\) and
\(t^{2j+1}(-\epsilon\Delta_y)^jB/(2j+1)!\), and source sum
\begin{equation}\protect\phantomsection\label{eq:gg-linear-source-wave-series}{
 -2\epsilon\sum_{l,j\ge0}
 \frac{t^{l+2j+2}}{(l+2j+2)!}
              (-\epsilon\Delta_y)^j\partial_t^lT(0,y).
}\end{equation} Cauchy bounds give numerator factors at most
\(C3^j(2j)!l!\) times the fixed-buffer inverse powers of \(R\). Since
\((2j)!l!\le(l+2j+2)!\), sufficiently small fixed \(|t|/R\) makes all
sums geometric. This gives the asserted bound and justifies termwise
differentiation and uniqueness. For a homogeneous polynomial source all
terms have degree \(d+2\) and only finitely many survive. The formulas
use real coefficients for \(\epsilon=-1\) and normally convergent
complex-linear operations, proving the remaining claims. ∎

\protect\phantomsection\label{sec:gg-same-symbol-inverse-budget}

\subsubsection{C.5. The same finite inverse and its analytic
budget}\label{c.5.-the-same-finite-inverse-and-its-analytic-budget}

For a compatible trace-free polynomial \(A\) of degree \(m-2\) with
\(\mathrm d\operatorname{div}A=0\), put
\begin{equation}\protect\phantomsection\label{eq:gg-compatible-schouten-source}{
 j=\frac43\int_0^1x^b(\operatorname{div}A)_b(tx)\,\mathrm dt,
 \qquad P_A=A+\frac14j\eta,\qquad S_A=2P_A+j\eta .
}\end{equation} Closedness and the Euler homotopy give
\(\operatorname{div}A=3\,\mathrm dj/4\), so
\(\operatorname{div}P_A=\mathrm dj\), \(\operatorname{tr}P_A=j\). The
source \(S_A\) satisfies the conservation equation of Lemma C.1. Its
potential \(h=GS_A\) has \(\mathsf P(h)=P_A\), and therefore
\([\mathsf P(h)^\circ]=[A]\) in the Hessian quotient. Remove its entire
free record with the radial vacuum inclusion: \(c=(1-\iota\Pi)J[h]\). It
has zero free record and optical image \([A]\).

If its optical image is zero, then
\(\mathsf P(h)^\circ=\operatorname{Hess}_0f\). Contracted linear Bianchi
makes the remaining pure trace in \(\mathsf P(h)-\partial^2f\) constant.
In degree \(m-2>0\) it vanishes; in degree zero a multiple of \(\eta\)
is itself a Hessian of a multiple of \(q\). A conformal change therefore
makes the same class Ricci-flat. Equation
(\ref{eq:gg-contact-cotton-kernel}) then says that the zero-record class
is zero. Thus the class just constructed depends only on \([A]\) and is
the unique inverse \(L_m\) of \(\mathfrak o_m\) on \(\ker\Pi\). This
proves the kernel and image assertions used in Appendix D. Surjectivity
of \(d_m[A]=\mathrm d\,\operatorname{div}A\) onto closed polynomial
two-forms of coefficient degree \(m-4\) is proved separately in §D.5;
for \(m<4\) this target is zero. No analytic bound for an inverse of
\(d_m\) is used in the estimate below.

This construction also bounds that same inverse, not an unrelated right
inverse. For an analytic compatible representative \(A(x)\), its scalar
Euler primitive has a common bound after a fixed derivative buffer;
Lemma C.1 gives an analytic \(h\) with the stated radius loss. Its whole
free record has the factorial Cauchy bound, and §C.3 supplies one
analytic vacuum removal. The ray primitive gives a bounded analytic
representative of the resulting contact class. More explicitly take
\(\widehat h=h-v(\Pi J[h])\); in degree \(m=n-1\),
\(b_{\widehat h}=\widehat h_{ab}k^ak^b/(2n)\) is a global quotient
representative with this analytic coefficient bound. Its minimal
quotient representative has no larger norm. The
polynomial-times-exponential bound of the fixed \(\ell_n\) in §C.1 then
sums on another fixed smaller tube. Hence the actual selected cochain
\(\sum_{m\ge2}\ell_{m+1}L_m[A_m]\) is analytic for this one correction.
Likewise \(\sum\ell_{m+1}\iota_{m+1}Z_{m-2}\) is analytic for every
factorially bounded record \(Z\).

The selection rule evaluates the finite matrices \(L_m\) on the optical
defect and uses these fixed cochains; it does not evaluate a nonlinear
Einstein solution or replace the selected maps by the auxiliary
potential. The potential proves consistency and a common spatial
estimate for the same unique class inverse. Repetition over a finite
parameter algebra makes each finite stage analytic. Common convergence
of all normalized stages is the separate coefficient comparison of §6.

\protect\phantomsection\label{app:gg-optical}

\subsection{Appendix D. The intrinsic projective optical bundle and its
obstruction}\label{appendix-d.-the-intrinsic-projective-optical-bundle-and-its-obstruction}

\protect\phantomsection\label{sec:gg-optical-density-line}

\subsubsection{D.1. The density line of compact contact
skies}\label{d.1.-the-density-line-of-compact-contact-skies}

Write \(\mathcal K=\ker\mathrm dq\) on the complete incidence, to
distinguish this line from the fixed spinor bundles. Its projection to
\(T_xM\) is the kernel of first-jet evaluation of
\(H^0(Q_x,\mathcal L)\). In bilinear coordinates on
\(\mathbb P^1\times\mathbb P^1\), a section whose value and two first
derivatives vanish at a point spans a line of type
\(\mathcal O(-1,-1)\). Thus \(\mathcal K|_{Q_x}=\mathcal O(-1,-1)\) as
an isomorphism type. This does not choose its varying constant factor.

The line \(\mathcal N=\mathcal K\otimes q^*\mathcal L\) is trivial on
every compact sky. Since \(H^0(\mathcal O_Q)=\mathbb C\) and
\(H^{>0}(\mathcal O_Q)=0\), proper base change and evaluation give
\begin{equation}\protect\phantomsection\label{eq:gg-intrinsic-density-line}{
 \mathcal B=p_*\mathcal N,\qquad
 p^*\mathcal B\xrightarrow{\sim}
               \mathcal K\otimes q^*\mathcal L .
}\end{equation} The evaluation is an isomorphism on the compact central
fiber and hence on a common nearby incidence. Over a finite local
algebra we use the proper flat compact family, a relatively smooth
moduli germ, and invertible \(\mathcal K,\mathcal L\); these follow from
the residue regularity and the finite construction. Vanishing of
\(H^1(\mathcal O_Q)\) lifts the fiber section through each highest
ideal. Its evaluation remains a unit, proving the same assertion by
induction.

On a small smooth contractible moduli representative take a root pair
\((\mathcal E[1],\beta)\) with \(\beta:\mathcal E[1]^2\to\mathcal B\).
The marked central fiber fixes its sign. Such pairs have a unique
comparison preserving \(\beta\) and that fiber lift: in local frames the
comparison is the square root of a nonzero holomorphic unit with its
prescribed central sign. It exists by a local logarithm; a nilpotent
unit has its finite binomial root. All comparisons commute with finite
parameter restriction. In particular
\begin{equation}\protect\phantomsection\label{eq:gg-intrinsic-root-line}{
 (p^*\mathcal E[1])^2
             =\mathcal K\otimes q^*\mathcal L .
}\end{equation} There is no use of a square root of \(\mathcal O(1,1)\).
The two sky degrees cancel before this root is taken.

The contact variation defines
\(\operatorname{ev}_z:T_xM\to\mathcal L_z\). Its polar hyperplane is
that of the ray line \(\mathcal K_z\) for the complete Segre null
quadric. For \(k\in\mathcal K_z\) and \(v\in T_xM\), evaluation in the
preceding line identifies \(k\otimes\operatorname{ev}_z(v)\) with an
element of \(\mathcal B_x\). Choose temporarily a scalar conformal
representative \(g_0\) with this same Segre null cone and a nonzero
\(b\in\mathcal B_x\). The complete-family and tangent-space
identifications used here are proved in §D.7. The two polar maps have
the same lines and hyperplanes, so
\begin{equation}\protect\phantomsection\label{eq:gg-intrinsic-metric-ratio}{
 \operatorname{ev}_{\mathcal B}^{-1}
       (k\otimes\operatorname{ev}_z(v))
          =a(z)g_0(k,v)b .
}\end{equation} The nonzero ratio \(a(z)\) is holomorphic on the compact
connected sky and is therefore constant. It defines a symmetric
nondegenerate \(\mathcal B\)-valued metric \(\boldsymbol g\) on \(M\),
independent of both temporary choices. Holomorphic dependence follows
from the same finite quadric coefficient equations. In a frame \(e\) of
\(\mathcal E[1]\), write \(\boldsymbol g=g_e e^2\). Replacing \(e\) by
\(\Omega^{-1}e\) gives \(\widehat g_e=\Omega^2g_e\) and
\(\widehat\sigma=\Omega\sigma\). Thus this root is exactly the density
line required in §5.

\protect\phantomsection\label{sec:gg-optical-projective-ray}

\subsubsection{D.2. A projective connection on each
ray}\label{d.2.-a-projective-connection-on-each-ray}

At fixed \(z\in Z\), the tangent spaces to the incident skies form a
Lagrangian curve \(\Lambda_u=T_zQ_{p(f(u))}\subset\mathcal D_z\). Here
\(u\) is any ray coordinate. A constant symplectic frame writes it as a
symmetric graph \(S(u)\), with \(S'\) invertible by regularity. Put
\begin{equation}\protect\phantomsection\label{eq:gg-intrinsic-matrix-schwarzian}{
 \mathcal S_u(S)=
       \tfrac12(S')^{-1}S'''
       -\tfrac34\big((S')^{-1}S''\big)^2,\qquad
 \kappa_u=\tfrac12\operatorname{tr}\mathcal S_u(S).
}\end{equation} We prove invariance without first introducing a
spacetime metric. Locally solve
\begin{equation}\protect\phantomsection\label{eq:gg-intrinsic-jacobi-normalization}{
 Y'=-\tfrac12YS''(S')^{-1},\qquad
 Y(u_0)^{\mathsf T}Y(u_0)=S'(u_0)^{-1},
 \qquad Z=YS .
}\end{equation} The initial factor exists locally over \(\mathbb C\) and
is only a proof witness. Differentiation and ODE uniqueness give
\(Y^{\mathsf T}Y=(S')^{-1}\), with \(Y^{\mathsf T}Y'\) symmetric. Set
\(R=-Y''Y^{-1}\). Differentiating the symmetry proves
\(R^{\mathsf T}=R\). The equation for \(Y\) and
\(S'=Y^{-1}Y^{-{\mathsf T}}\) give
\begin{equation}\protect\phantomsection\label{eq:gg-intrinsic-jacobi-identities}{\begin{gathered}
Y''+RY=0, \\
Z''+RZ=0, \\
Y^{\mathsf T}Z'-(Y')^{\mathsf T}Z=1, \\
Y^{\mathsf T}Y'-(Y')^{\mathsf T}Y=0 .
\end{gathered}}\end{equation} The remaining symplectic identity follows from
\(S=S^{\mathsf T}\). If \(A=(S')^{-1}S''\), then
\(A=-2(Y^{-1}Y')^{\mathsf T}\) and direct differentiation gives
\begin{equation}\protect\phantomsection\label{eq:gg-intrinsic-schwarzian-conjugacy}{
 \mathcal S_u(S)=(Y^{-1}RY)^{\mathsf T}.
}\end{equation} A constant symplectic graph change takes constant linear
combinations of the same normalized fundamental pair \(Y,Z\); it retains
\(R\) and its trace. Scaling the symplectic form by a nonzero raywise
constant also changes no trace. This covers contact frames, whose Levi
form is line-valued.

For \(v=f(u)\), direct chain differentiation, with all derivatives at
the same fixed ray, gives
\begin{equation}\protect\phantomsection\label{eq:gg-intrinsic-projective-covariance}{\begin{gathered}
\mathcal S_v(S)=f'^{-2}
       \left(\mathcal S_u(S)-\tfrac12\{f,u\}1_2\right), \\
\kappa_v=f'^{-2}\left(\kappa_u-\tfrac12\{f,u\}\right),
 \quad
 \{f,u\}=\frac{f'''}{f'}-\frac32\left(\frac{f''}{f'}\right)^2.
\end{gathered}}\end{equation} For verification,
\((S_v')^{-1}S_v''=f'^{-1}A-f''f'^{-2}1_2\). The third derivative
cancels its mixed \(A\) term in
\(\frac12(S_v')^{-1}S_v'''-\frac34A_v^2\); the remaining scalar is the
displayed Schwarzian. No simultaneous diagonalization or commutation of
matrix derivatives is used.

\protect\phantomsection\label{sec:gg-optical-actual-cauchy}

\subsubsection{D.3. The global operator and its transition
frames}\label{d.3.-the-global-operator-and-its-actual-transition-frames}

Let \(l\) be a local frame of \(\mathcal L\) pulled back from \(Z\).
Choose a local frame \(m_u\) of \(p^*\mathcal E[1]\) with
\(m_u^2=\partial_u\otimes l\) under the root-line isomorphism. For
\(v=f(u)\) and \(l_v=h\,l\), with \(h\) constant along rays, the same
root pair gives
\begin{equation}\protect\phantomsection\label{eq:gg-intrinsic-half-density-change}{
 m_v=\sqrt{h/f'}\,m_u,\qquad
 y_v=\sqrt{f'/h}\,y_u .
}\end{equation} These are comparisons in an existing root line, not new
independent roots of \(\mathcal L\). The preceding projective covariance
gives
\begin{equation}\protect\phantomsection\label{eq:gg-intrinsic-operator-change}{
 (\partial_v^2+\kappa_v)y_v
       =h^{-1/2}f'^{-3/2}
                     (\partial_u^2+\kappa_u)y_u.
}\end{equation} Hence the ray equations define one intrinsic operator
\begin{equation}\protect\phantomsection\label{eq:gg-intrinsic-ray-operator}{
 P_J:p^*\mathcal E[1]\longrightarrow
                  p^*\mathcal E[1]\otimes(\mathcal K^*)^2.
}\end{equation} Its solutions descend to a rank-two holomorphic bundle
\(\mathcal E_J\) on the short-ray space. Local transverse slices provide
the two Cauchy values; finite holomorphic ODE uniqueness supplies the
inverse transitions and exact cocycle. Connected simply connected short
fibers ensure that continuation defines one solution. The construction
is natural under contact/incidence comparisons and the retained root
mark.

For solutions \(s_a=y_a m_u\), the Wronskian is
\(W_u(s_1,s_2)=y_1\partial_uy_2-(\partial_uy_1)y_2\). It changes by
\(W_v=h^{-1}W_u\). Thus \(s_1\wedge s_2\mapsto W_u(s_1,s_2)l\) is the
specified, nondegenerate isomorphism
\begin{equation}\protect\phantomsection\label{eq:gg-intrinsic-wronskian-line}{
 \bigwedge^2\mathcal E_J\xrightarrow{\ \sim\ }\mathcal L .
}\end{equation} To obtain equality with the specified matrices in §5,
use its already fixed \(g_{\mathrm{aux}}\) only at this comparison step.
Write \(\boldsymbol g=g_{\mathrm{aux}}b_{\mathrm{aux}}\), and choose
\(e^2=b_{\mathrm{aux}}\) with the central sign. The incidence definition
of \(\boldsymbol g\) says exactly \(D_i\otimes l_i=p^*e^2\). Thus
\(m_i=p^*e\) in the original calibrated affine parameters. The Jacobi
trace comparison gives \(\kappa_i=P^{g_{\mathrm{aux}}}(k_i,k_i)\).

If \(\alpha_i=c_{ij}\alpha_j\), then \(l_i=c_{ij}^{-1}l_j\),
\(D_i=c_{ij}D_j\), and \(\mathrm du_i/\mathrm du_j=c_{ij}^{-1}\). The
half-density change therefore gives \(y_i=y_j\) and multiplication of
its slope by \(c_{ij}\). For the identical slices, moving base point and
hitting time,
\begin{equation}\protect\phantomsection\label{eq:gg-intrinsic-literal-optical-matrices}{
 B^J_{ij}=
       \begin{pmatrix}1&0\\0&c_{ij}\end{pmatrix}
                       F_j(\tau_{ij})=B_{ij}.
}\end{equation} Their inverses also coincide. With an arbitrary other
density frame only the stated \(C_\Omega\) comparison follows; literal
equality is asserted in the preceding specified frame. For
\(\widehat g=\Omega^2g\), \(\mathrm d\widehat u=\Omega^2\mathrm du\) and
the half-density formula gives \(\widehat y=\Omega y\). Its
differentiated first jet is exactly the matrix \(C_{\Omega,i}\) in §5.

For a controlled analytic incidence family, three-jet interior Cauchy
estimates and a common lower bound on the regular graph minor bound
\(\kappa\). The finite two-dimensional ODE and its inverse have common
positive slice domains and bounds. This observation does not provide a
new bound for an arbitrary unnormalized formal contact atlas. Section 6
supplies the required common witness for the specified stored
representatives.

\protect\phantomsection\label{sec:gg-optical-weighted-resolution}

\subsubsection{D.4. The weighted resolution and its intrinsic
obstruction
map}\label{d.4.-the-weighted-resolution-and-its-intrinsic-obstruction-map}

Let \(P_m=\mathcal P_m=\operatorname{Sym}^mV^*\), with both notations
zero for negative degrees, and let \(D=k^a\partial_a\) in the original
flat model. The weight-\(m\) Cauchy state has value degree \(m\) and
slope degree \(m-1\). For \(m\ge2\) its coefficient bundle is
\begin{equation}\protect\phantomsection\label{eq:gg-app-weighted-resolution}{
 0\longrightarrow\mathcal S_m
 \longrightarrow P_m\otimes\mathcal O_Q
 \xrightarrow{D^2}P_{m-2}\otimes\mathcal L^2
 \longrightarrow0 .
}\end{equation} On a direction patch choose \(\ell(k)=1\). For a
degree-\((m-2)\) polynomial \(q(x)\),
\begin{equation}\protect\phantomsection\label{eq:gg-app-double-primitive}{
 (\mathsf I_kq)(x)=
       \int_0^{\ell(x)}(\ell(x)-t)
                    q(x-\ell(x)k+tk)\,\mathrm dt
}\end{equation} is homogeneous of degree \(m\) and satisfies
\(D^2\mathsf I_kq=q\). Its kernel consists precisely of
\(a_m(b)+t\,b_{m-1}(b)\) with \(x=b+tk\). Slice and generator changes
are the flat intrinsic Cauchy transitions, proving the coefficient
identification, including its different value/slope degrees.

The kernel has rank \(\dim P_m-\dim P_{m-2}=(m+1)^2\). The two
right-hand bundles are sums of \(\mathcal O_Q\) and \(\mathcal O(2,2)\)
and have zero higher cohomology. The long exact sequence gives
\begin{equation}\protect\phantomsection\label{eq:gg-app-hessian-quotient}{
 H^1(\mathcal S_m)=\mathcal Q_m,\qquad H^2(\mathcal S_m)=0,
 \qquad
 \mathcal Q_m=
 \frac{P_{m-2}\otimes\operatorname{Sym}^2_0V^*}
      {\operatorname{Hess}_0P_m}.
}\end{equation} The null evaluation identifies \(H^0(\mathcal L^2)\)
with the trace-free symmetric quadratic forms. With
\(\delta u_{ij}=u_j-u_i\), a scalar lift \(\delta u=r\) of an optical
cocycle has agreeing \(D^2u_i=A(k,k)\) and connecting class \([A]\). A
different lift adds \(\operatorname{Hess}_0f\); conversely the local
double primitives of any \(A(k,k)\) realize its class. Thus the two Čech
constructions are inverse with the asserted sign.

If \(\partial^2f=\tau\eta\), symmetry of third derivatives forces
\(\mathrm d\tau=0\). Hence the global section kernels in weights
\(0,1,2\) are \(\mathbb C,V^*,\mathbb Cq\), and vanish in weights at
least three. Here \(q=\eta(x,x)\). Consequently \(\mathcal Q_2=0\); only
for \(m\ge3\) is its dimension \(9\binom{m+1}{3}-\binom{m+3}{3}\).

The intrinsic construction defines \(o_m\) by varying the actual
\(\mathcal E_J\) transitions and the residue section. Its identification
with the fixed symbol in §5 is
\begin{equation}\protect\phantomsection\label{eq:gg-app-intrinsic-contact-derivative}{
 o_mJ[h]=[P(h)^\circ].
}\end{equation} To prove it, choose the derived scalar representative
for comparison. Its highest ray residual on the lower density \(1\) is
\(P(h)(k,k)\). The local solution is exactly
\(1-\mathsf I_k(P(h)(k,k))\) over the highest square-zero ideal. Its
target-minus-source defect is the positive Čech difference of those
double primitives, giving the displayed class. Changing a local section
lift adds a coboundary. Changing a contact representative also gives a
coboundary by naturality of \(\mathcal E_J\). Thus \(o_m\) is the actual
derivative of the independently defined preservation problem; the
Schouten identity verifies that derivative rather than defining it by an
assigned curvature tensor.

At a lower finite stage, first lift its selected Einstein scale. Every
product of a highest coefficient with a positive parameter coefficient
then vanishes, so the same flat symbol applies. All lower moving bases
and endpoint variations are kept before taking this highest coefficient.
The Bianchi identity and the spatially constant lower trace give
\(\mathrm d\,\operatorname{div}a_m=0\) for the extracted trial class.

\protect\phantomsection\label{sec:gg-optical-correction-exactness}

\subsubsection{D.5. Exactness, correction and the numerical
scalar}\label{d.5.-exactness-correction-and-the-numerical-scalar}

The contact comparison and linear vacuum construction of Appendix C
identify \(\mathcal C_{m+1}\) with degree-\(m\) trace-free metric
polynomials modulo conformal Killing gradients, with both full curvature
projections. We verify the remaining exact-symbol assertions here.

The identity
\(\operatorname{div}\operatorname{Hess}_0f=\frac34\mathrm d\Delta f\)
makes \(d_m[A]=\mathrm d\,\operatorname{div}A\) well defined. If
\(\mathrm d\,\operatorname{div}A=0\), choose homogeneous \(j\) with
\(\operatorname{div}A=\frac34\mathrm dj\). For \(m\ge3\) take
\begin{equation}\protect\phantomsection\label{eq:gg-app-compatible-trace}{
 j=\frac4{3(m-2)}x^b\partial^aA_{ab};
}\end{equation} for \(m=2\) choose \(j=0\). Then \(S=2A+\frac32j\eta\)
has \(\operatorname{div}S=\frac12\mathrm d\operatorname{tr}S\). The
explicit conserved-source linear potential in Appendix C gives
homogeneous \(h\) with \(\operatorname{Ric}_{\mathrm{lin}}h=S\). It has
\(P(h)^\circ=A\). This proves \(\operatorname{im}o_m=\ker d_m\) for
these actual symbols.

If \([P(h)^\circ]=0\), choose \(f\) with
\(P(h)^\circ=\operatorname{Hess}_0f\). Then
\(P(h)-\partial^2f=\tau\eta\), and Bianchi gives \(\mathrm d\tau=0\).
Homogeneity makes \(\tau=0\) unless \(m=2\), where replacing \(f\) by
\(f+\frac12\tau q\) removes it. The metric polynomial \(h+2f\eta\) is
Ricci-flat and has the same conformal contact class. Its two complete
vacuum records give exactly the image of \(\iota\); the zero-record
ambiguity is a symmetric gradient by Appendix C's two-Poincaré
construction.

For surjectivity of \(d_m\), divergence from trace-free symmetric
polynomial tensors onto polynomial one-forms is onto. In a complex
orthogonal basis its Fischer adjoint is trace-free multiplication by
\(x\). If \(\operatorname{sym}_0(x\otimes v)=0\), contraction once and
twice gives
\begin{equation}\protect\phantomsection\label{eq:gg-app-divergence-adjoint-injectivity}{
 qv+\tfrac12(x\cdot v)x=0,\qquad
 \tfrac32q(x\cdot v)=0.
}\end{equation} The polynomial ring has no zero divisors, so \(v=0\);
finite-dimensional duality proves surjectivity. Exterior differentiation
is onto closed polynomial two-forms by the Euler homotopy
\(\omega=\mathrm d(\iota_E\omega/(j+2))\) for coefficient degree \(j\).
This establishes the full exact sequence in §5.

The fixed \(L_m\) is therefore the inverse of the derivative of the
intrinsic section obstruction on the zero-record complement. The
correction \(-L_ma_m+\iota Z_{m-2}\) cancels that actual obstruction and
the entire record defect. Recompute the incidence, projective potential
and optical matrices afterwards. For its exact residual solve
\(\delta u=r'\) and write \(D^2u=A(k,k)\) with
\(A=\operatorname{Hess}_0f\). Then \(u_i-f\) are local flat-ray
solutions with the same difference. Their Cauchy states cancel the
defect. For \(m\ge3\) an explicit primitive is
\begin{equation}\protect\phantomsection\label{eq:gg-app-exact-section-primitive}{
 j=\frac4{3(m-2)}x^b\partial^aA_{ab},\qquad
 P'=A+\tfrac14j\eta,\qquad
 f=\frac{x^ax^bP'_{ab}}{m(m-1)}.
}\end{equation} For \(m=2\) take \(\frac12A_{ab}x^ax^b\) and retain its
free multiple of \(q\); weights zero and one retain a constant and a
linear function. These are exactly the six flat whole-neighborhood
section freedoms. Central value and gradient remove five. With them
fixed, adding \(\gamma q\) changes the numerical scalar by \(12\gamma\),
since \(\Delta q=8\) and the central prolongation is the one in §5. Thus
\(\gamma=(\lambda_{\mathrm{target}}-\lambda_{\mathrm{actual}})/12\)
reaches the target with no division by \(\Lambda\). The actual scalar
comes from this same section's jet.

The invariance of the highest full Weyl record under this final scale
correction follows because its linear conformal variation is zero; all
action on lower nonzero records is a product in \(\mathfrak m_BJ\). This
argument applies also to their differentiated slots.

A closed change \(v\) of contact trial changes \((a,Z)\) to
\((a+ov,Z-\Pi v)\). The identities \(Lo=1-\iota\Pi\) and \(\Pi\iota=1\)
make the selected correction change by \(-v\). This proves the trial
independence claimed in §5 as a correction of the intrinsic
contact-plus-section problem.

\protect\phantomsection\label{sec:gg-optical-normalized-comparisons}

\subsubsection{D.6. Normalized formal
comparisons}\label{d.6.-normalized-formal-comparisons}

There is a direct Čech proof of uniqueness at the relevant highest
degrees. For \(n\ge3\),
\begin{equation}\protect\phantomsection\label{eq:gg-app-no-high-contact-corners}{
 H^0(Q,\mathcal F_n)=0.
}\end{equation} A global element of the middle contact bundle is
\(k\cdot\xi(x)\), with \(\xi\) homogeneous of degree \(n\). Its null
derivative vanishes exactly when
\begin{equation}\protect\phantomsection\label{eq:gg-app-conformal-killing-equation}{
 \partial_a\xi_b+\partial_b\xi_a=2\phi\eta_{ab},\qquad
 \phi=\tfrac14\operatorname{div}\xi.
}\end{equation} The finite-type calculation in §C.2, starting from
(\ref{eq:gg-contact-killing-finite-type}), gives \(\partial^3\xi=0\).
Hence a homogeneous solution of degree \(n\ge3\) vanishes.

Compare two corrected lifts with the same lower object, original record
and numerical scalar, after fixing the same normal central gauge and
actual spin lift. Translation, linear frame/dilation and quadratic
special-conformal freedom have already been fixed. The remaining highest
Hamiltonian corners have \(n\ge3\). Their contact difference \(v\) is
closed, with \(o[v]=0\) because both sections glue and \(\Pi[v]=0\)
because their full records agree. Exactness gives \([v]=0\), so
\(\delta b=v\); the preceding vanishing makes this highest corner
unique. Transport the section by the natural optical comparison. The
remaining difference lies in \(H^0(\mathcal S_m)\) and vanishes by its
kernel calculation and the same value, gradient and scalar
normalizations.

Induction gives compatible formal comparisons. The actual spin lift is
fixed in this statement; the distinct retained central spin arrows are
not collapsed by their identical tensor action. This formal proof
neither proves convergence of arbitrary raw gauges nor replaces the
literal normalized common-domain witness in §6.

\protect\phantomsection\label{sec:gg-affine-compact-lifts}

\subsubsection{D.7. Canonical lifts of the compact
skies}\label{d.7.-canonical-lifts-of-the-compact-skies}

The contact-quadric reconstruction and its deformation framework are
classical (LeBrun 1983, III.2--3 and III.6). We first prove the embedded
deformation statement used for the compact family.

\pandocstatementheading{\protect\phantomsection\label{lem:gg-compact-support-completeness}Lemma D.1.}

Let \(Q\) be a compact smooth complex submanifold of a separated
finite-dimensional complex manifold \(Y\). If \(H^1(Q,N_{Q/Y})=0\), then
its germ of embedded compact deformations is smooth of dimension
\(h^0(Q,N_{Q/Y})\), with tangent space \(H^0(Q,N_{Q/Y})\). It has an
actual universal proper smooth embedded family after restriction.

\emph{Proof.} Douady's theorem represents proper flat closed analytic
subspaces of \(S\times Y\) on finite-dimensional analytic bases,
including nonreduced bases, by a locally finite-dimensional analytic
space \(H(Y)\). Its tangent at \([Q]\) is
\(\operatorname{Hom}_Y(\mathcal I_Q,\mathcal O_Q)=H^0(Q,N_{Q/Y})\)
{[}Douady (1966), §9.1, Remarks 2--3; §9.7, Theorem 1; §9.8, Theorem
2{]}. For finite-dimensional bases the anaplat condition in the
construction is exactly coherence and ordinary base flatness (Douady
1966, sec. 8.3, Proposition 1). We apply this commutative-support
theorem only to the fixed ordinary ambient \(Y\). For a smooth embedded
compact submanifold with \(H^1(N)=0\), its unobstructed complete-family
consequence is also the classical completeness theorem (Kodaira 1962);
the small-extension argument below gives the form needed here.

We first prove the needed infinitesimal lifting directly. Let
\(A'\to A\) be a surjection of finite local Artin complex algebras with
kernel \(J\) and \(\mathfrak m_{A'}J=0\). A proper flat embedded
deformation \(Q_A\) of \(Q\) is locally a graph \(z=f_i(q)\) in
coordinates with \(Q=\{z=0\}\). Indeed flatness makes reduction of its
ideal exactly \(\mathcal I_Q\): the Tor sequence injects that reduction
into \(\mathcal O_Y\). Lift the central normal coordinate generators,
use Nakayama to generate the ideal, and solve the equations by finite
nilpotent recursion with invertible normal Jacobian. This gives the
graph and an \(A\)-flat quotient \(A\otimes_{\mathbb C}\mathcal O_q\).

Lift the finitely many coefficient functions of each \(f_i\) linearly to
\(A'\). On an overlap let \(T_{ij}\) be graph substitution and
\(n_{ij}\) the central normal-bundle transition. The failure
\(c_{ij}=f_i-T_{ij}(f_j)\) lies in \(N_{Q/Y}\otimes J\) and satisfies
\begin{equation}\protect\phantomsection\label{eq:gg-compact-support-small-extension}{
 c_{ik}=c_{ij}+n_{ij}c_{jk}.
}\end{equation} Terms involving positive base coefficients or graph
derivatives times \(J\) vanish since \(\mathfrak m_{A'}J=0\). Hence the
normal linearization and this cocycle equation are exact. Since
\(H^1(N_{Q/Y}\otimes J)=0\), after refinement write
\(c_{ij}=b_i-n_{ij}b_j\) and replace \(f_i\) by \(f_i-b_i\). The graphs
now glue exactly; \(J^2=0\) leaves no quadratic error. Their quotients
are coherent and \(A'\)-flat. Their underlying support is the compact
\(Q\), so the resulting embedded deformation is proper. This gives every
small-extension lift.

Put \(n=h^0(N_{Q/Y})\). A minimal analytic presentation of the
represented germ has local ring \(R=\mathbb C\{x_1,\ldots,x_n\}/I\),
\(I\subset\mathfrak m^2\); linear relations have been eliminated by the
ordinary analytic inverse theorem. If \(I\ne0\), choose a nonzero
relation \(f\) of least order \(d\ge2\) among all relations. Since
\(I\subset\mathfrak m^d\), the map \(R\to A=\mathbb C[x]/\mathfrak m^d\)
exists. Representability and the graph lift just proved give a lift to
\(A'=\mathbb C[x]/\mathfrak m^{d+1}\). Its images of \(x_i\) are
\(x_i+u_i\) with \(u_i\in\mathfrak m^d/\mathfrak m^{d+1}\). But
\begin{equation}\protect\phantomsection\label{eq:gg-compact-support-least-relation}{
 f(x+u)\equiv f_d(x)\not=0
                  \pmod{\mathfrak m^{d+1}},
}\end{equation} because \(2d-1\ge d+1\). This contradicts \(f=0\) in
\(R\). Thus \(I=0\) and the actual analytic germ is smooth.

Finally restrict the universal proper flat family to this smooth germ.
At each central point the same flat ideal reduction and Nakayama give
normal generators; their central normal Jacobian is invertible. The
holomorphic implicit theorem makes the family a smooth graph there. Let
\(V\) be a finite union of these good neighborhoods covering the central
fiber. Properness makes the image of the closed complement of \(V\)
closed and disjoint from the central base point. Removing that image
gives a base neighborhood on which the entire family is smooth.
Universal embeddedness, properness and the tangent identification
persist under this restriction. ∎

The compact family in Definition 5.6 is derived from its Legendrian seed
\(Q\). Its contact normal sequence is
\begin{equation}\protect\phantomsection\label{eq:gg-contact-sky-normal}{
 0\longrightarrow T^*Q\otimes\mathcal L
   \longrightarrow N_{Q/Z}\longrightarrow\mathcal L\longrightarrow0 .
}\end{equation} Here the left term is
\(\mathcal O(-1,1)\oplus\mathcal O(1,-1)\). Its \(H^0,H^1\) vanish,
while \(H^0(\mathcal L)\) has dimension four and \(H^1(\mathcal L)=0\).
Thus \(h^0(N)=4\), \(H^1(N)=0\). Lemma D.1 gives the universal smooth
complete four-dimensional embedded family. The quadric and polarization
stay of type \((\mathbb P^1\times\mathbb P^1,\mathcal O(1,1))\). Indeed
proper base change for \(\mathcal L\) and \(\mathcal L^2\) gives ranks
four and nine. Their complete linear system remains an embedding near
the compact central member: its derivative stays injective, distinct
limiting points remain separated, and a common local immersion
neighborhood excludes collisions of points with the same limit. The
multiplication map from the symmetric square of the four sections onto
the nine sections is surjective at the center and nearby. Its
one-dimensional kernel gives a nondegenerate quadric equation. The
embedded compact surface is that smooth quadric, proving the asserted
polarized type. Every nearby compact member is Legendrian because
\(H^0(\Omega_Q^1\otimes\mathcal L)=0\).

An infinitesimal Legendrian normal field is its contact scalar and its
first derivative. Consequently the complete normal evaluation is the
first-jet evaluation of \(H^0(\mathcal O(1,1))\), which is onto at every
point, as the local sections \(1,\varrho,\zeta,\varrho\zeta\) show. The
universal evaluation \(q:F\to Z\) is therefore a submersion near the
complete central sky, with one-dimensional kernel
\(\mathcal K=\mathcal O(-1,-1)\). Its projectivized contact-scalar
evaluation identifies \(F\) with the null-covector incidence of the
conformal form read from these four sections. The pullback of the
contact form is the tautological null-covector form; its characteristic
line is precisely \(\ker\mathrm dq\). This identifies the rays with the
null geodesics of any representative, without an Einstein equation. We
take short representatives with the buffered regularity used in §§4--5.

The affine-torsor construction below is the metric-cohomology
correspondence of (LeBrun 1983, Appendix, p.230). It selects a metric in
the reconstructed conformal class; the additional projective-extension
condition is treated in §D.8.

On any sky the affine torsor has translation \(\mathcal O(-1,-1)\),
whose \(H^0,H^1,H^2\) all vanish. Its Čech cocycle has exactly one
primitive. These unique sections vary holomorphically: the fixed
two-circle splittings give bounded relative primitives, or proper base
change gives \(p_*q^*\mathcal L^{-1}=R^1p_*q^*\mathcal L^{-1}=0\). The
zero kernel makes local primitives agree. Over a highest finite
extension the same \(H^1\) gives existence and \(H^0\) uniqueness. We
obtain the canonical \(j_A:F\to A\).

For a lifted sky,
\begin{equation}\protect\phantomsection\label{eq:gg-lifted-sky-normal}{
 0\longrightarrow\mathcal O(-1,-1)
 \longrightarrow N_{\widehat Q/A}
 \longrightarrow N_{Q/Z}\longrightarrow0 .
}\end{equation} The complete lifted family has the same four-dimensional
tangent space, because projection is an isomorphism on \(H^0\) and the
middle \(H^1\) vanishes. The finite inverse theorem identifies it with
the same \(M\). At a point of \(Q\), the kernel of complete normal
evaluation to \(N_{Q/Z}\) is \(\mathcal K\). Its lifted evaluation into
the vertical line is exactly
\begin{equation}\protect\phantomsection\label{eq:gg-affine-vertical-derivative}{
 \Theta=\mathrm d_qj_A:\mathcal K\longrightarrow q^*\mathcal L^{-1}.
}\end{equation} Thus \(\Theta\) is a unit if and only if complete
evaluation \(T_xM\otimes\mathcal O_Q\to N_{\widehat Q/A}\) is an
isomorphism, if and only if \(N_{\widehat Q/A}=\mathcal O_Q^4\). Central
triviality gives this on a common neighborhood by compactness. It also
makes \(j_A\) a local biholomorphism. Completeness alone does not do so:
for the trivial torsor the unique lifts all lie in its zero section and
\(\Theta=0\).

\protect\phantomsection\label{sec:gg-affine-projective-extension}

\subsubsection{D.8. The holomorphic extension
condition}\label{d.8.-the-holomorphic-extension-condition}

Evaluation of ray solutions defines
\(\iota:F\to\mathbb P(\mathcal E_J^*)\). Locally its fiber coordinate is
\(w=y_2/y_1\) and \(w'=W(y_1,y_2)/y_1^2\), a unit wherever \(y_1\ne0\).
It is a local biholomorphism over \(Z\). Uniform short injectivity is
explicit: normalized solutions with initial data \((1,0)\) and \((0,1)\)
obey \(\|y_1-1\|\le(K\ell^2/2)/(1-K\ell^2/2)\le1/19\) when
\(K\ell^2\le1/10\). Hence \(|w'-1|<1/2\) on that convex disc. Segment
integration proves injectivity. The finite parent cover gives the short
evaluation embedding used in Definition 5.6; no long-geodesic
compactification is asserted.

Let \(\tau\) be the scalar coordinate of \(j_A\) in an affine frame.
Since affine changes along each ray have constant coefficients, the
equality
\begin{equation}\protect\phantomsection\label{eq:gg-affine-schwarzian}{
 \{\tau,u\}=2\kappa_u
}\end{equation} is frame independent. By the Schwarzian chain rule it is
equivalent to \(\{\tau,w\}=0\). For clarity its fractional-linear
conclusion has a root-free proof. Put \(v=\tau_w\), \(h=v'/v\),
\(v_0=\tau_w(0)\) and \(d_0=\tau_{ww}(0)/(2\tau_w(0))\). The zero
Schwarzian equation is \(h'=h^2/2\). Its unique local holomorphic
solution with \(h(0)=2d_0\) is \(h=2d_0/(1-d_0w)\). Then
\(v=v_0/(1-d_0w)^2\), and integration with the value \(\tau(0)\) gives
\begin{equation}\protect\phantomsection\label{eq:gg-affine-projective-jet-formula}{
 \tau(w)=\tau(0)+\frac{v_0w}{1-d_0w},\qquad
 \begin{pmatrix}
 v_0-\tau(0)d_0&\tau(0)\\-d_0&1
 \end{pmatrix}
 \quad\text{has determinant }v_0 .
}\end{equation} This single formula also works when \(d_0\) is zero or
nilpotent. The ODE uniqueness is coefficientwise over a finite parameter
algebra. The formula uses only the unit \(v_0\), has holomorphic
base-dependent coefficients, and defines a projective isomorphism of the
full fiber. Its inverse is \(w=(\tau-\tau(0))/(v_0+d_0(\tau-\tau(0)))\),
with unit denominator at the slice. On overlapping base charts the
projective maps agree on the open physical ray and therefore everywhere.
They glue to the unique \(\overline J\). Conversely a fiberwise
projective extension has zero Schwarzian in \(w\) and proves the
displayed condition. Nonzero derivative on a short open set alone would
not imply degree one; the isomorphism requirement in the definition is
essential.

By §D.1, \(\Theta^{-1}\) is a section of
\(\mathcal K\otimes q^*\mathcal L=p^*\mathcal B\). Its coefficients are
constant on compact connected skies and descend to a unit \(b\) on
\(M\). Take the unique root \(\sigma\) with \(\beta(\sigma^2)=b\) and
its specified compatible central value. In a ray frame its coefficient
is \((\partial_u\tau)^{-1/2}\). Direct differentiation gives
\(\sigma_u''/\sigma_u=-\{\tau,u\}/2\), so the extension is exactly
\(P_J\sigma=0\). Its Cauchy values define the actual whole section \(s\)
of \(\mathcal E_J\). Proposition 5.4 now gives the Einstein scale and
the same numerical invariant. No metric was supplied to define the
extension.

Conversely let \(s\) be a whole section with unit evaluated density.
Define
\begin{equation}\protect\phantomsection\label{eq:gg-affine-from-section}{\begin{gathered}
A_s=\{\lambda\in\mathcal E_J^*:\lambda(s)=1\}, \\
j_s(f)=\epsilon_f/\sigma(p(f)), \\
\alpha_\nu(v)=\nu(s\wedge v),\quad \nu\in\mathcal L^{-1}.
\end{gathered}}\end{equation} The last formula fixes the translation identification
\(\operatorname{Ann}(s)=\mathcal L^{-1}\) by the Wronskian. With
\(W(s,t)=l\), its coordinate is \(\tau=t/\sigma\) and
\(\partial_u\tau=\sigma_u^{-2}\). Thus its lifted normal bundle is
trivial, and its completion is the same \(\mathbb P(\mathcal E_J^*)\)
with infinity \(\lambda(s)=0\).

For the inverse comparison, choose any transverse slice. Send its value
\(j_A\) to \(\epsilon/\sigma\) with identity translation map under the
specified \(\mathcal L^{-1}\) identifications. The two maps have the
same derivative \(\Theta\) along a ray and agree on the slice, hence
agree on each connected short fiber. The comparison is slice independent
and glues exactly. Unique compact lifts, the compatible marked root and
unique projective extension make these comparisons natural on the full
admitted domains. The actual spin lifts remain separate arrow data.

\protect\phantomsection\label{sec:gg-actual-affine-cocycles}

\subsubsection{D.9. Actual affine transitions and their
obstruction}\label{d.9.-actual-affine-transitions-and-their-obstruction}

For the local section states let
\begin{equation}\protect\phantomsection\label{eq:gg-affine-adapted-cauchy}{
 G_i=
 \begin{pmatrix}\sigma_i&0\\\dot\sigma_i&\sigma_i^{-1}\end{pmatrix},
 \qquad
 T_{ij}=(G_i\circ\Phi_{ij})^{-1}B_{ij}G_j .
}\end{equation} The first column is the selected solution; the second is
the companion with Wronskian \(l_i\). Its section equation and
\(\det B_{ij}=c_{ij}\) give
\begin{equation}\protect\phantomsection\label{eq:gg-affine-translation-from-source}{
 T_{ij}=\begin{pmatrix}1&b_{ij}\\0&c_{ij}\end{pmatrix},
 \qquad a_{ij}=-b_{ij}/c_{ij}.
}\end{equation} Indeed the transported companion is
\(t_j=b_{ij}s_i+c_{ij}t_i\), so \(\tau_j=b_{ij}+c_{ij}\tau_i\). The
exact matrix product has \(b_{ik}=b_{jk}+(b_{ij}\circ\Phi_{jk})c_{jk}\)
and therefore
\begin{equation}\protect\phantomsection\label{eq:gg-affine-translation-cocycle}{
 a_{ik}=a_{ij}\circ\Phi_{jk}
                  +(c_{ij}^{-1}\circ\Phi_{jk})a_{jk}.
}\end{equation} The inverse contact unit is \(c^{-1}\circ\Phi^{-1}\) and
the inverse translation is \(-(ca)\circ\Phi^{-1}\). Conversely these
affine maps recover \(T\), then the same \(B\) by the same \(G\),
including inverses and all moving arguments.

The specified calibration is stronger than an infinity line. For a
general \(T\) of determinant \(c\) the projective coordinate law is
\(\tau_i=(T_{11}\tau_j-T_{12})/(T_{22}-T_{21}\tau_j)\). Infinity
preservation gives \(T_{21}=0\), but its affine slope is \(T_{11}^2/c\).
The calibrated slope \(c^{-1}\) forces \(T_{11}^2=1\), and the fixed
branch forces \(T_{11}=1\); over finite algebras \(T_{11}+1\) is a unit.
This preserves the section.

For a trial highest target-minus-source defect \(r_{ij}\) the same
matrices give
\begin{equation}\protect\phantomsection\label{eq:gg-affine-obstruction-cochain}{
 (G_i\circ\Phi_{ij})^{-1}r_{ij}=e_1-T_{ij}e_1 .
}\end{equation} Its entries are exactly \(1-T_{11}\) and \(-T_{21}\).
Consequently the connecting class and symbol of §D.4 are the obstruction
and derivative of this calibrated affine reduction itself. The
correction of Theorem 5.5, followed by fresh incidence, fresh \(B\) and
the exact scalar correction, makes this cochain zero. There is no
independent affine correction or additional free translation.

Finally let \(|\sigma_i|\ge u_0\), \(|\sigma_i|\le M_0\),
\(|\dot\sigma_i|\le M_1\) on retained parents, and put
\(H=C(M_0+M_1+u_0^{-1})\) in a fixed submultiplicative matrix norm. Then
\(\|G_i^{\pm1}\|\le H\), \(\|T_{ij}^{\pm1}\|\le H^2\|B_{ij}^{\pm1}\|\)
and \(|a_{ij}|\le\|c_{ij}^{-1}\|H^2\|B_{ij}\|\). Nested source and
target fiber discs supply the actual inverse domains. The jet formula in
§D.8 requires a common lower bound on \(v_0\) and common parent jet
bounds for uniform independent-family estimates; pointwise invertibility
is insufficient. The generated families obtain these margins from §6.
All these finite formulas commute with truncation, root restriction on
the marked branch, and Lorentz conjugation. They divide by neither a
Weyl invariant nor \(\Lambda\).

\protect\phantomsection\label{app:gg-curved-penrose}

\subsection{Appendix E. A fixed representative inverse for the curved
Penrose
transform}\label{appendix-e.-a-fixed-representative-inverse-for-the-curved-penrose-transform}

Let \(g_+\) be a selected analytic Einstein background with vanishing
negative Weyl spinor. Write \(p_+:F_+=\mathbb P(S')\to M\) and
\(q_+:F_+\to Z_+\) for its complete-line correspondence. The
\(q_+\)-fibers are the short integral surfaces of
\(\pi^{A'}\nabla_{AA'}\). The degree-one line is the same calibrated
\(L\) as in §7. All frames below are induced by this selected physical
scale and the original spin mark.

Denote the curved negative contour by \(P_Z\). In these frames its
components are the fiber residues of
\(\pi_{A'}\pi_{B'}\pi_{C'}\pi_{D'}q_+^*b\). The relative flat line
connection defines a complex
\(\mathcal R^j=q_+^*L^{-6}\otimes\Omega^j_{F_+/Z_+}\). On a complete
\(p_+\)-fiber its three terms have types \(\mathcal O(-6)\),
\(S^*\otimes\mathcal O(-5)\) and \(\mathcal O(-4)\). Its degree-one
fiber-cohomology equation is precisely
\(\nabla^{AA'}\psi_{A'B'C'D'}=0\). This follows by applying the relative
derivative \(\pi^{A'}\nabla_{AA'}\) to the contour and contracting the
alternating primed indices; equivalently the four surviving
\(\mathcal O(-5)\) Laurent terms in each unprimed component are that
divergence.

\pandocstatementheading{\protect\phantomsection\label{thm:gg-fixed-curved-penrose-inverse}Theorem E.1.}

On the germ of the whole marked line there is a fixed linear
representative map \(K_Z\) from analytic negative massless fields to
actual ordinary \(L^{-6}\) cocycles, with
\begin{equation}\protect\phantomsection\label{eq:gg-curved-penrose-actual-inverses}{
 P_ZK_Z\psi=\psi,\qquad
 K_ZP_Zb-b=\delta v
}\end{equation} on common smaller representatives. The second equality
has an actual bounded holomorphic ordinary zero-cochain \(v\). For
conditioned bounded background families and their retained parent
charts,
\begin{equation}\protect\phantomsection\label{eq:gg-curved-penrose-representative-bound}{
 \|K_Z\psi\|\le C(\|\psi\|+\|\nabla\psi\|).
}\end{equation} A single choice of spatial and angular domains works at
all jet orders. The map preserves spatial/line-normal ideal filtrations,
holomorphic parameters and finite Taylor reductions. No inverse power of
the cosmological constant is used.

\emph{Proof.} We construct the domains and the operator together.
Supplied ordinary cocycles are first restricted to common parent charts
containing all complete lines over a sufficiently small source ball. The
choices below are made inside those parents. Uniform assertions require
the stated background, inverse-frame and parent margins; they do not
concern arbitrary unbounded coordinate changes.

\protect\phantomsection\label{sec:gg-penrose-complete-fiber-splitting}

\subsubsection{E.1. Complete fiber
splitting}\label{e.1.-complete-fiber-splitting}

Choose geodesic normal coordinates \(x=(u^A,v^A)\) and radial spin
frames. The coframe is the identity and the connection is zero at the
mark. On a common source ball \(B_R\) take complete fiber charts
\(|\zeta|<4\), \(|\xi|<4\), \(\xi=\zeta^{-1}\), with the original lower
spinor \(\pi=(1,\zeta)\) in the north frame. With \(u^A=x^{A0'}\) and
\(v^A=x^{A1'}\), the flat relative fields are
\(\pi^{A'}\partial_{AA'}=-\zeta\partial_{u^A}+\partial_{v^A}\). Their
\(dv\)-normalized graphs are exactly the north graphs in
eq.~\ref{eq:gg-penrose-flat-leaf-graphs}. In the south lift
\(\pi=(\xi,1)\), multiplication by \(-1\) gives the \(du\)-normalized
fields \(\partial_{u^A}-\xi\partial_{v^A}\) and the displayed south
graphs. Thus the leaf formula and the contour differential use the same
index convention eq.~\ref{eq:gg-spin-index-convention}. A massless field
has the fiber representative
\begin{equation}\protect\phantomsection\label{eq:gg-penrose-five-term-lift}{
 h(x,\zeta)=\sum_{j=0}^4\psi_j(x)\zeta^{-j-1}.
}\end{equation} The components are tensor components, not coefficients
of a binomially expanded spin polynomial. The relative derivative
\(d_\alpha h\) has zero fiber class by the massless equation. Split it
with the fixed Laurent circles \(|\zeta|=7/2\) and \(|\zeta|=2/7\),
using the ordinary north/south tail allocation, to obtain a zero-cochain
\(k\) in \(\mathcal R^1\) satisfying \(\delta k=d_\alpha h\). The output
angular radii can be taken to be three. A spatial Cauchy restriction
pays the one derivative, giving \(\|k\|\le C(\|\psi\|+\|\nabla\psi\|)\).

Now \(\delta(d_\alpha k)=d_\alpha^2h=0\). On every complete
\(p_+\)-fiber, \(d_\alpha k\) is therefore a global section of
\(\mathcal O(-4)\) and is zero. This proves
\begin{equation}\protect\phantomsection\label{eq:gg-penrose-complete-fiber-identities}{
 \delta k=d_\alpha h,\qquad d_\alpha k=0 .
}\end{equation} The vanishing is applied before any leaf refinement. It
would not follow from negative-degree cohomology on a disconnected or
incomplete part of a fiber.

\protect\phantomsection\label{sec:gg-penrose-whole-leaf-domains}

\subsubsection{E.2. Whole short fibers and quotient
overlaps}\label{e.2.-whole-short-fibers-and-actual-quotient-overlaps}

In the north chart use the transversal \(v=0\) and leaf coordinate
\(\lambda=v\). Normalize the two alpha fields by the inverse of their
\(dv\) matrix. It has its fixed unit value at the marked model. The
normalized fields \(E_A\) satisfy \(dv^B(E_A)=\delta_A^B\). Their
bracket is in the alpha distribution and has zero \(dv\) projection, so
\([E_0,E_1]=0\). Their initialized holomorphic flows give coordinates
\((\lambda;w,\zeta_0)\), where \((w,\zeta_0)\) is the actual hit on the
transversal. In the south chart use \(u=0\), \(\lambda=u\), with the
analogous normalization. At the flat model the spacetime graphs are
\begin{equation}\protect\phantomsection\label{eq:gg-penrose-flat-leaf-graphs}{
 X_N(\lambda;w,\zeta_0)=(w-\zeta_0\lambda,\lambda),\qquad
 X_S(\lambda;w,\xi_0)=(\lambda,w-\xi_0\lambda).
}\end{equation} The actual graphs have the same free coordinate block.
Finite ODE and implicit-inverse estimates give common bounds on their
first and second leaf derivatives. Fix these bounds on larger parents
before choosing smaller radii.

Here are the domain choices and the whole-fiber verification. Fix a
positive angular-footprint separation \(\eta\) between the closed
quotient direction-two patch and the complement of the corresponding
direction-three \(p_+\) patch on the compact central line. First choose
\(r\ll R\) so the flow parents stay in the source, the connection and
inverse-frame margins hold, and \(C_0C_1r<1/2\), \(C_qr<\eta/4\),
\(C_\pi r^2<1/10\). Next choose \(\delta\ll r\), with
\(C_\perp\delta<\eta/4\). The quotient patches \(V_i\) have normal
radius \(\delta\) and direction radius two. Their fixed local sections
\(s_i(z)=X_i(0;z)\) lie in \(B_{r/10}\). The leaf graphs extend to
\(|\lambda|<2r\) with a direction buffer up to three, and remain in the
original spatial/angular source.

For the auxiliary Hermitian coordinate norm, the free block gives
\(|D_\lambda X_i a|\ge|a|\). On the whole convex parent
\(|\lambda|<2r\), \(|w|<\delta\), the chosen bounds give
\(|X_i|\le C_0r\) and \(|D_\lambda^2X_i|\le C_1\). Thus
\begin{equation}\protect\phantomsection\label{eq:gg-penrose-whole-parent-convexity}{\begin{aligned}
\frac12\,\operatorname{Hess}_{\mathbb R}|X_i|^2(a,a)&=|D_\lambda X_i a|^2+
       \operatorname{Re}\langle X_i,D_\lambda^2X_i[a,a]\rangle \\
&\ge(1-C_0C_1r)|a|^2 \\
&\ge\frac12|a|^2 .
\end{aligned}}\end{equation} The sublevel \(|X_i|<r\) is convex, connected and
contains \(\lambda=0\). It lies in \(|\lambda|<r\) because \(\lambda\)
is an actual coordinate block. It cannot meet the parent leaf-chart
boundary.

This graph is the whole small-ball fiber, not merely one component. On
the complete central line, \(q_+(0,\pi)\) is its fixed line
identification. Uniform finite flow bounds give
\(|q_+(x,\pi)-q_+(0,\pi)|\le C_q|x|\) in the finite auxiliary charts.
The angular separation and the chosen \(\delta,r\) therefore imply
\begin{equation}\protect\phantomsection\label{eq:gg-penrose-angular-containment}{
 q_+^{-1}(V_i)\cap p_+^{-1}B_r
 \subset\{\text{the corresponding }p_+\text{ angle has modulus }<3\}.
}\end{equation} Every point of this preimage can be transported to the
chosen transversal. The connection is \(O(|x|)\) in the radial frame, so
a path of length \(O(r)\) changes its bounded spin ratio by \(O(r^2)\)
and stays inside the angular parent. Its hit has the same quotient image
\(z\). The quotient restricted to that transversal has a common unit
Jacobian and an injective neighborhood, so the hit is uniquely
\(s_i(z)\). Flow uniqueness identifies the point with \(X_i(\lambda;z)\)
for its actual free block. The displayed convex sublevel therefore
describes the full fiber.

Put \(\Omega_i=q_+^{-1}(V_i)\cap p_+^{-1}B_r\). For
\(z\in V_N\cap V_S\), both local sections lie in \(B_{r/10}\) and in the
same full fiber, which the containment argument puts in both middle
angular charts. Hence
\begin{equation}\protect\phantomsection\label{eq:gg-penrose-actual-overlap-surjectivity}{
 q_+(\Omega_N\cap\Omega_S)=V_N\cap V_S .
}\end{equation} The fibers there are connected and are the same whole
short fibers in both descriptions. The quotient transition is the
transverse hit map, with its line factor. No global section of \(q_+\)
is chosen.

Finally choose the output ball \(B_\varepsilon\) only after \(\delta\),
with \(\varepsilon\ll\delta\ll r\ll R\). Every complete projective fiber
has \(|\zeta|\le1\) or \(|\xi|\le1\). Its corresponding hit has normal
displacement at most \(C_{\mathrm{cov}}\varepsilon<\delta/10\) and
angular drift \(O(\varepsilon^2)\). The complete final \(p_+\) cover of
radii \(5/4\) thus lands inside quotient directions \(3/2\), with a
buffer to two. All complete output fibers are retained. The patches
\(V_i\) are coordinate products of a normal ball and a direction disc,
so ordinary two-chart cocycles modulo actual coboundaries compute the
required line-bundle \(H^1\) there.

\protect\phantomsection\label{sec:gg-penrose-typed-primitive}

\subsubsection{E.3. The line-valued primitive and
descent}\label{e.3.-the-line-valued-primitive-and-descent}

In the following, the superscript \(p\) on \(u_i^p\) and \(k_i^p\)
denotes components in the original \(p_+\) line frame. On \(\Omega_i\)
let \(t_i(\lambda,z)\) be the change from that frame to its parallel
frame along the alpha leaf, initialized by \(t_i(0,z)=1\). Solve its
radial linear ODE. Flatness of the relative line connection and
\([E_0,E_1]=0\) make this frame parallel in both leaf directions. The
same finite ODE bounds control \(t_i\) and its inverse on the parents.

In the original frame define
\begin{equation}\protect\phantomsection\label{eq:gg-penrose-typed-line-primitive}{
 u_i^p(\lambda,z)=t_i(\lambda,z)
 \int_0^1t_i(a\lambda,z)^{-1}
          \lambda^A k_i^p(E_A)(a\lambda,z)\,\mathrm da .
}\end{equation} The integrand contracts the relative one-form with the
coordinate-normalized vector fields \(E_A\). It is not an unconverted
spin-frame component. The path stays in the convex whole fiber. In the
parallel frame, \(d_\alpha k_i=0\) is the ordinary closed-one-form
identity, so the radial homotopy gives \(d_\alpha u_i=k_i\), with zero
value at the chosen transversal. The bound is \(\|u_i\|\le Cr\|k_i\|\).

The overlap section \(h^{\mathrm{bas}}=h-\delta u\) is now horizontal:
\(d_\alpha h^{\mathrm{bas}}=0\). The actual overlap surjectivity and
connected fibers proved above make it descend uniquely to an ordinary
cocycle on \(V_N\cap V_S\). Its quotient line frames are the frames
induced at the two actual local sections; their transition is computed
by these same parallel factors and the original \(p_+\) transition.
Evaluation on either section proves holomorphy and the stated bound.
Define that actual representative to be \(K_Z\psi\). On each complete
final \(p_+\)-fiber it differs from \(h\) by an actual coboundary,
proving \(P_ZK_Z\psi=\psi\).

\protect\phantomsection\label{sec:gg-penrose-actual-coboundary}

\subsubsection{E.4. The inverse on arbitrary ordinary
classes}\label{e.4.-the-inverse-on-arbitrary-ordinary-classes}

Suppose an analytic ordinary cocycle \(b\) has zero contour field. Pull
it back on a smaller complete \(p_+\) cover of radii \(3/2\), inside the
quotient direction-two parents. Its zero fiber class has a unique
splitting \(q_+^*b=\delta v\) in \(\mathcal R^0\), since
\(H^0(\mathcal O(-6))=0\). Use fixed circles \(7/5\) and \(5/7\) to
obtain bounded \(v_i\) on the complete radius-\(5/4\) cover.

Horizontality of \(q_+^*b\) makes \(d_\alpha v\) glue on every complete
fiber. Its type is \(S^*\otimes\mathcal O(-5)\), with no global
sections, so \(d_\alpha v=0\) there. Restrict to quotient directions
\(11/10\) and a normal radius sufficiently small relative to the
available moduli ball. The same angular separation, transversal-hit and
whole-fiber argument places the full preimage inside the \(5/4\) \(p_+\)
patch and makes \(v_i\) descend to bounded ordinary zero-cochains. Their
coboundary is \(b\). Apply this result to \(K_ZP_Zb-b\), whose contour
is zero by the first inverse, to obtain the second equality in Theorem
E.1. There are only finitely many fixed restrictions; none is indexed by
a Taylor order.

\protect\phantomsection\label{sec:gg-penrose-filtration-and-families}

\subsubsection{E.5. Filtration and analytic
families}\label{e.5.-filtration-and-analytic-families}

In each leaf chart the spatial ideal of \(p_+^{-1}(o)\) is
\(\mathfrak m=(\lambda,w)\), and \(q_+^*\mathcal I_\ell=(w)\). The
transversals meet the central line at \(\lambda=w=0\). Fixed
complete-fiber Laurent projections and analytic unit frames preserve
spatial order. For \(n\ge1\),
\(d_\alpha(\mathfrak m^n)\subset\mathfrak m^{n-1}\mathcal R^1\). The
primitive in §E.3 takes \(\mathfrak m^n\mathcal R^1\) to
\(\mathfrak m^{n+1}\): its path preserves total order and its prefactor
\(\lambda\) adds one. Consequently \(h\mapsto h-\delta u\) preserves
spatial order. For a horizontal output, its order is exactly line-normal
order, by evaluation at \(\lambda=0\). Thus its line-normal degree \(n\)
uses only finitely many field and background jets, with no field Taylor
order above \(n\).

All choices are fixed: parent covers, normalized commuting fields,
transversals, radial initial values, line frames, Laurent circles and
tail allocations. The auxiliary Hermitian inequalities only prove that
their paths and domains are valid. They introduce no conjugation of
holomorphic parameters into the operator. Finite ODE and implicit
inverses, contours and radial integrals therefore depend holomorphically
on the background and field parameters. Uniform parent bounds give a
single set of output domains.

For a finite coefficient algebra use the coefficients of these same
holomorphic formulas, with the same unit residues and initialized ODEs.
Uniqueness gives compatibility with an analytic family's Taylor
reduction. This does not postulate an analytic extension of arbitrary
finite algebra data or a norm uniform in the algebra's dimension.
Parameter derivatives of a supplied analytic family are taken before
Taylor reduction. These facts prove all remaining assertions of Theorem
E.1. ∎

\protect\phantomsection\label{sec:gg-initial-alpha-quotient}

\subsubsection{E.6. Constructing the initial separated
quotient}\label{e.6.-constructing-the-initial-separated-quotient}

The use of \(q_+\) in §§E.1--E.5 is legitimate for an independently
supplied ordinary background. For a source-generated pure chiral metric,
its initial quotient must first be constructed from the distribution
itself. We give that step separately.

\pandocstatementheading{\protect\phantomsection\label{lem:gg-initial-alpha-quotient}Lemma E.2.}

Let an analytic metric germ have an integrable projective alpha
distribution, with the standard marked tangent spin plane at the central
fiber. On sufficiently small common representatives it has a Hausdorff
three-dimensional short-leaf quotient along the whole embedded central
\(\mathbb P^1\). It contains all complete fibers over a smaller source
ball. Its charts and transitions are actual initialized finite flows and
hit maps, uniformly in conditioned bounded metric families. This
assertion does not assume an already existing quotient map.

\emph{Proof.} Take finitely many projective-spin cores covering the
compact central fiber, with successively larger angular buffers. In a
geodesic-normal spin frame the connection is \(O(|x|)\). The alpha
distribution has its fixed injective spacetime projection. On each
angular parent choose two free spacetime coordinates \(y\); their minor
is the model unit and remains invertible after shrinking. Normalize the
fields to
\begin{equation}\protect\phantomsection\label{eq:gg-initial-quotient-commuting-fields}{
 E_A=\partial_{y^A}
        +V_A^\mu(y,t)\partial_{t^\mu},\qquad [E_0,E_1]=0,
}\end{equation} where \(t\) consists of the remaining two spacetime
coordinates and the spin direction. The bracket vanishes because
involutivity puts it in the distribution and its free-coordinate
projection is zero.

The initialized flows give actual graph coordinates \(X(y;t_0)\) and
their inverse hit on \(y=0\). First fix all derivative bounds on larger
spatial and angular parents. Then choose \(r\) small. Any point with
\(|x|<r\) and direction in a core hits \(y=0\) in a still larger spatial
parent with normal transverse value at most \(C_wr\). This follows by
reverse finite flows of free-coordinate length at most \(Cr\). For these
hits the graphs extend throughout the convex parent \(|y|<2r\), have
\(|X|\le C_0r\), and satisfy \(|D_yX a|\ge|a|\), \(|D_y^2X|\le C_1\).
Here the free coordinate is a block of \(X\). Taking \(C_0C_1r<1/2\)
gives the strict convexity estimate
(\ref{eq:gg-penrose-whole-parent-convexity}). Hence a graph's
intersection with \(|X|<r\) is convex and connected. Its free
coordinates lie in \(|y|<r\), away from the parent boundary.

The spin component of \(E_A\) is \(O(r)\) on these spatial parents. Over
any one retained graph its direction moves by at most \(C_\pi r^2\), by
the initialized path of length \(O(r)\). Choose \(r\) so this is smaller
than every angular buffer. A connected leaf inside the working domain
\(U=p_+^{-1}B_r\) cannot escape the graph of its first core while
remaining in \(U\). Along that graph its free coordinates stay within
\(|y|<r\), its transverse hit is constant, and its spin direction stays
inside the larger angular parent. Unique continuation of the same
commuting flows then keeps any path in the leaf on this graph. This does
not estimate the length of an arbitrary winding path: it uses one graph
already defined on the whole convex free-coordinate parent. Conversely
any two graph points in \(B_r\) are joined by its convex free-coordinate
segment. Thus the connected leaf in \(U\) is exactly that whole graph
sublevel.

Define \(R\) to mean membership in the same such connected leaf in
\(U\). It is closed in \(U\times U\). Indeed for a convergent sequence
of related pairs, select a subsequence with the first directions in one
fixed core. The whole related leaves stay in its larger parent by the
preceding no-exit argument. Their two transverse hits are equal. The
buffered hit maps extend at both limiting points, so their hits remain
equal. The whole graph sublevel is connected, and the limiting pair is
again related.

The quotient projection is open: the saturation of an open set is a
union of its images under the local holomorphic leaf flows and is open.
For an unrelated pair, closedness of \(R\) supplies a product
neighborhood disjoint from \(R\); its two open quotient images are
disjoint. The quotient is therefore Hausdorff.

To exhibit its complex charts near the central line, choose a smaller
normal transverse radius \(\delta\ll r\) and buffered angular slices.
Their zero sections lie in \(B_{r/10}\). Every corresponding graph
sublevel is nonempty and connected and remains in its parent, so the
transverse hit gives a bijective local quotient chart. Its projection is
a holomorphic submersion and is open. Overlapping quotient charts are
related by the actual finite hit on the other transversal. Three routes
hit the same point on the same connected short leaf, proving exact
transition and inverse laws with moving arguments. No prior quotient or
its angular coordinates were used.

The central projective fiber is transverse to the alpha distribution.
Two different points of that fiber cannot be related: in a parent graph
both have free spacetime coordinate zero, whose hit is unique. The
compact fiber is consequently an embedded \(\mathbb P^1\) in this
separated threefold. Finally choose \(\varepsilon\ll\delta\). The finite
inverse-hit bound puts every point of every complete \(p_+\)-fiber over
\(B_\varepsilon\) inside these transverse chart images, using the finite
core cover of the projective line. Thus the initial quotient contains
the complete incidence needed in the preceding sections.

All constants came from bounded finite flows, unit free minors and fixed
angular margins. They can be chosen on the retained parameter parents
and do not involve an inverse scalar curvature. This proves the uniform
and holomorphic family assertions. ∎

\protect\phantomsection\label{app:gg-source-operators}

\subsection{Appendix F. Source operators and their first negative-order
units}\label{appendix-f.-source-operators-and-their-first-negative-order-units}

\protect\phantomsection\label{sec:gg-op-local-model}

\subsubsection{F.1. The fixed local operator
model}\label{f.1.-the-fixed-local-operator-model}

Let \(\mathcal E=\mathcal E_{T_0}^{1/2}\) be the analytic
microdifferential algebra on the retained open part of
\(\mathbb P(T^*T_0)\), acting on half-densities of the original
threefold. Its order filtration has
\(\operatorname{gr}_m\mathcal E=\mathcal L^m\). In local coordinates
\(y=(w^0,w^1,\zeta)\) and momenta \(p\), its product and principal
commutator are
\begin{equation}\protect\phantomsection\label{eq:gg-op-symbol-product}{
 P\circ Q=\sum_\alpha\frac1{\alpha!}
                  (\partial_p^\alpha P)(\partial_y^\alpha Q),
 \qquad
 \sigma_{m+n-1}[P,Q]
       =\{\sigma_mP,\sigma_nQ\}_{\mathrm{can}} .
}\end{equation} The product formula is for the total symbol
\(\sigma_{\mathrm{tot}}(P)=\sum_{r\ge0}P_{m-r}\) of an order-\(m\)
operator. On every strict conic compact its homogeneous terms satisfy
\(\|P_{m-r}\|\le MA^rr!\) for suitable \(M,A>0\). These are the
analytic, rather than unrestricted formal, microdifferential operators.
The complex-linear formal adjoint reverses products and satisfies
\(\sigma_m(P^*)=(-1)^m\sigma_m(P)\). It is distinct from Lorentz
conjugation. These are the local models of (Kashiwara 1996, Lemma 1 and
Theorem 2; Polesello 2007, sec. 1).

This model uses the projective reduction of the original twistor
Heisenberg operators. Set \(r=\pi_{0'}\), \(y=(Z_i/r)_{i=1}^3\) and
\(\Upsilon=r\partial_r\) in homogeneous coordinates. The operators
\(y^i\) and \(r\partial_{Z_i}=\partial_{y^i}\) commute with \(\Upsilon\)
and obey \([\partial_{y^i},y^j]=\delta_i^j\). Euler reduction at
\(\Upsilon=\nu\) gives the local twisted differential operators on
\(\mathcal O(\nu)\). Since \(K_{T_0}=\mathcal O(-4)\), the half-density
choice is \(\nu=-2\). Microlocalization gives \(\mathcal E\). This fixed
Euler sector does not retain every generator of the unreduced
four-variable algebra.

For \(\alpha=\mathrm dw^0+\varrho\,\mathrm dw^1+s\,\mathrm d\zeta\),
write \((p_0,p_1,p_2)=t(1,\varrho,s)\). A contact Hamiltonian \(h\)
becomes the degree-one symbol \(H=th\). Choose the local normal symbol
lift \(Q_0(H)\) with only the term \(H\) and put
\(B(H)=(Q_0(H)-Q_0(H)^*)/2\). Then \(B(H)^*=-B(H)\) and
\(\sigma_1 B(H)=H\). Its negative total-symbol terms satisfy the
required analytic bounds: in
\begin{equation}\protect\phantomsection\label{eq:gg-op-adjoint-bound}{
 \sigma_{\mathrm{tot}}(Q_0(H)^*)
   =\sum_\alpha\frac1{\alpha!}
       \partial_p^\alpha\partial_y^\alpha[H(y,-p)],
}\end{equation} the derivatives act on the displayed composite function.
On fixed strict coordinate and conic parents, the term of degree \(1-k\)
is bounded by \(M C^k k!\). Indeed two Cauchy derivative estimates give
\((\alpha!)^2\), division by \(\alpha!\) leaves one factorial, and the
number of multiindices of length \(k\) is polynomial in \(k\). The
constants are uniform for bounded holomorphic families of \(H\). This
estimates an individual skew lift, not an integrated operator flow.

Both original source operations have such lifts in the same model. The
positive one is the divergence-free half-density operator
\(F_{w^1}\partial_{w^0}-F_{w^0}\partial_{w^1}\), with symbol
\(p_0F_{w^1}-p_1F_{w^0}\). For the negative source, set
\(K=w^0p_0+w^1p_1+\zeta p_2\). With \(U^-_{r;a,j}\) as in Theorem 3.1, a
direct formula is
\begin{equation}\protect\phantomsection\label{eq:gg-op-general-negative}{
 \begin{aligned}
 H_-={}&\sum_{r\ge0}\sum_{a=0}^r\sum_{j=0}^{r+4}
 \frac{(-1)^rU^-_{r;a,j}}
 {(r+4-j)!\,j!\,p_0^{r+1-a}p_1^{a+1}}\\
 &\quad{}\times
 \left[-j(-K)^{r+4-j}p_2^{j-1}
       +\zeta(r+4-j)(-K)^{r+3-j}p_2^j\right].
 \end{aligned}
}\end{equation} A summand with zero prefactor is omitted before
evaluating its power. Differentiate the two monomials of the derived
dual class before making the incidence substitution \(\eta^0=-K/p_0\),
\(\eta^1=p_2/p_0\), \(\varrho=p_1/p_0\), and multiply by \(t=p_0\); this
gives the displayed formula. Every summand has degree one. Its analytic
sum is the already convergent Hamiltonian symbol of §5 on the same
smaller sectors. Quantize that analytic symbol as one order-one
operator; no convergence of an increasing-order sum of operator
representatives is needed.

For example, \(F_{-6}=24/\zeta\) gives
\begin{equation}\protect\phantomsection\label{eq:gg-op-negative-example}{
 H_-=-4\zeta K^3/(p_0p_1),\qquad
 Q_-=-4\zeta
 (w^0\partial_{w^0}+w^1\partial_{w^1}
                         +\zeta\partial_\zeta)^3
 \partial_{w^0}^{-1}\partial_{w^1}^{-1}.
}\end{equation} The inverse derivatives exist microlocally on
\(p_0p_1\ne0\). The skew part of \(Q_-\) has the displayed symbol. Thus
Proposition 9.5's nonprojectible field is represented by an operator in
the original twistor coordinates, although it is not a finite
differential operator or an ordinary point transformation. For any two
source symbols, \(\sigma_1[B(h),B(k)]=t\{h,k\}\). This includes their
mixed operation; lower-order quantum terms remain.

\protect\phantomsection\label{sec:gg-op-finite-descent}

\subsubsection{F.2. Coherent finite amplitude
overlaps}\label{f.2.-coherent-finite-amplitude-overlaps}

\pandocstatementheading{\protect\phantomsection\label{thm:gg-op-finite-descent}Theorem F.1.}

At every amplitude order \(N\), the corrected contact overlaps admit
filtered, adjoint-preserving operator comparisons with their actual
classical symbols and coherent normalized inner triple comparisons.
These constructions commute with amplitude truncation.

\emph{Proof.} Work over
\(A_N=\mathbb C[\varepsilon]/(\varepsilon^{N+1})\). Let
\(\varphi^0_{ij}\) be the exact reference half-density substitutions. In
one orientation, take the Hamiltonian logarithm of the actual operator
\(E_{ij}=(\Phi_{ij}^{-1})^*\) relative to this reference, and lift it by
the fixed linear skew section of F.1 to
\(B_{ij}\in\varepsilon F_1\mathcal E\otimes A_N\). Use exact inverse
comparisons in the other orientation. Set
\begin{equation}\protect\phantomsection\label{eq:gg-op-finite-comparisons}{
 U_{ij}=\exp B_{ij},\qquad
 \varphi_{ij}=\operatorname{Ad}(U_{ij})\varphi^0_{ij},
 \qquad
 C_{ijk}=U_{ij}\varphi^0_{ij}(U_{jk})U_{ik}^{-1}.
}\end{equation} The exponentials are finite in \(\varepsilon\). Their
coefficients can have positive order, but \(\operatorname{ad}B_{ij}\)
preserves every filtration piece, since \([F_1,F_m]\subset F_m\). The
associated graded action is the actual homogeneous contact action of
\(E_{ij}\). Skewness gives \(U_{ij}^*=U_{ij}^{-1}\) and makes
\(\varphi_{ij}\) commute with \(*\).

The finite Baker--Campbell--Hausdorff series \(R=\log C_{ijk}\) has
order at most one. Its order-one symbol is the Hamiltonian logarithm of
the classical triple defect and is zero. Thus \(R\) has order at most
zero. Since \(C_{ijk}\) is unitary, \(R^*=-R\); the adjoint fixes
degree-zero symbols, so that symbol is also zero. Therefore
\begin{equation}\protect\phantomsection\label{eq:gg-op-coherence}{
 \begin{gathered}
 C_{ijk}\in1+F_{-1}\mathcal E\otimes A_N,\qquad
 C_{ijk}^*C_{ijk}=1,\\
 \varphi_{ij}\varphi_{jk}
        =\operatorname{Ad}(C_{ijk})\varphi_{ik},\qquad
 C_{ijk}C_{ikl}
        =\varphi_{ij}(C_{jkl})C_{ijl}.
 \end{gathered}
}\end{equation} The last identity follows by expanding the products and
using the exact reference cocycle. These identities are the descent data
for the local module categories. All operations use fixed sections and
finite polynomials in the amplitude ideal, so reduction from \(N+1\) to
\(N\) gives precisely the order-\(N\) construction. ∎

The filtration and its symbol identification are part of this assertion.
Over a nonreduced amplitude algebra, conjugation by the positive-order
nilpotent exponentials is legitimate, but forgetting the filtration
would erase the classical deformation. Faithfulness of an unfiltered
module category is not asserted.

\protect\phantomsection\label{sec:gg-op-first-units}

\subsubsection{F.3. The calibrated torsor and the additional operator
cocycle}\label{f.3.-the-calibrated-torsor-and-the-additional-operator-cocycle}

Inside the units of the truncated order-zero algebra define
\begin{equation}\protect\phantomsection\label{eq:gg-op-first-units}{
 \mathscr U_1=
 \{1+\bar A:\bar A\in F_{-1}\mathcal E/F_{-2}\mathcal E\}
       \subset(F_0\mathcal E/F_{-2}\mathcal E)^\times .
}\end{equation} Here \(F_{-2}\) is an ideal of \(F_0\), not of the whole
operator algebra. Products of two negative-one classes vanish, so
\(\mathscr U_1\) is the additive sheaf \(\mathcal L^{-1}\). Any
normalized order-zero inner change acts trivially on it, by the
commutator estimate. Thus the quantized comparisons induce its strict
classical contact transport despite their inner triple defects.

\pandocstatementheading{\protect\phantomsection\label{prop:gg-op-torsor-transport}Proposition F.2.}

The calibrated affine torsor of Proposition 5.7 is canonically a torsor
of the derived sheaf \(\mathscr U_1\). Its transported contact action,
transition coefficients, inverses and compact lifts agree with the
existing affine data. This is not a strict embedding of the affine Lie
algebra into the raw microdifferential commutator algebra.

\emph{Proof.} The symbol isomorphism identifies the underlying
translation sheaves. If \(\Phi_{ij}^*\alpha_i=c_{ij}\alpha_j\), the
homogeneous scale satisfies \(t_i\circ\Phi_{ij}=c_{ij}^{-1}t_j\).
Transport by \(E_{ij}\) on degree-minus-one symbols therefore has the
coefficient multiplier \(c_{ij}^{-1}\) in the affine point transition.
The same torsor is consequently written
\begin{equation}\protect\phantomsection\label{eq:gg-op-affine-transport}{
 \tau_i=c_{ij}^{-1}\tau_j+a_{ij},\qquad
 a_{ik}=a_{ij}\circ\Phi_{jk}
            +(c_{ij}^{-1}\circ\Phi_{jk})a_{jk}.
}\end{equation} The first equality uses the same local point arguments
as in §5. Since the identification is an isomorphism of sheaves with
their transport, it preserves the actual cochains and their inverses,
rather than only their cohomology classes.

Infinitesimally, the homogeneous lift is
\(\widetilde X_H=X_h-r_ht\partial_t\), whence
\begin{equation}\protect\phantomsection\label{eq:gg-op-unit-action}{
 \{th,t^{-1}a\}=t^{-1}(X_ha+r_ha).
}\end{equation} For affine generators \((h,\chi)\) and \((k,\psi)\) as
in §8.1, this gives the external semidirect bracket
\((\{h,k\},X_h\psi-X_k\chi+r_h\psi-r_k\chi)\) in
eq.~\ref{eq:gg-affine-operator-bracket}, with the existing negative
stored point velocity. It is independent of the skew principal section,
because the difference of two such sections has order at most minus one
and acts trivially on \(\mathscr U_1\). This proves the asserted
transport, not a section of the full operator extension. ∎

The distinction can be tested on the mixed example of Proposition 9.1.
Use Weyl symbols, for which the homogeneous order-one symbols themselves
give skew lifts. If \(\Pi\) is the canonical Poisson bidifferential, the
degree-minus-one term of their commutator beyond the principal bracket
is \(\Pi^3(H,K)/24\). Put
\begin{equation}\protect\phantomsection\label{eq:gg-op-moyal-example}{
 H=-4a^3p_1/\zeta,\qquad
 K=-4\zeta(ap_0+bp_1+\zeta p_2)^3/(p_0p_1).
}\end{equation} Differentiate before restricting to \(b=p_2=0\). The
contributing paired derivatives on that slice are
\begin{equation}\protect\phantomsection\label{eq:gg-op-moyal-derivatives}{
 \begin{aligned}
 H_{\zeta\zeta\zeta}&=24a^3p_1/\zeta^4,&
 K_{p_2p_2p_2}&=-24\zeta^4/(p_0p_1),\\
 H_{p_1\zeta\zeta}&=-8a^3/\zeta^3,&
 K_{bp_2p_2}&=-24\zeta^3/p_0 .
 \end{aligned}
}\end{equation} All other paired third derivatives vanish on this slice:
those of \(H\) containing \(b\) or two momentum derivatives vanish
identically, and the remaining paired derivatives of \(K\) containing
\(p_0\) have a positive power of \(bp_1+\zeta p_2\). Hence
\begin{equation}\protect\phantomsection\label{eq:gg-op-moyal-defect}{
 \left.\frac{\Pi^3(H,K)}{24}\right|_{b=p_2=0}
 =\frac{-H_{\zeta\zeta\zeta}K_{p_2p_2p_2}
          +3H_{p_1\zeta\zeta}K_{bp_2p_2}}{24}
 =48a^3/p_0\ne0 .
}\end{equation} Thus even an allowed skew section need not identify the
raw operator bracket with the external affine bracket. This example does
not prove that every other section has a nontrivial cohomology class.

Likewise an external \(\mathscr U_1\)-torsor must not be confused with
neutralizing the whole algebroid by units. For normalized units
\(g_{ij}\) that stronger problem is
\begin{equation}\protect\phantomsection\label{eq:gg-op-neutralization}{
 g_{ij}\varphi_{ij}(g_{jk})C_{ijk}=g_{ik}.
}\end{equation} Modulo \(F_{-2}\), with \(g_{ij}=1+A_{ij}\) and
\(c^{(-1)}_{ijk}=\sigma_{-1}(C_{ijk}-1)\), it becomes
\(\delta a=-c^{(-1)}\). Equation (\ref{eq:gg-op-affine-transport}) gives
the homogeneous equation \(\delta a=0\) and does not supply such a
primitive. The quadruple identity makes \(c^{(-1)}\) closed; its
vanishing class has not been established here. This additional problem
concerns an invertible local object used to neutralize the algebroid. It
is not required for the existence of the supported simple modules below.

\protect\phantomsection\label{sec:gg-op-fixed-enrichment}

\subsubsection{F.4. Canonical enrichment of a fixed contact-affine
object}\label{f.4.-canonical-enrichment-of-a-fixed-contact-affine-object}

Kashiwara's Theorem 2 gives the canonical module stack of a complex
contact manifold \(Y\). The relevant uniqueness statement is stronger
than unfiltered Morita equivalence: (Polesello 2007, Theorem 3.3) gives
uniqueness up to an equivalence unique up to a unique isomorphism, when
the filtration, full graded-symbol identification and compatible
anti-involution are retained. The anti-involution includes its square
and degree-zero duality coherences. A global operator ring is not
supplied by this statement.

For fixed \(Y\), its commutator recovers the contact Jacobi bracket on
\(\operatorname{gr}_1\), and its first-unit sheaf recovers
\(\mathcal L^{-1}\). Proposition F.2 therefore transports the calibrated
torsor and the independent admission of Definition 5.6 to the structured
quantization. Conversely the full symbols and torsor recover the same
contact-affine object. Structured uniqueness removes the extra
quantization choice; Einstein recognition still uses that holomorphic
admission and Proposition 5.4.

The relative analytic construction is proved in §§F.6--F.9 and Appendix
G. Proposition F.7 compares its finite jets with the auxiliary
construction of §F.2.

\protect\phantomsection\label{sec:gg-op-static-points}

\subsubsection{F.5. Supported point
objects}\label{f.5.-supported-point-objects}

\pandocstatementheading{\protect\phantomsection\label{prop:gg-op-static-points}Proposition F.3.}

On each actual complete sky \(Q\simeq\mathbb P^1\times\mathbb P^1\), the
marked half-form choice gives a unique isomorphism class of rank-one
simple supported module. Framing its rank-one local-system fiber at a
calibrated cone lift of the marked direction removes its scalar
automorphisms.

\emph{Proof.} The contact restriction is
\(\mathcal L|_Q=\mathcal O(1,1)\) and the chosen root is
\(K_Q^{1/2}=\mathcal O(-1,-1)=\mathcal L^{-1}|_Q\). The precise
regular-module classification is (D'Agnolo and Schapira 2007, Theorem
6.1 and Corollary 6.4), which corrects the earlier formulation in
(Kashiwara 1996, Proposition 4). After its half-form twist, local
systems live on the symplectization cone
\(\widehat Q=\mathcal O(-1,-1)^\times\), not merely on \(Q\). The
displayed root neutralizes that twist.

The homotopy sequence of this punctured line bundle contains
\begin{equation}\protect\phantomsection\label{eq:gg-op-sky-cone-homotopy}{
 \pi_2(Q)=\mathbb Z^2
   \longrightarrow\pi_1(\mathbb C^\times)=\mathbb Z
   \longrightarrow\pi_1(\widehat Q)
   \longrightarrow\pi_1(Q)=0.
}\end{equation} The first map is \((a,b)\mapsto-a-b\), by the Chern
class of \(\mathcal O(-1,-1)\), and is onto. Thus \(\widehat Q\) is
connected and simply connected. There is exactly one rank-one
local-system isomorphism class, with automorphisms \(\mathbb C^\times\).
A nonzero framed value at the calibrated cone point leaves only the
identity. The cited equivalence gives the asserted supported object and
its automorphisms. ∎

Simple connectedness of a Legendrian base alone would not suffice:
(D'Agnolo and Schapira 2007, Example 6.5) has a point support with cone
\(\mathbb C^\times\) and nontrivial monodromy choices. The complete
compact sky and its actual cone are essential here. The physical full
spin lift is retained as in §8, rather than inferred from a projective
support.

Sections F.6--F.9 construct the relative operator overlaps and
comparisons. Appendix G constructs the moving modules, relative
half-form calibration and whole-parameter framing directly, rather than
obtaining them pointwise from the static classification.

\protect\phantomsection\label{sec:gg-op-analytic-transport}

\subsubsection{F.6. Analytic transport on the actual classical
tube}\label{f.6.-analytic-transport-on-the-actual-classical-tube}

Quantized Hamiltonian transport is classical; compare (Kashiwara 1986,
Example 7.3.2). We give the coefficient estimate and locality argument
needed for the relative construction. All symbol algebras in this
argument are scalar.

\pandocstatementheading{\protect\phantomsection\label{thm:gg-op-analytic-transport}Theorem F.4.}

Let \(B(t,\mu)\) be a holomorphic order-one microdifferential family
with uniform factorial symbol bounds on fixed larger cotangent-sector
parents. Suppose the homogeneous Hamiltonian \(H=\sigma_1B\) has an
actual holomorphic classical transport \(T_t\), with \(T_0=\mathrm{id}\)
and \(\dot T_t=X_HT_t\), on a common tube with bounded inverse maps and
fixed Cauchy buffers. Then
\(\dot{\mathcal Q}_t=\operatorname{ad}(B(t))\mathcal Q_t\),
\(\mathcal Q_0=\mathrm{id}\), has a filtered quantized-contact sheaf
solution on smaller tube parents. It is jointly holomorphic in all
passive parameters \(\mu\), with uniform factorial negative-symbol
bounds. If \(B^*=-B\), it preserves the formal adjoint. Its
principal-symbol action is exactly \(T_t\).

\emph{Proof.} Normal and Weyl total symbols determine the same analytic
microdifferential class. Their conversion is the formal differential
transform \(\exp(\pm\frac12\sum_i\partial_{y_i}\partial_{p_i})\). On
strict parents, a degree-drop-\(k\) term has two Cauchy factorials
divided by \(\alpha!\), leaving \(\alpha!\le k!\). Combining an input
drop \(r_0\) with \(k\) uses \(r_0!k!\le(r_0+k)!\). Thus this conversion
preserves factorial symbol bounds on smaller parents, uniformly in the
parameters.

Use Weyl symbols. For an input of order \(m\) and the generator write
\begin{equation}\protect\phantomsection\label{eq:gg-op-flow-input-bounds}{
 \|P_{m-r}\|\le M_PA_P^r r!,\qquad
 \|B_{1-l}\|\le M_BA_B^l l!,\qquad r,l\ge0.
}\end{equation} The Weyl commutator contribution involving \(B_{1-l}\)
and an odd total derivative order \(k\ge1\) is a differential operator
\(D_{l,k}\) of order \(k\) on the other symbol. It increases its
homogeneous drop by \(\Delta=l+k-1\). Cauchy estimates on the generator
factor cancel the factorial divisors \(\alpha!\beta!\) of the Weyl
product. After one fixed buffer its coefficients therefore satisfy
\begin{equation}\protect\phantomsection\label{eq:gg-op-flow-differential-bound}{
 \|{\rm coeff}(D_{l,k})\|
       \le M_BA_B^l l!C_B^k .
}\end{equation} The multiindex count is absorbed into \(C_B^k\). The
only zero-drop term is \(D_{0,1}=X_H\).

For a fixed input \(P\), put \(A(t)=\mathcal Q_t(P)\). Set
\(A(t)=T_tV(t)\) and conjugate the remaining operators by \(T_t\). An
analytic change of variables turns an order-\(k\) derivative into
derivatives of orders \(j\le k\) with coefficient bounds
\(C_{\rm flow}^k k!/j!\). To verify this bound, expand \(f(\chi(z+h))\)
in powers of \(\chi(z+h)-\chi(z)\). The \(j\)-th derivative has its
Taylor divisor \(j!\), while the analytic composition coefficients have
an exponential bound times \(k!\); several-variable counts are absorbed
into the exponential. Since \(k!/j!\le2^k(k-j)!\), and \(l+k=\Delta+1\),
the conjugated operators obey
\begin{equation}\protect\phantomsection\label{eq:gg-op-flow-interaction-bound}{
 \begin{gathered}
 \dot V=\sum_{\Delta\ge1}\widetilde D_\Delta(t)V,\qquad
 \operatorname{ord}_{\rm diff}\widetilde D_\Delta\le\Delta+1,\\
 \|\widetilde d_{\Delta,j}\|
   \le M_1C_1^{\Delta+1}(\Delta+1-j)!,
       \qquad 0\le j\le\Delta+1 .
 \end{gathered}
}\end{equation} These constants depend only on the supplied tube and its
buffers, the classical maps and their inverses, and the generator
bounds. They are independent of the input drop and the parameters.

In the time-ordered series, \(n\) steps from input drop \(r_0\) to
output drop \(r\) have positive drops \(\Delta_i\) with sum
\(\Delta=r-r_0\). Thus \(n\le\Delta\). For a derivative order \(j_i\)
put \(s_i=\Delta_i+1-j_i\) and \(J=\sum_i j_i\). On nested initial
polydisks allocate a total shrink \(\delta\) proportionally to the
positive \(j_i\), omitting zero-derivative steps. Cauchy's formula
bounds the derivative composition by \((e/\delta)^J J!\), with value one
when \(J=0\). This follows from \(j_i!\le j_i^{j_i}\) and
\(J^J\le e^JJ!\). Derivatives hitting earlier coefficients are included
in this operator-composition bound. The remaining factorials satisfy
\begin{equation}\protect\phantomsection\label{eq:gg-op-flow-factorial-budget}{
 r_0!\left(\prod_i s_i!\right)J!\le(r+n)!,
 \qquad
 r_0+\sum_i s_i+J=r+n .
}\end{equation} After increasing \(C_1\) to at least one, put
\(D=C_1\max(1,e/\delta)\). A chosen sequence is bounded before
integration by \(M_PA_P^{r_0}M_1^nD^{\Delta+n}(r+n)!\). There are at
most \(2^{\Delta+3n}\) allowed pairs \((\Delta_i,j_i)\): the nonnegative
\(2n\)-tuple \((s_i,j_i)\) has sum \(\Delta+n\), and discarding its
additional restrictions gives this bound. The time simplex contributes
\(|t|^n/n!\), while
\begin{equation}\protect\phantomsection\label{eq:gg-op-flow-time-budget}{
 \frac{(r+n)!}{n!}
   =r!\binom{r+n}{r}\le r!2^{r+n}.
}\end{equation} For \(|t|\le T\) set \(q=16M_1DT\). Since
\(n\le\Delta\), the finite sum is bounded by
\(\sum_{n=0}^{\Delta}q^n\le(1+q)^\Delta\). Summing also in \(r_0\)
gives, for example,
\begin{equation}\protect\phantomsection\label{eq:gg-op-flow-output-bound}{
 \|V_{m-r}(t)\|
 \le M_Pr!\left[2\{A_P+2D(1+q)\}\right]^r .
}\end{equation} The initial term \(n=0\), \(\Delta=0\) is included.
There is no condition \(q<1\). Pullback by the classical maps gives the
same type of bound on the corresponding smaller final tube.

At each fixed drop the series has finitely many positive-drop steps and
uses finitely many input derivatives at the classical endpoint. Hence it
applies to arbitrary input germs near that endpoint; the input need not
extend along the whole trajectory. The tube carries the coefficient
operators. This proves locality as a sheaf map, as well as
coefficientwise holomorphy and the analytic-microdifferential growth
estimate.

Uniqueness is recursive in the drop. The highest symbol has the
prescribed classical transport; every lower symbol has that transport
with source determined by earlier drops. The commutator derivation rule
makes products solve the same equation, so uniqueness gives
multiplicativity. Reverse evolution gives the inverse after finite
parent restrictions. Applying the formal adjoint to the equation gives
adjoint preservation when \(B\) is skew. All estimates are uniform on
the stated parameter parents. ∎

The theorem constructs a transport automorphism. It does not place a
positive-order exponential inside the operator ring, and it does not
assert scalar convergence of the negative homogeneous symbol series.

\protect\phantomsection\label{sec:gg-op-analytic-coherence}

\subsubsection{F.7. Analytic inner
coherence}\label{f.7.-analytic-inner-coherence}

\pandocstatementheading{\protect\phantomsection\label{lem:gg-op-negative-unit-evolution}Lemma F.5.}

Let \(R(t,\mu)\) have order at most minus one and uniform analytic
symbol bounds on a fixed larger parent and bounded time disk. Then
\(C'=RC\), \(C(0)=1\), has a jointly holomorphic normalized order-zero
unit solution on a smaller parent. If \(R^*=-R\), then \(C^*C=1\).

\emph{Proof.} In normal symbols assume \(\|R_{-l}\|\le M_RA_R^ll!\),
\(l\ge1\). A product term with component drop \(l\) and derivative order
\(k\) increases the drop by \(l+k\). Cauchy on the \(R\) factor cancels
the derivative factorial divisor. At output drop \(r\), an \(n\)-fold
product has \(n\le r\) and \(L+K=r\). The nested Cauchy estimate gives
factorial factor \((\prod_i l_i!)K!\le r!\). Counting the choices by
\(2^{r+n}\) and integrating over the simplex yields
\begin{equation}\protect\phantomsection\label{eq:gg-op-negative-unit-bound}{
 \|C_{-r}\|\le e^{2M_RT}(2D_R)^r r!,
 \qquad
 D_R\ge\max(1,A_R,eC_R/\delta).
}\end{equation} At each drop the series is finite. Right multiplication
similarly constructs the inverse, and the principal symbol is one. If
\(R\) is skew, differentiation gives \((C^*C)'=0\). The same estimates
and finite coefficient formulas prove holomorphy in every passive
parameter. ∎

\pandocstatementheading{\protect\phantomsection\label{thm:gg-op-analytic-source-descent}Theorem F.6.}

The actual corrected contact source overlaps admit analytic relative
filtered operator comparisons and normalized unitary triple units on
common bounded-family parents, including amplitude one. Their symbols
are the original actual contact maps, and their induced transport on
\(\mathscr U_1\) is the original calibrated affine transport.

\emph{Proof.} Theorem 6.3 provides the contact maps, multipliers and
inverses on a common complex amplitude disk containing one. Relative to
the exact reference atlas, their homogeneous symplectizations supply the
classical tubes of Theorem F.4. Fixed finite cover refinements and
parent restrictions give uniform coordinate and inverse bounds. Cauchy
in a larger amplitude disk reads the holomorphic Hamiltonian of the
actual logarithmic time derivative. No logarithm of an infinite
diffeomorphism series is used. Section F.1 lifts this Hamiltonian to a
skew order-one family \(B_{ij}(t)\) with the required uniform symbol
bounds.

Integrate only one independent orientation of each overlap, starting at
\(\varphi^0_{ij}\). Define the reverse literally by
\begin{equation}\protect\phantomsection\label{eq:gg-op-analytic-inverse-generator}{
 \varphi_{ji}=\varphi_{ij}^{-1},\qquad
 B_{ji}=-\varphi_{ji}(B_{ij}).
}\end{equation} The reverse generator is the displayed actual generator,
rather than a separately reselected principal lift. Theorem F.4 gives
the analytic filtered comparisons with their actual classical symbols.
Proposition F.2 then gives the old affine cochains and moving arguments.

On a triple put \(D_{ijk}=\varphi_{ij}\varphi_{jk}\varphi_{ik}^{-1}\).
Direct differentiation gives
\begin{equation}\protect\phantomsection\label{eq:gg-op-analytic-triple-generator}{
 \dot D_{ijk}D_{ijk}^{-1}=\operatorname{ad}R_{ijk},\qquad
 R_{ijk}=B_{ij}+\varphi_{ij}(B_{jk})-D_{ijk}(B_{ik}).
}\end{equation} The classical symbol of \(D_{ijk}\) is the identity for
every amplitude. Thus the degree-one symbol of \(R_{ijk}\) is Poisson
central and is zero, since a homogeneous degree-one central function
vanishes. Skewness then kills its degree-zero symbol, so \(R_{ijk}\) has
order at most minus one. Theorem F.4 applied with input families
\(B_{jk}\) and \(B_{ik}\) gives uniform factorial bounds for this
generator after finitely many extra parent restrictions.

Lemma F.5 constructs the actual holomorphic unit \(C_{ijk}\) with
\(C'_{ijk}=R_{ijk}C_{ijk}\) and \(C_{ijk}(0)=1\). Its adjoint action
equals \(D_{ijk}\) by uniqueness. The two sides of the quadruple
relation have the same inner action, are unitary, and have degree-zero
symbol one. A normalized order-zero unit with trivial inner action is
one, so the quadruple relation is exact. This also follows from the
normalized-inner injectivity of the local model in (Kashiwara 1996,
Lemma 1). This constructs relative analytic coherence. ∎

The construction uses the source family's classical parents and finitely
many extra buffer restrictions.

\protect\phantomsection\label{sec:gg-op-analytic-finite-comparison}

\subsubsection{F.8. Comparison with the auxiliary finite
operators}\label{f.8.-comparison-with-the-auxiliary-finite-operators}

\pandocstatementheading{\protect\phantomsection\label{prop:gg-op-analytic-finite-comparison}Proposition F.7.}

The amplitude jets of Theorem F.6 are compatible with Theorem F.1 up to
normalized filtered inner equivalence. The contact and calibrated affine
tuples are unchanged.

\emph{Proof.} For an oriented overlap write \(B(a)\) for the generator
of Theorem F.6 in its amplitude variable \(a\). Over
\(A_N=\mathbb C[\varepsilon]/(\varepsilon^{N+1})\), use a separate path
parameter \(t\) and set
\begin{equation}\protect\phantomsection\label{eq:gg-op-finite-fundamental-unit}{\begin{gathered}
B_N(t,\varepsilon)=\varepsilon B(t\varepsilon)\pmod{\varepsilon^{N+1}}, \\
\partial_t U_N(t,\varepsilon)=B_N(t,\varepsilon)U_N(t,\varepsilon), \\
U_N(0,\varepsilon)=1.
\end{gathered}}\end{equation} Since \(B_N\in\varepsilon F_1\mathcal E\otimes A_N\),
its time-ordered series has at most \(N\) factors and defines a finite
polynomial with operator coefficients. Put \(U_N=U_N(1,\varepsilon)\).
Its adjoint action solves the amplitude-jet equation of Theorem F.6.
This construction takes the finite Taylor coefficients of \(B\) and does
not exponentiate a positive-order operator over the ordinary analytic
coefficient field. Its logarithm belongs to \(\varepsilon F_1\) by the
finite Magnus recursion, since \(F_1\) is closed under commutators. Let
\(U_N^{\rm old}\) be the auxiliary exponential of Theorem F.1. The
logarithm of \(W_N=U_N(U_N^{\rm old})^{-1}\) has order at most one and
zero degree-one symbol: both have the same classical contact action.
Unitarity makes its logarithm skew and kills the degree-zero symbol.
Hence \(W_N\in1+F_{-1}\mathcal E\otimes A_N\).

These compatible units relate the finite comparisons by conjugation.
Their triple units satisfy the exact change law
\begin{equation}\protect\phantomsection\label{eq:gg-op-analytic-finite-triple-change}{
 C^{\rm an}_{ijk}
 =W_{ij}\varphi^{\rm old}_{ij}(W_{jk})
             C^{\rm old}_{ijk}W_{ik}^{-1}
       \pmod{\varepsilon^{N+1}} .
}\end{equation} Both sides are the normalized unit with the indicated
inner action, so equality follows also by injectivity. Conjugation by
\(W_N\) is trivial on \(\mathscr U_1\). Thus the affine cochains, their
inverses and compact lifts remain the same. ∎

This comparison does not assert literal equality or convergence of the
previously chosen auxiliary infinite operator section. The physical
source tuples retain their classical gauge.

\protect\phantomsection\label{sec:gg-op-relative-inner-comparisons}

\subsubsection{F.9. Relative normalized inner
comparisons}\label{f.9.-relative-normalized-inner-comparisons}

\pandocstatementheading{\protect\phantomsection\label{lem:gg-op-relative-inner-comparisons}Lemma F.8.}

Let \(\psi_\mu\) be a holomorphic family of filtered, adjoint-preserving
automorphisms of the local half-density algebra, inducing the identity
on its full graded-symbol algebra. On a common parent assume the
coordinate differences \(\Delta_i=\psi_\mu(y_i)-y_i\) have uniform
analytic factorial symbol bounds. There is a unique holomorphic
normalized unit \(P_\mu\) with
\begin{equation}\protect\phantomsection\label{eq:gg-op-relative-inner-normalization}{
 \psi_\mu=\operatorname{Ad}P_\mu,\qquad
 \operatorname{ord}P_\mu=0,\qquad
 \sigma_0(P_\mu)=1,\qquad P_\mu^*P_\mu=1.
}\end{equation} It and its inverse have common factorial estimates after
fixed parent restrictions. The construction commutes with restrictions
and finite jets of the holomorphic parameter family.

\emph{Proof.} The fixed-manifold normalized-inner theorem (Kashiwara
1996, Lemma 1) supplies existence and uniqueness at each fixed parameter
and spatial germ. We prove the relative assertion directly. Since
\(y_i^*=y_i\), the difference \(\Delta_i\) is self-adjoint. Identity on
the graded algebra gives order at most minus one, and self-adjointness
kills that odd degree. Thus \(\operatorname{ord}\Delta_i\le-2\).

Use normal total symbols, with momenta \(p_i\). The orientation
\(\psi(A)=PAP^{-1}\) gives
\begin{equation}\protect\phantomsection\label{eq:gg-op-relative-inner-euler}{
 [P,y_i]=\partial_{p_i}P=\Delta_i\circ P,\qquad
 P_{-r}=-\frac1r
  \left[\sum_i p_i(\Delta_i\circ P_{<r})\right]_{-r},
 \quad P_0=1.
}\end{equation} Here \(p_i\) outside the parentheses is ordinary
multiplication of total-symbol functions, not another operator product.
A term from \(\Delta_{i,-l}\), \(P_{-s}\) and \(k\) paired derivatives
has output drop \(r=s+l+k-1\). Because \(l\ge2\), it uses only \(s<r\).
Each recursive coefficient is therefore a finite holomorphic
differential polynomial. At every fixed parameter the actual normalized
unit obeys this recursion; triangular uniqueness identifies it with the
recursive coefficients. This establishes all intertwining and adjoint
equations, rather than inferring them from the Euler-contracted equation
alone.

For completeness, suppose \(\|\Delta_{i,-l}\|\le MA^l l!\). Reserve one
fixed momentum buffer first. Normal composition and Cauchy estimates
give
\begin{equation}\protect\phantomsection\label{eq:gg-op-relative-inner-coefficient-bound}{
 \left\|\frac{p_i}{\alpha!}
           \partial_p^\alpha\Delta_{i,-l}\right\|
 \le M_1 A^l l! C^k,\qquad k=|\alpha|.
}\end{equation} The finite coordinate sum, momentum factor and
multiindex count are included in \(M_1,C\). These coefficient functions
remain holomorphic on a fixed larger \(y\)-parent. In an iterated
recursive term with \(n\) steps write \(L=\sum l_i\), \(K=\sum k_i\).
Its positive drops \(d_i=l_i+k_i-1\) give \(r=L+K-n\) and \(n\le r\).
The successive Euler divisors are the strictly increasing cumulative
drops, so their product is at least \(n!\).

Allocate a total \(y\)-buffer \(\delta\) in proportions \(\delta k_i/K\)
to steps with \(k_i>0\). Nested Cauchy estimates, applied to each entire
preceding expression, include derivatives of its intermediate
coefficients and give \((e/\delta)^K K!\). For \(K=0\) the factor is
one. Moreover,
\begin{equation}\protect\phantomsection\label{eq:gg-op-relative-inner-factorial-budget}{\begin{gathered}
\left(\prod_i l_i!\right)K!\le(r+n)!, \\
\frac{(r+n)!}{n!}\le r!2^{r+n}, \\
\#\{(l_i,k_i)\}\le
       \binom{r+n-1}{2n-1}\le2^{r+n}.
\end{gathered}}\end{equation} The count uses \(l_i-2,k_i\ge0\) with total \(r-n\);
counting terms that differentiate the constant \(P_0\) only enlarges the
bound. With \(D=\max(1,A,eC/\delta)\),
\begin{equation}\protect\phantomsection\label{eq:gg-op-relative-inner-final-bound}{
 \|P_{-r}\|
 \le r!(4D)^r\sum_{n=1}^r(4M_1D)^n
 \le r!\,[4D(1+4M_1D)]^r.
}\end{equation} No smallness condition is needed. These are analytic
microdifferential symbol bounds, not scalar convergence of the
negative-homogeneous series. The inverse is \(P^*\); its normal adjoint
expansion has drop \(r=s+k\), and the two Cauchy factorials divided by
\(\alpha!\) leave at most \(s!k!\le r!\), with exponentially many
indices. It therefore has the same type of bound. Finite coefficient
formulas establish joint holomorphy, restrictions and parameter-jet
compatibility. This does not assert a theorem about arbitrary nonlifted
nonreduced bases.

Finally, two such units differ by a central passive-parameter function.
Its degree-zero symbol is itself, so symbol-one normalization makes it
one. The locally constructed units thus glue uniquely. ∎

For two actual structured quantizations covering the same full classical
symbol map and satisfying the common-parent bounds, their quotient
satisfies this lemma. Its unique normalized unit gives the actual
relative comparison and the coherent finite comparison diagrams.

\protect\phantomsection\label{app:gg-relative-sky-modules}

\subsection{Appendix G. Relative simple modules on complete
skies}\label{appendix-g.-relative-simple-modules-on-complete-skies}

\protect\phantomsection\label{sec:gg-module-standard-model}

\subsubsection{G.1. An analytic order-zero local
model}\label{g.1.-an-analytic-order-zero-local-model}

Use contact coordinates \((z,q^1,q^2;p_1,p_2)\) with
\(\alpha=dz+p_i\,dq^i\), and homogeneous cotangent coordinates
\((t,\xi_1,\xi_2)\) with \(t\ne0\), \(p_i=\xi_i/t\). The reference
Legendrian has \(z=\xi=0\). All algebras below are the scalar analytic
half-density algebras of Appendix F.

\pandocstatementheading{\protect\phantomsection\label{lem:gg-module-standard-model}Lemma G.1.}

The finite presentation
\begin{equation}\protect\phantomsection\label{eq:gg-module-standard-presentation}{
 P_0=zD_z+\tfrac12,\qquad
 \mathcal M_0=\mathcal E/
       (\mathcal E P_0+\mathcal E D_{q^1}+\mathcal E D_{q^2})
}\end{equation} defines a simple analytic module of order zero on this
Legendrian. It has coefficient normal form \(f(q,t)u_0\), with the usual
factorial bounds, and actions
\begin{equation}\protect\phantomsection\label{eq:gg-module-standard-action}{
 D_z=t,\qquad D_{q^i}=\partial_{q^i},\qquad
 z=Z=-\partial_t+\frac1{2t}.
}\end{equation} Its endomorphisms are \(\mathbb C\); in a passive
holomorphic family they are \(\mathcal O_S\). All normal-form and action
formulas are analytic and relative.

\emph{Proof.} These operators satisfy the Weyl relations. In particular
\((Zt+\tfrac12)f=-t\partial_t f\), so \(P_0u_0=0\) and
\(zD_zu_0=-u_0/2\).

We give analytic division. Expand a normal symbol of order \(m\) in
\(z\) and \(\theta=\xi/t\):
\begin{equation}\protect\phantomsection\label{eq:gg-module-symbol-taylor}{
 P=\sum_{l,h,\beta}
 a_{l,h,\beta}(q)\,z^hD_q^\beta D_z^{m-l-|\beta|},
 \qquad
 |a_{l,h,\beta}|\le MA^l l!C^{h+|\beta|}.
}\end{equation} For \(\beta\ne0\), the division uses finitely many
actual analytic generators. If \(F_l(z,q,\theta)\) is the analytic
Taylor coefficient at order \(m-l\), write
\begin{equation}\protect\phantomsection\label{eq:gg-module-finite-taylor-division}{
 F_l(z,q,\theta)-F_l(z,q,0)
 =\sum_i\theta_i\int_0^1
           \partial_{\theta_i}F_l(z,q,a\theta)\,da .
}\end{equation} On a fixed strict \(\theta\)-parent, Cauchy gives the
integrals a bound \(M'A^l l!\). Since \(\theta_i=\xi_i/t\), their
order-\((m-1)\) symbols define the two analytic rightmost-\(D_{q^i}\)
quotients. No \(q\)-dependent coefficient is moved through a \(D_q\).
Thus the entire infinite \(\beta\) tail is in the displayed finite left
ideal with analytic quotient bounds. For \(\beta=0\), put \(s=m-l\) and
use
\begin{equation}\protect\phantomsection\label{eq:gg-module-ore-division}{
 zD_z^s=D_z^{s-1}P_0+(\tfrac12-s)D_z^{s-1},\qquad
 c_{h,s}=(-1)^h(s-\tfrac12)_{\underline h}.
}\end{equation} The remainder of \(z^hD_z^s\) is \(c_{h,s}D_z^{s-h}\);
its quotient by the rightmost \(P_0\) is
\(\sum_{v=0}^{h-1}z^{h-1-v}c_{v,s}D_z^{s-1-v}\). At quotient drop
\(r=l+v\), only finitely many \((l,v)\) occur. The remaining tail is the
convergent analytic sum \(\sum_{h\ge v+1}a_{l,h,0}(q)z^{h-1-v}\). On
\(C|z|\le1/2\) its norm is at most \(2MA^l l!C^{v+1}\). The elementary
falling-factorial bound
\begin{equation}\protect\phantomsection\label{eq:gg-module-division-bound}{
 l!\,|c_{v,m-l}|\le C_m^{l+v}(l+v)!
}\end{equation} therefore gives factorial bounds for the actual
quotient. The remainder has \(r=l+h\) and satisfies the same estimate.
Consequently there are actual analytic operators \(Q_0,Q_1,Q_2\) and an
analytic normal form such that \(P=Q_0P_0+Q_iD_{q^i}+f(q,D_z)\). Its
action on \(u_0\) retains exactly \(f\), proving that the kernel is the
displayed finite left ideal, not only its formal closure.

For the full action one has
\begin{equation}\protect\phantomsection\label{eq:gg-module-z-action}{
 Z^h(t^sf)=(-1)^h
   \sum_{a=0}^h\binom ha
    (s-\tfrac12)_{\underline{h-a}}\,
       t^{s-h+a}\partial_t^a f .
}\end{equation} Here \(s=m-l-|\beta|\). Put \(\Delta=l+h+|\beta|\) and
\(j=|\beta|+a\). The falling factorial, together with \(l!\), is bounded
by \(C_m^{l+|\beta|+h-a}(l+h-a)!\). For example, bound it by the ratio
of factorials with arguments \(l+|\beta|+c+h-a-1\) and
\(l+|\beta|+c-1\), where \(c\) depends only on \(m\), and absorb its
binomial ratio into an exponential constant. Since \(l+h-a=\Delta-j\),
the order-\(j\) differential coefficient is bounded by
\begin{equation}\protect\phantomsection\label{eq:gg-module-action-bound}{
 M C_m^\Delta(\Delta-j)!,\qquad j\le\Delta.
}\end{equation} The fixed radial annulus and Taylor counts only alter
the exponential constants. An input coefficient of drop \(r_0\),
followed by \(j\) Cauchy derivatives on one buffer, obeys
\(r_0!(\Delta-j)!j!\le(r_0+\Delta)!\). Thus the full action preserves
the analytic normal-form class. At each drop only finitely many jets
enter. The Weyl relations give coefficientwise associativity, and the
bounds restrict that same representation to the analytic class.

The graded annihilator is the reduced ideal \((z,\xi)\) and the simple
symbol line has rank one. Its order-zero local half-form is
\begin{equation}\protect\phantomsection\label{eq:gg-module-order-zero-half-form}{
 \rho_0=\sqrt{dq^1\,dq^2\,dt/t}.
}\end{equation} Indeed \(P_0\) has zero subprincipal symbol and
principal Hamilton field \(-t\partial_t\) on the cone; its Lie action on
this half-form is zero. This agrees with the simple-symbol calculus of
(D'Agnolo and Schapira 2007, sec. 4). It differs from the order-one-half
delta generator. An endomorphism sends \(u_0\) to \(f(q,t)u_0\) killed
by \(D_q\) and \(P_0\), whence \(\partial_qf=t\partial_tf=0\). This
proves the endomorphism assertion. Uniform Taylor convergence and the
finite drop formulas prove all relative assertions. ∎

\protect\phantomsection\label{sec:gg-module-tangent-transport}

\subsubsection{G.2. Actual analytic tangent
transport}\label{g.2.-actual-analytic-tangent-transport}

\pandocstatementheading{\protect\phantomsection\label{lem:gg-module-tangent-transport}Lemma G.2.}

Let \(B(a,\mu)\) be a skew-adjoint order-one family with uniform
analytic symbol bounds whose principal Hamiltonian vanishes on the
reference Legendrian. Suppose its actual tangent cone flow and inverse
have common holomorphic tubes. The equation \(S'=BS\), \(S_0=1\), has an
analytic filtered module evolution with inverse and uniform factorial
bounds. If \(\varphi'=\operatorname{ad}(B)\varphi\) and
\(\varphi_0=\mathrm{id}\), then
\begin{equation}\protect\phantomsection\label{eq:gg-module-transport-semilinearity}{
 S_a(Pm)=\varphi_a(P)S_a(m).
}\end{equation} Its leading symbol is the actual half-form transport.
The construction is local and holomorphic in passive parameters.

\emph{Proof.} The first transverse Taylor terms of the principal symbol
are \(H=tzb(q)+\xi_i v^i(q)+\text{higher transverse terms}\). There is
no degree-raising action on \(\mathcal M_0\). In normal symbols
skew-adjointness gives \(B_0|_\Lambda=(b+\operatorname{div}_qv)/2\). The
action of \(bzD_z\) is \(-bt\partial_t-b/2\), so the zero-drop part is
exactly
\begin{equation}\protect\phantomsection\label{eq:gg-module-leading-lie-transport}{
 L_0=v^i\partial_{q^i}-bt\partial_t+
                         \tfrac12\operatorname{div}_qv .
}\end{equation} For the volume \(dq^1dq^2dt/t\), the divergence of
\(v\partial_q-bt\partial_t\) is \(\operatorname{div}_qv\). Thus \(L_0\)
is its half-form Lie derivative, including the scalar term. Its actual
evolution is the tangent pullback with the initialized nonzero
half-density Jacobian amplitude.

For all higher terms put \(\Delta=l+h+|\beta|-1\ge1\). Equation
(\ref{eq:gg-module-action-bound}), with \(m=1\), gives differential
order \(j\le\Delta+1\) and coefficient bound
\(MC^{\Delta+1}(\Delta+1-j)!\). Conjugate by the actual \(L_0\)
evolution, including its amplitude and inverse. Analytic coordinate
changes and multiplication by these amplitudes change an order-\(j\)
derivative to order \(i\le j\) at cost at most \(C_1^j j!/i!\). The
inequality
\begin{equation}\protect\phantomsection\label{eq:gg-module-conjugated-action-bound}{
 (\Delta+1-j)!\frac{j!}{i!}
 \le 2^j(\Delta+1-i)!
}\end{equation} retains exactly the same factorial type.

The interaction-picture iteration now has positive drops, as in Theorem
F.4. At output drop \(r\), an input of drop \(r_0\) permits only
\(n\le r-r_0\) steps. If \(s_i=\Delta_i+1-j_i\) and \(J=\sum j_i\), the
single-total-buffer estimate is
\begin{equation}\protect\phantomsection\label{eq:gg-module-transport-factorial-budget}{
 r_0!\left(\prod_i s_i!\right)J!\le(r+n)!.
}\end{equation} The ordered time simplex supplies \(1/n!\). The finite
composition counts and \((r+n)!/n!\le2^{r+n}r!\) give a uniform
exponential times \(r!\), by the finite \(n\) sum of Theorem F.4. There
is no loss of a new domain at each normal order. This constructs the
analytic evolution. Reverse evolution gives the inverse. Its leading
term is the \(L_0\) half-form transport.

Both sides of (\ref{eq:gg-module-transport-semilinearity}) have
derivative \(B\) times themselves: on the right the commutator term and
the module derivative combine to \(B\varphi(P)S(m)\). They have the same
initial value, so coefficientwise uniqueness proves equality. Each
output coefficient uses finitely many endpoint jets; this proves
locality, passive holomorphy and compatibility with restrictions and
finite parameter jets. ∎

\protect\phantomsection\label{sec:gg-module-relative-descent}

\subsubsection{G.3. Moving modules and exact algebroid
descent}\label{g.3.-moving-modules-and-exact-algebroid-descent}

Write \(B\) for the passive parameter base and \(S\to B\) for the total
physical point family; at a fixed passive parameter its fiber is the
four-germ \(M\). The sky projection is \(p:F\to S\). For the source
family \(F=Q\times S\), \(Q=\mathbb P^1\times\mathbb P^1\), with its
actual parameterized embeddings. Retain an actual relative
half-canonical root \(R^2=K_{F/S}\) and its marked calibration. The
relative cotangent sequence gives the specified canonical-line
comparison
\begin{equation}\protect\phantomsection\label{eq:gg-module-relative-canonical-base}{
 K_{F/B}\simeq K_{F/S}\otimes p^*K_{S/B}.
}\end{equation} Thus the canonical line relative only to the passive
base has an additional point-base factor. At a fixed passive parameter
\(S=M\), so the sky-relative notation becomes \(K_{F/M}\). In the ruled
parameterization one may take \(R=\mathcal O(-1,-1)\) with its actual
root frames. Any base line factor in comparing \(R\) to the inverse
contact line is retained with its calibrated frame; fiberwise bidegree
alone does not identify relative lines.

\pandocstatementheading{\protect\phantomsection\label{thm:gg-module-relative-descent}Theorem G.3.}

The analytic operator descent of Theorem F.6 has an actual holomorphic
family of locally standard order-zero regular simple modules supported
on its complete skies. Their local semilinear maps obey
\begin{equation}\protect\phantomsection\label{eq:gg-module-actual-triple}{
 \tau_{ij}\tau_{jk}=C_{ijk}\tau_{ik}.
}\end{equation} They have common bounded-family representatives,
inverses and finite parameter-jet compatibility. Their chosen order-zero
lattices have degree-zero symbol line \(R\).

\emph{Proof.} Locally choose a fixed contact polarization adapted to the
central sky. A fixed partial Legendre quantized chart change is
permitted. Nearby skies are graphs \(z=f_s(q)\), \(p=-df_s\). The base
substitution \((z,q)\mapsto(z-f_s(q),q)\) and its cotangent half-density
pullback flatten them analytically. Pull back \(\mathcal M_0\) to obtain
the local modules and their finite presentations.

Write \(L_i\) for the flattening algebra maps and
\(\psi_{ij}=L_i\varphi_{ij}L_j^{-1}\). If
\(\dot L_i=\operatorname{ad}(A_i)L_i\) and
\(\dot\varphi_{ij}=\operatorname{ad}(B_{ij})\varphi_{ij}\), the
flattened generator is
\begin{equation}\protect\phantomsection\label{eq:gg-module-actual-flat-generator}{
 B^{\rm flat}_{ij}
     =A_i+L_i(B_{ij})-\psi_{ij}(A_j).
}\end{equation} Every \(\psi_{ij}\) preserves the reference Legendrian,
so this generator is tangent. It is the generator of the given
comparison, not a separately reselected principal lift.

We construct joint Hom generators in two stages. At source amplitude
zero, move the physical point along \(x=ax\), starting at \(x=0\). The
reference \(\varphi_{ij}\) is fixed, so the displayed formula has
\(B_{ij}=0\). Start from a single fixed same-order simple intertwiner
\(\tau_0\), supplied by the local simple-symbol calculus (D'Agnolo and
Schapira 2007, sec. 4). Apply Lemma G.2 to \(\psi_a\psi_0^{-1}\) and
take \(S_a\tau_0\); this gives an actual holomorphic family in \(x\).
Next follow the source amplitude path, with \(x\) and the remaining
source parameters passive, using the same initialized relative evolution
with the full generator. It gives holomorphic semilinear intertwiners on
the whole retained family. Thus local Hom is actually free of rank one
over \(\mathcal O_S\), rather than merely one-dimensional at each
parameter.

Choose local simple generators \(u_i\) with prescribed projective
half-form symbols \(\rho_i\) of \(R\), by an elliptic order-zero symbol
lift. The cone half-form in (\ref{eq:gg-module-order-zero-half-form})
descends to this relative half-root: changing the radial frame adds a
base \(d\log\) term, which wedges to zero with the top \(q\)-form. Let
\(\rho_j=s_{ij}\rho_i\) at the actual moving arguments. For an
intertwiner \(u_j\mapsto b_{ij}u_i\) let \(e_{ij}\) be its leading
coefficient. The fixed initial simple-symbol calculus and Lemma G.2's
half-form transport imply
\begin{equation}\protect\phantomsection\label{eq:gg-module-symbol-normalization}{
 e_{ij}/s_{ij}\in\mathcal O_S^\times .
}\end{equation} This is an actual holomorphic unit, evaluated on a local
calibrated cone section; its holomorphy follows from the Hom
construction just given. Divide by it. The resulting \(\tau_{ij}\) has
symbol ratio exactly \(s_{ij}\) and is unique with that normalization.

Both \(\tau_{ij}\tau_{jk}\) and \(C_{ijk}\tau_{ik}\) are semilinear for
\(\varphi_{ij}\varphi_{jk}=\operatorname{Ad}(C_{ijk})\varphi_{ik}\).
Their symbols agree because \(\sigma_0(C_{ijk})=1\) and
\(s_{ij}s_{jk}=s_{ik}\) is the actual root cocycle. Their ratio is an
\(\mathcal O_S\) scalar by Lemma G.1 and is one by its actual symbol.
This proves the full module equation. On generators it reads
\begin{equation}\protect\phantomsection\label{eq:gg-module-generator-triple}{
 \varphi_{ij}(b_{jk})b_{ij}u_i=C_{ijk}b_{ik}u_i .
}\end{equation} The unit acts from the left; this is a module identity,
not a ring identity for representatives \(b_{ij}\). The actual
associator has therefore been retained. Inverses are literal inverses
and have the inverse root symbol.

A finite cover of compact \(Q\), its actual graph tubes and fixed parent
restrictions give the common estimates. The filtered comparisons and
elliptic order-zero intertwiners preserve the chosen
\(\mathcal E_i(0)u_i\) lattices with their module descent. The
associator has degree-zero action one, so their degree-zero lines glue
honestly with transitions \(s_{ij}\) to \(R\). All coefficient
constructions commute with restrictions and finite jets of holomorphic
parameter families. ∎

\protect\phantomsection\label{sec:gg-module-relative-framing}

\subsubsection{G.4. Calibration on the whole parameter
family}\label{g.4.-calibration-on-the-whole-parameter-family}

\pandocstatementheading{\protect\phantomsection\label{prop:gg-module-relative-framing}Proposition G.4.}

In the locally standard family category of Theorem G.3, retain an actual
nonzero holomorphic vector section of the degree-zero line along the
calibrated cone section \(\sigma:S\to\widehat F\). The supported module
is then unique up to unique framed isomorphism, locally on the parameter
base. The source construction has this calibration, and auxiliary local
choices have canonical relative comparisons.

\emph{Proof.} Local Hom is the rank-one \(\mathcal O_S\) local system
constructed above. Both \(Q\) and its actual symplectization
\(\mathcal O(-1,-1)^\times\) are simply connected, by Proposition F.3.
Thus its global Hom along each complete sky/cone is a rank-one
\(\mathcal O_S\) line locally on \(S\). Evaluation at the calibrated
cone section identifies its scalar with the action on the actual
degree-zero line. The prescribed nonzero vector section fixes that
scalar holomorphically. An automorphism acts by
\(f(s)\in\mathcal O_S^\times\); preserving the entire section gives
\(f=1\). The marked ruled-root calibration supplies this vector section
for the source family. This proves existence and uniqueness of
comparisons of its auxiliary realizations. ∎

Fixing a vector only at one parameter \(s_0\) would leave nonconstant
units with \(f(s_0)=1\) and would not prove this assertion. The
proposition concerns the stated locally standard analytic family
category; it does not classify arbitrary nonflat relative regular
modules or arbitrary good lattices.

\protect\phantomsection\label{sec:gg-op-physical-comparisons}

\subsubsection{G.5. Actual physical comparisons and coupled
reality}\label{g.5.-actual-physical-comparisons-and-coupled-reality}

Physical admission is the enriched admission of §8: preserve the
complete incidence, calibrated affine/projective data, canonical section
and actual full spin lift. It is stronger than preserving a bare contact
distribution. Comparisons below are between the same actual source
presentations and use their conditioned common parents.

\pandocstatementheading{\protect\phantomsection\label{thm:gg-op-physical-comparisons}Theorem G.5.}

Every actual physical source comparison of Theorem 8.3 has an analytic
structured operator lift and a comparison of its calibrated supported
modules. The lifts have coherent composition, identity and inverse
comparisons. Their first negative-order unit action is the original
affine action, and their physical Hom sets after normalized comparison
are precisely the full-spin Hom sets of §8, with the same map
directions. The coupled Lorentz operation of Proposition 8.4 has an
antilinear structured lift with its actual square coherence and
Real-compatible framed supported modules.

\emph{Proof.} Put \(V=R_{a,A}U\). The displayed metric map is
\(\Xi_{a,A}:g[V]\to g[U]\). Use its ray point map \(K\) in that
direction and quantize \((K^{-1})^*\). Thus local algebra maps have
direction \(k_i:\mathcal E_i^V\to\mathcal E_i^U\). The inverse realizes
the covariant source instruction \(U\to V\); the directions are not
interchanged.

Use the actual nonlinear physical amplitude family
\begin{equation}\protect\phantomsection\label{eq:gg-op-physical-amplitude}{
 U_\varepsilon=(\varepsilon U;\varepsilon\Lambda),
 \qquad V_\varepsilon=R_{a,A}(U_\varepsilon).
}\end{equation} In general \(V_\varepsilon\ne\varepsilon R_{a,A}U\). The
record-flow bounds and the actual exponential, transport and hit
formulas of §8 give the analytic \(K_\varepsilon\) and its inverse on
fixed finite conditioned diagrams. Evaluate Theorem F.6 at the passive
source family \(V_\varepsilon\); its own internal amplitude is an
independent variable, set to one after construction. Joint bounds
justify this evaluation.

At \(\varepsilon=0\) the transformation is the flat translation and the
two determinant-one spin actions. Its homogeneous twistor substitution
has determinant one and the fixed half-density action on
\(\mathcal O(-2)\). It supplies the reference quantum map, including its
actual parameter dependence. Read the logarithmic derivative of
\((K_\varepsilon^{-1})^*\) with the point/operator sign convention of
§8.1. After reference identification, Theorem F.4 constructs the
initialized local \(k_i\). Extra parent restrictions are finite and
retain the whole compact sky.

On overlaps the two maps \(k_i\varphi^V_{ij}\) and \(\varphi^U_{ij}k_j\)
have the same full classical symbol by Proposition 8.1. Their quotient
satisfies Lemma F.8, with bounds from the actual constructed maps and
inverses. It gives a unique holomorphic normalized unit \(T_{ij}\) such
that
\begin{equation}\protect\phantomsection\label{eq:gg-op-physical-descent}{
 k_i\varphi^V_{ij}
       =\operatorname{Ad}(T_{ij})\varphi^U_{ij}k_j,\qquad
 T_{ij}\varphi^U_{ij}(T_{jk})C^U_{ijk}
       =k_i(C^V_{ijk})T_{ik}.
}\end{equation} For the second equation, compute the inner action of
\(k_i\varphi^V_{ij}\varphi^V_{jk}\) in the two ways. All implementers
are unitary with symbol one, so Lemma F.8 makes them equal. Both
original associators occur in the equation. A composite physical lift
and a selected lift of the composite have the same actual symbol and
full spin label; the same lemma compares them. Its uniqueness proves all
finite coherence equations, including identities and inverses.

Normalized inner conjugation is trivial on \(\mathscr U_1\). Hence the
affine map remains exactly \((K,d^{-1}\tau+u)\) of Proposition 8.1, with
its moving arguments and product law. Keep the actual full \(A\),
including central elements whose scalar or tensor action has a kernel.
Over each fixed admitted physical map and full label, the structured
quantum choices form a contractible comparison groupoid. Conversely,
physical symbol admission is exactly that of §8, so Theorem 8.3 gives
its unique \((a,A)\). Thus the resulting physical Hom classes are the
same, without forgetting the label or passing to unfiltered Morita
equivalence.

For modules, pull the source and target sky families to a common
parameter base by the actual physical map. With their actual flatteners,
\(\psi=L_UkL_V^{-1}\) has generator \(A_U+L_U(B^k)-\psi(A_V)\). Its
principal vanishes on the reference Legendrian. Apply Lemma G.2 to
\(\psi_\varepsilon\psi_0^{-1}\), preceded by the fixed reference module
map. Theorem G.3's root-symbol normalization and Proposition G.4's
entire framing section give actual module comparisons \(F_i\) satisfying
\begin{equation}\protect\phantomsection\label{eq:gg-op-physical-module-descent}{
 F_i\tau^V_{ij}=T_{ij}\tau^U_{ij}F_j .
}\end{equation} Both sides are semilinear for \(k_i\varphi^V_{ij}\),
have the same root symbol and whole frame, and their ratio is a scalar;
hence they agree. The full triple equations above and in Theorem G.3
give descent and coherence. The framing is transported with the actual
sky, cone lift, full spin mark and any calibrated base factor, rather
than fixed in an old coordinate direction.

For reality, conjugate coefficients and use the actual coupled operation
of Proposition 8.4, including both chiral records, their different
residue weights and conjugation of \(\Lambda\). On conjugate complex
manifolds the remaining contact map is holomorphic. The preceding
construction quantizes it; a fixed quantized contact chart change
handles its reference polarization. Write the resulting maps as
\(J_U:\overline{\mathcal E_U}\to\mathcal E_{\mathfrak cU}\). Their
square has identity full classical symbol. Lemma F.8 gives its
normalized unit \(Z_U\) and exact coherence
\begin{equation}\protect\phantomsection\label{eq:gg-op-real-square}{
 J_{\mathfrak cU}\overline{J_U}
       =\operatorname{Ad}(Z_U),\qquad
 J_U(\overline{Z_U})=Z_{\mathfrak cU}.
}\end{equation} For the second equation, transport the first square
through \(J_U\); both normalized implementers have the same inner
action. The same uniqueness proves the descent and finite Real
comparison diagrams. This is a Real structure in the structured
algebroid category with its specified square isomorphism, rather than an
assertion that every selected local map squares strictly to the
identity.

The Lorentz sky operation exchanges the two conjugate spin factors.
Their product \(\mathcal O(-1,-1)\) has an antilinear lift of square
one, compatible with the half-canonical root after fixing its calibrated
phase. In ordered local coordinates the factor exchange reverses the
top-form wedge; the corresponding fixed half-form phase compensates its
minus sign. Multiplying an antilinear lift by a unit phase does not
change its antilinear square. Fix this compatible phase once with the
Lorentz marking. Choose the calibrated cone section on its real
null-direction locus and a nonzero compatible real vector along it, and
extend holomorphically on the source germ. These compatible choices
exist locally and are transported over the source family by the fixed
real spin identification and initialized transports.

With these actual compatible root and whole-parameter frames, Theorems
G.3--G.4 give the unique framed module comparison
\(j_U:\overline{\mathcal M_U}\to\mathcal M_{\mathfrak cU}\). Its square
is
\begin{equation}\protect\phantomsection\label{eq:gg-op-real-module-square}{
 j_{\mathfrak cU}\overline{j_U}=Z_U\cdot(-).
}\end{equation} Both sides have the same semilinearity, actual root
symbol and frame, so scalar injectivity proves equality. Higher
Real/module diagrams follow in the same way. The antilinear operation is
distinct from the complex-linear formal adjoint, and from a one-factor
antipodal involution. On real numerical \(\Lambda\) and fixed original
data its classical metric and physical arrows are exactly the Lorentz
ones of Proposition 8.4.

All constructions use the existing finite conditioned physical diagrams
and finite buffer restrictions; they do not claim uniformity over
unbounded point/spin histories. The actual chiral cotangent and
source/ordinary counit contact comparisons can be lifted by the same
argument, with their original directions and classical symbols. No point
action or point reality on one fixed mixed ordinary background is
asserted. ∎

The lifted maps are the physically admitted classical comparisons of §8.

\protect\phantomsection\label{app:gg-independent-operator-domain}

\subsection{Appendix H. Independently admitted operator
data}\label{appendix-h.-independently-admitted-operator-data}

\protect\phantomsection\label{sec:gg-independent-op-definition}

\subsubsection{H.1. The independent
domain}\label{h.1.-the-independent-domain}

\pandocstatementheading{\protect\phantomsection\label{def:gg-independent-op-domain}Definition H.1.}

A calibrated operator presentation consists of local copies of the
analytic half-density algebra of F.1, their order filtration, full
graded-symbol identifications and complex-linear formal adjoints, with
the square and degree-zero duality coherences. Its overlaps are actual
filtered adjoint-preserving algebra maps \(\varphi_{ij}\) and normalized
units \(C_{ijk}\) satisfying
\begin{equation}\protect\phantomsection\label{eq:gg-independent-op-algebroid}{\begin{gathered}
\varphi_{ij}\varphi_{jk}
   =\operatorname{Ad}(C_{ijk})\varphi_{ik}, \\
C_{ijk}C_{ikl}=\varphi_{ij}(C_{jkl})C_{ijl}, \\
C_{ijk}^*C_{ijk}=1,\quad \sigma_0(C_{ijk})=1.
\end{gathered}}\end{equation} Presentations and their comparison maps may be
restricted to a common finite refinement covering the whole compact
characteristic family. In a holomorphic family this may also refine
parameter charts over a retained compact parameter set. All algebra
maps, associators, torsors, module maps, root square identifications and
calibrated frames are restricted together. Restriction is an equivalence
of the represented descent data: sections and morphisms glue by their
sheaf or module descent equations. Thus two such refinements represent
the same object and the same physical arrows; a full spin label is
retained under restriction. No global algebra map on an old overlap is
required merely because its germs have local quantizations.

The graded-symbol identifications must glue their contact sectors to a
separated germ \(Y\) along a compact characteristic. Thus \(Y\) and its
contact line \(\mathcal L=\operatorname{gr}_1\mathcal E\) are read from
the presentation; the local operator total symbols satisfy the factorial
bounds of F.1. Retain a torsor of the intrinsic sheaf \(\mathscr U_1\)
of F.3, with its translation identification. The symbol identification
\(\mathscr U_1\simeq\mathcal L^{-1}\) gives its actual affine bundle
\(A\).

Retain one compact, locally standard order-zero simple module with
reduced multiplicity-one Legendrian characteristic
\(Q\simeq\mathbb P^1\times\mathbb P^1\), with
\(\mathcal L|_Q=\mathcal O(1,1)\) and ordered rulings. Its complete
point family is derived as in H.3. The resulting contact-affine data
must satisfy Definition 5.6, including the lifted normal-bundle and
projective-extension conditions, the optical root pair
\(\beta:\mathcal E[1]^2\to\mathcal B\), the marked value \(e_o\) and the
derived numerical label \(\Lambda\).

For a passive parameter base \(B\), let \(S\to B\) be the total derived
point family and \(p:F\to S\) its sky projection, as in G.3; for a fixed
parameter \(S=M\). The module half-canonical root is \(R^2=K_{F/S}\),
with the canonical-line comparison
eq.~\ref{eq:gg-module-relative-canonical-base}. It is distinct from the
optical density root. Retain its actual square identification and
projective simple-symbol calibration, with any parameter-base factor and
its frame. Module framing means a nonzero vector of the rank-one
local-system factor on the actual cone for the whole parameter family.
Module morphisms are all semilinear module isomorphisms carrying that
vector. They are not additionally required to preserve a fixed linear
identification \(\operatorname{gr}_0\mathcal M=R\) with principal scalar
one. Normalized algebra comparison units, separately, have symbol one.

A bounded family has the conditioned finite-diagram margins of
Definition 2.1 on common analytic coordinate and operator parents,
uniform factorial bounds on the actual operator coordinate images and
inverse comparisons, actual locally standard presentations, and upper
and lower root/frame margins. Parameter bases are the finite-dimensional
germs of Definition 2.1. On a relatively compact smaller representative,
a required compact parameter set is kept in full and covered by finitely
many parameter charts. In the source construction that set includes the
retained amplitude disk containing one. Finite refinement changes the
presentation over this same set, rather than selecting a subset of its
inputs. The intrinsic incidence and projective-extension inverses have
the conditioned margins of D.7--D.9. These are geometric and analytic
conditions on supplied data, not membership in a source image.

Physical comparisons preserve the full symbolic incidence, torsor,
projective extension and optical calibration, and retain an actual full
spin lift of the induced calibrated tangent map. They can carry a
central sky to a nearby sky; point displacement is allowed. Their
operator and module comparisons obey the actual descent equations below.
Normalized comparison isomorphisms lie over the same actual physical map
and the same full spin label. Auxiliary cone framings are transported
with their local-system identification rather than fixed in an old
coordinate direction. Write \(\mathfrak O_\Lambda\) for this structured
comparison category, with physical Hom sets taken modulo normalized
comparison isomorphisms.

No metric, independent dual-twistor class or four-dimensional spacetime
solution is input in this definition. The torsor and projective
condition are substantive structure; they are not claimed to follow from
the bare algebra or a characteristic support.

\protect\phantomsection\label{sec:gg-general-relative-qct}

\subsubsection{H.2. Actual local contact
comparisons}\label{h.2.-actual-local-contact-comparisons}

\pandocstatementheading{\protect\phantomsection\label{lem:gg-general-relative-qct}Lemma H.2.}

A holomorphic family of actual contact chart maps with actual inverses
and conditioned common parents has local analytic filtered
adjoint-preserving quantizations of its inverse pullback, with uniform
factorial estimates after fixed finite restrictions. Their normalized
comparison groupoid is contractible. Supplied structured comparisons
with bounded coordinate images have actual skew logarithmic generators
of order at most one along holomorphic parameter paths. All finite
composition, identity and inverse comparison diagrams are coherent.

\emph{Proof.} Work with \(\alpha=dz+p_i\,dq^i\) and center a source
contact point at zero. If its image is \(a=(z_0,q_0,p_0)\), remove it by
\begin{equation}\protect\phantomsection\label{eq:gg-qct-contact-translation}{
 T_a(z,q,p)=(z+z_0-p_0\cdot q,\ q+q_0,\ p+p_0).
}\end{equation} This preserves \(\alpha\) and is the contact cotangent
lift of an ordinary base substitution, with its exact half-density
action.

For a contact germ fixing zero, use the dilations
\(\delta_s(z,q,p)=(s^2z,sq,sp)\). Its weighted leading map has the form
\begin{equation}\protect\phantomsection\label{eq:gg-qct-weighted-leading}{\begin{gathered}
S(z,q,p)=(\lambda z+Q(q,p),L(q,p)), \\
L^*d(p\,dq)=\lambda\,d(p\,dq), \\
dQ=\lambda p\,dq-L^*(p\,dq).
\end{gathered}}\end{equation} The last one-form is closed and linear in the
coordinates, so its quadratic primitive with zero constant is unique. In
particular \(L\) is conformal symplectic. Removing \(S\) gives a
remainder \(R\) with weighted leading identity.

The conformal symplectic group is connected. At one fixed reference
matrix choose a finite product of Lie-algebra exponentials. Near it, the
remaining near-identity matrix has a holomorphic logarithm. Their
contact paths are generated by Hamiltonians quadratic in \((q,p)\) and
linear in \(z\). A finite path and its inverse have common small spatial
tubes. Symplectization gives degree-one Hamiltonians; F.1 and Theorem
F.4 quantize these paths with the required bounds.

For the nonlinear remainder take the actual path
\begin{equation}\protect\phantomsection\label{eq:gg-qct-weighted-path}{
 R_s=\delta_s^{-1}R\delta_s,\qquad R_0=\mathrm{id}.
}\end{equation} The remainders in \((q,p)\) have weight at least two and
that in \(z\) has weight at least three. Their scaled Taylor terms
therefore have strictly positive powers of \(s\). After one weighted
spatial restriction, both this series and its actual inverse converge
normally on an \(s\)-disk containing \([0,1]\). The contact multiplier
is one at \(s=0\) and has a common unit margin after restriction. The
given map and inverse Cauchy bounds thus supply an actual classical
tube. Read the Hamiltonian of the inverse-pullback logarithmic
derivative, with the sign convention of §5, and apply F.1 and Theorem
F.4. Composing the three factors quantizes exactly the supplied contact
map.

Bounded weighted first jets and inverse jets have compact closure in the
finite-dimensional conformal symplectic group. Finitely many reference
neighborhoods suffice, with maximum constants and a common finite
spatial shrink. The other constants depend on the stipulated map,
inverse and parent margins. This proves uniformity on the conditioned
diagrams, without a global matrix logarithm or a strict global
quantizer.

For a supplied comparison \(\varphi\) and a constructed controlled
comparison \(k\) with the same full symbols, first use the left quotient
\begin{equation}\protect\phantomsection\label{eq:gg-qct-controlled-left-quotient}{
 \psi=k^{-1}\varphi,\qquad
 \psi(y_i)-y_i=k^{-1}\bigl(\varphi(y_i)\bigr)-y_i .
}\end{equation} The known bounds for \(k^{-1}\) act on the stipulated
bounded coordinate-image families of \(\varphi\). Lemma F.8 gives a
normalized source unit \(Q\) with \(\psi=\operatorname{Ad}(Q)\). Hence
\(\varphi=\operatorname{Ad}(P)k\) with \(P=k(Q)\), including uniform
bounds for \(P\) and its inverse. No uniform action of the unknown
\(\varphi\) on an arbitrary symbol family was presumed. Two supplied
quantizers are each compared with the same constructed \(k\); their
target units \(P,P'\) give the comparison \(P'P^{-1}\). Lemma F.8's
uniqueness proves all finite coherence equations and glues the local
parameter choices up to that unique comparison.

For the generator assertion, let \(\dot k_t=\operatorname{ad}(B_t)k_t\),
initialized by \(k_0=\mathrm{id}\) independently of a passive parameter
\(\mu\). Differentiating \(k_t^{-1}\partial_\mu k_t\) gives
\begin{equation}\protect\phantomsection\label{eq:gg-qct-passive-generator}{
 \partial_\mu k_t\,k_t^{-1}=\operatorname{ad}(A_\mu(t)),
 \qquad
 A_\mu(t)=k_t\int_0^t k_u^{-1}(\partial_\mu B_u)\,du .
}\end{equation} Theorem F.4 and one parameter Cauchy gap give its
order-one skew symbol bounds. Ordinary parameter-dependent substitutions
have their explicit first-order half-density generators; fixed reference
factors are independent of these parameters. Finite products use the
transported sum. Finally, from \(\varphi_s=\operatorname{Ad}(P_s)k_s\)
its actual generator is
\begin{equation}\protect\phantomsection\label{eq:gg-qct-given-generator}{
 A_s^\varphi=P_s'P_s^{-1}
                  +\operatorname{Ad}(P_s)A_s^k .
}\end{equation} The first term is skew of order at most minus one, and
the second is skew of order at most one. This is the generator of the
supplied comparison, rather than a new principal lift. ∎

\protect\phantomsection\label{sec:gg-independent-module-points}

\subsubsection{H.3. Complete point objects from a compact
characteristic}\label{h.3.-complete-point-objects-from-a-compact-characteristic}

\pandocstatementheading{\protect\phantomsection\label{prop:gg-independent-module-points}Proposition H.3.}

The compact characteristic in Definition H.1 determines a complete
smooth four-dimensional support germ and its actual calibrated
supported-module family. With proper smooth supports and actual locally
standard holomorphic presentations, its framed module comparisons are
unique. The construction has the same conditioned analytic bounds and
finite jets of holomorphic families as Appendix G. It does not classify
arbitrary flat relative regular modules.

\emph{Proof.} The contact normal sequence is
\begin{equation}\protect\phantomsection\label{eq:gg-independent-op-normal-sequence}{
 0\longrightarrow\Omega_Q^1\otimes\mathcal L
   \longrightarrow N_{Q/Y}\longrightarrow\mathcal L
   \longrightarrow0,\qquad
 \Omega_Q^1\otimes\mathcal L
     =\mathcal O(-1,1)\oplus\mathcal O(1,-1).
}\end{equation} Its left term has zero cohomology in every degree. Thus
\(H^0N=H^0\mathcal O(1,1)\) has dimension four and \(H^1N=0\). Lemma D.1
gives the complete support germ; D.7 also proves polarized-quadric
preservation and Legendrianness. This derives the family from the single
compact characteristic.

For a smooth parameter germ \(B\) of dimension \(d\), the relative
contact atlas has an ordinary smooth total ambient \(Y_{\mathrm{tot}}\).
The central total normal sequence has quotient \(\mathcal O_Q^d\) and
kernel \(N_{Q/Y_0}\). Its \(H^1\) vanishes and its \(H^0\) has dimension
\(4+d\), so Lemma D.1 gives the total complete family. The same
polarization argument of D.7 makes nearby members compact quadrics.
Their holomorphic functions are constant; proper base change for
\(\mathcal O\) therefore makes the ambient projection descend to a
holomorphic map of the complete base to \(B\). Its tangent is the
surjection \(H^0(N_{Q/Y_{\mathrm{tot}}})\to\mathbb C^d\). The finite
implicit theorem gives the actual relative four-dimensional family.
Denote its total point base by \(S\); thus \(S\to B\) has relative
dimension four and the sky projection is \(p:F\to S\), consistently with
G.3 and H.1. This uses the fixed ordinary total ambient, not a general
relative theorem over nilpotent ambients.

The full normal evaluation is the first-jet evaluation. In affine
coordinates \((\varrho,\zeta)\) its matrix and kernel are
\begin{equation}\protect\phantomsection\label{eq:gg-independent-op-jet-evaluation}{
 \begin{pmatrix}
 1&\varrho&\zeta&\varrho\zeta\\
 0&1&0&\zeta\\
 0&0&1&\varrho
 \end{pmatrix},\qquad
 \ker=\mathbb C(\varrho\zeta,-\zeta,-\varrho,1).
}\end{equation} It has rank three and its kernel lines form the entire
nondegenerate Segre quadric. The universal incidence evaluation
\(q:F\to Y_{\mathrm{tot}}\) has relative rank five over \(B\) and total
rank \(5+d\); both differentials have a line kernel. D.7 identifies it
with the same conformal null incidence and short rays, before Einstein
recognition.

For a fixed actual algebroid, flatten the moving skies by \(L_i(x)\).
Along \(x=ax\), the comparison \(\psi_{ij}=L_i\varphi_{ij}L_j^{-1}\) has
actual generator \(A_i-\psi_{ij}(A_j)\) because \(\varphi_{ij}\) is
point independent. It preserves the reference Legendrian, so its
principal generator vanishes there. Apply Lemma G.2 to
\(\psi_a\psi_0^{-1}\) and the given initial intertwiner. This is the
construction of G.3 with \(B_{ij}=0\), with no flat-source assumption.

For parameterized algebroids, Lemma H.2 supplies their actual
\(B_{ij}\), so the flattened generator is
\(A_i+L_i(B_{ij})-\psi_{ij}(A_j)\). Contraction of the small parameter
coordinate germ, with the remaining variables passive, gives joint Hom
generators from one fixed reference. Root-symbol normalization gives
exactly \(\tau_{ij}\tau_{jk}=C_{ijk}\tau_{ik}\). The finite compact
cover, tangent transport and inverse estimates are those of G.1--G.3.

For a holomorphic family of supports, its classifying map is unique
locally. In the complete family choose four nearby grid points of one
affine sky chart. Evaluation of \(H^0\mathcal O(1,1)\) there is an
isomorphism, so their four contact graph values give coordinates by the
finite inverse theorem. Equal embedded supports have the same four
values and hence the same point parameter. This uses actual graph
evaluations in addition to the tangent dimension.

Over a fixed support family the Hom line is the actual rank-one
\(\mathcal O_S\) line of G.4. Root frames first construct and glue it
and a local-symbol-normalized comparison. Locally on the parameter base,
a single global unit along the complete sky/cone then adjusts its
section to carry the entire local-system vector. This scalar is not
constrained again by a fixed linear root identification: frames \(e\)
and \(2e\) on the same module are compared by \(2\,\mathrm{id}\).

The auxiliary cone section is fixed on the derived family and pulled
back, or its vector is read by the canonical transport in the simply
connected cone local system. It is not an independent sky-direction
modulus. Relative root choices retain their actual square
identifications and base factors; on the parameter germ their ratio is a
base two-torsion line, whose calibrated root comparison is unique. These
facts give the stated framed comparisons and parameter restrictions. No
general microdifferential-module flatness or nonlifted Artin-family
classification has entered the proof. ∎

\protect\phantomsection\label{sec:gg-independent-op-equivalence}

\subsubsection{H.4. Both inverse comparisons on the independent
domain}\label{h.4.-both-inverse-comparisons-on-the-independent-domain}

\pandocstatementheading{\protect\phantomsection\label{thm:gg-independent-op-equivalence}Theorem H.4.}

Taking full symbols, compact characteristics and the intrinsic
first-unit torsor gives an equivalence between the independently
admitted operator data of Definition H.1 and the calibrated affine data
of Definition 5.6. It has actual inverse comparisons on the whole
independent domains, coherent on conditioned bounded holomorphic
families and finite physical diagrams. It preserves the selected optical
section, numerical scalar, all physical point/full-spin Hom classes and
the coupled Lorentz operation.

\emph{Proof.} Write \(\mathfrak C_\Lambda\) for the calibrated affine
category. Taking the classical data while forgetting the remaining
operator choices defines
\(F:\mathfrak O_\Lambda\to\mathfrak C_\Lambda\). Proposition H.3
supplies its complete incidence. Definition H.1 supplies the
affine/projective admission of Definition 5.6.

Conversely take any \(\mathfrak C_\Lambda\) object. We first choose a
finite refinement on which every overlap has a single quantization. In a
bounded family retain the required compact passive-parameter set
\(B_0\). At every vertex \(v\) of the fixed finite physical diagram
choose a compact point core \(S_v^0\Subset S_v\) containing all the
required endpoints and composite images over \(B_0\). Let
\begin{equation}\protect\phantomsection\label{eq:gg-quantization-whole-diagram-core}{
 K=\bigcup_v q_v\bigl(p_v^{-1}(S_v^0)\bigr).
}\end{equation} Each sky projection \(p_v\) is proper, so \(K\) is
compact; take the finite disjoint union of the vertex ambients when
necessary. The point cores keep the entire stipulated diagram and
passive-parameter set. All later restrictions concern neighborhoods of
this same \(K\), including its moving skies, rather than selecting a
subset of the original inputs. Enlarge it to a compact neighborhood
\(K^+\) inside the supplied parents, with \(K\) in its interior. Choose
finitely many original chart parents \(U_i\) and closed subsets
\(K_i\Subset U_i\) covering \(K^+\). At \(y\in K^+\) put
\(I_y=\{i:y\in K_i\}\). Lemma H.2 gives, near \(y\), quantizers and
inverses for all pairs of original charts with indices in this finite
set. Choose nested neighborhoods \(V_y\Subset V_y^+\) inside all these
quantization parents, and disjoint from every \(K_j\) with
\(j\notin I_y\). Finite intersections reserve the parents for each map
and its inverse. All choices are in the total space, so they include the
required local parameter neighborhoods.

A finite closed-star refinement \(\{W_a\}\) of \(\{V_y\}\) covers \(K\)
and has each closed star contained in one \(V_y\). To obtain it, take a
finite subcover on the compact metrizable set \(K^+\) and a Lebesgue
number \(\lambda>0\). Put
\(d=\operatorname{dist}(K,Y_{\mathrm{tot}}\setminus K^+)>0\) and use
ambient metric balls centered on \(K\) of radius
\(\rho<\min(\lambda/8,d/4)\). Each closed star has diameter at most
\(6\rho<\lambda\) and lies in \(K^+\). A finite subcover of \(K\)
therefore has the required property. Choose \(x_a\in W_a\cap K\) and an
original index \(i(a)\) with \(x_a\in K_{i(a)}\). If
\(W_a\cap W_b\ne\varnothing\), a single \(V_y\) contains both sets and
their closures. Since it contains \(x_a,x_b\) and avoids \(K_j\) for
\(j\notin I_y\), both \(i(a),i(b)\) belong to \(I_y\). Its chosen
quantizer for that pair is therefore defined on a parent of the entire
new overlap. On that overlap restrict this quantizer to obtain
\(\varphi_{ab}\), and define the reverse as its literal inverse. The
closed-star buffer, taken before restriction, is a positive Cauchy
buffer for these maps. For the specified finite composite diagram,
intersect each source neighborhood with the preimages of the assigned
target parents along its composites, and take a finite closed shrinking
once more. The compact cores are retained because each stipulated
composite already has its image in its target parent with positive
buffer. This constructs maps on the new overlaps; coherence is the
following separate step.

On a triple the classical cocycle makes
\(D_{abc}=\varphi_{ab}\varphi_{bc}\varphi_{ac}^{-1}\) have identity full
symbols. Lemma F.8 supplies local normalized units for it. Their
uniqueness makes them agree on intersections, so they give one
\(C_{abc}\) on the entire triple, with uniform bounds from the finite
compact cover. The two units \(C_{abc}C_{acd}\) and
\(\varphi_{ab}(C_{bcd})C_{abd}\) implement the same inner automorphism
on a quadruple and have symbol one and unitary adjoint; injectivity in
Lemma F.8 makes them equal. Thus the full algebroid equation, not only
the pairwise maps, holds on this refinement. The constants are maxima
over finitely many parents and the buffers and unit margins are their
positive minima. They do not depend on homogeneous order. This uses
local quantization and normalized-inner uniqueness, not a claim that the
twisted-algebra classification of (Polesello 2007, Corollary 1.6)
trivializes every old overlap.

Keep the same torsor under \(\mathscr U_1\simeq\mathcal L^{-1}\), with
its restricted cochains. Pull this finite operator cover to the complete
incidence and choose H.3's finite graph/module charts over it. Carry any
extra chart indices in all pulled-back algebra data as well. Restrict
the relative root, its square identification and all base factors to
those same charts. Proposition F.3 supplies the static central module
and H.3 constructs its family there. The root transitions obey their
exact cocycle; G.3's symbol normalization then gives
\(\tau_{ab}\tau_{bc}=C_{abc}\tau_{ac}\) with these same \(C_{abc}\). Its
tangent transports and their inverses have the retained buffers, so
taking the finite maxima gives the common module bounds. The
whole-parameter calibrated vector is restricted, not reselected at each
parameter. This defines \(Q:\mathfrak C_\Lambda\to\mathfrak O_\Lambda\)
without first decoding an Einstein metric. The composition \(FQ\) is the
original contact-affine object under the restriction equivalence of H.1,
with its tautological complete-family identifications.

For an independently supplied operator object compare it to \(QF\) on a
common chart refinement obtained by applying the same closed-star
construction to the union of their finite chart families and identity
comparisons. This retains the same compact parameter set and admits all
maps in the finite comparison diagram on buffered parents. Lemma H.2
quantizes the actual identity contact comparison in these possibly
different chart coordinates. Its normalized overlap units satisfy
\begin{equation}\protect\phantomsection\label{eq:gg-independent-op-two-inverse-descent}{
 \begin{aligned}
 k_i\varphi^{s}_{ij}
   &=\operatorname{Ad}(T_{ij})\varphi^{t}_{ij}k_j,\\
 T_{ij}\varphi^{t}_{ij}(T_{jk})C^{t}_{ijk}
   &=k_i(C^{s}_{ijk})T_{ik}.
 \end{aligned}
}\end{equation} Indeed the two expansions of
\(k_i\varphi^s_{ij}\varphi^s_{jk}\) have the same inner action. Both
implementers are normalized, so Lemma F.8 makes them equal. Both
associators are retained.

The module comparisons \(F_i\) obey
\begin{equation}\protect\phantomsection\label{eq:gg-independent-op-module-counit}{
 F_i\tau^s_{ij}=T_{ij}\tau^t_{ij}F_j .
}\end{equation} Construct their Hom line by the actual flattened tangent
transport of H.3, normalize its projective root symbols and then its
whole local-system vector. The two sides have the same semilinearity and
global normalization, so scalar injectivity proves equality. The
\(\mathscr U_1\)-torsor action is literally the original affine
comparison, since normalized inner changes act trivially on this sheaf.

If \(k'_i=\operatorname{Ad}(P_i)k_i\), the two normalized overlap
implementers give
\begin{equation}\protect\phantomsection\label{eq:gg-independent-op-counit-uniqueness}{
 T'_{ij}\varphi^t_{ij}(P_j)=P_iT_{ij}.
}\end{equation} Their module maps are compared by \(P_iF_i\) and framed
Hom uniqueness. Hence \(QF\) is canonically equivalent to the supplied
object. On a further common refinement all comparison equations restrict
literally, and any two choices are compared by the unique normalized
units and framed Hom maps just constructed. Consequently these equations
descend independently of the auxiliary cover. The unit/counit triangle,
composition and finite naturality diagrams follow from these exact
uniqueness statements. This compares actual representatives, rather than
just their original records.

The symbols of an admitted physical comparison preserve the entire
affine/projective calibration. D.8 and Proposition 5.4 recover the
selected metric naturally, so that map is an isometry. Its full spin
lift gives the distinct label of Theorem 8.3. Conversely Lemma H.2 and
H.3 lift each such physical comparison. Thus the physical Hom classes
agree, including spin-center distinctions not detectable by the scalar
action.

The coupled antilinear construction and normalized square of Theorem G.5
use the same comparisons and compatible root/local-system frames. The
preceding uniqueness transports its square and finite Real diagrams. On
the declared Lorentz slice it gives Proposition 8.4's physical real
structure. Every parameter comparison used the given common parents,
finite refinements and unit margins; finite jets are taken from those
holomorphic families. ∎

\protect\phantomsection\label{sec:gg-independent-op-active-comparison}

\subsubsection{H.5. The same source operation and preservation
obstruction}\label{h.5.-the-same-source-operation-and-preservation-obstruction}

\pandocstatementheading{\protect\phantomsection\label{prop:gg-independent-op-active-comparison}Proposition H.5.}

The equivalence of Theorem H.4 identifies the actual operator source
construction with the same coupled contact-affine law, including its
higher preservation correction. Its independently read holomorphic
obstruction is the cochain of D.9, with the same derivative \(o_m\) and
the same inverse \(L_m\) on the zero-record complement. Its Einstein
consequence is Proposition 5.4.

\emph{Proof.} Use F.1's positive differential operator and its complete
negative symbol formula, with the fixed skew lift. The latter is a
direct formula in the original \(U^-_{r;a,j}\) and the original twistor
coordinates and derivative momenta. The dual class \(G\) used to derive
it is not a separate execution primitive.

At a trial operator stage, the regular compact seed gives the
characteristic-family construction of H.3 and F.3 gives the intrinsic
first-unit translation sheaf. Retain the calibrated affine reduction of
the preceding corrected stage and choose provisional section-state lifts
at the next order. Lemma I.1 reads their actual defect on the freshly
recomputed incidence; its vanishing supplies the next affine reduction.
A particular torsor is not manufactured by the bare algebra.
Projective-extension admission is not assumed before this defect
vanishes. Its states and matrices are exactly those of D.9. In
particular,
\begin{equation}\protect\phantomsection\label{eq:gg-independent-op-same-obstruction}{
 T_{ij}=(G_i\circ\Phi_{ij})^{-1}B_{ij}G_j,\qquad
 (G_i\circ\Phi_{ij})^{-1}r_{ij}=e_1-T_{ij}e_1.
}\end{equation} Thus its entries are \(1-T_{11}\) and \(-T_{21}\), the
same actual connecting cochain. Its derivative on the fixed zero-record
complement is \(o_m\). The same \(L_m\) and original record correction
remove it, followed by fresh incidence, root, matrices and scalar
recomputation. Lemma 5.1's cubic-gauge induction applies to this actual
source execution; it is not inferred from a first point jet.

At finite source stages the incidence graphs are the explicit
finite-coefficient construction of Proposition 5.2. Lift each corrected
Hamiltonian by F.1. Theorems F.1 and F.6 give finite and analytic
structured comparisons, Lemma F.8 supplies their actual normalized
comparisons, and Appendix G supplies the full module families. Their
full symbols and first-unit torsor are the same corrected tuple.
Proposition F.7 compares the finite active recursion with their
amplitude jets by compatible normalized filtered inner equivalence. This
requires neither equality nor convergence of another auxiliary infinite
quantum section.

The mixed principal commutator is \(t\{h,k\}\), while the raw
order-minus-one Weyl correction of F.3 remains. No strict affine Lie
embedding has replaced that correction. These are actual operations in
the same fixed projective Heisenberg half-density model, with its
intrinsic derived sheaves and point objects.

The projective extension and calibrated root of Definition H.1 supply
the whole optical section. Proposition 5.4 gives Einsteinity with the
same numerical scalar, and \(L_m\) is the same inverse of the optical
obstruction on the zero-record complement. Proposition 9.2 exhibits a
nonzero obstruction before correction; the flat-scale calculation of
Proposition I.3 identifies the role of the calibrated scale data. ∎

Theorem H.4 and Proposition H.5 give both inverse comparisons and
identify the active preservation law on the independent operator domain.
Appendix I identifies its infinity structure with the normal tractor and
rank-four local-twistor forms, and assembles the primitive, admission
and point comparisons used in Theorem 9.11.

\protect\phantomsection\label{app:gg-original-positive-comparison}

\subsection{Appendix I. Infinity structure and the original twistor
model}\label{appendix-i.-infinity-structure-and-the-original-twistor-model}

\protect\phantomsection\label{sec:gg-preadmission-reduction}

\subsubsection{I.1. The calibrated reduction before
admission}\label{i.1.-the-calibrated-reduction-before-admission}

\pandocstatementheading{\protect\phantomsection\label{lem:gg-preadmission-reduction}Lemma I.1.}

The intrinsic first-unit translation sheaf does not select a particular
affine torsor on an uncorrected trial. At that stage the obstruction to
extending the previously calibrated optical reduction is the cochain of
D.9. Its vanishing supplies the actual next calibrated affine torsor.

\emph{Proof.} Recompute the trial compact incidence and its optical
transitions \(B_{ij}\). Choose provisional local section states
\(s_i=(\sigma_i,\dot\sigma_i)^T\), lifting the preceding corrected
reduction, in the retained Cauchy charts where \(\sigma_i\) is a unit.
Use exactly
\(G_i=\left(\begin{smallmatrix}\sigma_i&0\\\dot\sigma_i&\sigma_i^{-1}\end{smallmatrix}\right)\),
so \(\det G_i=1\). The fixed unit margins persist under these
highest-order lifts. Set
\begin{equation}\protect\phantomsection\label{eq:gg-preadmission-actual-states}{
 r_{ij}=s_i\circ\Phi_{ij}-B_{ij}s_j,\qquad
 T_{ij}=(G_i\circ\Phi_{ij})^{-1}B_{ij}G_j .
}\end{equation} Then direct multiplication gives
\begin{equation}\protect\phantomsection\label{eq:gg-preadmission-exact-defect}{
 (G_i\circ\Phi_{ij})^{-1}r_{ij}
       =e_1-T_{ij}e_1
       =\begin{pmatrix}1-T_{11}\\-T_{21}\end{pmatrix}.
}\end{equation} When this cochain is zero, \(T_{ij}\) fixes \(e_1\). Its
determinant is the contact-line factor \(c_{ij}\), so it has the form
\(\left(\begin{smallmatrix}1&b_{ij}\\0&c_{ij}\end{smallmatrix}\right)\).
With the companion and ratio convention of D.9,
\(a_{ij}=-b_{ij}/c_{ij}\) and \(\tau_i=c_{ij}^{-1}\tau_j+a_{ij}\). The
matrix cocycle is exactly the affine cocycle with all moving arguments.
Before vanishing, a general \(T_{ij}\) is not such an affine transition.
F.3 provides its translation sheaf
\(\mathscr U_1\simeq\mathcal L^{-1}\); it does not remove this reduction
obstruction or manufacture a distinguished torsor from the bare ring.
The same calculation applies at a highest finite stage. ∎

\protect\phantomsection\label{sec:gg-projective-infinity-structure}

\subsubsection{I.2. The intrinsic infinity-structure
predicate}\label{i.2.-the-intrinsic-infinity-structure-predicate}

\pandocstatementheading{\protect\phantomsection\label{def:gg-projective-infinity-structure}Definition I.2.}

The projective infinity structure on an operator presentation is its
actual first-unit affine torsor, the induced projective optical
extension, optical density root and central calibration. Its admission
predicate is the whole-neighborhood holomorphic extension condition of
Definition H.1, with its derived numerical scalar. The predicate is
defined on these primitive holomorphic data before Einstein recognition.
A compact point object means the calibrated locally standard supported
simple module of H.1--H.3.

This definition specifies the projective alternative used here. It does
not assert a literal equivalence with every presentation of the
speculative unreduced Palatial algebra.

\pandocstatementheading{\protect\phantomsection\label{prop:gg-geometric-infinity-structure}Proposition I.3.}

The admitted projective infinity structure is equivalent to the
calibrated nonvanishing parallel normal tractor of norm \(-\Lambda/3\).
Its local-twistor spin realization has the infinity tensors and
numerical normalization of the original gravitational construction,
including their degeneration at \(\Lambda=0\). The correspondence is
natural for the actual physical operator comparisons and their full spin
lifts. On the Lorentz slice it has the original local phase-scaled ray
and point fibers.

\emph{Proof.} D.8 and Proposition 5.4 identify the holomorphic
projective extension and root with a nonvanishing density \(\sigma\)
satisfying
\begin{equation}\protect\phantomsection\label{eq:gg-original-infinity-almost-einstein}{
 (\nabla_a\nabla_b\sigma+P_{ab}\sigma)_0=0 .
}\end{equation} Its normal prolongation is parallel, has norm
\(-\Lambda/3\), and selects \(g=\sigma^{-2}g_{\mathrm{aux}}\).
Conversely the parallel section with the retained unit scale component
gives this equation and D.8's projective extension. The two
constructions are inverse comparisons preserving the actual root and
central value.

Here is an explicit spin and sign comparison. Put \(c=\Lambda/6\) and
use the selected Einstein scale, where the standard Schouten tensor is
\(P=cg\). Let
\(J=\left(\begin{smallmatrix}0&1\\-1&0\end{smallmatrix}\right)\) in the
marked spin bases. Write \(\theta_v\) for the tangent soldering matrix,
with \(g(v,w)=\operatorname{tr}(\theta_v\theta_w^\sharp)\) and
\(\theta^\sharp=\operatorname{adj}(\theta)\). Represent the raised
middle tractor component by the matrix \(\mu\). Define
\begin{equation}\protect\phantomsection\label{eq:gg-original-infinity-spin-connection}{
 \Omega_v=
 \begin{pmatrix}0&i\theta_v\\ic\theta_v^\sharp&0\end{pmatrix},
 \qquad
 \mathcal F(\sigma,\mu,\rho)=
 \begin{pmatrix}
 \rho J&-iJ\mu\\-i\mu^TJ&\sigma J
 \end{pmatrix}.
}\end{equation} The matrix \(\mathcal F\) is a covariant alternating
form on the rank-four local-twistor bundle \(\mathcal T_{\mathrm{LT}}\),
with marked conformal splitting \(S\oplus S'\). Its Spin representation
differs from that of the middle bundle \(W_{\mathrm{scr}}\simeq TM\),
which has type \(S\otimes S'\). For its induced covariant connection
\(D^{\wedge}K=\nabla K-\Omega^TK-K\Omega\), ordinary two-by-two
multiplication gives
\begin{equation}\protect\phantomsection\label{eq:gg-original-infinity-exterior-square}{
 \begin{aligned}
 D_v^\wedge\mathcal F(\sigma,\mu,\rho)
   =\mathcal F\bigl(&d_v\sigma-g(v,\mu),\\
                   &\nabla_v\mu+(c\sigma+\rho)\theta_v,\\
                   &d_v\rho-cg(v,\mu)\bigr),\\
 2\operatorname{Pf}\mathcal F(\sigma,\mu,\rho)
   &=2\sigma\rho+g(\mu,\mu).
 \end{aligned}
}\end{equation} For example \((\theta^\sharp)^TJ=J\theta\) and
\(JX+X^TJ=\operatorname{tr}(X)J\) prove the connection identity; the
Pfaffian is \(\sigma\rho+\det\mu\). Thus this is exactly the normal
tractor connection and its quadratic form in §4, including its
soldering, scale and norm signs.

The selected tractor is \((1,0,-c)\) in this scale. Its covariant spin
form and volume dual are therefore
\begin{equation}\protect\phantomsection\label{eq:gg-original-infinity-product}{
 I_{\alpha\beta}=\operatorname{diag}(-cJ,J),\qquad
 I^{\alpha\beta}=\operatorname{diag}(J,-cJ),\qquad
 I_{\alpha\beta}I^{\beta\gamma}
                   =c\,\delta_\alpha{}^\gamma .
}\end{equation} Both forms are parallel, and
\(2\operatorname{Pf}(I)=-2c=-\Lambda/3\). The initial full spin marking
fixes the volume used in the dual and the entire identification. Their
ranks are four for \(c\ne0\) and two at zero.

The curvature-component dictionary is
(\ref{eq:gg-original-infinity-curvature-dictionary}); in particular,
\(\Lambda_{\mathrm P}=-\Lambda\) and \(P^{\mathrm P}=P\), while the
prescribed Einstein coefficient \(\Lambda\) is unchanged. The spinor
component dictionary requires a separate declaration. For this
comparison only, use the numerical arrays
\begin{equation}\protect\phantomsection\label{eq:gg-original-infinity-external-spinor-convention}{
 \begin{gathered}
 [\varepsilon^{\mathrm P}_{AB}]
 =[\varepsilon_{\mathrm P}^{AB}]
 =[\varepsilon^{\mathrm P}_{A'B'}]
 =[\varepsilon_{\mathrm P}^{A'B'}]=J,\\
 u^A=\varepsilon_{\mathrm P}^{AB}u_B,\qquad
 u_A=u^B\varepsilon^{\mathrm P}_{BA},
 \end{gathered}
}\end{equation} with the same order in the primed operations. Raising
acts on the left and lowering on the right. On coefficient columns these
are respectively \(J\) and \(-J\), so they are inverse operations. In
particular,
\(\varepsilon^{\mathrm P}_{AB}\varepsilon_{\mathrm P}^{CB}=\delta_A{}^C\),
whereas contraction in the order \(AB,BC\) gives \(-\delta_A{}^C\).
These comparison arrays are distinct from the arrays in
(\ref{eq:gg-spin-index-convention}): the lower arrays agree numerically,
the upper arrays have opposite signs, and the lowering operation uses
the opposite side. Thus an upper \(\varepsilon_{\mathrm P}\) in a
primary block formula cannot be replaced by the upper \(\epsilon\) of
§2.

This reading agrees with the primary coordinate normalization. Its
unnumbered formulas for \(\delta\pi\) and \(\tfrac12\mathrm d\delta\pi\)
give
\(\delta\pi=\varepsilon_{\mathrm P}^{F'G'}\pi_{F'}\mathrm d\pi_{G'}\)
and
\(\tfrac12\mathrm d\delta\pi=\mathrm d\pi_{0'}\wedge\mathrm d\pi_{1'}\),
which fix \(\varepsilon_{\mathrm P}^{0'1'}=1\) in those coordinates
(Penrose 2015, 5). The primary article does not explicitly prescribe the
side of the two-spinor raising and lowering operations. The operations
just declared specify our component comparison rather than attribute an
unstated historical convention to it.

Here is the calibrated basis map. If \((b_0,b_1,b_{0'},b_{1'})\) is the
ordered rank-four frame used above, identify the primary ordered twistor
column \((\omega^0,\omega^1,\pi_{0'},\pi_{1'})^T\) with
\begin{equation}\protect\phantomsection\label{eq:gg-original-infinity-external-basis}{
 C(\omega,\pi)
   =\omega^0b_0+\omega^1b_1
                         +\pi_{0'}b_{0'}+\pi_{1'}b_{1'} .
}\end{equation} Thus \([C]=\operatorname{Id}_4\) in the indicated
frames. The second slot is identified using the covariant primed
coefficient column itself, not a column obtained by applying the raising
rule of §2. The map on covectors is the inverse transpose, and the map
on bivectors is the exterior square of \(C\); these are not independent
choices. Calibrate the two twistor volumes by
\(\varepsilon^{\mathrm P}_{0123}=1=\varepsilon_{\mathrm P}^{0123}\), as
in the primary unnumbered duality formulas (Penrose 2015, 7). Since
\(\det C=1\), this is the volume fixed by the existing full spin
marking.

Write \(\lambda_{\mathrm P}=\Lambda_{\mathrm P}/6=-c\). In these ordered
coordinates the primary displays become
\begin{equation}\protect\phantomsection\label{eq:gg-original-infinity-external-volume-forms}{
 \begin{aligned}
 I^{\mathrm P}_{\alpha\beta}
   &=\operatorname{diag}
          (\lambda_{\mathrm P}\varepsilon^{\mathrm P}_{AB},
                                  \varepsilon_{\mathrm P}^{A'B'})
     =\operatorname{diag}(-cJ,J),\\
 I_{\mathrm P}^{\alpha\beta}
   &=\operatorname{diag}
          (\varepsilon_{\mathrm P}^{AB},
                        \lambda_{\mathrm P}\varepsilon^{\mathrm P}_{A'B'})
     =\operatorname{diag}(J,-cJ),\\
 I_{\mathrm P}^{\alpha\beta}
   &=\tfrac12\varepsilon_{\mathrm P}^{\alpha\beta\rho\sigma}
                                    I^{\mathrm P}_{\rho\sigma},
 \qquad
 I^{\mathrm P}_{\alpha\beta}I_{\mathrm P}^{\beta\gamma}
      =-\lambda_{\mathrm P}\delta_\alpha{}^\gamma .
 \end{aligned}
}\end{equation} Indeed the volume dual has \(01\) entry
\(I^{\mathrm P}_{23}=1\) and \(23\) entry
\(I^{\mathrm P}_{01}=\lambda_{\mathrm P}\); all mixed entries vanish.
Its product with the covariant matrix is
\(-\lambda_{\mathrm P}\operatorname{Id}_4\), since
\(J^2=-\operatorname{Id}_2\). Conversely the same lower volume recovers
the covariant form. Starting with the primary normalization
\(\varepsilon_{\mathrm P}^{0'1'}=1\), these duality and product displays
force the other three alternating arrays to be \(J\) when
\(\lambda_{\mathrm P}\ne0\); the same fixed arrays give their rank-two
degeneration at zero (Penrose 2015, 7--8). Consequently \(C\) carries
both tensors together, their volume duality, their product and their
ranks to (\ref{eq:gg-original-infinity-product}).

The transport comparison uses the same basis map. Read the primary
tangent matrix \([t^{AA'}]\) as \(\theta_v\). Since
\(P^{\mathrm P}=cg\), its lower spinor components are
\(P^{\mathrm P}_{AA'BB'}=c\varepsilon^{\mathrm P}_{AB}\varepsilon^{\mathrm P}_{A'B'}\).
For an arbitrary
\(\theta_v=\left(\begin{smallmatrix}a&b\\d&e\end{smallmatrix}\right)\),
contraction gives
\begin{equation}\protect\phantomsection\label{eq:gg-original-infinity-external-transport}{
 \left[
  \sum_{A,A'}t^{AA'}P^{\mathrm P}_{AA'BB'}
 \right]_{B'B}
   =c\begin{pmatrix}e&-b\\-d&a\end{pmatrix}
   =c\theta_v^\sharp .
}\end{equation} Thus the primary equations
\[\nabla_v\omega^B=-it^{BB'}\pi_{B'}\] and
\(\nabla_v\pi_{B'}=\allowbreak-it^{AA'}P^{\mathrm P}_{AA'BB'}\omega^B\) are
\(\nabla_v(\omega,\pi)^T=\allowbreak-\Omega_v(\omega,\pi)^T\) in this calibration
(Penrose 2015, 12). The initialized map \(C\) therefore intertwines the
connections as well as both infinity tensors. This comparison uses no
metric rescaling, no inverse infinity form and no division by
\(\Lambda\).

Along every short null ray the initialized local-twistor ODE has a
four-dimensional solution space and carries these same forms. Conversely
constancy of a calibrated tractor along all null directions implies its
full parallelism, since the null vectors span each tangent space. This
explains the original raywise infinity structure without imposing any
Weyl-type restriction.

For a normalization test put \(q=\eta_{ab}x^ax^b\) and \(\sigma_C=1+Cq\)
on a flat conformal carrier. Its Hessian is pure trace, its prolongation
has \(\rho=-2C\), and
\begin{equation}\protect\phantomsection\label{eq:gg-original-infinity-scale-test}{
 |d\sigma_C|_\eta^2+2\sigma_C\rho=-4C,
 \qquad \Lambda_C=12C .
}\end{equation} Every member has central value one and central gradient
zero. Thus bare contact geometry and a bare standard local operator
algebra do not determine this infinity structure or selected scale. The
retained torsor, whole projective extension and calibration are
substantive.

An actual structured comparison carries the full symbols, torsor,
projective extension and root, and hence their actual optical section.
The derived metric comparison is an isometry; its retained full spin
lift carries the normal and rank-four spin connections and the displayed
forms. Normalized inner comparison units act trivially on
\(\mathscr U_1\) and its classical data, while their actual point-module
and associator actions remain in H.4's descent. This gives coherent
naturality on the full structured category.

Finally on the Lorentz slice a nonzero primed spinor gives
\(p_{AA'}=\overline\pi_A\pi_{A'}\). Its parallel transport along short
null geodesics defines the local phase-scaled ray space: four real
spacetime variables and four real nonzero-spinor variables, modulo one
geodesic parameter, give dimension seven. The phase quotient is
dimension six and its positive momentum-scaling quotient is dimension
five. At a fixed point the fibers are respectively
\(\mathbb C^2\setminus0\) and \(S^2\), with the intermediate positive
null-momentum cone. The real locus of the coupled sky operation is this
same \(S^2\) inside the complex quadric. The actual full spin marking
retains its phase and central distinctions. Only local short
representatives are used; no global hyperbolicity or global topology is
assumed. ∎

\protect\phantomsection\label{sec:gg-positive-original-adequacy}

\subsubsection{I.3. Operator, infinity-structure and point
comparisons}\label{i.3.-operator-infinity-structure-and-point-comparisons}

\pandocstatementheading{\protect\phantomsection\label{thm:gg-positive-original-adequacy}

\protect\phantomsection\label{thm:gg-primitive-infinity-point-structure}Theorem I.4.}

In the standard projective Heisenberg/Euler-minus-two analytic model,
the source construction and the independent domain of Definition H.1
have these properties:

\begin{enumerate}
\tightlist
\item
  Both complete original helicity classes act in the same analytic
  half-density operator primitive. Their actual nonlinear gluing has
  filtered adjoint-preserving overlaps, normalized inner associators and
  the original coupled source corrections.
\item
  The primitive holomorphic infinity predicate of Definition I.2 is
  equivalent, with its actual root and central calibration, to the
  parallel normal tractor and rank-four local-twistor infinity forms of
  Proposition I.3. Its failure before affine admission is the reduction
  defect of Lemma I.1, and Proposition H.5 removes it with the same
  source rule.
\item
  Compact supported simple modules determine the complete physical point
  family and null directions by Proposition H.3. The constructor and
  decoder have both inverse comparisons on the whole independent
  domains, with the physical point/full-spin Hom, numerical scalar,
  bounded-family and coupled-real assertions of Theorem 9.10.
\end{enumerate}

\emph{Proof.} For the first assertion, §F.1 gives the positive
differential operator and the direct analytic negative symbol in the
original twistor variables and their derivative momenta. The negative
operation introduces no new dual input and need not project to a vector
field on \(T_0\). Theorems F.1 and F.6 give actual finite and analytic
algebra gluing with normalized triple units. Proposition 9.1 and §F.3
retain the mixed principal bracket and lower quantum terms. Theorem 3.1
identifies both complete source-to-curvature lists with their original
factorials. The primitive's Euler-minus-two ket sector does not replace
either gravitational deformation class by a scalar-field state.

For the second assertion, the torsor, projective extension and density
root are holomorphic data specified before Einstein recognition. Section
D.8 and Proposition 5.4 recover the unit almost-Einstein density and its
canonical parallel tractor. Proposition I.3's explicit exterior-square
map intertwines the normal and rank-four spin connections, with the norm
and numerical dictionary fixed in §2. The reverse construction recovers
the same projective extension and root. The operator, torsor and module
comparisons of Theorem H.4 make this correspondence natural; a principal
symbol alone is insufficient. The flat scale test in Proposition I.3
shows why the retained infinity data are substantive.

Lemma I.1 reads the pre-admission failure on the actual provisional
Cauchy states. Proposition H.5 identifies its cochain, derivative
\(o_m\), zero-record inverse \(L_m\) and six normalizations with Theorem
5.5. Each contact correction is followed by fresh incidence, root,
optical-state, matrix and scalar computation. Finite-prefix execution is
Proposition 5.9. The bounded comparison of §6.1 and Theorem 6.3 gives
its common-domain analytic limit, together with the operator and module
estimates of Appendices F--G.

For the third assertion, Proposition H.3 derives the complete
four-dimensional family from one compact characteristic; the first-jet
kernel is its full null Segre quadric. Sections G.1--G.4 construct the
actual simple modules, Hom generators, root-symbol comparisons,
whole-family framing and exact associator descent. Thus points are
nonlocal supported objects rather than an input four-manifold.
Proposition I.3 supplies their Lorentz sky and geometric phase-scaled
fibers. Finally Theorem H.4 compares every independently admitted
operator object with its constructed classical quantization, retaining
both associators and all module maps. Theorem 9.10 composes those
counits with the source/affine correspondence and supplies all stated
physical and analytic quantifiers, including every Weyl type and
numerical \(\Lambda=0\). ∎

These are the operator, infinity-structure and point comparisons
collected in Theorem 9.11. Appendix J identifies their same-data raywise
local-twistor, optical and affine diagrams; Theorem 9.12 states the
resulting unified correspondence.

\protect\phantomsection\label{app:gg-simultaneous-projective-comparison}

\subsection{Appendix J. The same source, infinity structure and point
objects}\label{appendix-j.-the-same-source-infinity-structure-and-point-objects}

\protect\phantomsection\label{sec:gg-derived-ward-construction}

\subsubsection{J.1. Normal spin transport before Einstein
admission}\label{j.1.-normal-spin-transport-before-einstein-admission}

We use the independent operator domain \(\mathfrak O_\Lambda\) of
Definition H.1 and the symbol and quantization functors \(F\) and \(Q\)
of Theorem H.4. For the constructions before Einstein admission, retain
its filtered embedded compact characteristic, complete support family,
global first-unit torsor \(A\), compact lifts, unit optical derivative,
root and central calibration, and temporarily omit its
projective-extension predicate and numerical readback. Call this the
regular affine domain. A highest source trial with a nonzero affine
obstruction is treated separately in J.4.

Let \(M\xleftarrow{p}\mathcal F\xrightarrow{q}Y\) be the incidence
derived from the compact characteristic. Write \(\mathcal L\) for the
contact line on \(Y\), \(\mathcal E[1]^2\xrightarrow{\beta}\mathcal B\)
for the optical root, and \(\sigma\) for the density selected by the
same torsor. Appendix D gives these objects from the complete embedded
support family and its actual affine data. They are retained in all
comparisons below.

For clarity, the contact normal sequence of a sky is
\begin{equation}\protect\phantomsection\label{eq:gg-ward-contact-normal-sequence}{
 0\longrightarrow\Omega^1_{Q_x}\otimes\mathcal L
 \longrightarrow N_{Q_x/Y}\longrightarrow\mathcal L
 \longrightarrow0,
 \qquad
 \mathcal L|_{Q_x}=\mathcal O(1,1).
}\end{equation} The left term is
\(\mathcal O(-1,1)\oplus\mathcal O(1,-1)\) and is acyclic. Completeness
and proper base change identify \(TM=p_*q^*\mathcal L\), of rank four.
First-jet evaluation gives the Segre null quadric and the conformal
structure. The line \(K=\ker dq\) obeys
\(K\otimes q^*\mathcal L=p^*\mathcal B\). The derivative
\(\Theta=d_qj_A\) of the unique compact lifts is constant on each
compact sky and gives the same unit metric section as in D.8. The
density is characterized by
\begin{equation}\protect\phantomsection\label{eq:gg-ward-selected-root}{
 \beta(\sigma^2)=\Theta^{-1},\qquad \sigma(o)=e_o.
}\end{equation} The rank-six standard tractor is the jet quotient
\begin{equation}\protect\phantomsection\label{eq:gg-ward-tractor-jet-quotient}{
 \mathcal T=J^2\mathcal E[1]\big/
 \bigl(\operatorname{Sym}^2_0T^*M\otimes\mathcal E[1]\bigr).
}\end{equation} In a conformal representative \(g\), set
\(J=\operatorname{Scal}(g)/6\) and write
\begin{equation}\protect\phantomsection\label{eq:gg-ward-normal-tractor}{\begin{aligned}
L\sigma&=(\sigma,\mu,\rho), \\
\mu&=\nabla\sigma, \\
\rho&=-\tfrac14(\Delta\sigma+J\sigma), \\
H(\sigma,\mu,\rho)&=2\sigma\rho+g^{-1}(\mu,\mu), \\
D_a(\sigma,\mu_b,\rho)
 &=\bigl(\nabla_a\sigma-\mu_a,
 \nabla_a\mu_b+P_{ab}\sigma+g_{ab}\rho,
 \nabla_a\rho-P_a{}^b\mu_b\bigr).
\end{aligned}}\end{equation} These formulas apply before Einstein admission. Under
\(\widehat g=e^{2\omega}g\), \(\Upsilon=d\omega\), their splitting
change is
\begin{equation}\protect\phantomsection\label{eq:gg-ward-conformal-splitting}{
 \widehat\sigma=e^\omega\sigma,\qquad
 \widehat\mu=e^\omega(\mu+\sigma\Upsilon),\qquad
 \widehat\rho=e^{-\omega}
 \bigl(\rho-\Upsilon\cdot\mu-\tfrac12|\Upsilon|^2\sigma\bigr).
}\end{equation} Substitution, together with
\(\widehat P=P-\nabla\Upsilon+\Upsilon\otimes\Upsilon-|\Upsilon|^2g/2\),
preserves \(H\) and intertwines \(D\). The full marked spin lift gives
the rank-four local-twistor bundle \(\mathcal T_{\mathrm{LT}}\) with
volume and the exterior-square identification
\(\mathcal F_{\mathrm{spin}}:\mathcal T\simeq\Lambda^2\mathcal T_{\mathrm{LT}}^*\)
of Proposition I.3. Before admission its connection has upper-right
block \(i\theta_v\) and lower-left block
\(i\theta_{P^\sharp(v)}^\sharp\). The same two-by-two identities in I.3
give the middle tractor term \(\nabla_v\mu+\sigma P^\sharp(v)+\rho v\)
and last term \(d_v\rho-P(v,\mu)\). Thus the normal spin connection is
defined for the derived conformal geometry, with no Einstein assumption.

The bundle \(\mathcal T_{\mathrm{LT}}\) is a geometric functor of the
structured operator object. It is distinct from \(TM=p_*q^*\mathcal L\)
and carries no asserted fiberwise finite-dimensional representation of
the Heisenberg algebra: taking a trace of \([d_a,y^a]=1\) would exclude
such a representation.

\pandocstatementheading{\protect\phantomsection\label{prop:gg-exact-ward-defect}Proposition J.1.}

For \(I=L\sigma\) selected by the same torsor and root, put
\(E_{ab}=(\nabla_a\nabla_b\sigma+P_{ab}\sigma)_0\). On the regular
affine domain,
\begin{equation}\protect\phantomsection\label{eq:gg-exact-ward-defect}{\begin{gathered}
D_aI=\bigl(0,E_{ab},-\tfrac13\nabla^bE_{ab}\bigr). \\
DI=(0,P^\circ,-\tfrac14dJ)
 \quad\text{in the selected scale }\sigma=1.
\end{gathered}}\end{equation} In particular \(DI=0\) if and only if \(E=0\).

\emph{Proof.} The first component is zero and the middle component has
zero trace by the definition of \(\rho\). Put
\(\zeta_a=\nabla_a\rho-P_a{}^b\mu_b\). The gradient commutator,
contracted Bianchi identity and \(\operatorname{Ric}=2P+Jg\) give
\begin{equation}\protect\phantomsection\label{eq:gg-ward-defect-divergence}{
 \begin{aligned}
 \nabla^bE_{ab}
 &=\nabla_a\Delta\sigma+\operatorname{Ric}_a{}^b\mu_b
   +(\nabla_aJ)\sigma+P_a{}^b\mu_b+\nabla_a\rho\\
 &=3P_a{}^b\mu_b-3\nabla_a\rho=-3\zeta_a.
 \end{aligned}
}\end{equation} In the selected scale \(\mu=0\) and \(\rho=-J/4\),
giving the remaining formula. ∎

\protect\phantomsection\label{sec:gg-relative-ward-descent}

\subsubsection{J.2. Raywise descent and the infinity
predicate}\label{j.2.-raywise-descent-and-the-infinity-predicate}

Here raywise descent is descent of the normal rank-four local-twistor
bundle by its initialized flat partial connection along the connected
short null rays. The bundle, its flag and its tensor comparisons are
constructed by that transport.

\pandocstatementheading{\protect\phantomsection\label{thm:gg-derived-ward-datum}Theorem J.2.}

On the regular affine domain, the normal rank-four spin bundle descends
along the short null-ray incidence to a holomorphic bundle
\(\mathcal T_{\mathrm{LT,ray}}\) on \(Y\), with volume and a flag
\(\ell\subset h\) of ranks one and three. The alternating form of the
selected \(I=L\sigma\) descends if and only if the projective-extension
predicate of Definition H.1 holds. On that admitted subdomain its
covariant and contravariant forms have the numerical label, ranks and
zero-constant limit of Proposition I.3.

\emph{Proof.} Restrict the connection on \(p^*\mathcal T_{\mathrm{LT}}\)
to \(K=\ker dq\). Its one-dimensional fibers have a flat partial
connection. On the retained short flow-box cover, transverse sections
and the initialized finite matrix ODE give
\begin{equation}\protect\phantomsection\label{eq:gg-ward-descent}{
 \mathcal T_{\mathrm{LT,ray}}(U)
 =\{s\in\Gamma(q^{-1}U,p^*\mathcal T_{\mathrm{LT}}):D_qs=0\},
 \qquad q^*\mathcal T_{\mathrm{LT,ray}}\simeq p^*\mathcal T_{\mathrm{LT}}.
}\end{equation} The fibers are connected and contractible. Uniqueness
gives exact inverse and triple laws, and holomorphic dependence gives
local freeness. Trace zero of the connection descends the volume.

Write a null tangent as \(k^{AA'}=\kappa^A\pi^{A'}\) and choose parallel
spin factors along its affine geodesic. In the splitting
\(\mathcal T_{\mathrm{LT}}=S\oplus S'\), the vector \((0,\pi)\) is
parallel: its upper derivative is proportional to
\(\kappa(\pi\cdot\pi)=0\) and its lower Schouten term multiplies zero.
The dual covector \((\kappa_A,0)\in \mathcal T_{\mathrm{LT}}^*\), with
\(\kappa_A=\epsilon_{AB}\kappa^B\), is parallel by
\(\kappa_A\kappa^A=0\). Their line and kernel descend to \(\ell\) and
\(h\), with \(\ell\subset h\). Evaluation at \(x\) identifies
\begin{equation}\protect\phantomsection\label{eq:gg-ward-flag-lines}{\begin{gathered}
\mathcal T_{\mathrm{LT,ray}}|_{Q_x}\simeq (\mathcal T_{\mathrm{LT}})_x\otimes\mathcal O_{Q_x}, \\
\ell|_{Q_x}=\mathcal O(0,-1), \\
\operatorname{Ann}(h)|_{Q_x}=\mathcal O(-1,0), \\
\ell\otimes\operatorname{Ann}(h)\simeq\mathcal L^{-1}.
\end{gathered}}\end{equation} The last identification is obtained by pulling back to
the incidence: \(\kappa_A\pi_{A'}\) and \(q^*\mathcal L^{-1}\) are the
same tautological null-covector line, with the same parallel transport
and root calibration. It descends along \(q\). Triviality of the
rank-four bundle on a sky leaves these tautological sublines intact.

Under \(\mathcal F_{\mathrm{spin}}\), the same \(I\) is an alternating
form on \(p^*\mathcal T_{\mathrm{LT}}\). Descent requires \(D_qI=0\) on
the whole incidence. Every null direction occurs on its complete compact
sky; null vectors span \(T_xM\). Consequently this condition is
equivalent to \(DI=0\). Proposition J.1, D.8 and Proposition 5.4
identify it with the original projective-extension predicate for the
same unit density. The tensor being tested is specified by the torsor on
the whole incidence; parallel extension of arbitrary initial tensor
values is a different operation.

For an admitted object \(\Lambda=-3H(I,I)\) is constant. In the selected
Einstein splitting, with \(c=\Lambda/6\) and the alternating matrix
\(J_2\), Proposition I.3 gives
\begin{equation}\protect\phantomsection\label{eq:gg-ward-infinity-pair}{\begin{gathered}
I_{\mathrm{cov}}=\operatorname{diag}(-cJ_2,J_2), \\
I_{\mathrm{contra}}=\operatorname{diag}(J_2,-cJ_2), \\
I_{\mathrm{cov}}I_{\mathrm{contra}}=c\,1, \\
2\operatorname{Pf}(I_{\mathrm{cov}})=-\Lambda/3.
\end{gathered}}\end{equation} The contravariant form is the volume dual. Their ranks
are four for \(\Lambda\ne0\) and two for \(\Lambda=0\); the construction
uses their values and not their inverses. The marked root fixes the sign
of \(I\). ∎

\pandocstatementheading{\protect\phantomsection\label{prop:gg-ward-naturality}Proposition J.3.}

The raywise local-twistor descent is natural for the physical
comparisons of H.1/H.4. A normalized comparison 2-isomorphism over the
same classical data and Spin label has identity finite-rank geometric
comparison. The actual associator actions on point modules remain those
of H.4. The physical central Spin arrow acts by \(-1\) on
\(\mathcal T_{\mathrm{LT,ray}}\).

\emph{Proof.} Write the descent as
\begin{equation}\protect\phantomsection\label{eq:gg-ward-actual-associators}{
 \varphi_{ij}\varphi_{jk}
   =\operatorname{Ad}(C_{ijk})\varphi_{ik},\qquad
 C_{ijk}C_{ikl}=\varphi_{ij}(C_{jkl})C_{ijl},\qquad
 \tau_{ij}\tau_{jk}=C_{ijk}\tau_{ik}.
}\end{equation} Since \(\sigma_0(C)=1\), the associated graded glues
strictly. Normalized inner action on
\(\mathscr U_1=(1+F_{-1})/(1+F_{-2})\simeq\mathcal L^{-1}\) is identity.
The torsor maps are supplied in H.1 and the normalized comparisons lie
over that same classical torsor and root. Thus the normalized associator
comparison fixes the embedded support, torsor, root, their compact lifts
and derived jets. Its specified same Spin label fixes the normal spin
comparison. Uniqueness of the initialized ray ODE makes its raywise
local-twistor comparison identity. This uses both the strict translation
sheaf and the supplied torsor maps; identity on a translation sheaf
alone would not determine an affine map.

The last equation in the display continues to describe the actual module
action. The finite-rank construction reads classical characteristic and
affine data and imposes no neutralization equation for \(C\). A physical
comparison carries those data, hence \(\Theta,\sigma,H,D\) and the
retained spin lift. Proper base change and ray transport commute with
it. This proves naturality and its finite composition diagrams. Finally
the central spin element acts by \(-1\) on \(S\oplus S'\), so by \(-1\)
on \(\mathcal T_{\mathrm{LT,ray}}\). It acts trivially on exterior
squares and projective flags, and on the even Euler-minus-two sector;
the full physical spin label nevertheless remains distinct. ∎

\protect\phantomsection\label{sec:gg-optical-point-pencil}

\subsubsection{J.3. The optical evaluation, point plane and affine
slice}\label{j.3.-the-optical-evaluation-point-plane-and-affine-slice}

The classical Einstein-bundle recognition is described in §1 (LeBrun
1985; Wolf 2007, sec. 3.8). For the intrinsic ray operator used here,
the following Cauchy/tractor identification is explicit.

\pandocstatementheading{\protect\phantomsection\label{prop:gg-same-optical-point-pencil}Proposition J.4.}

Let \(\gamma\) be a retained short null ray. Its optical solution space
\(\mathcal E_{J,\gamma}\) is canonically \(V_\gamma^*\), where
\(V_\gamma\) is the parallel totally null two-plane spanned by its
position tractor and derivative. For \(x\in\gamma\), its evaluation
gives the point plane
\begin{equation}\protect\phantomsection\label{eq:gg-same-point-plane}{
 K_{x,\gamma}=\ker\mathcal F_{\mathrm{spin}}(X(x))
       \subset(\mathcal T_{\mathrm{LT,ray}})_\gamma.
}\end{equation} The completed optical pencil is the projective line of
planes \(K\) with \(\ell_\gamma\subset K\subset h_\gamma\). On an
admitted object, the same torsor is the affine slice \(H(I,v)=1\) in
\(V_\gamma\), with compact evaluation \(X/\sigma\) and the original
moving affine transitions.

\emph{Proof.} In an auxiliary conformal scale, along an affinely
parametrized null geodesic,
\begin{equation}\protect\phantomsection\label{eq:gg-same-optical-plane}{
 D_kX=(0,k^\flat,0),\qquad
 D_k^2X=-P(k,k)X,\qquad
 V_\gamma=\operatorname{span}(X,D_kX).
}\end{equation} The two vectors are independent, and their three
pairings under \(H\) vanish. A parallel \(\lambda\in V_\gamma^*\) gives
the optical solution \(\lambda(X(u))\). Value and slope on a transverse
slice give its inverse, so this is the Cauchy identification
\begin{equation}\protect\phantomsection\label{eq:gg-same-optical-cauchy-comparison}{
 V_\gamma^*\simeq\mathcal E_{J,\gamma},\qquad
 X(u)\longleftrightarrow\epsilon_u\in\mathcal E_{J,\gamma}^*.
}\end{equation} At a point \(x\), in the marked spin splitting and with
\(\operatorname{rank}\mu=1\), the exterior-square map gives
\begin{equation}\protect\phantomsection\label{eq:gg-same-point-pencil-matrices}{
 \mathcal F_{\mathrm{spin}}(X)=
 \begin{pmatrix}J_2&0\\0&0\end{pmatrix},\qquad
 \mathcal F_{\mathrm{spin}}(D_kX)=
 \begin{pmatrix}0&-iJ_2\mu\\-i\mu^TJ_2&0\end{pmatrix}.
}\end{equation} The first kernel is the point plane \(S'_x\). For a
finite pencil coordinate \(z\), the kernel is
\(\{(iz\mu v,v):v\in S'_x\}\), since \(\mu^TJ_2\mu=0\). Its common line
is \(0\oplus\ker\mu=\ell_\gamma\), and the span of all kernels is
\(\operatorname{im}\mu\oplus S'_x=h_\gamma\). The other projective chart
completes the pencil between this line and hyperplane. The short
physical evaluations select the open subset of D.8, rather than all
points in the completed projective line. For a compact point module, H.3
supplies its support \(Q_x\); evaluation of
\(\mathcal T_{\mathrm{LT,ray}}\) on that support supplies these same
point planes.

This argument commutes with actual finite Taylor restriction. Over a
finite Artin base one uses primitive homogeneous pairs in
\(\mathbb P(V_\gamma)\). On its two standard charts a rank-two minor is
a unit by the reduced computation, while the Pfaffian vanishes by total
nullity. The kernel is therefore a rank-two direct summand. The unit
spinor charts likewise give its common line and hyperplane. This
verifies the assertion on finite jets of the admitted holomorphic
families.

For an admitted object \(H(I,X)=\sigma\) and \(I\) is parallel. Its
restriction to \(V_\gamma\) is exactly the optical section \(s\), with
the same Cauchy values. Its affine slice and evaluation are
\begin{equation}\protect\phantomsection\label{eq:gg-same-infinity-affine-slice}{
 A_s=\{v\in V_\gamma:H(I,v)=1\},\qquad
 j_s=X/\sigma\longleftrightarrow\epsilon/\sigma.
}\end{equation} With the Wronskian
\(\Lambda^2\mathcal E_J\simeq\mathcal L\), a translation
\(\nu\in\mathcal L^{-1}\) is the functional \(v\mapsto\nu(s\wedge v)\)
in \(\operatorname{Ann}(s)\subset\mathcal E_J^*\). In the calibrated
Cauchy trivialization, choose the companion with \(W(s,t)=l\). Its
affine coordinate is \(t/\sigma\), with
\(\partial_u(t/\sigma)=\sigma^{-2}\). It is the same \(\Theta=d_qj_A\)
used in D.8. On a transverse slice identify the supplied \(j_A\) with
\(\epsilon/\sigma\), with identity translation map. The evaluations have
the same derivative on each connected short fiber, so they agree
throughout that fiber. Uniqueness of the compact lifts glues these
identifications over the complete sky. The comparison preserves the
origins, moving arguments and affine transitions, and is independent of
the chosen slice. ∎

\protect\phantomsection\label{sec:gg-source-ward-defect}

\subsubsection{J.4. The source obstruction and six scale
modes}\label{j.4.-the-source-obstruction-and-six-scale-modes}

\pandocstatementheading{\protect\phantomsection\label{prop:gg-source-ward-defect}Proposition J.5.}

In the flat-residue source filtration of Theorem 5.5, the actual highest
optical mismatch is the Hessian quotient of the middle normal-tractor
defect of Proposition J.1. Its derivative and compatible zero-record
inverse are the same \(o_m\) and \(L_m\). The six normalizations are the
corresponding six flat tractor initial values.

\emph{Proof.} Scale the original records and \(\Lambda\) by
\(\varepsilon\), as in §6. The residue at \(\varepsilon=0\) is flat. At
a highest square-zero ideal \(\mathfrak j\), with
\(\mathfrak m\mathfrak j=0\), lift the lower selected scale
provisionally and keep its moving incidence and Cauchy frames. The trial
already has conformal incidence, but need not have a global affine
reduction. In that scale the provisional density \(1\) has ray residual
\(P^{\mathrm{trial}}(k,k)\). Use the fixed double primitives of D.4 to
set
\begin{equation}\protect\phantomsection\label{eq:gg-source-ward-actual-mismatch}{\begin{gathered}
\widetilde\sigma_i
 =1-\mathsf I_i(P^{\mathrm{trial}}(k,k)), \\
r_{ij}=\widetilde s_i\circ\Phi_{ij}-B_{ij}\widetilde s_j, \\
(G_i\circ\Phi_{ij})^{-1}r_{ij}=e_1-T_{ij}e_1.
\end{gathered}}\end{equation} Thus the entries are \(1-T_{11}\) and \(-T_{21}\). With
\(\delta u_{ij}=u_j-u_i\), the fixed splitting has \(D^2u_i=A(k,k)\) and
\begin{equation}\protect\phantomsection\label{eq:gg-source-ward-hessian-class}{
 a_m=[A]=[P^{\mathrm{trial}\,\circ}]
 \in\frac{\mathcal P_{m-2}\otimes\operatorname{Sym}^2_0V^*}
              {\operatorname{Hess}_0\mathcal P_m},
 \qquad o_mJ[h]=[P(h)^\circ].
}\end{equation} Proposition J.1 gives
\(DI=(0,P^{\mathrm{trial}\,\circ},-dJ_{\mathrm{trial}}/4)\) for the
provisional density. This identifies the class and its sign. The
contracted Bianchi identity gives \(d\operatorname{div}a_m=0\), so D.5
applies its fixed inverse and the correction is
\(\Delta c=-L_ma_m+\iota Z\). Recompute the incidence, optical states,
adapted matrices, root and scalar. Once the class vanishes, the same
double-primitive argument gives \(v_i=u_i-f\) with \(\delta v=r\) and
\(D^2v=0\); adding it cancels the target-minus-source mismatch. This is
the stage at which the next global affine reduction is obtained.

The remaining modes \(a+b\cdot x+\gamma q\) have flat tractor initial
values \((a,b,-2\gamma)\). Marked value, four gradients and the
numerical invariant fix them. Since \(\Delta q=8\), the scalar changes
by \(12\gamma\) and \(\gamma=(\lambda_{\mathrm{target}}-\lambda_*)/12\).
These highest-coefficient formulas are used in this flat-residue
filtration. The analytic output at \(\varepsilon=1\) is the unrestricted
prescribed representative of Theorem 6.3. ∎

\protect\phantomsection\label{sec:gg-simultaneous-comparison-proof}

\subsubsection{J.5. Simultaneous comparison on the independent
domain}\label{j.5.-simultaneous-comparison-on-the-independent-domain}

For \(O\in\mathfrak O_\Lambda\), let \(\mathcal R(O)\) be the marked
Einstein germ read from \(F(O)\) by D.8, Proposition 5.4 and Proposition
4.3. Let \(\mathcal P(O)\) be the diagram consisting of its supplied
operator algebroid and complete compact point modules, their actual
associators, the derived incidence and affine torsor, optical evaluation
and Wronskian, local-twistor volume and flags, infinity pair and
point-plane map. It is an output diagram, not additional admission data.
For an Einstein germ \(g\), define \(\operatorname{Can}(g)\) from its
canonical ray geometry, tractor and spin transport, point planes and
calibrated affine slice, enriched by the specified \(Q\) of Theorem H.4
and the framed modules of Appendix G. These definitions use the
canonical geometry of \(g\) independently of \(O\).

\pandocstatementheading{\protect\phantomsection\label{thm:gg-simultaneous-projective-comparison}Theorem J.6.}

On the whole independent domain there are natural actual diagram
comparisons
\begin{equation}\protect\phantomsection\label{eq:gg-simultaneous-diagram-comparisons}{
 \beta_O:\mathcal P(O)\xrightarrow{\sim}
       \operatorname{Can}(\mathcal R(O)),\qquad
 \alpha_A:\mathcal P(\mathcal S_{\mathrm{op}}A)
       \xrightarrow{\sim}\operatorname{Can}(g_A),
}\end{equation} where \(g_A\) is the canonical Einstein germ of
Proposition 4.1 for the original records. These preserve the same source
operations, torsor, optical evaluation, infinity pair and compact point
objects. Both inverse comparisons of Theorem 9.10 act on these diagrams,
with their physical point/full-spin Hom classes, bounded holomorphic
families, actual finite jets and conditioned finite diagrams. Their
coupled Lorentz comparison has the type of Proposition J.7.

\emph{Proof.} H.4 supplies an actual counit \(\rho_O:O\simeq QF(O)\).
D.8 and Proposition 4.3 supply the calibrated classical comparison
\(\theta_O\) from \(F(O)\) to the canonical ray geometry of
\(\mathcal R(O)\). The operator and module component of \(\beta_O\) is
\(Q(\theta_O)\rho_O\). Propositions J.1, J.3 and J.4, together with
Theorem J.2, identify all its finite-rank components with the
independently defined \(\operatorname{Can}(\mathcal R(O))\): the same
\(\Theta\) selects \(\sigma\) and \(I\), the same spin connection gives
raywise local-twistor transport, and the same optical evaluation gives
the point pencil and affine slice.

For actual comparison units \(U_{ij}\) the algebra and module equations
are
\begin{equation}\protect\phantomsection\label{eq:gg-simultaneous-actual-comparison-equations}{
 \begin{aligned}
 k_i\varphi^s_{ij}
   &=\operatorname{Ad}(U_{ij})\varphi^t_{ij}k_j,\\
 U_{ij}\varphi^t_{ij}(U_{jk})C^t_{ijk}
   &=k_i(C^s_{ijk})U_{ik},\\
 F_i\tau^s_{ij}&=U_{ij}\tau^t_{ij}F_j.
 \end{aligned}
}\end{equation} Both associators and their module actions occur here.
Proposition J.3 gives the identity finite-rank shadow for normalized
2-isomorphisms over the same physical map and Spin label. For a physical
arrow \(f:O\to O'\), complete embedded support classification carries
the point parameter and incidence. Unique compact lifts carry
\(j_A,\Theta\); the calibrated root carries \(\sigma,I\); jet
functoriality carries \(H,D\). Its spin lift and initialized transport
then carry the flags, point planes and affine slice. Framed Hom
uniqueness in H.4 gives
\begin{equation}\protect\phantomsection\label{eq:gg-simultaneous-naturality}{
 \beta_{O'}\mathcal P(f)
   =\operatorname{Can}(\mathcal R(f))\beta_O
}\end{equation} on physical Hom classes, with the coherent normalized
comparison of H.4 on actual representatives. This proves naturality and
its finite composition laws on every independent object.

For an original input \(A\), H.5 identifies
\(F(\mathcal S_{\mathrm{op}}A)\) with the full corrected contact-affine
tuple of Theorems 5.5--6.3. Proposition J.5 verifies its same source
obstruction and six normalizations. The operators in F.1, the mixed
principal bracket and lower quantum term of F.3, and Proposition 9.2's
shared optical defect remain its source operations. F.6 and Appendix G
supply their analytic operator and module comparisons; F.7 identifies
the amplitude jets with the normalized finite operator construction.
Theorem 6.3 gives the canonical marked comparison
\(\eta_A:\mathcal R(\mathcal S_{\mathrm{op}}A)\simeq g_A\). Define
\(\alpha_A=\operatorname{Can}(\eta_A)\beta_{\mathcal S_{\mathrm{op}}A}\).
Its components are the identical constructions just described.
Naturality of \(\eta\) and \(\beta\) proves the source diagrams.

The actual counit
\(\mathcal S_{\mathrm{op}}\mathcal D_{\mathrm{op}}O\simeq O\) of Theorem
9.10, including its H.4 module maps, transports through \(\mathcal P\).
The opposite composite is its initialized equality on original classes
and scalar. H.4's uniqueness gives the unit/counit and finite
composition coherences on these diagrams. The physical central Spin
label acts by \(-1\) on \(\mathcal T_{\mathrm{LT,ray}}\) and remains
distinct, as in J.3.

All operations use the existing complete families, fixed finite covers,
unit root and frame margins, and initialized finite matrix transports.
Finitely many buffered restrictions remain within the common parents of
§§3--6 and F--H. Proper base change, finite jet operations and the
primitive-pair argument of J.4 commute with the parameter restrictions
of Definition 2.1. Proposition J.7 proves compatibility with the typed
coupled Lorentz operation and its retained square. ∎

\protect\phantomsection\label{sec:gg-ward-lorentz-duality}

\subsubsection{J.6. The type of the Lorentz
comparison}\label{j.6.-the-type-of-the-lorentz-comparison}

\pandocstatementheading{\protect\phantomsection\label{prop:gg-ward-lorentz-duality}Proposition J.7.}

For the coupled antiholomorphic operation \(\iota\) of §8 and G.5, the
raywise local-twistor comparison is an antilinear duality
\begin{equation}\protect\phantomsection\label{eq:gg-typed-ward-lorentz-duality}{
 J_{\mathrm{LT}}:(\mathcal T_{\mathrm{LT,ray}})_y\longrightarrow (\mathcal T_{\mathrm{LT,ray}}^*)_{\iota(y)}.
}\end{equation} Its companion on the dual and its composite have the
retained normalized square. It exchanges \(\ell\) with
\(\operatorname{Ann}(h)\), sends a point plane to the annihilator of its
paired point plane, and exchanges the covariant and contravariant
infinity forms. These operations are natural for Theorem J.6 and the
actual operator/module comparison of G.5/H.4.

\emph{Proof.} On a fixed Lorentz object, use compatible marked spin
frames and
\begin{equation}\protect\phantomsection\label{eq:gg-ward-hermitian-duality-matrix}{
 H_{\mathrm{LT}}=\begin{pmatrix}0&1_2\\1_2&0\end{pmatrix},\qquad
 J_{\mathrm{LT}}(Z)=Z^\dagger H_{\mathrm{LT}},\qquad
 \Omega_v=\begin{pmatrix}0&i\theta_v\\ic\theta_v^\sharp&0\end{pmatrix}.
}\end{equation} The row \(Z^\dagger H_{\mathrm{LT}}\) is a covector.
Since \(H_{\mathrm{LT}}^\dagger=H_{\mathrm{LT}}\) and
\(H_{\mathrm{LT}}^2=1\), its companion has square one in these
normalized frames; globally use the actual root/frame square of G.5. For
real tangent \(v\), \(\theta_v\) and its adjugate are Hermitian and
\(c\) is real, so
\begin{equation}\protect\phantomsection\label{eq:gg-ward-duality-connection}{
 \Omega_v^\dagger H_{\mathrm{LT}}+H_{\mathrm{LT}}\Omega_v=0.
}\end{equation} Before admission replace \(c\theta_v\) by
\(\theta_{P^\sharp(v)}\); the same identity follows from reality of
\(P\). Including the Lorentz spin connection proves parallelism. It
therefore descends through the initialized ray ODE. On paired complex
objects the same identity conjugates the parameter and \(c\) to their
values in the paired object.

For a real ray, \((0,\pi)\) maps to \((\pi^\dagger,0)\), giving the
asserted flag exchange. In a spin splitting at \(x\) the point plane is
\(0\oplus S'_x\), and its image is its annihilator \(S_x^*\oplus0\).
Equivalently, in a flat marked frame with Hermitian solder \(X\),
\begin{equation}\protect\phantomsection\label{eq:gg-ward-point-plane-duality}{
 U_x=\{(iX\pi,\pi)\},\qquad
 \operatorname{Ann}(U_x)=\{(\kappa,-i\kappa X)\}.
}\end{equation} The same statement is tensorial for curved splittings
and natural under the spin connection and raywise local-twistor
transport. For complex points it pairs \(x\) with \(\iota x\). Finally
the matrices in J.2 satisfy
\begin{equation}\protect\phantomsection\label{eq:gg-ward-infinity-duality}{
 H_{\mathrm{LT}}^{-1}\overline{I_{\mathrm{cov}}}H_{\mathrm{LT}}^{-T}
       =I_{\mathrm{contra}}
}\end{equation} on the real slice, with conjugated \(c\) on the paired
complex object. This uses the same fixed volume and is valid at \(c=0\).
The central Spin arrow is \(-1\) on both
\(\mathcal T_{\mathrm{LT,ray}}\) and \(\mathcal T_{\mathrm{LT,ray}}^*\)
and commutes with this duality. The operation on the operator ring
called formal adjoint is complex linear; its type is distinct from the
coupled antiholomorphic comparison just proved. ∎

\protect\phantomsection\label{app:sh-schwarzschild}

\subsection{Appendix K. Finite Schwarzschild comparisons and screen
transport}\label{appendix-k.-finite-schwarzschild-comparisons-and-screen-transport}

Throughout this appendix, \(m>0\) is fixed and \(g_E=g_m\) denotes the
explicitly displayed Schwarzschild metric. Its conformal representative
is \(g_c=r^{-2}g_m\).

\protect\phantomsection\label{sec:sh-integrated-scale}

\subsubsection{K.1. Einstein scale and calibrated full-spin
frame}\label{k.1.-einstein-scale-and-calibrated-full-spin-frame}

\protect\phantomsection\label{sec:sh-scale-1}

\paragraph{K.1.1. Metric, conventions, and
domain}\label{k.1.1.-metric-conventions-and-domain}

Fix \(m>0\). On a finite connected real-analytic future-horizon
neighborhood \(U\), choose ingoing coordinates \((v,r,\theta,\phi)\)
with \begin{equation}\protect\phantomsection\label{eq:sh-scale-a1}{
0<r_-<r<r_+<\infty,\qquad r_-<2m<r_+,
\qquad f(r)=1-\frac{2m}{r}.
}\end{equation} An angular coordinate patch is understood in each
formula; different regular angular patches give the same tensors. Set
\begin{equation}\protect\phantomsection\label{eq:sh-scale-a2}{
g_E=f\,dv^2-2\,dv\,dr-r^2d\Omega_2^2,
\qquad g_c=r^{-2}g_E,
\qquad \sigma=r^{-1}.
}\end{equation} Here \(g_c\) denotes the analytically continued
conformal tensor; no real orthogonal complex structure for a real
Lorentz metric is asserted. Thus \(g_E\) has the signature of §2
\(+---\). Its Levi-Civita connection is the connection of the negative
of the usual \(-+++\) Schwarzschild metric. The curvature and Ricci
conventions are
\begin{equation}\protect\phantomsection\label{eq:sh-scale-a3}{
R(X,Y)=[\nabla_X,\nabla_Y]-\nabla_{[X,Y]},\qquad
\operatorname{Ric}(Y,Z)=\operatorname{tr}(X\mapsto R(X,Y)Z),
}\end{equation} and
\begin{equation}\protect\phantomsection\label{eq:sh-scale-a4}{
P=\frac12\left(\operatorname{Ric}-\frac{\operatorname{Scal}}6g\right),
\qquad J=\frac{\operatorname{Scal}}6.
}\end{equation} Every statement is local on \(U\) or on a sufficiently
small common complexification with \(r\ne0\). There is no limit
\(m\to0\), infinite-time assertion, or global quotient assertion.

For reference, the vacuum property can be checked before extending
across the horizon. In static \(-+++\) coordinates, direct substitution
into the Christoffel formula gives
\begin{equation}\protect\phantomsection\label{eq:sh-scale-a5}{
\begin{aligned}
\operatorname{Ric}_{tt}&=f\left(\frac{f''}{2}+\frac{f'}r\right),&
\operatorname{Ric}_{rr}&=-f^{-1}\left(\frac{f''}{2}+\frac{f'}r\right),\\
\operatorname{Ric}_{\theta\theta}&=1-f-rf',&
\operatorname{Ric}_{\phi\phi}&=(1-f-rf')\sin^2\theta.
\end{aligned}
}\end{equation} All mixed components vanish. For \(f=1-2m/r\) these
expressions are zero. Metric negation preserves the connection and
covariant Ricci tensor. The ingoing formula eq.~\ref{eq:sh-scale-a2} is
analytic and nondegenerate at \(f=0\), so the zero Ricci tensor extends
there by continuity. Thus \(\operatorname{Ric}(g_E)=0\) everywhere under
consideration.

\protect\phantomsection\label{sec:sh-scale-2}

\paragraph{K.1.2. Exact scale
statement}\label{k.1.2.-exact-scale-statement}

\pandocstatementheading{\protect\phantomsection\label{prop:sh-scale-1}Proposition K.1.}

In the representative \(g_c\), the nonvanishing density whose component
is \(\sigma=1/r\) selects precisely \(g_E\):
\begin{equation}\protect\phantomsection\label{eq:sh-scale-a6}{
g_E=\sigma^{-2}g_c.
}\end{equation} It satisfies the almost-Einstein equation of Proposition
5.4. Its normal prolongation is
\begin{equation}\protect\phantomsection\label{eq:sh-scale-a7}{
I=L\sigma=(\sigma,\mu,\rho)
=\left(\frac1r,-\frac{dr}{r^2},\frac{f}{2r}\right),
\qquad D I=0,\qquad H(I,I)=0.
}\end{equation} It consequently has prescribed numerical Einstein label
\(\Lambda=-3H(I,I)=0\), including at the future horizon.

\emph{Proof.} The algebraic identity eq.~\ref{eq:sh-scale-a6} is
immediate from eq.~\ref{eq:sh-scale-a2}. We verify the differential and
tractor statements directly. Write \(g=g_E\), \(\omega=-\log r\), and
\(h=g_c=e^{2\omega}g\). The basic radial norm is
\begin{equation}\protect\phantomsection\label{eq:sh-scale-a8}{
|dr|_g^2=-f,\qquad |d\omega|_g^2=-\frac f{r^2}.
}\end{equation} The conformal connection formula, obtained by
substituting \(e^{2\omega}g\) in the Koszul formula, is
\begin{equation}\protect\phantomsection\label{eq:sh-scale-a9}{
\nabla^h_XY=\nabla^g_XY+d\omega(X)Y+d\omega(Y)X
-g(X,Y)\operatorname{grad}_g\omega.
}\end{equation} Since \(P^g=0\), its Schouten consequence is
\begin{equation}\protect\phantomsection\label{eq:sh-scale-a10}{
P^h=-\nabla_g^2\omega+d\omega\otimes d\omega
-\frac12|d\omega|_g^2g
=\frac1r\nabla_g^2r+\frac f{2r^2}g.
}\end{equation} For \(\sigma=e^\omega\), eq.~\ref{eq:sh-scale-a9} gives
\begin{equation}\protect\phantomsection\label{eq:sh-scale-a11}{
\begin{aligned}
\nabla_h^2\sigma
&=\nabla_g^2\sigma-d\omega\otimes d\sigma-d\sigma\otimes d\omega
+\langle d\omega,d\sigma\rangle_g g\\
&=\sigma\left(\nabla_g^2\omega-d\omega\otimes d\omega
+|d\omega|_g^2g\right).
\end{aligned}
}\end{equation} Adding \(P^h\sigma\) from eq.~\ref{eq:sh-scale-a10}
cancels both trace-free terms and leaves
\begin{equation}\protect\phantomsection\label{eq:sh-scale-a12}{
\nabla_h^2\sigma+P^h\sigma
=\frac12\sigma|d\omega|_g^2g
=-\frac f{2r}h.
}\end{equation} Taking the trace yields
\(\rho=-\frac14(\Delta_h\sigma+J_h\sigma)=f/(2r)\), with the exact
convention of Proposition 5.4 and Appendix J.1 of the correspondence.
The first two tractor equations are therefore
\begin{equation}\protect\phantomsection\label{eq:sh-scale-a13}{
d\sigma-\mu=0,\qquad
\nabla_h\mu+P^h\sigma+h\rho=0.
}\end{equation} We check the last equation separately. Raising
\(\mu=-dr/r^2\) by \(h^{-1}=r^2g^{-1}\) gives
\(\mu^{\sharp_h}=-\operatorname{grad}_g r\). Also \[
\nabla_g^2r(\operatorname{grad}_g r,\cdot)
=\frac12d(|dr|_g^2)=-\frac{f'}2dr.
\] It follows from eq.~\ref{eq:sh-scale-a10} that
\begin{equation}\protect\phantomsection\label{eq:sh-scale-a14}{
P^h(\mu^{\sharp_h},\cdot)
=\left(\frac{f'}{2r}-\frac f{2r^2}\right)dr
=d\left(\frac f{2r}\right)=d\rho.
}\end{equation} Thus every component of \(DI\) vanishes. Finally,
\begin{equation}\protect\phantomsection\label{eq:sh-scale-a15}{
H(I,I)=2\sigma\rho+|\mu|_h^2
=\frac f{r^2}-\frac f{r^2}=0.
}\end{equation} All the displayed expressions are analytic at \(r=2m\).
At the horizon \(\rho=0\), while \(\sigma=1/(2m)\ne0\) and
\(\mu=-dr/(4m^2)\) remains a nonzero null covector. The tractor itself
is nonzero there. No inverse power of \(f\) or of \(\Lambda\) appears. ∎

In the selected Einstein splitting the same tractor has the simpler
components \((1,0,0)\). The calculation above specifies its components
in the different auxiliary splitting, rather than equating these two
triples term by term.

\protect\phantomsection\label{sec:sh-scale-3}

\paragraph{K.1.3. Scalar curvature of the conformal
representative}\label{k.1.3.-scalar-curvature-of-the-conformal-representative}

\pandocstatementheading{\protect\phantomsection\label{prop:sh-scale-2}Proposition K.2.}

With eq.~\ref{eq:sh-scale-a3}, the scalar curvature of the auxiliary
metric in eq.~\ref{eq:sh-scale-a2} is
\begin{equation}\protect\phantomsection\label{eq:sh-scale-a16}{
\operatorname{Scal}(g_c)=-\frac{12m}{r},\qquad
\sigma=-\frac{\operatorname{Scal}(g_c)}{12m}.
}\end{equation} If the metric is negated to signature \(-+++\), its
scalar is \(12m/r\) and its scale component is the scalar divided by
\(12m\).

\emph{Proof.} For the static \(-+++\) metric \(g_-=-g_E\), the radial
divergence identity gives
\begin{equation}\protect\phantomsection\label{eq:sh-scale-a17}{
\Box_{g_-}(r^{-1})
=\frac1{r^2}\partial_r\left(r^2f\partial_r(r^{-1})\right)
=-\frac{2m}{r^4}.
}\end{equation} In dimension four, the scalar transformation is
\begin{equation}\protect\phantomsection\label{eq:sh-scale-a18}{
\operatorname{Scal}(\Omega^2g)
=\Omega^{-2}\left[\operatorname{Scal}(g)-6\Omega^{-1}\Box_g\Omega\right].
}\end{equation} Use \(\Omega=r^{-1}\) and \(\operatorname{Scal}(g_-)=0\)
to obtain \(12m/r\). Negating the metric changes the inverse metric and
hence the scalar sign, while fixing the covariant Ricci tensor. This
proves eq.~\ref{eq:sh-scale-a16}. An independent check is
\(\Box_{r^{-2}g_-}r=2m\), obtained directly from
\(\sqrt{|\det(r^{-2}g_-)|}=r^{-2}\sin\theta\) and
\((r^{-2}g_-)^{rr}=r^2f\). For \(r^{-2}g_-\), the reciprocal of its
scalar, \(r/(12m)\), has Laplacian \(1/6\), the value required for the
scalar-flat conformal metric. ∎

\protect\phantomsection\label{sec:sh-scale-4}

\paragraph{K.1.4. Point normalization and the auxiliary
representative}\label{k.1.4.-point-normalization-and-the-auxiliary-representative}

Fix a marked point \(o\in U\) and put \(r_0=r(o)>0\). Densities are
represented by functions only after specifying a representative metric.
We use the density convention of §5: under \(\widehat h=c^2h\) with a
positive constant \(c\), a weight-one density component changes by
\(\widehat\sigma=c\sigma\).

\pandocstatementheading{\protect\phantomsection\label{prop:sh-scale-3}Proposition K.3.}

The auxiliary representative and scale
\begin{equation}\protect\phantomsection\label{eq:sh-scale-a20}{
g_{\rm aux}=r_0^2g_c,\qquad
\sigma_o=\frac{r_0}{r}
}\end{equation} satisfy
\begin{equation}\protect\phantomsection\label{eq:sh-scale-a21}{
\sigma_o(o)=1,\qquad
\sigma_o^{-2}g_{\rm aux}=g_E.
}\end{equation} The same parallel tractor in this constant-rescaled
splitting has components
\begin{equation}\protect\phantomsection\label{eq:sh-scale-a22}{
I=\left(\frac{r_0}{r},-\frac{r_0}{r^2}dr,\frac{f}{2r_0r}\right).
}\end{equation} In particular this change of auxiliary representative
preserves the selected physical metric, the tractor itself, its zero
norm, and its calibrated central value.

\emph{Proof.} Equation eq.~\ref{eq:sh-scale-a21} is direct
multiplication. Constant rescaling fixes the connection and covariant
Schouten tensor, and scales the inverse metric by \(r_0^{-2}\). Hence
\(\mu_o=d\sigma_o=r_0\mu\) and \[
\rho_o=-\frac14(\Delta_{r_0^2h}\sigma_o+J_{r_0^2h}\sigma_o)
=r_0^{-1}\rho.
\] All three equations
eq.~\ref{eq:sh-scale-a13}--eq.~\ref{eq:sh-scale-a14}, and the norm
eq.~\ref{eq:sh-scale-a15}, are preserved. These are the component
transformation laws of the same normal tractor. ∎

If instead one keeps \(g_c\) fixed and replaces only \(\sigma\) by
\(r_0\sigma\), the selected metric becomes \(r_0^{-2}g_E\), and the
tractor becomes \(r_0 I\). That is a different constant physical
calibration. At \(\Lambda=0\) the numerical Einstein label remains zero
and cannot detect this error. Proposition K.3 avoids it by changing the
auxiliary representative at the same time.

The auxiliary scalar after eq.~\ref{eq:sh-scale-a20} is
\begin{equation}\protect\phantomsection\label{eq:sh-scale-a23}{
\operatorname{Scal}(g_{\rm aux})=-\frac{12m}{r_0^2r},\qquad
\sigma_o=-\frac{r_0^3}{12m}\operatorname{Scal}(g_{\rm aux}).
}\end{equation} Consequently even a correct scalar/scale formula must
track the auxiliary constant: it is not invariant as an equality of
numerical component functions.

\protect\phantomsection\label{sec:sh-scale-5}

\paragraph{K.1.5. Full spin marking and the optical
root}\label{k.1.5.-full-spin-marking-and-the-optical-root}

The auxiliary normalization above also admits a regular, fully marked
implementation at the horizon. Put
\begin{equation}\protect\phantomsection\label{eq:sh-scale-a24}{\begin{gathered}
a=dv, \\
b=\frac f2dv-dr, \\
\vartheta^0=\frac{a+b}{\sqrt2},\quad
\vartheta^1=\frac{a-b}{\sqrt2},\quad
\vartheta^2=r\,d\theta,\quad
\vartheta^3=r\sin\theta\,d\phi.
\end{gathered}}\end{equation} Then
\begin{equation}\protect\phantomsection\label{eq:sh-scale-a25}{
g_E=(\vartheta^0)^2-(\vartheta^1)^2-(\vartheta^2)^2-(\vartheta^3)^2.
}\end{equation} The forms \(a,b\) are independent even at \(f=0\). On an
angular coordinate patch away from its coordinate poles,
eq.~\ref{eq:sh-scale-a24} is an analytic orthonormal coframe throughout
the crossing neighborhood. A regular angular coframe replaces it on the
other finite patches.

Choose its orientation and time orientation once. Fix a spin lift at
\(o\) and retain its actual label, including the sign distinguishing the
two lifts of the same orthonormal frame. In the lower-index conventions
of §2 \(\epsilon_{01}=\epsilon_{0'1'}=1\), a compatible solder is
\begin{equation}\protect\phantomsection\label{eq:sh-scale-a26}{
\vartheta^{AA'}=\frac1{\sqrt2}
\begin{pmatrix}
\vartheta^0+\vartheta^3&\vartheta^1-i\vartheta^2\\
\vartheta^1+i\vartheta^2&\vartheta^0-\vartheta^3
\end{pmatrix},\qquad
2\det(\vartheta(v))=g_E(v,v).
}\end{equation} The conformal frame identification from \(g_E\) to
\(g_{\rm aux}\) sends its dual vector frame \(e_a\) to
\(\sigma_o^{-1}e_a\), or equivalently its solder coframe to
\(\sigma_o\vartheta^{AA'}\). It preserves the Spin factor of the
conformal frame and changes only the positive dilation. At \(o\) its
dilation is one, so the two marked spin frames coincide under this
specified identification. A conformal weight-half convention can be
lifted using the unique positive square root \(\sqrt{\sigma_o}\) with
value one at \(o\); on a sufficiently small simply connected
complexification, the same root is its analytic continuation.
Multiplication of an internal symplectic spin basis by
\(\sqrt{\sigma_o}\) is not being asserted to be an \(SL(2)\) gauge
change.

Let \(\beta:\mathcal E[1]^{\otimes2}\to\mathcal B\) be the retained
density-square identification and \(\Theta\) the compact-lift optical
derivative of Appendix J.1. Its root condition is
\(\beta(\sigma_o^2)=\Theta^{-1}\) with marked value \(\sigma_o(o)=+1\);
these choices fix the sign of the scale root. Changing it to
\(-\sigma_o\) would keep \(g_E\) but send \(I\) to \(-I\) and change the
selected affine slice, and is therefore excluded by the marked root. The
central spin element instead changes the local-twistor spin lift by
\(-1\), while fixing its exterior-square tractor; it is retained as a
distinct physical spin label.

To compare with marked normal coordinates in the manuscript, apply the
exponential map of the \emph{selected} \(g_E\) using the dual of
eq.~\ref{eq:sh-scale-a24} at \(o\) and the specified spin lift. This
gives \(g_E(o)=\eta\) and its initialized radial spin frame. This same
choice is used on both sides of the metric-mediated comparison. The
Levi-Civita radial frames of \(g_E\) and of \(g_{\rm aux}\) are not
identified away from \(o\). Their difference is
eq.~\ref{eq:sh-scale-a9}, and the normal tractor comparison has already
accounted for it.

\protect\phantomsection\label{sec:sh-scale-6}

\paragraph{K.1.6. Optical section and local-twistor infinity
tensors}\label{k.1.6.-optical-section-and-local-twistor-infinity-tensors}

The normal connection and its norm used above are exactly
\begin{equation}\protect\phantomsection\label{eq:sh-scale-a27}{
\begin{aligned}
D_a(s,\nu_b,q)
&=(\nabla_as-\nu_a,
\nabla_a\nu_b+P_{ab}s+h_{ab}q,
\nabla_aq-P_a{}^b\nu_b),\\
H(s,\nu,q)&=2sq+h^{-1}(\nu,\nu).
\end{aligned}
}\end{equation} Let \(k\) be any \(h\)-affinely parametrized null
tangent on a short ray and put \(\kappa=P^h(k,k)\). Equation
eq.~\ref{eq:sh-scale-a12} implies
\begin{equation}\protect\phantomsection\label{eq:sh-scale-a28}{
\frac{d^2\sigma}{dt^2}+\kappa\sigma=0.
}\end{equation} Thus the same weight-one density gives the optical
section of Proposition 5.3 on every null direction of the complete sky,
rather than only on one chosen principal ray. Conversely, testing this
equation on all null directions forces the trace-free symmetric tensor
\((\nabla_h^2\sigma+P^h\sigma)_0\) to vanish: a quadratic form vanishing
on the nondegenerate null cone is proportional to \(h\). This is the
almost-Einstein recognition step in Proposition 5.4. The raywise
section, its Cauchy derivative, and the parallel prolongation are
consequently those of the very same \(I\).

The selected affine parameter obeys
\begin{equation}\protect\phantomsection\label{eq:sh-scale-a29}{
\frac{d\widehat t}{dt}=\sigma^{-2},
\qquad\text{or}\qquad
\frac{d\widehat t}{dt_o}=\sigma_o^{-2}
\text{ in the representative }g_{\rm aux}.
}\end{equation} More explicitly, if \(t\) is the \(g_c\) affine
parameter calibrated by \(dt/d\widehat t=\sigma^2\), choose the
\(g_{\rm aux}\) affine parameter \(t_o\) with \(dt_o=r_0^2dt\) and the
same calibrated origin. Then \(dt_o/d\widehat t=\sigma_o^2\). Constant
metric rescaling alone does not force an affine-parameter rescaling; it
is this stated Cauchy calibration that selects it. These formulas
therefore describe the same physical affine parameter under Proposition
K.3. A bare numerical comparison of affine parameters without their
initial Cauchy scale would discard retained data.

In the selected Einstein splitting \(P^{g_E}=0\), the rank-four normal
local-twistor bundle has splitting \(S\oplus S'\) and the connection of
Proposition I.3 with off-diagonal solder block \(i\vartheta\) and zero
lower off-diagonal block. The exterior-square tractor map is
\begin{equation}\protect\phantomsection\label{eq:sh-scale-a30}{
\mathcal F(s,\nu,q)=
\begin{pmatrix}qJ_2&-iJ_2\nu\\-i\nu^TJ_2&sJ_2\end{pmatrix},
\qquad
J_2=\begin{pmatrix}0&1\\-1&0\end{pmatrix}.
}\end{equation} Here \(\nu\) is the raised middle component in the
specified solder frame. For \(I=(1,0,0)\), the covariant infinity form
and its volume dual are
\begin{equation}\protect\phantomsection\label{eq:sh-scale-a31}{
I_{\alpha\beta}=\operatorname{diag}(0,J_2),\qquad
I^{\alpha\beta}=\operatorname{diag}(J_2,0),\qquad
I_{\alpha\beta}I^{\beta\gamma}=0.
}\end{equation} They are parallel, each has rank two, and
\(2\operatorname{Pf}(I)=0\). The chosen full spin marking fixes the
volume dual and the exterior-square identification; there is no division
by a cosmological constant. In an auxiliary conformal splitting the
nonzero middle component and the \(\rho\) of eq.~\ref{eq:sh-scale-a7} or
eq.~\ref{eq:sh-scale-a22} must be inserted in eq.~\ref{eq:sh-scale-a30}
with the corresponding solder. One cannot substitute their numerical
triples into the selected-scale block formula eq.~\ref{eq:sh-scale-a31}
without that change of splitting. Appendix J.1 gives this same
normal-spin naturality. The tensors therefore agree on the two sides of
the metric-mediated comparison, including their exact constant
calibration and root.

\protect\phantomsection\label{sec:sh-scale-7}

\paragraph{K.1.7. Uniqueness of the calibrated scale on this
Schwarzschild
neighborhood}\label{k.1.7.-uniqueness-of-the-calibrated-scale-on-this-schwarzschild-neighborhood}

For this fixed nonzero mass, the scale selecting an Einstein metric in
the conformal class is determined up to a constant. This is stronger
than the assertion that arbitrary null-cone data determine a scale.

\pandocstatementheading{\protect\phantomsection\label{prop:sh-scale-4}Proposition K.4.}

Any holomorphic parallel normal tractor on a connected complexification
of \(U\) has, in the Ricci-flat \(g_E\) splitting, the form \((c,0,0)\)
with a constant \(c\). Thus any nonvanishing almost-Einstein density is
a constant multiple of the density in Proposition K.1. In the
representative eq.~\ref{eq:sh-scale-a20}, its value one at \(o\) forces
it to equal \(\sigma_o\).

\emph{Proof.} In the Ricci-flat splitting the parallel equations are
\begin{equation}\protect\phantomsection\label{eq:sh-scale-a32}{
\nabla s=\nu,\qquad \nabla\nu=-qg_E,\qquad dq=0.
}\end{equation} Commuting the middle equation gives
\(R(X,Y)\nu^\sharp=0\) for every \(X,Y\). We verify directly that
Schwarzschild curvature has no vector nullity for \(m\ne0\), rather than
deducing this from a scalar invariant alone. In the exterior static
region use the orthonormal frame \[
e_0=f^{-1/2}\partial_t,\quad e_1=f^{1/2}\partial_r,\quad
e_2=r^{-1}\partial_\theta,\quad e_3=(r\sin\theta)^{-1}\partial_\phi.
\] Let \(a=f'/(2\sqrt f)\) and \(b=\sqrt f/r\). The Koszul formula gives
\(\nabla_{e_0}e_0=ae_1\), \(\nabla_{e_0}e_1=ae_0\),
\(\nabla_{e_1}e_j=0\), and \(\nabla_{e_A}e_1=be_A\),
\(\nabla_{e_A}e_A=-be_1\) for \(A=2\), with the usual additional sphere
connection for \(A=3\). The relevant brackets are \([e_0,e_1]=ae_0\) and
\([e_0,e_A]=0\). Substitution in eq.~\ref{eq:sh-scale-a3} yields
\begin{equation}\protect\phantomsection\label{eq:sh-scale-a33}{
\begin{aligned}
R(e_0,e_1)e_0&=-\frac{f''}{2}e_1=\frac{2m}{r^3}e_1,&
R(e_0,e_1)e_1&=\frac{2m}{r^3}e_0,\\
R(e_0,e_A)e_0&=-\frac{f'}{2r}e_A=-\frac m{r^3}e_A,&
R(e_0,e_A)e_A&=-\frac m{r^3}e_0.
\end{aligned}
}\end{equation} Each displayed curvature operator vanishes on the
remaining frame directions. Thus \(R(e_0,e_1)v=0\) forces \(v^0=v^1=0\),
and the equations with \(A=2,3\) force \(v^2=v^3=0\). Therefore
\(\nu=0\) in the nonempty exterior open subset. Analyticity and
connectedness extend this equality throughout the chosen
horizon-crossing neighborhood and its connected sufficiently small
complexification. Equation eq.~\ref{eq:sh-scale-a32} then gives \(q=0\)
and constant \(s=c\). In the auxiliary representative the component is
\(c/r\), or \(cr_0/r\) after eq.~\ref{eq:sh-scale-a20}. The marked value
one fixes \(c=1\) in the latter representative. ∎

No assertion of this kind is made for the zero-mass limit, where
additional parallel tractors occur. The role of the complete sky in the
optical scale construction and the role of the Einstein equation in
Proposition K.4 are substantive: a point value plus unparametrized null
rays alone would not determine a scale field.

\protect\phantomsection\label{sec:sh-integrated-comparison}

\subsubsection{K.2. Complete records and finite horizon
comparisons}\label{k.2.-complete-records-and-finite-horizon-comparisons}

\protect\phantomsection\label{sec:schw-bc-regular-metric}

\protect\phantomsection\label{sec:sh-comparison-1}

\paragraph{K.2.1. The regular metric and the exact
scope}\label{k.2.1.-the-regular-metric-and-the-exact-scope}

Put \(f(r)=1-2m/r\). On an ingoing Eddington--Finkelstein neighborhood
let \begin{equation}\protect\phantomsection\label{eq:schw-bc-ef-metric}{
 g_m=f(r)\,dv^2-2\,dv\,dr-r^2\gamma,
 \qquad \gamma=g_{S^2}.
}\end{equation} Fix \(0<\delta<2m\) and finite \(v_-<v_+\). The compact
physical core is
\begin{equation}\protect\phantomsection\label{eq:schw-bc-physical-core}{
 K=[v_-,v_+]\times[2m-\delta,2m+\delta]\times S^2.
}\end{equation} Choose an open real-analytic parent of \(K\) with
\(r>r_->0\) and finite angular charts with compact subcores. Complexify
these charts on smaller parents when using the holomorphic results of
the manuscript. All comparisons in this appendix are short
representatives near points of this parent, or finite compositions with
explicitly retained common parents. The horizon is the hypersurface
\(r=2m\) in this given Schwarzschild region.

\pandocstatementheading{\protect\phantomsection\label{lem:schw-bc-ef-vacuum}Lemma K.5 (regular vacuum and regular full-spin frames).}

The metric in eq.~\ref{eq:schw-bc-ef-metric} is nondegenerate and real
analytic across \(r=2m\), has Lorentz signature \((+---)\), and has zero
Ricci tensor. There is a finite system of real-analytic orthonormal
frames over the parent of \(K\) with compatible full-spin lifts. The
induced complexified frames, density calibration and sky-root
comparisons are regular after finite restrictions of their parents.

\emph{Proof.} The \((v,r)\) block and its inverse are
\begin{equation}\protect\phantomsection\label{eq:schw-bc-ef-blocks}{
 \begin{pmatrix}f&-1\\-1&0\end{pmatrix},\qquad
 \begin{pmatrix}0&-1\\-1&-f\end{pmatrix}.
}\end{equation} Its determinant is \(-1\) for every \(r>0\). Thus no
factor \(f^{-1}\) is required. The angular block is negative definite on
the real sphere. For a direct curvature check, temporarily allow an
arbitrary function \(f(r)\). In angular coordinates the non-spherical
Christoffel symbols, together with their lower-index symmetries, are
\begin{equation}\protect\phantomsection\label{eq:schw-bc-ef-connection}{
 \begin{aligned}
 \Gamma^v_{vv}&=\tfrac12f',&
 \Gamma^r_{vv}&=\tfrac12ff',&
 \Gamma^r_{vr}&=-\tfrac12f',\\
 \Gamma^v_{AB}&=-r\gamma_{AB},&
 \Gamma^r_{AB}&=-fr\gamma_{AB},&
 \Gamma^A_{rB}&=r^{-1}\delta^A_B.
 \end{aligned}
}\end{equation} The remaining angular symbols are those of \(\gamma\).
Substituting these into the curvature convention above gives, with
\(a=f''/2+f'/r\),
\begin{equation}\protect\phantomsection\label{eq:schw-bc-ef-ricci}{
 \operatorname{Ric}_{vv}=fa,\quad
 \operatorname{Ric}_{vr}=-a,\quad
 \operatorname{Ric}_{rr}=0,\quad
 \operatorname{Ric}_{AB}=(1-f-rf')\gamma_{AB},
}\end{equation} and zero mixed angular components. For \(f=1-2m/r\),
both \(a\) and \(1-f-rf'\) vanish. These formulas also fix the sign of
the check: replacing \(f\) by \(1-cr^2\) gives
\(\operatorname{Ric}=-3c\,g\) in the convention of §2.

For regular frames set \(a_0=(f-1)/2\), \(b_0=(f+1)/2\), and choose an
oriented orthonormal sphere coframe \((\eta^2,\eta^3)\) on each angular
chart. One explicit choice on a stereographic chart is
\(\eta^2=2\,dx/(1+x^2+y^2)\), \(\eta^3=2\,dy/(1+x^2+y^2)\). Finitely
many smaller oriented stereographic disks with compact subcores cover
the sphere; their complex thickenings can be chosen to keep this
denominator nonzero. Then
\begin{equation}\protect\phantomsection\label{eq:schw-bc-regular-coframe}{
 e^0=b_0\,dv-dr,\qquad
 e^1=dr-a_0\,dv,\qquad
 e^2=r\eta^2,\qquad e^3=r\eta^3
}\end{equation} satisfies \(g_m=(e^0)^2-(e^1)^2-(e^2)^2-(e^3)^2\).
Indeed \(b_0-a_0=1\), \(a_0+b_0=f\), and the determinant of the first
two coframe rows is \(1\). Its dual timelike and radial spacelike
vectors are \[
 T=\partial_v+a_0\partial_r,\qquad
 S=\partial_v+b_0\partial_r.
\] They obey \(g_m(T,T)=1\), \(g_m(S,S)=-1\), \(g_m(T,S)=0\) also at
\(f=0\). Choose the time orientation in which \(T\) is future directed.
These are regular analytic frames on the finite angular parents, without
the static square root \(\sqrt f\).

On a common angular chart with the same angular frame, let \(e_0,e_1\)
be dual to the first two coframe elements in eq.~\ref{eq:sh-scale-a24}.
Direct substitution gives
\begin{equation}\protect\phantomsection\label{eq:sh-r29-radial-frame-boost}{
 \begin{pmatrix}T\\S\end{pmatrix}
 =\frac1{2\sqrt2}
  \begin{pmatrix}3&1\\1&3\end{pmatrix}
  \begin{pmatrix}e_0\\e_1\end{pmatrix}.
}\end{equation} This is the proper time-oriented boost of rapidity
\(\log\sqrt2\). Its two full-spin lifts retain the distinction fixed by
the specified marks.

The oriented time-oriented product region has a spin structure: its
tangent bundle is the sum of a trivial real rank-two bundle and
\(TS^2\), and \(w_2(TS^2)\) is the Euler class \(2\) modulo \(2\). Fix
that spin structure and take local analytic sections of its full
principal spin bundle on a finite good cover. Their transition functions
satisfy the spin cocycle. This constructs compatible lifts first; it
does not choose unrelated signs on different double overlaps. Their
local complexifications extend after finite shrinkage. The unit physical
density for \(g_m\) fixes the density calibration. The spinor-line
construction of the sky roots, with the base factor retained as in G.5
and H.4, is functorial for these full-spin transitions. Its root
comparisons therefore extend on the same finite system after a further
finite restriction. ∎

For orientation with the usual exterior chart, on any retained overlap
with \(r\ge 2m+\varepsilon\) one may use \(v=t+r_*(r)\),
\(r_*'=f^{-1}\). Direct substitution gives
\(g_m=f\,dt^2-f^{-1}dr^2-r^2\gamma\). Both directions of this coordinate
comparison are analytic on such an exterior overlap. The EF charts,
rather than this exterior change, cover \(r=2m\) and the inside part of
\(K\).

\protect\phantomsection\label{sec:schw-bc-all-records}

\protect\phantomsection\label{sec:sh-comparison-2}

\paragraph{K.2.2. The complete records and the initialized metric
map}\label{k.2.2.-the-complete-records-and-the-initialized-metric-map}

Let \(p\) belong to the parent of \(K\) and let \(s\) be a chosen
full-spin frame above the orthonormal frame at \(p\). In that frame
extract \(\Psi^\pm\) from the covariant Weyl tensor by the fixed
manuscript's readout
\begin{equation}\protect\phantomsection\label{eq:schw-bc-spin-readout}{
 C_{AA'BB'CC'DD'}
 =-\Psi^+_{ABCD}\epsilon_{A'B'}\epsilon_{C'D'}
  -\Psi^-_{A'B'C'D'}\epsilon_{AB}\epsilon_{CD}.
}\end{equation} For every \(r\ge0\), define
\begin{equation}\protect\phantomsection\label{eq:schw-bc-all-record-extraction}{
 U_r^+(p,s)=\operatorname{Sym}_{A_1\cdots A_{r+4}}
 \operatorname{Sym}_{A'_1\cdots A'_r}
 \nabla_{A_1A'_1}\cdots\nabla_{A_rA'_r}
 \Psi^+_{A_{r+1}\cdots A_{r+4}}\big|_{(p,s)}.
}\end{equation} The negative formula exchanges the two spin factors.
Each symmetrization is an average, exactly as in Theorem 3.1. This is
one definition of every order using the connection of \(g_m\); it
introduces no separate guessed formula for higher derivatives. On the
physical real slice it gives \(U^-_r=\overline{U^+_r}\) in the coupled
spin identification.

\pandocstatementheading{\protect\phantomsection\label{lem:schw-bc-record-family}Lemma K.6 (analytic all-record extraction).}

The records in eq.~\ref{eq:schw-bc-all-record-extraction} satisfy the
factorial evaluation bounds of Theorem 3.1. They do so uniformly on each
fixed finite family of buffered complex point/frame parents with bounded
frame and inverse. Their exact Theorem 3.1 encoders define analytic
classes \(A_{p,s}\in\mathcal A_0\), holomorphically in those parameters.
An actual isometry with a specified full-spin lift transports every one
of these records by the same point/full-spin rule.

\emph{Proof.} Solve the analytic geodesic and parallel-spin-frame
equations starting at \((p,s)\) with initial direction
\(k=\xi\otimes\pi\), where \(|\xi|,|\pi|\le1\). For a compact parameter
core of this type, holomorphic ODE existence gives a common complex
affine disk \(|t|<R\) and a common bound \(C\) on the Weyl components
evaluated in the transported frame. Such a disk exists because the
initial data form a compact set inside finitely many analytic parents,
the metric and inverse are bounded there, and the initial values have
positive distance to their boundaries. The \(r\)th derivative at \(t=0\)
is \(\nabla_k^r\Psi^\pm\) in that parallel frame. Its coherent
contraction is exactly \(U_r^+(\xi^{r+4},\pi^r)\) or its negative
counterpart: repeated identical spinors make the averaged
symmetrizations invisible. Cauchy's formula gives
\begin{equation}\protect\phantomsection\label{eq:schw-bc-record-bound}{
 \|U_r^\pm(p,s)\|_{\mathrm{ev}}\le C\,r!R^{-r}.
}\end{equation} This also follows by the metric-to-record argument in §4
immediately after Proposition 4.1, now on the explicit finite EF
parents. Theorem 3.1, including its fixed factorial multipliers and
Laurent projection, produces the two classes on common smaller fiber
disks and annuli. Holomorphic dependence follows from the same ODE and
uniform Cauchy estimates. All orders use the same contraction convention
eq.~\ref{eq:schw-bc-spin-readout}. An isometry intertwines Levi--Civita
connections, curvature and every iterated covariant derivative; its
specified spin lift intertwines the two spinor contractions and averaged
projections. This proves the last assertion for all \(r\), not just a
finite prefix. ∎

Write \(g[A_{p,s}]\) for the generated marked metric and let
\(s:\mathbb C^4\to T_pM\) denote the soldering identification supplied
by the mark. Here and below, \(g[A]\) is identified with the canonical
marked germ \(g_A\) of Proposition 4.1 through the initialized marked
isometry of Proposition A.4 and Theorem 6.3; the canonical bundle and
operator constructions are transported through that isometry. Define the
physical normal map and the generated normal map by \[
 n_{p,s}(x)=\operatorname{Exp}^{g_m}_p(sx),\qquad
 n_A(x)=\operatorname{Exp}^{g[A]}_o(x).
\] They are holomorphic embeddings on sufficiently small common
representatives. Theorem 2.4 gives the exact entire records of \(g[A]\),
with scalar \(\Lambda=0\). Proposition A.4 therefore gives equality of
the two pulled-back analytic normal metrics and the actual initialized
isometry
\begin{equation}\protect\phantomsection\label{eq:schw-bc-normal-isometry}{
 f_{p,s}=n_{p,s}\circ n_A^{-1}:
 (g[A_{p,s}],o,\text{full spin})\xrightarrow{\sim}(g_m,p,s).
}\end{equation} For specificity, A.4 obtains the normal metric from the
unique radial Jacobi equations with curvature Taylor coefficients given
by the completed ordered record. Equality of all the extracted free
records gives equality of those ordered coefficients by A.3, equality of
the radial Jacobi Taylor series, and equality of convergent metric
series. Thus eq.~\ref{eq:schw-bc-normal-isometry} is an actual analytic
map on common short representatives with its initialized spin lift and
inverse.

\protect\phantomsection\label{sec:schw-bc-typed-comparison}

\protect\phantomsection\label{sec:sh-comparison-3}

\paragraph{K.2.3. The actual enriched
comparison}\label{k.2.3.-the-actual-enriched-comparison}

\pandocstatementheading{\protect\phantomsection\label{thm:schw-bc-metric-mediated-comparison}Theorem K.7 (marked metric-mediated comparison).}

For the extracted input \(A=A_{p,s}\), the source-generated contact,
calibrated affine and projective operator diagram has an actual
comparison with \(\operatorname{Can}(g_m,p,s)\). It preserves the full
compact point family, its modules, the optical evaluation and Wronskian,
the density root and selected scale, the normal tractor and
local-twistor transport, the calibrated infinity pair and point planes.
The comparison and its inverse retain the specified physical full-spin
label. They are holomorphic on conditioned common parents and satisfy
the finite coherence laws on conditioned diagrams of Definition 2.1.

\emph{Proof.} We describe the comparison induced by the initialized
isometry \(f=f_{p,s}\), and then its operator representative.

\textbf{Contact geometry.} On the nonzero null cotangent cone use the
cotangent lift
\begin{equation}\protect\phantomsection\label{eq:schw-bc-cotangent-map}{
 (x,\ell)\longmapsto
 \bigl(f(x),(df_x^{-1})^*\ell\bigr).
}\end{equation} It preserves the tautological one-form, the metric null
Hamiltonian, homogeneous scaling and its Hamiltonian characteristic
curves. On short ray flow boxes it therefore descends to a holomorphic
contact map \(C_f\) with inverse \(C_{f^{-1}}\). It sends the whole
compact sky \(Q_x\) to \(Q_{f(x)}\). In contact charts set
\(K_i=\psi_i^t C_f(\psi_i^s)^{-1}\) and
\(K_i^*\alpha_i^t=d_i\alpha_i^s\). Then \(d_i\) is a holomorphic unit
and
\begin{equation}\protect\phantomsection\label{eq:schw-bc-contact-comparison}{
 \Phi^t_{ij}=K_i\Phi^s_{ij}K_j^{-1}.
}\end{equation} The map is induced on every null direction, so its
definition includes the complete compact sky rather than only one affine
direction chart.

\textbf{Optical solution space and affine calibration.} The optical
bundle \(\mathcal E_J\) here is the rank-two bundle of scalar solutions
of \(D^2y+P(k,k)y=0\), with the contact-line-valued Wronskian specified
in §5 and Appendix D. It is distinct from the full transverse Jacobi
phase space used in the transverse screen calculation. The metric
isometry and its ray-parameter/line identifications carry this
differential equation and its Cauchy evaluation to the same equation.
Denote the resulting solution map by
\(\mathsf S:\mathcal E_J^s\to\mathcal E_J^t\). In transported Cauchy
frames this is the identity; in arbitrary retained frames it is the
unique map intertwining value and derivative at a transversal.
Uniqueness for the scalar ODE proves intertwining throughout the short
connected ray. Wronskians agree under the same contact-line map, and
\(\mathsf S\) carries the selected density solution \(s_s\) to \(s_t\)
because \(f^*g_m=g[A]\) with the retained root calibration.

The affine slice has the intrinsic description
\begin{equation}\protect\phantomsection\label{eq:schw-bc-affine-solution-map}{
 A_{s_s}=\{\lambda\in(\mathcal E_J^s)^*: \lambda(s_s)=1\},
 \qquad
 \lambda\longmapsto\lambda\circ\mathsf S^{-1}\in A_{s_t}.
}\end{equation} This map carries the evaluation
\(\epsilon_x/\sigma_s(x)\) to \(\epsilon_{f(x)}/\sigma_t(f(x))\). By
Proposition J.4 it is the same affine comparison as that obtained from
the compact incidence. In §8.1's affine charts its point map is
\begin{equation}\protect\phantomsection\label{eq:schw-bc-affine-chart-map}{
 \widehat K_i(z,\tau)
  =\bigl(K_i(z),d_i(z)^{-1}\tau+u_i(z)\bigr),\qquad
 u_i=\tau_i^t\circ f^{\mathcal F}-d_i^{-1}\tau_i^s.
}\end{equation} The latter expression is evaluated at any one retained
transversal hit. It is independent of that hit: the identities
\(f^{\mathcal F}_*(\sigma_s^2D_s)=d_i^{-1}\sigma_t^2D_t\) and
\(D_a\tau_i^a=\sigma_a^{-2}\) give \(D_su_i=0\), with target quantities
evaluated at the image. Consequently \[
 \widehat\Phi^t_{ij}
 =\widehat K_i\widehat\Phi^s_{ij}\widehat K_j^{-1}.
\] The compact lifts, affine origins and projective extensions are
carried by this same map. The selected root is transported as a
line/root map with its central calibration, not recovered merely from
the square of a function. If two such root comparisons cover the same
squared comparison and have the retained central value, their ratio is a
holomorphic unit \(q\) with \(q^2=1\) and \(q(o)=1\). On a smaller
connected parent \(q+1\) is a unit, hence \(q=1\). Thus the calibration
removes a possible extra sign.

\textbf{Normal transport, infinity and point planes.} The isometry and
its specified full-spin lift induce the connection isomorphisms on the
normal tractor and local-twistor bundles. In a metric splitting the
tractor map is the identity on the scalar slots and the induced
cotangent map on the middle slot. Substitution in J.1's formula for
\(D\) verifies the intertwining, since \(f\) preserves \(\nabla\), \(P\)
and \(g\). It preserves \(H\), \(X\), the selected density and
\(I=L\sigma\). In the physical vacuum scale \(\sigma=1\), \(P=0\), so
\(I=(1,0,0)\) and \(H(I,I)=0\). No division by this zero norm is used.
Proposition I.3's spin exterior-square identification is natural for the
full-spin lift, hence carries both infinity tensors, including their
normalization, and all raywise local-twistor flags. The point-plane
kernels in J.4 are carried by the same linear map. Evaluation and
\(X/\sigma\) agree with eq.~\ref{eq:schw-bc-affine-solution-map}, so
there is no separate unmatched choice of affine lift.

\textbf{Operators, associators and point modules.} Form the comparison
\begin{equation}\protect\phantomsection\label{eq:schw-bc-enriched-map}{
 \mathcal P(\mathcal S_{\mathrm{op}}A)
 \xrightarrow{\alpha_A}\operatorname{Can}(g[A])
 \xrightarrow{\operatorname{Can}(f)}\operatorname{Can}(g_m,p,s),
}\end{equation} where \(\alpha_A\) is the actual comparison of Theorem
J.6. For the second arrow, H.4 quantizes the symbol map \((K_i^{-1})^*\)
in the calibrated projective Heisenberg charts, retaining the chosen
root, and G.5 supplies the framed module maps. More explicitly, on a
common finite refinement there are chart algebra maps \(k_i\),
normalized units \(U_{ij}\) and point-module maps \(F_i\) satisfying
\begin{equation}\protect\phantomsection\label{eq:schw-bc-actual-operator-comparison}{
 \begin{aligned}
 k_i\varphi^s_{ij}
    &=\operatorname{Ad}(U_{ij})\varphi^t_{ij}k_j,\\
 U_{ij}\varphi^t_{ij}(U_{jk})C^t_{ijk}
    &=k_i(C^s_{ijk})U_{ik},\\
 F_i\tau^s_{ij}&=U_{ij}\tau^t_{ij}F_j.
 \end{aligned}
}\end{equation} Here \(\varphi\) and \(C\) are the overlap maps and
associators, and \(\tau\) are the module overlaps. All base changes and
moving arguments are those of H.4 and Theorem J.6. The second equation
retains both associators; replacing it by a principal-symbol equality
would be insufficient. Compact embedded support classification and the
framed module construction identify the module supported on \(Q_x\) with
the module on \(Q_{f(x)}\) over its whole compact support. G.5 and H.4's
normalized-inner and framed-Hom uniqueness give inverses and coherent
finite composites. The geometric components are those already described,
by Theorem J.6. In particular, a central full-spin label is retained
even if some tensor components factor through its orthogonal image.

All constituent operations are the existing analytic normal maps, ray
hit maps, root maps, finite matrix transports and H.4 finite
refinements. Under the conditioned-parent hypotheses their finitely many
restrictions retain common positive domains. This proves the asserted
parameter and finite-coherence statements. ∎

\protect\phantomsection\label{sec:schw-bc-finite-coherence}

\protect\phantomsection\label{sec:sh-comparison-4}

\paragraph{K.2.4. Finite representatives across the future
horizon}\label{k.2.4.-finite-representatives-across-the-future-horizon}

\pandocstatementheading{\protect\phantomsection\label{thm:schw-bc-finite-horizon}Theorem K.8 (finite horizon system).}

The core \(K\) in eq.~\ref{eq:schw-bc-physical-core} admits a finite
system of short analytic normal representatives, their generated
contact/affine/operator objects, and actual physical comparisons
covering its exterior, horizon and interior portions. The system can be
chosen with common positive buffers and unit margins on its retained
compact subcores, including the required inverse maps and every
specified finite double/triple comparison. It carries compatible
full-spin lifts, roots and density calibration, compact point modules
and actual operator associators. Its finite composition diagrams have
the coherence of Definition 2.1, Theorem 8.3 and H.4. All assertions are
relative to this finite system; no global ray quotient or global
operator ring is constructed.

\emph{Proof.} We give the parent-domain construction before the overlap
identities.

\textbf{Short common parents.} Enlarge \(K\) slightly inside \(r>0\) and
use the finite real-analytic angular frame parents of Lemma K.5. On
their compact subcores the metric, inverse, connection, chosen frames
and inverse frames are bounded. The analytic exponential and
parallel-spin-transport ODEs depend holomorphically on their point/frame
parameters. At zero displacement the exponential differential is the
initialized isomorphism. The analytic inverse theorem with compact
parameters therefore gives a common strictly positive short radius after
finite restriction. More concretely, the differential remains uniformly
close to the initial identity in the finite coordinate trivializations
on a small complex polydisk; shrinking that disk yields injective normal
maps and bounded inverses on smaller disks with a reserved buffer.

Cover \(K\) by finitely many small cores inside these normal parents.
Choose the cores small enough relative to the retained parent radius
that if two cores meet, each center and the relevant double/triple
overlap lie within the larger normal parents needed for the comparison.
Restricting an overlap at any of its points then gives an actual short
physical comparison. Its normal displacement, spin transport and inverse
stay inside the retained parents. This is a finite covering
construction, not a lower bound on the injectivity radius of every
possible long null segment.

The bundle of whole complex skies over each compact point core is
compact: its fiber is \(\mathbb P^1\times\mathbb P^1\). Cover it by
finitely many bounded direction charts, normalizing a nonzero direction
component in each. The null-geodesic vector field is nonzero there. At
each central incidence point choose a transversal coordinate with
nonzero directional derivative; shrink so this derivative has a positive
lower bound on the chosen compact subcore. The implicit-function theorem
for the short hit time gives a ray flow box and its inverse with a
positive buffer. Compactness provides finitely many such boxes and a
common positive short bound. Their transitions describe the entire
compact sky, including directions missing in one affine spinor chart.
This is the geometric hypothesis used by Proposition 4.2 and Appendix D,
now verified for the EF parents.

Lemma K.16 below is local in the physical displacement. Its application
to the finite core \(K\) of Theorem K.8 is as follows. Enlarge \(K\)
slightly within \(r>0\) and keep its finite analytic point/full-spin
frame parents. Lemma K.6 gives uniformly bounded record families
\(U(p,s)\) on their compact subcores. Apply the preceding construction
also with \((p,s)\) as passive parameters. Compactness and finitely many
frame charts give a positive common local displacement radius for the
fixed finite comparison shapes, including the double/triple and required
higher coherence diagrams. This radius controls small local stars of
point charts; it is not a normal-coordinate radius from one chosen point
that contains all of \(K\).

Now choose the finite physical cover of \(K\) with cores small relative
to these local radii and with larger buffered normal parents. It may be
refined before selecting the final finite list of charts. On every
nonempty local comparison star, take one center as a local reference
mark and express the other marks and frames by small displacements and
actual initialized spin transports from that center. At physical
amplitude one these are the prescribed Schwarzschild point/frame
records. The compatible family in that star is
\begin{equation}\protect\phantomsection\label{eq:schw-bc-local-star-families}{
 U^{(\alpha)}_\varepsilon=\varepsilon U_\alpha,\qquad
 V^{(\alpha)}_\beta(\varepsilon)
      =R_{a_{\beta\alpha},A_{\beta\alpha}}
                                    (U^{(\alpha)}_\varepsilon).
}\end{equation} All finite-star parents are chosen before their passive
marks are specialized. The local cover may be chosen sufficiently fine
that the centers in a common double/triple comparison lie inside the
relevant passive displacement balls. The same construction with finitely
many extra passive vertices handles any additional specified finite
local composite. There is no quantifier over arbitrary unbounded words.

At \(\varepsilon=1\), a vertex obtained from either incident star has
exactly the same complete marked Schwarzschild records. Its prescribed
generated object is the one for those records. If auxiliary
representative choices differ, H.4 and Theorem J.6 give their actual
normalized comparisons. Thus the physical operator/module lifts
furnished by G.5 at amplitude one compare the original vertex objects.
The selected composite and the composite of selected lifts have the same
actual physical symbol and full-spin label. Lemma F.8's normalized-inner
uniqueness and the whole-family framed-Hom uniqueness then give the
actual comparison units and all finite coherence equations. This
argument does not require the local star families to glue to a single
spacetime metric for \(\varepsilon\ne1\).

The compact point cores still cover \(K\): for any required point, the
sufficiently small chart cores and their larger parents can be chosen
around that point before the finite subcover is taken. Subsequent
simultaneous closed-star refinements retain those whole compact cores
and the complete compact sky fibers. All norms have finite maxima and
all actually inverted units and composite buffers have positive minima
over the finite system. This proves the amplitude and
structured-comparison part of Theorem K.8 on its original finite core.
It does not replace the exact physical transition maps by gluing
independently truncated charts.

In particular, the older ray family from \(r=5m/2\) to \(r=3m/2\) can
remain in that finite physical cover. It is not thereby contained in one
root germ's common physical amplitude parent. Nor are independent
families \(\varepsilon U_\alpha\) at different marks automatically
physical comparisons of one another.

\textbf{Physical double and triple maps.} Write
\(f_\alpha:g[A_\alpha]\to U_\alpha\subset(M,g_m)\) for
eq.~\ref{eq:schw-bc-normal-isometry}. On the retained overlaps the
comparison from chart \(\alpha\) to chart \(\beta\) is
\begin{equation}\protect\phantomsection\label{eq:schw-bc-physical-chart-coherence}{
 t_{\beta\alpha}=f_\beta^{-1}f_\alpha,\qquad
 t_{\gamma\beta}t_{\beta\alpha}=t_{\gamma\alpha},\qquad
 t_{\alpha\beta}=t_{\beta\alpha}^{-1}.
}\end{equation} Use the full-spin frame transitions of the already fixed
spin bundle for their lifts. Those lifts satisfy the same identities,
because they are changes of its actual local sections. To describe an
overlap by Theorem 8.3's instructions, take \(a\) to be the short normal
displacement from the \(\beta\) mark to the \(\alpha\) mark and take
\(A\) so that radial spin transport followed by \(A\) is the actual
\(\alpha\) frame. Then \(U_\alpha=R_{a,A}U_\beta\) and
\(t_{\beta\alpha}=\Xi_{a,A}\) on their common representative. Composite
instructions use
\begin{equation}\protect\phantomsection\label{eq:schw-bc-full-instruction-composition}{
 (b,B)\circ(a,A)
   =\bigl(h_U(a,Ab),K_U(a,Ab)AB\bigr),
}\end{equation} with the source-dependent radial/broken-path spin
correction \(K_U\) of §8.3. The physical point-map direction is the one
specified there. In particular, we have not replaced this law by
addition of displacements or by multiplication \(AB\) alone. The spin
holonomy around a curved path is allowed to be nontrivial. The strict
chart-transition identity eq.~\ref{eq:schw-bc-physical-chart-coherence}
follows from changes of sections, not from path independence of the
connection.

The cotangent lifts and descended contact maps compose exactly with
eq.~\ref{eq:schw-bc-physical-chart-coherence}. For successive
source-to-target affine maps \((K_{12},d_{12},u_{12})\) and
\((K_{23},d_{23},u_{23})\), direct substitution in
eq.~\ref{eq:schw-bc-affine-chart-map} gives
\begin{equation}\protect\phantomsection\label{eq:schw-bc-affine-coherence}{
 \begin{aligned}
 K_{13}&=K_{23}K_{12},\\
 d_{13}&=(d_{23}\circ K_{12})d_{12},\\
 u_{13}&=(d_{23}^{-1}\circ K_{12})u_{12}
                   +u_{23}\circ K_{12}.
 \end{aligned}
}\end{equation} These equations include the moving arguments. The
inverse is the inverse affine map, with translation
\(-(d_{12}u_{12})\circ K_{12}^{-1}\). The full-spin maps carry the root
and its calibration, hence the optical section, infinity tensors and
affine slice agree on triples by functoriality, with no sign selected
afresh on the triple. The compact support \(Q_x\) and its lifted
evaluation are carried to the same compact support and evaluation under
both routes.

\textbf{Finite operator refinement and module coherence.} Take the
compact union of whole skies over the enlarged point cores and all
stipulated finite comparison images. It is compact. Apply H.4's finite
buffered closed-star refinement to the contact and root charts and to
the comparison maps and their inverses simultaneously. Intersect the
finitely many required preimages of parent neighborhoods so every
stipulated composite remains inside its assigned parent with a positive
buffer. This does not remove points from the required physical core: at
any retained physical point all chart-transition routes represent that
same point in charts whose larger parents contain it. Their finitely
many preimages contain a neighborhood of that point; a finite refinement
of those neighborhoods still covers \(K\) and its whole skies. This is
precisely the finite-conditioned-diagram procedure of Definition 2.1. It
does not assert that a quantizer exists on an arbitrary previously
chosen full overlap without refinement.

On this finite system each \(t_{\beta\alpha}\) has actual
operator/module comparisons as in
eq.~\ref{eq:schw-bc-actual-operator-comparison}. On a triple choose
their normalized comparison between the two routes. At the chart-algebra
level denote the equivalences by \(L_{\beta\alpha}\) and the actual
implementer by \(P_{\gamma\beta\alpha}\). On common refined charts they
satisfy
\begin{equation}\protect\phantomsection\label{eq:schw-bc-operator-triple}{
 L_{\gamma\beta}L_{\beta\alpha}
   =\operatorname{Ad}(P_{\gamma\beta\alpha})L_{\gamma\alpha},
 \qquad
 M_{\gamma\beta}M_{\beta\alpha}
   =P_{\gamma\beta\alpha}M_{\gamma\alpha}
}\end{equation} for the framed point-module maps \(M\). These equations
are in addition to the internal chart-overlap associators \(C_{ijk}\)
retained in eq.~\ref{eq:schw-bc-actual-operator-comparison}. The two
routes through a fourfold composite have the equality
\begin{equation}\protect\phantomsection\label{eq:schw-bc-operator-quadruple}{
 P_{\delta\gamma\beta}P_{\delta\beta\alpha}
 =L_{\delta\gamma}(P_{\gamma\beta\alpha})P_{\delta\gamma\alpha}.
}\end{equation} To see this equality, associate
\(L_{\delta\gamma}L_{\gamma\beta}L_{\beta\alpha}\) in its two ways. The
displayed sides implement the same inner action over the same physical
map and full-spin label. Both have the normalization of Lemma F.8. Their
equality follows from normalized-inner uniqueness; the framed module
equality then follows from the framed-Hom uniqueness in H.4. The same
uniqueness gives identities, inverse comparisons and every retained
finite composition diagram. Whole compact point modules, the
local-twistor flags and calibrated infinity forms are preserved by
Theorem J.6, including their module actions. The full-spin center
remains part of the label in these equalities.

\textbf{Margins and reality.} The geometric maps above are analytic
isomorphisms on their parents. On smaller compact subcores every
coordinate or frame determinant, contact multiplier, scalar
root-transition coefficient and hit-time derivative that is used as a
unit is nonzero. Thus the absolute values of these units have positive
minima, the inverse frame matrices are bounded, and the core images have
positive distance from the finite parent boundaries. Individual entries
of an invertible frame matrix need not be nonzero. The explicit \(r\)
denominators stay bounded away from zero by \(r_->0\); \(f=0\) is not a
denominator. There are finitely many stipulated maps, inverses and
composites, so the minima of their positive buffers and unit margins are
positive and the maxima of their bounds are finite. The source common
parents follow from Lemma K.6 and Theorem 3.1; the prescribed generated
representatives use §6; the finite operator refinement uses H.4. These
observations verify every type of conditioning item in Definition 2.1
for this constructed finite diagram.

The real analytic metric and chosen real spin structure give the coupled
Lorentz operation. Its induced operator comparison is the structured
Real comparison of G.5 and J.7, with its retained square isomorphism. We
do not replace that square by a claim that every local representative is
strictly involutive. Refining the finite cover with its conjugate
carries the same maps and compact cores and preserves the positive
margins. This proves the theorem. ∎

\protect\phantomsection\label{sec:sh-integrated-optics}

\subsubsection{K.3. Nonradial screen transport and stationary
relations}\label{k.3.-nonradial-screen-transport-and-stationary-relations}

\protect\phantomsection\label{sec:sh-optics-1}

\paragraph{K.3.1. Conventions and the two optical
systems}\label{k.3.1.-conventions-and-the-two-optical-systems}

Use the auxiliary metric \[
 h=-f(r)\,dv^2+2\,dv\,dr+r^2(d\theta^2+\sin^2\theta\,d\phi^2),
 \qquad f(r)=1-\frac{2m}{r},
\] of signature \((-,+,+,+)\). The metric in the signature of the
correspondence is \(g=-h\). We use \[
 R(X,Y)=[\nabla_X,\nabla_Y]-\nabla_{[X,Y]},\qquad
 \operatorname{Ric}(Y,Z)=\operatorname{tr}(X\mapsto R(X,Y)Z).
\] Metric negation leaves the connection, the curvature operator, the
null geodesics with their affine parameters, and the covariant Ricci
tensor unchanged. On the screen \(k^\perp/\langle k\rangle\), its
positive definite inner product is \(h=-g\). If \(e_1,e_2\) are parallel
and \(h\)-orthonormal screen vectors, put \[
 \mathcal R_{ij}=h(e_i,R(e_j,k)k)=-g(e_i,R(e_j,k)k).
\] Then the transverse Jacobi equation is \(x''+\mathcal R x=0\). This
convention agrees with the curvature-operator Jacobi equation of the
correspondence.

The correspondence's rank-two scalar optical bundle \(\mathcal E_J\)
consists of solutions of \(y''+\kappa y=0\), where
\(\kappa=\tfrac12\operatorname{Ric}(k,k)\) in an affine metric
calibration; see Proposition 5.3 and its proof. It is a two-dimensional
scalar solution space. In this Ricci-flat example its affine equation is
\(y''=0\).

The transverse screen has dimension two at a point. Its Jacobi solution
space, or phase space of \((x,x')\), has dimension four. It retains the
trace-free Weyl tidal term even when \(\kappa=0\). These three spaces
are distinct. The calculation below concerns the four-dimensional screen
phase space.

\protect\phantomsection\label{sec:sh-optics-2}

\paragraph{K.3.2. An explicit compact family of finite crossing
rays}\label{k.3.2.-an-explicit-compact-family-of-finite-crossing-rays}

Let \(u\) be affine parameter, and impose \[
 \begin{gathered}
u_0=0,\quad r(0)=\frac52m,\quad v(0)=0,\quad
 \theta=\frac\pi2,\quad \phi(0)=0, \\
1\le E\le2,\quad \frac14\le\ell:=\frac Lm\le\frac12.
\end{gathered}
\] Choose the future ingoing branch. The conserved quantities and the
null equation give \[
 E=f\dot v-\dot r,\qquad L=r^2\dot\phi,\qquad
 q(r)=\sqrt{E^2-f(r)\frac{L^2}{r^2}},
\] \begin{equation}\protect\phantomsection\label{eq:sh-optics-d-1}{
 \dot r=-q,\qquad
 \dot v=\frac{L^2}{r^2(E+q)},\qquad
 \dot\phi=\frac L{r^2}.
}\end{equation} The expression for \(\dot v\) is rationalized before
crossing the horizon. It is regular at \(r=2m\) and gives
\(\dot v=L^2/(8Em^2)\) there. Moreover \(E=f\dot v+q\) holds, so the
stated energy has the same value on both sides. These equations solve
the geodesic equation: conservation of the cyclic momenta gives the
\(v,\phi\) Euler--Lagrange equations, and differentiation of the null
first integral gives the radial equation because \(\dot r\ne0\). The
equatorial plane is totally geodesic.

Put \(t=u/m\), \(\rho=r/m\), and \(V=v/m\). On \(3/2\le\rho\le5/2\), \[
 a(\rho)=\frac{\rho-2}{\rho^3},\qquad
 a'(\rho)=\frac{6-2\rho}{\rho^4}>0,\qquad
 -\frac4{27}\le a\le\frac4{125}.
\] Consequently
\begin{equation}\protect\phantomsection\label{eq:sh-optics-d-2}{
 \frac{124}{125}\le q^2=E^2-\ell^2a(\rho)\le\frac{109}{27},
 \qquad \frac{99}{100}<q<\frac{21}{10}.
}\end{equation} There is a unique endpoint \(T\) at \(r(T)=3m/2\). The
dimensionless length \(\tau=T/m\) satisfies
\begin{equation}\protect\phantomsection\label{eq:sh-optics-d-3}{
 \frac{10}{21}<\tau<\frac{100}{99}.
}\end{equation} The horizon is met exactly once, transversely, with
\(\dot r=-E\le-1\). More strongly, \(-\dot r>99/100\) throughout. The
ray is nonradial everywhere: \(L\ge m/4\).

The complete coordinate solution is specified by monotone quadratures:
\begin{equation}\protect\phantomsection\label{eq:sh-optics-d-4}{\begin{gathered}
t=\int_{\rho(t)}^{5/2}\frac{ds}{q(s)}, \\
\phi(t)=\ell\int_{\rho(t)}^{5/2}\frac{ds}{s^2q(s)}, \\
V(t)=\ell^2\int_{\rho(t)}^{5/2}\frac{ds}{s^2q(s)(E+q(s))}.
\end{gathered}}\end{equation} They define real-analytic functions of the parameters
and the affine variable on a common neighborhood of the compact family:
the radicand has the uniform positive lower bound in
eq.~\ref{eq:sh-optics-d-2}, and \(dt/d\rho=-1/q\) is a unit. No
elliptic-integral inversion convention is needed.

The coordinate derivatives obey the explicit bounds \[
 \frac1{410}\le\frac{dV}{dt}=\dot v\le\frac1{16},\qquad
 \frac1{25}\le\frac{d\phi}{dt}=\frac\ell{\rho^2}\le\frac29.
\] Indeed the lower bound for \(\dot v\) follows from \(\ell^2\ge1/16\),
\(\rho^2\le25/4\), and \(E+q\le41/10\); its upper bound follows from
\(E+q\ge199/100\). Hence
\begin{equation}\protect\phantomsection\label{eq:sh-optics-d-5}{
 0\le V(t)<\frac{25}{396}<\frac1{15},\qquad
 0\le\phi(t)<\frac{200}{891}<\frac14.
}\end{equation} In particular the whole family lies in one regular
angular and advanced-time chart inside the parent domain
\begin{equation}\protect\phantomsection\label{eq:sh-optics-d-6}{
 \frac43<\rho<\frac83,\quad |V|<\frac14,\quad
 |\theta-\pi/2|<\frac18,\quad -\frac18<\phi<\frac38.
}\end{equation} The radial buffer is at least \(1/6\), the angular
\(\phi\) buffer is at least \(1/8\), and the advanced-time buffer is
greater than \(1/4-1/15=11/60\), in the displayed dimensionless
coordinates. The equator has a \(1/8\) latitude buffer. The metric
determinant is \(-r^4\sin^2\theta\), so its nondegeneracy on the closure
of a slightly smaller parent is explicit, including the horizon. These
are finite coordinate buffers, not statements about the global null-ray
quotient.

\protect\phantomsection\label{sec:sh-optics-3}

\paragraph{K.3.3. A genuinely parallel Sachs
basis}\label{k.3.3.-a-genuinely-parallel-sachs-basis}

Along each ray define the regular null frame \[
 l=\partial_v+\frac f2\partial_r,\quad n=-\partial_r,
 \quad e_\theta=\frac1r\partial_\theta,\quad
 e_\phi=\frac1r\partial_\phi.
\] On the equator, \(h(l,n)=-1\) and \(e_\theta,e_\phi\) are unit
spacelike. Write \[
 A=\dot v>0,\qquad c=\frac Lr,\qquad
 b=q+\frac{fA}{2},\qquad
 k=A l+b n+c e_\phi,\qquad 2Ab=c^2.
\] Set \[
 z=\frac cA=\frac{r(E+q)}L,\qquad
 H(u)=\frac{Eu+r(0)-r(u)}L,
\] \begin{equation}\protect\phantomsection\label{eq:sh-optics-d-7}{
 e_1=e_\theta,\qquad e_2=e_\phi+z n+Hk.
}\end{equation} These satisfy \(h(e_i,e_j)=\delta_{ij}\) and
\(h(e_i,k)=0\). In particular \(e_2\) is not obtained by deleting the
angular momentum of the ray. The nonprincipal character of the ray
remains in both \(k\) and \(e_2\).

Here is a direct parallelism calculation. The EF connection on the
equator gives \[
 \nabla_k e_\theta=0,\quad
 \nabla_k e_\phi=-\frac crl+\frac{fc}{2r}n,\quad
 \nabla_k n=-\frac{Af'}2 n-\frac cr e_\phi.
\] Also \(\dot A=-f'A^2/2+c^2/r\), \(\dot c=qc/r\), and thus \[
 \dot z=-\frac{Ez}{r}+\frac{cf'}2,\qquad
 \nabla_k(e_\phi+z n)=-\frac zr k=-\frac{E+q}L k,
 \qquad \dot H=\frac{E+q}L.
\] Equation eq.~\ref{eq:sh-optics-d-7} is therefore an exactly parallel
orthonormal Sachs basis, not merely a choice parallel modulo \(k\). It
is regular at the horizon. Uniform coefficient bounds useful for
comparison are
\begin{equation}\protect\phantomsection\label{eq:sh-optics-d-8}{
 0<z\le41,\qquad 0\le H<\frac{1196}{99}<13.
}\end{equation} The second bound uses \(u\le100m/99\) and
\(r(0)-r\le m\). These are coefficient bounds in the regular frame
above, not Lorentzian positive norms of arbitrary spacetime vectors.

\protect\phantomsection\label{sec:sh-optics-4}

\paragraph{K.3.4. Actual curvature and the two nonzero tidal
eigenvalues}\label{k.3.4.-actual-curvature-and-the-two-nonzero-tidal-eigenvalues}

For completeness the nonzero connection coefficients, besides symmetry
in the lower slots, are
\begin{equation}\protect\phantomsection\label{eq:sh-optics-d-9}{
 \begin{gathered}
 \Gamma^v_{vv}=f'/2,\quad \Gamma^r_{vv}=ff'/2,
 \quad \Gamma^r_{vr}=-f'/2,\quad
 \Gamma^v_{\theta\theta}=-r,\quad
 \Gamma^v_{\phi\phi}=-r\sin^2\theta,\\
 \Gamma^r_{\theta\theta}=-fr,\quad
 \Gamma^r_{\phi\phi}=-fr\sin^2\theta,\quad
 \Gamma^\theta_{r\theta}=\Gamma^\phi_{r\phi}=1/r,\\
 \Gamma^\theta_{\phi\phi}=-\sin\theta\cos\theta,
 \qquad\Gamma^\phi_{\theta\phi}=\cot\theta.
 \end{gathered}
}\end{equation} Substitution in the stated curvature convention yields
\[
 \begin{gathered}
 \operatorname{Ric}_{vv}=f(f''/2+f'/r),\qquad
 \operatorname{Ric}_{vr}=-(f''/2+f'/r),\qquad
 \operatorname{Ric}_{rr}=0,\\
 \operatorname{Ric}_{\theta\theta}=1-f-rf',\qquad
 \operatorname{Ric}_{\phi\phi}=\sin^2\theta(1-f-rf').
 \end{gathered}
\] All other components vanish; \(f=1-2m/r\) makes this tensor zero.

The out-of-plane tidal entry can be calculated without taking a Ricci
trace. For the two-dimensional base metric \(-f\,dv^2+2\,dv\,dr\), the
Hessian of \(r\) has entries \(-ff'/2,f'/2,0\) in its \(vv,vr,rr\)
slots. If \(k_B=A\partial_v-q\partial_r\), the null identity
\(fA^2+2Aq=c^2\) gives \[
 (\nabla_B^2r)(k_B,k_B)=-\frac{f'c^2}{2}.
\] The mixed and spherical curvature terms from
eq.~\ref{eq:sh-optics-d-9} are respectively
\(-(\nabla_B^2r)(k_B,k_B)/r\) and \(c^2(1-f)/r^2\). Their sum is
\begin{equation}\protect\phantomsection\label{eq:sh-optics-d-10}{
 h(e_\theta,R(e_\theta,k)k)
 =c^2\left(\frac{f'}{2r}+\frac{1-f}{r^2}\right)
 =\frac{3mL^2}{r^5}=:d_u>0.
}\end{equation} The equatorial reflection fixes \(k\) and \(e_2\) and
reverses \(e_1\), so the off-diagonal tidal entries vanish. The screen
trace is \(\operatorname{Ric}(k,k)=0\); the two missing null-frame
contributions vanish by curvature antisymmetry. It follows that
\begin{equation}\protect\phantomsection\label{eq:sh-optics-d-11}{
 \mathcal R=\begin{pmatrix}d_u&0\\0&-d_u\end{pmatrix},
 \qquad x_1''=-d_u x_1,\quad x_2''=d_u x_2.
}\end{equation} Adding \(Hk\) in eq.~\ref{eq:sh-optics-d-7} changes
neither screen curvature entry. The ordering is fixed: the out-of-plane
direction focuses, while the in-plane direction has the opposite Weyl
tide.

There is a separate sign check using a geometric Jacobi field. Rotation
of the orbital plane gives an out-of-plane amplitude \(y=r\sin\phi\);
this follows by differentiating a spatial rotational isometry and
pairing with \(e_\theta\), up to an immaterial constant sign. The
geodesic first integrals give \[
 \ddot r=\frac{L^2}{r^3}-\frac{3mL^2}{r^4},\qquad
 \dot\phi=\frac L{r^2},\qquad
 2\dot r\dot\phi+r\ddot\phi=0.
\] Therefore \(y''=-(3mL^2/r^5)y\). Using the second rotational field
gives \(r\cos\phi\) and removes any issue at zeros of \(\sin\phi\). This
independently checks the sign and the nonzero coefficient of
eq.~\ref{eq:sh-optics-d-10}. Ricci-flatness alone would not give
eq.~\ref{eq:sh-optics-d-10}.

\protect\phantomsection\label{sec:sh-optics-5}

\paragraph{K.3.5. Full phase transport and an exact
representation}\label{k.3.5.-full-phase-transport-and-an-exact-representation}

Use \(t=u/m\) and the dimensionally consistent phase vector \[
 \begin{gathered}
Z=\binom{x}{p}, \\
p=m\frac{dx}{du}=\frac{dx}{dt}, \\
\omega(t)=\frac{3\ell^2}{\rho(t)^5}, \\
K(t)=\operatorname{diag}(-\omega(t),+\omega(t)).
\end{gathered}
\] Then
\begin{equation}\protect\phantomsection\label{eq:sh-optics-d-12}{
 \frac{dZ}{dt}=\mathsf A(t)Z,\qquad
 \mathsf A(t)=\begin{pmatrix}0&I_2\\K(t)&0\end{pmatrix},
 \qquad\Phi(s,s)=I_4.
}\end{equation} All matrix norms below are Euclidean operator norms in
this phase basis. Direct evaluation gives
\begin{equation}\protect\phantomsection\label{eq:sh-optics-d-13}{
 \frac6{3125}\le\omega\le\frac8{81},\qquad
 \|\mathsf A\|=1,\qquad
 |\omega'|=\frac{15\ell^2q}{\rho^6}\le\frac{56}{81}.
}\end{equation} The exact fundamental matrix is the uniformly convergent
Peano--Baker series
\begin{equation}\protect\phantomsection\label{eq:sh-optics-d-14}{
 \Phi(t,s)=I_4+
 \sum_{n=1}^\infty
 \int_{s<t_n<\cdots<t_1<t}
 \mathsf A(t_1)\cdots\mathsf A(t_n)\,dt_n\cdots dt_1.
}\end{equation} This, together with the explicit quadratures
eq.~\ref{eq:sh-optics-d-4}, is a specification of every entry of the
full transport. It does not replace the matrix by its trace. With
\(h=t-s\le100/99\) and the series truncated after degree \(N\),
\begin{equation}\protect\phantomsection\label{eq:sh-optics-d-15}{
 \|\Phi(t,s)\|\le e^h<3,\qquad
 \|\Phi(t,s)-\Phi_N(t,s)\|
 \le e^h\frac{h^{N+1}}{(N+1)!}
 \le3\frac{(100/99)^{N+1}}{(N+1)!}.
}\end{equation} For example the \(N=8\) upper bound is less than
\(1/40320\), using \((100/99)^9<(11/10)^9<3\). This is an analytic
finite error certificate, not a numerical integration result. The loose
elementary bound \(e^{100/99}<3\) follows from \(e<11/4\) and
\(e^{1/99}\le99/98\).

More explicitly let \(c_\pm,s_\pm\) solve \(y''=\pm\omega y\) with
initial data \((1,0)\) and \((0,1)\) at \(t=0\). In the ordered basis
\((x_1,x_2,p_1,p_2)\) the exact matrix is
\begin{equation}\protect\phantomsection\label{eq:sh-optics-d-16}{
 \Phi(t,0)=
 \begin{pmatrix}
 c_-&0&s_-&0\\
 0&c_+&0&s_+\\
 c_-'&0&s_-'&0\\
 0&c_+'&0&s_+'
 \end{pmatrix}.
}\end{equation} Two entries are available directly from the rotational
Jacobi fields. With \(\rho_0=5/2\) and \(q_0=q(\rho_0)\),
\begin{equation}\protect\phantomsection\label{eq:sh-optics-d-17}{
 s_-(t)=\frac{\rho_0}{\ell}\rho(t)\sin\phi(t),\qquad
 c_-(t)=\frac{\rho(t)}{\rho_0}\cos\phi(t)
                 +\frac{q_0}{\rho_0}s_-(t).
}\end{equation} The lower bound on \(\ell\) makes these formulas regular
for the whole family. The other two solutions are specified by their
convergent Volterra equations or eq.~\ref{eq:sh-optics-d-14}; no
elementary closed form is asserted.

There is also a finite approximation preserving the exact symplectic
structure. For a partition \(0=t_0<\cdots<t_n=\tau\) of mesh \(\delta\),
use the exact matrix exponential of the frozen generator
\(\mathsf A(t_j)\) on each subinterval and compose them in time order.
If \(h_j=t_{j+1}-t_j\) and \(w=\omega(t_j)\), that factor has blocks \[
 \begin{pmatrix}C&S\\K(t_j)S&C\end{pmatrix},\quad
 C=\operatorname{diag}(\cos(\sqrt w h_j),\cosh(\sqrt w h_j)),
\] \[
 S=\operatorname{diag}\left(
 \frac{\sin(\sqrt w h_j)}{\sqrt w},
 \frac{\sinh(\sqrt w h_j)}{\sqrt w}\right).
\] By eq.~\ref{eq:sh-optics-d-13} the piecewise coefficient error is at
most \((56/81)\delta\). Duhamel's formula gives the global bound
\begin{equation}\protect\phantomsection\label{eq:sh-optics-d-18}{
 \|\Phi_{\mathrm{frozen}}(\tau,0)-\Phi(\tau,0)\|
 \le\frac{5600}{2673}\delta<3\delta.
}\end{equation} For a uniform partition it is in particular less than
\(3/n\). Exact values of the sampled coefficients in this construction
are defined by eq.~\ref{eq:sh-optics-d-4}. Computation or certified
enclosure of those samples would be an additional numerical step; none
has been performed here. Further sample errors are covered by the
coefficient-stability estimate below. An unmodified finite Peano--Baker
polynomial need not itself be symplectic, whereas these exact frozen
factors are symplectic.

\protect\phantomsection\label{sec:sh-optics-6}

\paragraph{K.3.6. Symplecticity, finite composition, and stationary
generating
families}\label{k.3.6.-symplecticity-finite-composition-and-stationary-generating-families}

Let \[
 \mathsf J=\begin{pmatrix}0&I_2\\-I_2&0\end{pmatrix},\qquad
 \Omega(Z,\widehat Z)=Z^{\mathsf T}\mathsf J\widehat Z.
\] Since \(K=K^{\mathsf T}\),
\(\mathsf A^{\mathsf T}\mathsf J+\mathsf J\mathsf A=0\). Differentiating
\(\Phi^{\mathsf T}\mathsf J\Phi\) proves
\begin{equation}\protect\phantomsection\label{eq:sh-optics-d-19}{
 \Phi(t,s)^{\mathsf T}\mathsf J\Phi(t,s)=\mathsf J,
 \qquad \Phi(t,s)^{-1}=\Phi(s,t),\qquad
 \Phi(t,s)\Phi(s,r)=\Phi(t,r).
}\end{equation} The composition identity follows from the initial-value
theorem and holds for every finite partition. The Hamiltonian is \[
 H_t(x,p)=\frac12|p|^2+\frac12\omega(t)(x_1^2-x_2^2).
\] Its potential is indefinite; no positivity, minimizing principle, or
positive-definite closed-system theorem is used. With the metric
\(g=-h\), the Jacobi Wronskian of Proposition 5.3 is \(-m^{-1}\Omega\)
in this basis. Multiplication by the constant \(-m^{-1}\) changes
neither symplectic maps nor Lagrangian subspaces.

Write a short-segment matrix as
\(U_j=\begin{pmatrix}a_j&b_j\\c_j&d_j\end{pmatrix}\). When \(b_j\) is
invertible, symplecticity makes \(b_j^{-1}a_j\) and \(d_jb_j^{-1}\)
symmetric and gives \(c_j-d_jb_j^{-1}a_j=-b_j^{-\mathsf T}\). Thus
\begin{equation}\protect\phantomsection\label{eq:sh-optics-d-20}{
 S_j(x_j,x_{j+1})=
 \frac12 x_j^{\mathsf T}b_j^{-1}a_jx_j
 -x_j^{\mathsf T}b_j^{-1}x_{j+1}
 +\frac12x_{j+1}^{\mathsf T}d_jb_j^{-1}x_{j+1}
}\end{equation} generates the exact graph by \(p_j=-\partial_{x_j}S_j\)
and \(p_{j+1}=\partial_{x_{j+1}}S_j\).

For a finite chain retain all intermediate positions and put
\(S=\sum_j S_j\). Stationarity in \(x_j\) for \(0<j<n\) is exactly
equality of the outgoing and incoming momenta at that vertex. The
critical relation with endpoint momenta defined by the endpoint
derivatives is therefore the graph of \(U_{n-1}\cdots U_0\). This proof
does not require a nondegenerate Hessian in the intermediate positions,
nor does it eliminate those positions.

More generally, for any continuous bounded symmetric tidal matrix on a
finite interval, sufficiently short segments have invertible \(b_j\).
Indeed, if \(\|K\|\le M\) and \(a=\max(1,M)\), the initial solution with
position zero and momentum \(p_0\) satisfies
\(|x(s)|\le s e^{as}|p_0|\). The Volterra equation then gives
\(\|b_j-h_jI\|\le M e^{ah_j}h_j^3/6\). A finite partition with
\(M e^{ah_j}h_j^2/6<1\) therefore suffices. This is a direct finite
Hamiltonian argument.

Consequently the graph of the full transport is always a Lagrangian
subspace of \((\mathbb R^4,-\Omega)\times(\mathbb R^4,\Omega)\) and
remains of dimension four when its endpoint-position projection becomes
singular. A propagated initial Lagrangian is likewise still a Lagrangian
subspace at a caustic of that projection. The retained finite generating
family describes that phase relation. A singular projection need not
determine a unique phase point from its projected endpoints, and no
inverse projection or inverse-Hessian estimate follows there.

\protect\phantomsection\label{sec:sh-optics-7}

\paragraph{K.3.7. An evaluated endpoint-block margin for this
family}\label{k.3.7.-an-evaluated-endpoint-block-margin-for-this-family}

For a segment of dimensionless length \(0<h\le100/99\), the
initial-position-zero block \(b(h)\) obeys \[
 \|b(h)-hI_2\|
 \le\frac8{81}e^h\frac{h^3}{6}
 \le\frac4{81}\left(\frac{100}{99}\right)^2 h
 <\frac1{10}h.
\] Hence
\begin{equation}\protect\phantomsection\label{eq:sh-optics-d-21}{
 \sigma_{\min}(b(h))\ge\frac9{10}h,\qquad
 \|b(h)\|\le h e^h.
}\end{equation} At the full endpoint, eq.~\ref{eq:sh-optics-d-3} gives
\begin{equation}\protect\phantomsection\label{eq:sh-optics-d-22}{
 \sigma_{\min}(b(\tau))>\frac37,\qquad
 \|b(\tau)^{-1}\|<\frac73,\qquad
 \|b(\tau)\|<\frac{100}{33},\qquad
 \det b(\tau)>\frac9{49}.
}\end{equation} Here \(b=\operatorname{diag}(s_-,s_+)\). Its entries are
positive: eq.~\ref{eq:sh-optics-d-21} prevents a zero, and their initial
derivatives are positive. The Volterra equations then give \(s_-<h<s_+\)
at every positive segment length, since \(\omega>0\). Thus the optical
effect is not the scalar free propagation of \(\mathcal E_J\).

There is no point-source conjugate point within this particular short
family. The finite stationary-composition argument proves what remains
meaningful if a projection of a general initial Lagrangian becomes
singular; it does not assert that the point-source family constructed
here passes through a caustic. Nor is eq.~\ref{eq:sh-optics-d-22}
uniform as a segment length tends to zero. For initial physical
transverse rate \(\dot x(0)\), the physical endpoint map is \(m b\) and
its inverse norm is at most \(7/(3m)\).

A concrete singular Lagrangian projection can nevertheless be exhibited
on these same intervals. Choose \(t_c=\tau/2\), write
\(\Phi(t_c,0)=\begin{pmatrix}a_c&b_c\\c_c&d_c\end{pmatrix}\), and impose
the initial screen condition \[
 \Lambda_0=\{(x,P_0x):x\in\mathbb R^2\},\qquad
 P_0=-b_c^{-1}a_c=P_0^{\mathsf T}.
\] Equation eq.~\ref{eq:sh-optics-d-21} gives
\(\sigma_{\min}(b_c)>3/14\), so \(\|P_0\|<14\). At \(t_c\) the position
is identically \((a_c+b_cP_0)x=0\), while invertibility and
symplecticity of \(\Phi\) imply that the propagated two-plane is exactly
the vertical Lagrangian \(\{(0,p)\}\). The projection is therefore
singular there, although the full transport remains invertible. A
retained generating family for this propagated condition is obtained by
adding \(\tfrac12x_0^{\mathsf T}P_0x_0\) to \(\sum_jS_j\) and imposing
stationarity also in \(x_0\). It remains an exact phase relation at
\(t_c\); solving it as a single graph over the final position is
unavailable there. This example is a caustic of the specified linear
Lagrangian screen condition. No additional point-source conjugate point
or nonlinear congruence construction is claimed.

\protect\phantomsection\label{sec:sh-optics-8}

\paragraph{K.3.8. Finite coefficient and observation
error}\label{k.3.8.-finite-coefficient-and-observation-error}

This section is the ODE part of the metric-to-optics error bridge. It
assumes that the comparison has already placed the exact and approximate
systems on the same affine interval and in identified phase bases. The
derivation of that identification and its coefficient bound from finite
metric records is a separate input.

Suppose \[
 \sup_{0\le t\le\tau}\|\widetilde{\mathsf A}(t)-\mathsf A(t)\|
 \le\varepsilon,\qquad 0\le\varepsilon\le\frac1{10}.
\] The approximate coefficient need not be Hamiltonian for this
stability assertion. Both initial fundamental matrices are the identity.
Duhamel's formula gives
\begin{equation}\protect\phantomsection\label{eq:sh-optics-d-23}{
 \|\widetilde\Phi(\tau,0)-\Phi(\tau,0)\|
 \le\varepsilon\tau e^{(1+\varepsilon)\tau}\le4\varepsilon,
 \qquad \|\widetilde\Phi\|<4.
}\end{equation} For an elementary evaluated bound,
\((1+\varepsilon)\tau\le10/9\) and \(e^{10/9}<(11/4)(9/8)=99/32<4\); the
first bound is actually at most \((25/8)\varepsilon\). An initial phase
error \(\eta_0\) therefore contributes at most \(4\eta_0\), and the
coefficient error contributes at most \(4\varepsilon\|Z_0\|\) to
endpoint phase error.

If an exact-system Peano--Baker truncation error \(r_N\) is also
retained, a common total transfer tolerance is
\(\delta=4\varepsilon+r_N\). If truncating the approximate system
instead, its own series bound with
\(\|\widetilde{\mathsf A}\|\le1+\varepsilon\) must be used for \(r_N\).
The exact transfer blocks and any surrogate then satisfy
\(\|\widetilde b-b\|\le\delta\).

One finite optical observation is the oriented transverse area factor \[
 \mathcal O=m^2\det b(\tau),
\] with initial angular/rate normalization fixed as above. For
\(\delta\le1\) the \(2\times2\) determinant identity and
eq.~\ref{eq:sh-optics-d-22} give
\begin{equation}\protect\phantomsection\label{eq:sh-optics-d-24}{
 |\widetilde{\mathcal O}-\mathcal O|
 \le m^2(2\|b\|\delta+\delta^2)\le9m^2\delta.
}\end{equation} This is a forward observation estimate; it requires no
inverse projection. If in addition \(\delta\le3/14\), the evaluated
endpoint margin survives:
\begin{equation}\protect\phantomsection\label{eq:sh-optics-d-25}{
 \sigma_{\min}(\widetilde b)\ge\frac3{14},\qquad
 \|\widetilde b^{-1}\|\le\frac{14}{3},\qquad
 \|\widetilde b^{-1}-b^{-1}\|\le\frac{98}{9}\delta.
}\end{equation} The last estimate is the matrix inverse identity. For
the inverse of the physical map \(m b\), multiply its right side by
\(m^{-1}\). At a caustic the hypothesis \(\sigma_{\min}(b)>0\) is false,
so eq.~\ref{eq:sh-optics-d-25} is unavailable;
eq.~\ref{eq:sh-optics-d-23}, finite symplectic composition, and the full
phase relation still make sense wherever the coefficients remain bounded
on a finite interval.

\protect\phantomsection\label{sec:sh-integrated-error}

\subsubsection{K.4. Metric truncation and finite optical
errors}\label{k.4.-metric-truncation-and-finite-optical-errors}

Fix one of the finite coordinate representatives constructed in the
finite horizon construction and an inner real tube containing the chosen
ray segments. Coordinate and tensor bounds below use fixed auxiliary
Euclidean operator norms, including multilinear derivative norms; they
are not norms of an indefinite metric. Write \(g\) for the Schwarzschild
metric with the signature of §2 and use the same coordinates for every
approximation. All constants refer to this one fixed positive mass and
this finite tube.

\pandocstatementheading{\protect\phantomsection\label{lem:sh-holomorphic-metric-tail}Lemma K.9.}

Suppose the prescribed stored representative of Proposition 6.2 is
defined on complex spatial parents with positive buffer, holomorphically
for \(|\varepsilon|<T_b\), with \(T_b>1\), and that the finite metric
reconstruction and its inverse units stay in the fixed margins of that
proposition. Choose \(1<T<T_b\) and an inner spatial polydisc whose
coordinatewise gap from the retained parent is \(\rho>0\). There is a
finite \(H_g\) such that, for the degree-\(N\) amplitude Taylor
polynomial \(g_N\) of the reconstructed physical metric in the fixed
coordinates,
\begin{equation}\protect\phantomsection\label{eq:sh-metric-tail-cauchy}{\begin{gathered}
\|g-g_N\|\le\delta_N:=\frac{H_gT^{-N-1}}{1-T^{-1}}, \\
\|D^j(g-g_N)\|\le e_j:=4^j j!\rho^{-j}\delta_N
 \quad(0\le j\le3).
\end{gathered}}\end{equation} The metric coefficients are real on the real tube when
the source and retained frames obey the Lorentz reality of §8.4. The
same conclusion and bound hold on a fixed finite cover after taking the
largest \(H_g\) and smallest gap.

\protect\phantomsection\label{sec:sh-r1-finite-readout-proof}

\emph{Proof.} Fix the presentation choices and a bounded holomorphic
original family of the construction of §§5--6. Write \(Y(\varepsilon)\)
for its prescribed augmented stored tuple, and let \(q_N\) denote
reduction modulo \(\varepsilon^{N+1}\). Theorem 6.3 identifies the
actual finite output with \(X_N=q_NY\). In this proof all finite readout
operations are performed over
\(A_N=\mathbb C[\varepsilon]/(\varepsilon^{N+1})\). Evaluation at one is
applied only to the resulting coefficient polynomial.

Choose \(T<T_c<T_b\) and keep the compact amplitude disk
\(|\varepsilon|\le T_c\). Retain a smaller point-coordinate parent with
a positive complex Cauchy gap inside the stored graph parent. At each of
the fixed nine buffered directions put \(\ell_\nu=\partial_xf(q_\nu)\).
Proposition 5.2 reconstructs the symmetric inverse conformal form by the
ten-by-ten system
\begin{equation}\protect\phantomsection\label{eq:sh-r1-nine-direction-readout}{
 A(x,\varepsilon)G(x,\varepsilon)=b,\qquad
 G^{ab}(\ell_\nu)_a(\ell_\nu)_b=0\quad(1\le\nu\le9),\qquad
 \eta_{ab}G^{ab}=4.
}\end{equation} The system uses the ten independent entries of \(G\).
The fixed central sky jets and four-evaluation normalization give the
same invertible \(A(0,\varepsilon)=A_0\) for every amplitude. Put
\(a_0=\|A_0^{-1}\|\). Spatial Cauchy estimates on the first smaller
parent control the evaluated derivatives. A further uniform restriction,
keeping the entire compact amplitude disk, gives
\begin{equation}\protect\phantomsection\label{eq:sh-r1-readout-inverse-margin}{
 \|A-A_0\|\le\frac1{2a_0},\qquad
 \|A^{-1}\|\le2a_0.
}\end{equation} Indeed \(A=A_0(I+A_0^{-1}(A-A_0))\) and the Neumann
series has ratio at most \(1/2\). At the mark \(G=\eta^{-1}\). A further
common restriction bounds \(G^{-1}\). The retained selected density has
\(\sigma(0,\varepsilon)=1\), so we can also require \(|\sigma|\ge1/2\).
Consequently
\begin{equation}\protect\phantomsection\label{eq:sh-r1-physical-metric-readout}{
 \widetilde g_\varepsilon=\sigma_\varepsilon^{-2}G_\varepsilon^{-1}
}\end{equation} is jointly holomorphic, uniformly bounded and uniformly
invertible on a common point-coordinate parent. These restrictions exist
by compactness of the amplitude disk and the fixed central units, rather
than by smallness of the Schwarzschild records. For a fixed finite
passive family take the largest inverse and derivative bounds and the
smallest positive radii.

We spell out the uniform normal development used in this readout. Choose
a further positive complex spatial gap, and a convex complex ball
\(B_R\) on which \[
 \|\Gamma_\varepsilon\|\le K_\Gamma,\qquad
 \|D\Gamma_\varepsilon\|\le L_\Gamma
 \quad(|\varepsilon|\le T_c).
\] These are finite bounds from the metric, its inverse and interior
Cauchy estimates. For initial velocity \(y\), use the complex time disk
\(|s|\le1\) and write
\begin{equation}\protect\phantomsection\label{eq:sh-r1-normal-picard-map}{
 q_v(s)=\int_0^s v(u)\,du,\qquad
 (\mathcal T_\varepsilon v)(s)
 =y-\int_0^s\Gamma_\varepsilon(q_v(u))(v(u),v(u))\,du.
}\end{equation} The integrals are holomorphic primitives. Choose \(a>0\)
with
\begin{equation}\protect\phantomsection\label{eq:sh-r1-picard-margins}{
 a\le R/4,\qquad
 4K_\Gamma a\le1/4,\qquad
 4L_\Gamma a^2\le1/4,
}\end{equation} omitting a constraint whose coefficient is zero. If
\(\|y\|\le a\) and \(\|v-y\|_\infty\le a\), then \(\|v\|_\infty\le2a\)
and \(\|q_v\|_\infty\le2a<R\). The Picard map stays in this ball, since
\(\|\mathcal T_\varepsilon v-y\|_\infty\le4K_\Gamma a^2\le a/4\). For
two velocities in it,
\begin{equation}\protect\phantomsection\label{eq:sh-r1-picard-contraction}{
 \|\mathcal T_\varepsilon v-\mathcal T_\varepsilon w\|_\infty
 \le(4L_\Gamma a^2+4K_\Gamma a)\|v-w\|_\infty
 \le\tfrac12\|v-w\|_\infty.
}\end{equation} Thus the solution is unique and jointly holomorphic in
\(y,\varepsilon\), by uniformly convergent Picard iteration. Let
\(E_\varepsilon(y)=q(1)\). Its integral equation gives
\begin{equation}\protect\phantomsection\label{eq:sh-r1-normal-endpoint-margin}{\begin{gathered}
\|E_\varepsilon(y)-y\|\le2K_\Gamma a^2\quad(\|y\|\le a), \\
\|DE_\varepsilon-I\|\le4K_\Gamma a\le1/4
 \quad(\|y\|\le a/2).
\end{gathered}}\end{equation} The derivative estimate follows by vector-valued
directional Cauchy estimates, with radius \(a/2\) in each unit complex
direction. Hence \(\|DE_\varepsilon\|\le5/4\). Writing
\(E_\varepsilon=I+e_\varepsilon\), the contraction
\(y\mapsto x-e_\varepsilon(y)\) preserves \(\|y\|\le a/2\) for
\(\|x\|\le a/4\), because \(e_\varepsilon(0)=0\) and
\(\|De_\varepsilon\|\le1/4\). Its fixed point is the actual initialized
endpoint inverse, jointly holomorphic on that common target ball, with
derivative norm at most \(4/3\).

The normal-coordinate physical metric is therefore
\begin{equation}\protect\phantomsection\label{eq:sh-r1-normal-metric-readout}{
 g^{\mathrm{nor}}_\varepsilon(y)
 =DE_\varepsilon(y)^{\mathsf T}
   \widetilde g_\varepsilon(E_\varepsilon(y))DE_\varepsilon(y).
}\end{equation} On a common smaller normal parent \(P_1\), a bound
\(B_{\widetilde g}\) for the physical point-coordinate readout gives the
possible bound \(B_g=(25/16)B_{\widetilde g}\) for this normal metric.
This \(B_g\) is a metric-readout bound. It is not the original
stored-tuple bound \(M_C\); its dependencies include the first spatial
gaps, evaluation and density inverse margins, and the initialized
endpoint bounds. The normal initial equations give
\(g^{\mathrm{nor}}_\varepsilon(0)=\eta\) and
\(\partial g^{\mathrm{nor}}_\varepsilon(0)=0\) at every amplitude. The
fixed mark and tangent/full-spin identification identify it with the
marked normal geometry in Proposition 6.1.

We next verify compatibility with the finite computation. Let
\(A^{(0)}(x)\) be the amplitude-zero matrix on the common parent. With
\(\Delta_N=(A^{(0)})^{-1}(q_NA-A^{(0)})\), the matrix \(\Delta_N\)
belongs to the nilpotent amplitude ideal. The inverse of \(q_NA\) is the
finite expression
\begin{equation}\protect\phantomsection\label{eq:sh-r1-finite-inverse-recursion}{
 (q_NA)^{-1}
 =\sum_{j=0}^{N}(-\Delta_N)^j(A^{(0)})^{-1}
 =q_N(A^{-1}).
}\end{equation} The same argument applies to the density and every other
retained unit. Spatial derivatives, finite products and sums commute
with \(q_N\). Composition commutes as well: expand around the
amplitude-zero argument, and a coefficient through degree \(N\) uses
only finitely many derivatives and nilpotent Taylor terms through degree
\(N\).

For the initialized ODE, its amplitude-zero solution is the solution of
the residue equation. At each degree \(n\ge1\), the combined state
coefficient satisfies
\begin{equation}\protect\phantomsection\label{eq:sh-r1-finite-ode-recursion}{
 z_n'=D_z\mathcal V_0(z_0)z_n
          +F_n(s,z_0,\ldots,z_{n-1}),
}\end{equation} with the prescribed initial coefficient. This is a
uniquely initialized finite linear inhomogeneous ODE once the lower
coefficients are fixed. Induction identifies the finite-algebra ODE with
the reduction of the jointly holomorphic solution. Differentiation gives
the same conclusion for \(DE_\varepsilon\); the initialized endpoint
inverse is likewise compatible, either by coefficient recursion or its
unique inverse identity. It follows that the complete specified readout,
including the pullback in the preceding display, satisfies
\begin{equation}\protect\phantomsection\label{eq:sh-r1-finite-readout-identity}{
 \mathcal R_N(X_N)=q_N g^{\mathrm{nor}}_\varepsilon.
}\end{equation} This is an equality of analytic coefficient germs over
the finite algebra for every \(N\). If an original finite-stage
coefficient was constructed on a smaller parent, its germ has the common
representative given by the corresponding coefficient of this single
holomorphic metric. No estimate for norms of growing quotient algebras
is used. The common parents and the bounds were chosen from the
holomorphic family before \(N\) was fixed.

Now use the exact initialized coordinate isometry at amplitude one to
express the metric in the fixed physical chart of the lemma. Write
\(\kappa\) for its inverse, mapping that chart to the normal parent, and
define
\(g^{\mathrm{fix}}_\varepsilon=\kappa^*g^{\mathrm{nor}}_\varepsilon\).
The same \(\kappa\) is used for every amplitude and every \(N\); it is
not an independently recomputed approximate coordinate map. On a smaller
fixed complex parent whose closure maps into \(P_1\), this family has a
finite holomorphic bound \(H_g\), for example \(\sup\|D\kappa\|^2B_g\).
The fixed-coordinate pullback commutes with reduction. Only after these
operations set
\begin{equation}\protect\phantomsection\label{eq:sh-r1-final-polynomial-evaluation}{
 g_N=\operatorname{ev}^{\mathrm{poly}}_1
          (q_Ng^{\mathrm{fix}}_\varepsilon),\qquad
 g=g^{\mathrm{fix}}_1.
}\end{equation} Here \(\operatorname{ev}^{\mathrm{poly}}_1\) means
evaluation of the unique coefficient polynomial of degree at most \(N\),
not a homomorphism \(A_N\to\mathbb C\). In particular this is not the
nonlinear readout \(\mathcal R(X_N(1))\).

Cauchy's formula on \(|\varepsilon|=T\) gives the coefficient bound
\(H_gT^{-n}\). Summing the tail at one gives exactly \(\delta_N\) in the
statement of Lemma K.9. Keep the second, independent spatial gap
\(\rho\) between this fixed parent and the inner polydisc. For each
multi-index \(\alpha\) with \(|\alpha|=j\le3\), spatial Cauchy gives
\begin{equation}\protect\phantomsection\label{eq:sh-r1-second-gap-cauchy}{
 \|\partial^\alpha(g-g_N)\|
 \le\alpha!\rho^{-j}\delta_N.
}\end{equation} Expanding a \(j\)-linear derivative into the four
coordinate directions bounds its Euclidean operator norm by the safe
factor \(4^j j!\rho^{-j}\delta_N\) stated in the lemma. The first loss
supplies bounded differentiated readout, inverses and initialized
developments; the second loss differentiates the already bounded
amplitude tail. Neither depends on \(N\).

For real amplitude the coupled Lorentz operation of §8.4 gives a real
physical metric, and initialized normal development in the fixed real
tangent frame is real. Holomorphy and uniqueness give
\(g^{\mathrm{fix}}_{\bar\varepsilon}(\bar x)=\overline{g^{\mathrm{fix}}_\varepsilon(x)}\).
All coefficient metrics and their polynomial sum are therefore real
symmetric. On a fixed finite cover the same proof applies chartwise,
taking the largest bound and smallest positive gap. This proves the
lemma. ∎

\protect\phantomsection\label{sec:sh-r1-ef-tube-and-initialization}

\textbf{Common EF tube and initialization.} The \(C^3\) estimate must
hold on the whole inflated position tube used in the ray/frame
first-exit argument. Choose its closure inside the inner
fixed-coordinate polydisc on which the stated derivative estimate holds,
with a positive additional position buffer; containment only in an
earlier readout parent is insufficient. If a normal-coordinate \(C^3\)
estimate is transported instead of applying spatial Cauchy directly in
the fixed chart as above, the coordinate-change constant depends on
derivatives of \(\kappa\) through order four, because three derivatives
can hit either Jacobian factor. Fixed powers of the mass from the
dimensionless phase/time dictionary are included explicitly in that
constant. A numerical approximation of \(\kappa\) would require a
separate certified \(C^4\) error.

For the earlier \(r:5m/2\to3m/2\), \(E\in[1,2]\), \(L/m\in[1/4,1/2]\)
family, apply Proposition 10.3 under its stated single-readout
position-tube inclusion hypothesis. Its initial clock, screen order and
exact optical constants are unchanged. The finite-cover conclusion of
Lemma K.9 is chartwise; independently evaluated vertex polynomials are
not asserted to glue as exact metrics, and using several readouts would
also require compatible physical comparisons and finite state matching.
Section K.8 constructs those finite state and Cauchy matchings on the
whole original crossing family. For each approximation use the same
specified reference initial data and affine-clock normalization rule. At
a noncentral entrance apply Lemma K.11: the actual null tangent is
\(k_N=k+an\) and the normalized initial frame differs from the reference
frame by at most \(C_{\rm init}e_0\), as retained in Theorem K.12. At
the common marked point, coefficientwise calibration gives
\(g_N(o)=g(o)\), so the tangent and screen can actually be chosen
identical. The initialized screen component phase is the same; no
unmarked affine rescaling or screen rotation is introduced. The
approximate metric need not conserve the exact Schwarzschild energy and
angular momentum. Bounds such as \(B_g,H_g\) and the resulting
coordinate/stability constants are finite theoretical bounds until their
inputs and parent radii are numerically evaluated; no numerical
sufficient \(N\) is asserted here.

\pandocstatementheading{\protect\phantomsection\label{lem:sh-inverse-curvature-stability}Lemma K.10.}

On the real tube suppose \(g\) is symmetric with signature \((+,-,-,-)\)
and \(\|g^{-1}\|\le G\), \(\|D^jg\|\le K_j\) for \(0\le j\le3\). If
\(g_N\) is real symmetric, \(e_0\le(2G)^{-1}\) and \(e_0,e_1,e_2\le1\),
then \(g_N\) has the same signature and
\begin{equation}\protect\phantomsection\label{eq:sh-inverse-margin}{
 \|g_N^{-1}\|\le2G,\qquad
 \|g_N^{-1}-g^{-1}\|\le2G^2e_0.
}\end{equation} Put \(\Gamma\) for the Christoffel bilinear map,
\(R(a,b)c\) for the curvature operator in the convention of §2, and
\begin{equation}\protect\phantomsection\label{eq:sh-connection-error-constants}{
 \begin{aligned}
 \gamma_0&=3G(K_1+1),\\
 \gamma_1&=6G^2(K_1+1)^2+3G(K_2+1),\\
 d_0&=6G^2(K_1+1)e_0+3Ge_1,\\
 d_1&=\{24G^3(K_1+1)^2+6G^2(K_2+1)\}e_0
       +12G^2(K_1+1)e_1+3Ge_2,\\
 d_R&=2d_1+4\gamma_0d_0.
 \end{aligned}
}\end{equation} Then \(\|\Gamma\|,\|\Gamma_N\|\le\gamma_0\),
\(\|D\Gamma\|,\|D\Gamma_N\|\le\gamma_1\),
\(\|\Gamma_N-\Gamma\|\le d_0\), \(\|D\Gamma_N-D\Gamma\|\le d_1\) and
\(\|R_N-R\|\le d_R\). All these are coordinate tensor bounds on the
common tube.

\emph{Proof.} For \(g_s=g+s(g_N-g)\), \(\|g^{-1}(g_s-g)\|\le1/2\). The
Neumann series gives inverse bound \(2G\) and the inverse difference
identity gives the stated estimate. The path \(g_s\) consists of real
invertible symmetric matrices, so inertia is constant. No indefinite
positive norm is used.

Write \(\Gamma=\tfrac12g^{-1}A(Dg)\), where
\(A(Dg)(u,v,w)=Dg[u](v,w)+Dg[v](u,w)-Dg[w](u,v)\). Its norm is at most
\(3\|Dg\|\). Differentiate \(g^{-1}\) by \(Dg^{-1}=-g^{-1}(Dg)g^{-1}\).
The displayed bounds for \(\gamma_0,\gamma_1\) follow uniformly along
\(g_s\). The derivative in \(s\) of \(\Gamma_s\) is bounded by
\(6G^2(K_1+1)e_0+3Ge_1\). For \(D\Gamma_s\) its two inverse factors and
two first derivatives contribute respectively \(24G^3(K_1+1)^2e_0\) and
\(12G^2(K_1+1)e_1\); its second-derivative term contributes
\(6G^2(K_2+1)e_0+3Ge_2\). Integrating in \(s\) proves \(d_0,d_1\).

Finally
\(R(a,b)c=D\Gamma[a](b,c)-D\Gamma[b](a,c)+\Gamma(a,\Gamma(b,c))-\Gamma(b,\Gamma(a,c))\).
Subtract the two formulas and use the uniform \(\gamma_0\) bound,
obtaining \(2d_1+4\gamma_0d_0\). ∎

\pandocstatementheading{\protect\phantomsection\label{lem:sh-null-frame-initialization}Lemma K.11.}

A real initial null frame \((k,n,e_1,e_2)\) with \(g(k,n)=1\),
\(g(e_A,e_B)=-\delta_{AB}\), and every other pairing zero has a real,
locally analytic normalization for every sufficiently small symmetric
perturbation of \(g\). On a compact family with these initial vectors
bounded, its difference from the original frame is at most
\(C_{\rm init}e_0\) for a finite constant depending only on the
displayed frame bounds and the prescribed nonzero denominator margins.

\emph{Proof.} First choose the small root \(a\) of \(g_N(k+an,k+an)=0\),
with \(a=0\) at \(g_N=g\). The derivative in \(a\) there is
\(2g(k,n)=2\), so the analytic implicit theorem gives the root,
uniformly on the compact family. It can also be found by the quadratic
formula using the branch tending to zero. Put \(k_N=k+an\),
\begin{equation}\protect\phantomsection\label{eq:sh-null-frame-normalization}{\begin{gathered}
b=g_N(k_N,n), \\
n_N=b^{-1}\left(n-\frac{g_N(n,n)}{2b}k_N\right), \\
p_A=e_A-g_N(e_A,n_N)k_N-g_N(e_A,k_N)n_N.
\end{gathered}}\end{equation} These satisfy \(g_N(k_N,n_N)=1\) and both null
conditions, and \(p_A\) is orthogonal to \(k_N,n_N\). The matrix
\(H_{AB}=-g_N(p_A,p_B)\) tends to the identity. Require \(|b|\ge1/2\)
and \(\|H-I\|\le1/2\); its positive inverse square root is analytic.
Then \((e_{1,N},e_{2,N})=(p_1,p_2)H^{-1/2}\) is the desired screen
frame. All operations are analytic on these fixed margins. The supremum
of their first derivatives on a smaller compact matrix neighborhood is a
finite \(C_{\rm init}\). The mean-value formula gives the estimate. This
constant is specified by a finite supremum, not numerically evaluated
here. ∎

\pandocstatementheading{\protect\phantomsection\label{thm:sh-finite-optical-stability}Theorem K.12.}

Fix a compact family of the finite Schwarzschild null segments in the
explicit screen calculation, a coordinate/state tube with positive
buffer \(\eta_Z\), and individual affine-time intervals
\([0,S_\lambda]\) with a common bound \(S_\lambda\le S\). Each ray is
retained only on its own interval, with a fixed small extension when a
terminal hitting time is compared. Let \(Z=(x,k,n,e_1,e_2)\) consist of
the position, tangent, null companion and parallel screen frame. On a
larger compact tube assume their vector norms are at most \(V\ge1\), and
put
\begin{equation}\protect\phantomsection\label{eq:sh-geodesic-frame-error}{
 L_Z=1+\gamma_1 V^2+2\gamma_0V,\qquad
 \epsilon_Z=e^{L_ZS}\{C_{\rm init}e_0+SV^2d_0\}.
}\end{equation} Require \(\epsilon_Z<\eta_Z\), as well as the metric and
frame margins of Lemmas K.10--K.11. Then, for each parameter
\(\lambda\), the normalized \(g_N\)-null geodesic and its parallel frame
exist on that same individual finite time interval, remain on the common
parent and differ from \(Z\) by at most \(\epsilon_Z\).

Let \(R_0=\sup\|R\|\), \(R_1=\sup\|DR\|\) on that tube. These are
finite, with \(R_1\) controlled by the third metric derivative and
inverse bounds. Define
\begin{equation}\protect\phantomsection\label{eq:sh-tidal-error}{
 \epsilon_{\mathcal T}=2V^4\{R_0e_0+(K_0+1)d_R\}
 +2\{V^4(K_1R_0+K_0R_1)+4K_0R_0V^3\}\epsilon_Z.
}\end{equation} The \(2\times2\) screen tidal matrices differ by at most
\(\epsilon_{\mathcal T}\). If their first-order Jacobi generators have
norm at most \(B\) on the common interval, their full transfer matrices
obey
\begin{equation}\protect\phantomsection\label{eq:sh-jacobi-transfer-error}{
 \|U_N(t,s)-U(t,s)\|
 \le (t-s)e^{B(t-s)}\epsilon_{\mathcal T},\qquad s\le t\le S_\lambda\le S.
}\end{equation} One may choose
\(B=1+\sup\|\mathcal T\|+\epsilon_{\mathcal T}\) in the induced
Euclidean phase norm. The estimate also holds for a fixed finite
subdivision, with the same total-time factor and the largest tidal error
in the finitely many fixed screen charts.

\emph{Proof.} In the fixed coordinates the joint equations are
\(\dot x=k\), \(\dot k=-\Gamma(k,k)\), and \(\dot z=-\Gamma(k,z)\) for
\(z=n,e_1,e_2\). Use the maximum of the Euclidean norms of these five
blocks as state norm. Their coordinate derivative has row sums at most
\(L_Z\) on the stated larger tube. At the same state the two right-hand
sides differ by at most \(V^2d_0\). The initial difference has the bound
of Lemma K.11. The integral inequality and Gronwall give the displayed
\(\epsilon_Z\) until the first exit. Its strict buffer bound excludes
that exit, proving existence on the full fixed interval. Nullity, frame
pairings and signature follow from metric compatibility and the exact
normalized initial values.

The tidal entries are \(\mathcal T_{AB}=-g(R(e_A,k)k,e_B)\) in the
signature of §2. At a fixed state the metric and curvature perturbations
contribute at most \(V^4\{R_0e_0+(K_0+1)d_R\}\). Moving the position
contributes \(V^4(K_1R_0+K_0R_1)\epsilon_Z\), and moving the four vector
arguments contributes \(4K_0R_0V^3\epsilon_Z\). The factor two changes
an entrywise bound into a safe \(2\times2\) operator bound. This proves
the tidal estimate.

The first-order phase generator is
\(A=\left(\begin{smallmatrix}0&I\\-\mathcal T&0\end{smallmatrix}\right)\).
The variation-of-constants identity gives
\(U_N(t,s)-U(t,s)=\int_s^tU_N(t,u)(A_N(u)-A(u))U(u,s)\,du\). Both
fundamental matrices have bound \(e^{B|t-s|}\); their product inside the
integral has bound \(e^{B(t-s)}\). This proves the estimate. On a
subdivision, exact matching of the full intermediate phase state gives
matrix multiplication; the same Duhamel identity applies piecewise.
Orthogonal changes of screen basis preserve the norms. ∎

For the explicit screen calculation, the initial companion can be chosen
as \(n_0=(-\partial_r)/\dot v(0)\). At \(u=0\) its \(H\) coefficient is
zero, so \(g(k,n_0)=1\) and it is orthogonal to both initial screen
vectors. The lower bound \(\dot v(0)\ge1/410\) gives the required
compact companion bound.

The comparison is at fixed affine parameter in a common coordinate
domain. If terminal data are instead read on the fixed surface
\(r=r_-\), a separate hitting-time margin is needed. The exact family in
the explicit screen calculation has \(-\dot r\ge a>0\). If
\(\epsilon_Z<a/2\), the approximate family is also transverse, and the
difference of hitting times is at most \(2\epsilon_Z/a\), provided both
are extended by this amount within the retained buffers. Here \(S\) is
chosen to include those individual small extensions, and \(B\),
\(\epsilon_Z\) and all tube bounds are taken on the same extended
intervals; no member is required to remain in the radial tube up to the
maximum endpoint time of a different member. The transfer at the moved
endpoint acquires at most \(2Be^{BS}\epsilon_Z/a\). No trapped or
infinitely long ray estimate is involved.

To use the dimensionless screen equations, distinguish physical affine
time \(u\) from \(t=u/m\). The phase bases are \(Z_u=(x,dx/du)\) and
\(Z_t=(x,m\,dx/du)\), with
\begin{equation}\protect\phantomsection\label{eq:sh-phase-time-dictionary}{
 C_m=\operatorname{diag}(I_2,mI_2),\qquad
 Z_t=C_mZ_u,\qquad U_t=C_mU_uC_m^{-1},\qquad
 \mathcal T_t=m^2\mathcal T_u.
}\end{equation} Thus the coefficient error for the evaluated
dimensionless system is at most \(m^2\epsilon_{\mathcal T,u}\). A
physical hitting-time error \(\Delta u\) becomes
\(\Delta t=\Delta u/m\). If one uses the already proved physical
transfer estimate instead, its norm must be multiplied by
\(\|C_m\|\,\|C_m^{-1}\|=\max(m,m^{-1})\). The physical point-source
endpoint block is \(mB_t\), and the expansion traces satisfy
\(\theta_t=m\theta_u\). The mass is fixed, so these are finite
constants, not uniform zero-mass limits. To apply the evaluated
\(4\varepsilon\) and endpoint-inverse margins of the explicit screen
calculation, impose the dimensionless condition
\(m^2\epsilon_{\mathcal T,u}\le1/10\) (including any separate
coefficient approximation error).

\pandocstatementheading{\protect\phantomsection\label{cor:sh-optical-observables}Corollary K.13.}

Write \(U=\left(\begin{smallmatrix}A&B_0\\ C&D\end{smallmatrix}\right)\)
and let \(\epsilon_U\) denote the right-hand side of the transfer
estimate, including a chosen hitting-time term when applicable. For an
initial phase of norm at most \(J\), the outgoing phase error is at most
\(J\epsilon_U\). The signed point-source beam determinant obeys
\begin{equation}\protect\phantomsection\label{eq:sh-beam-determinant-error}{
 |\det B_{0,N}-\det B_0|
 \le2\|B_0\|\epsilon_U+\epsilon_U^2.
}\end{equation} If additionally the least singular value of \(B_0\) is
at least \(b>0\) and \(\epsilon_U\le b/2\), then
\(\|B_{0,N}^{-1}\|\le2/b\) and
\begin{equation}\protect\phantomsection\label{eq:sh-expansion-error}{
 |\operatorname{tr}(D_NB_{0,N}^{-1})-\operatorname{tr}(DB_0^{-1})|
 \le4\epsilon_U\{b^{-1}+\|D\|b^{-2}\}.
}\end{equation} \emph{Proof.} Block projections have operator norm one.
Expand the two-dimensional determinant by its columns; the two linear
terms have magnitude at most \(\|B_0\|\epsilon_U\) and the quadratic
term at most \(\epsilon_U^2\). The inverse estimates follow from the
Neumann series and
\(B_{0,N}^{-1}-B_0^{-1}=-B_{0,N}^{-1}(B_{0,N}-B_0)B_0^{-1}\). Subtract
the two products and use \(|\operatorname{tr}Q|\le2\|Q\|\). ∎

The phase and determinant estimates continue through a caustic. The
expansion estimate is restricted to its positive singular-value margin.
The full transfer relation, rather than an inverse endpoint block, is
the observable retained in the stationary representation of the explicit
screen calculation. Constants
\(H_g,T,\rho,G,K_j,V,R_0,R_1,\eta_Z,C_{\rm init}\) come from specified
analytic parents and compact-frame margins; they are finite theoretical
constants. The dimensionless Schwarzschild tidal and crossing bounds of
the explicit screen calculation are separately evaluated constants. No
floating-point experiment is used as a proof of these estimates.

\protect\phantomsection\label{sec:sh-r1-primitive-comparison}

\subsubsection{K.5. A calibrated Schwarzschild source
primitive}\label{k.5.-a-calibrated-schwarzschild-source-primitive}

Fix the physical mass \(m>0\) of Section 10 and put
\(\mu_{\rm src}=m/2\). We compute a specified local primitive of AASS
v2: its explicit principal spinors, the first line of its restricted
Dolbeault operator (3.12), and the contour on the left of (2.32) (Adamo
et al. 2026a, equations (2.5), (2.10), (2.15), (2.26), (2.32) and
(3.12)). The defining differential operator and integral will be
evaluated directly.

This construction is on one local source chart. Put \(u_s=x+iy\), choose
a signed square root \(\varrho^2=x^2+y^2+z^2\), and require
\(u_s\varrho\ne0\). Choose the square root \(q\) on a simply connected
parent and set
\begin{equation}\protect\phantomsection\label{eq:sh-r1-source-coordinates}{\begin{gathered}
q^2=-\frac{iu_s}{2}, \\
z_-=\frac{z+\varrho}{u_s}, \\
z_+=\frac{z-\varrho}{u_s}, \\
X^{aa'}=\frac1{\sqrt2}
 \begin{pmatrix}\tau-iz&ix-y\\ix+y&\tau+iz\end{pmatrix}.
\end{gathered}}\end{equation} For this source calculation use the raising and lowering
order of (Adamo 2018, arXiv v2, §1.2, equations (1.5)--(1.10)). The
upper and lower epsilon \emph{symbols} have the same numerical array
\(E\), whereas their left-acting matrices differ:
\begin{equation}\protect\phantomsection\label{eq:sh-r1-source-index-order}{\begin{gathered}
E=\begin{pmatrix}0&1\\-1&0\end{pmatrix}, \\
L=E^T=-E, \\
v^a=E^{ab}v_b, \\
v_a=v^bE_{ba}=L_{ab}v^b.
\end{gathered}}\end{equation} Apply the same rule in the primed factor. In particular
\(\langle ab\rangle=a^cb_c\), \(dX^b{}_{a'}=L_{a'c'}dX^{bc'}\), and
\(\varphi_{a'}{}^{b'}=E^{b'c'}\varphi_{a'c'}\). These are local source
conventions, distinct from the fixed component convention of §2. The
flat tensor is nevertheless
\begin{equation}\protect\phantomsection\label{eq:sh-r1-source-flat-metric}{
 \delta=L_{ab}L_{a'b'}dX^{aa'}dX^{bb'}
       =d\tau^2+dx^2+dy^2+dz^2,
}\end{equation} since changing both alternating factors from \(E\) to
\(L=-E\) has no effect on their product.

The required arrays, with lower dyad \(o_a=(1,0)\), are
\begin{equation}\protect\phantomsection\label{eq:sh-r1-source-arrays}{
 \begin{aligned}
 \alpha^{a'}&=q(z_-,-1),&\beta^{a'}&=q(z_+,-1),\\
 \alpha_{a'}&=q(1,z_-),&\beta_{a'}&=q(1,z_+),\\
 A_a&=q(-z_-,1),&B_a&=q(-z_+,1),\\
 \varphi_{a'b'}&=-\frac{4\mu_{\rm src}}{\varrho q^2}
                       \alpha_{a'}\alpha_{b'}.
 \end{aligned}
}\end{equation} They give \(s=\langle AB\rangle=-i\varrho\),
\([\beta\alpha]=i\varrho\) and \(c=\langle oA\rangle^2=q^2\). In
particular \(A,B\) are independent. Define an endomorphism of the
complexified tangent bundle by
\begin{equation}\protect\phantomsection\label{eq:sh-r1-primitive-projector}{
 \mathsf P^{bb'}{}_{aa'}=
 \frac{A^bB_a}{s}\delta_{a'}^{b'}
 +\frac{c}{s^2}\varphi_{a'}{}^{b'}B_aB^b,
 \qquad J=i(1-2\mathsf P).
}\end{equation} The image of \(\mathsf P\) is designated the \(-i\)
eigenspace. Define
\(Q(\lambda)=\langle A\lambda\rangle\langle B\lambda\rangle\),
\(D\lambda=\lambda^a d\lambda_a\), and, with the base point fixed, take
a positively oriented fiber contour around \([A]\) excluding \([B]\):
\begin{equation}\protect\phantomsection\label{eq:sh-r1-primitive-contour}{
 \omega_{\rm int}=\frac1{2\pi}\oint_{[A]}
 \frac{\lambda_a\lambda_b}{Q(\lambda)^2}
 D\lambda\wedge dX^{aa'}\wedge dX^b{}_{a'},
 \qquad
 G_{\rm int}(U,V)=\omega_{\rm int}(U,JV).
}\end{equation} We compare the two sides of (Adamo et al. 2026a,
equation (2.32)) with the displayed source index rule, a contour making
one positive turn around \([A]\), and the ordinary exterior-product
convention
\begin{equation}\protect\phantomsection\label{eq:sh-r27-exterior-evaluation}{
 (\alpha\wedge\beta)(U,V)=\alpha(U)\beta(V)-\alpha(V)\beta(U).
}\end{equation} For a two-form we use
\(\omega=\tfrac12\omega_{IJ}\,dX^I\wedge dX^J\), so
\(\omega(U,V)=\omega_{IJ}U^IV^J\), where \(I=(a,a')\). With these
conventions the left side of the cited (2.32), including its printed
\(1/(2\pi)\) and \(D\lambda=\langle\lambda,d\lambda\rangle\), is exactly
\(\omega_{\rm int}\) in eq.~\ref{eq:sh-r1-primitive-contour}. Its
printed right side is exactly \(\omega_0\) in
eq.~\ref{eq:sh-r1-contour-value}. Proposition K.14 evaluates their
relation directly.

\pandocstatementheading{\protect\phantomsection\label{prop:sh-r1-primitive-tensor}Proposition K.14 (Primitive tensor calculation).}

On this source chart, \(\mathsf P\) is a rank-two projector and its two
eigendistributions are involutive. The form and symmetric metric defined
above are
\begin{equation}\protect\phantomsection\label{eq:sh-r1-contour-value}{
 \omega_{\rm int}=2\omega_0,\qquad
 \omega_0=\frac{i}{s^3}A_aB_b\,dX^{aa'}\wedge dX^b{}_{a'},
}\end{equation}
\begin{equation}\protect\phantomsection\label{eq:sh-r1-general-source-metric}{
 G_{\rm int}=2G_0,\qquad
 G_0=-\frac1{\varrho^2}
 \left[\delta+\frac{4\mu_{\rm src}}{\varrho}
                  (d\tau+i\,d\varrho)^2\right].
}\end{equation} The form is closed and nondegenerate, and
\(\omega_0(JU,JV)=\omega_0(U,V)\). These identities hold on the whole
chart with either signed radial branch.

\emph{Proof.} Write \(\mathsf P=\mathsf P_0+\mathsf N\) for the two
terms in eq.~\ref{eq:sh-r1-primitive-projector}. The contractions
\(B_aA^a=s\), \(A_aB^a=-s\) and \(B_aB^a=0\) give
\begin{equation}\protect\phantomsection\label{eq:sh-r1-projector-identities}{
 \mathsf P_0^2=\mathsf P_0,\qquad
 \mathsf P_0\mathsf N=0,\qquad
 \mathsf N\mathsf P_0=\mathsf N,\qquad \mathsf N^2=0.
}\end{equation} Thus \(\mathsf P^2=\mathsf P\) and
\(\operatorname{tr}\mathsf P=2\). Its kernel is spanned by
\(B^a\partial_{aa'}\) and its image by
\begin{equation}\protect\phantomsection\label{eq:sh-r1-projector-image}{
 A^a\partial_{aa'}+\frac c s\varphi_{a'}{}^{b'}B^a\partial_{ab'}.
}\end{equation} Write \(V_{a'}\) for the displayed image generator. The
first line of (Adamo et al. 2026a, equation (3.12)) gives
\(\mathsf P(A^a\partial_{aa'})=V_{a'}\). Symmetry of \(\varphi_{a'b'}\)
and the source index rule give
\(\varphi^{b'}{}_{a'}=\varphi_{a'}{}^{b'}\). The annihilator of this
image is spanned by
\begin{equation}\protect\phantomsection\label{eq:sh-r29-projector-annihilator}{
 \begin{aligned}
 \vartheta^{b'}&=A_a\,dX^{ab'}
       +\frac cs\varphi^{b'}{}_{c'}B_a\,dX^{ac'},\\
 \vartheta^{b'}(V_{a'})
       &=-c\varphi_{a'}{}^{b'}+c\varphi^{b'}{}_{a'}=0,\\
 \vartheta^{b'}(B^a\partial_{aa'})&=-s\delta_{a'}^{b'}.
 \end{aligned}
}\end{equation} The last identity proves independence. The vectors and
forms printed in (Adamo et al. 2026a, equations (3.14)--(3.15)) instead
have minus signs in both correction terms; the displayed forms
annihilate the displayed vectors. Here the plus signs are fixed by the
adopted first line of (3.12). Also \(J^2=-1\).

Use the local homogeneous lift \(\lambda=A+wB\) of the fiber contour.
Directly,
\begin{equation}\protect\phantomsection\label{eq:sh-r1-residue}{
 Q=-s^2w,\qquad D\lambda=s\,dw,\qquad
 \operatorname*{Res}_{w=0}\frac{\lambda_a\lambda_bD\lambda}{Q^2}
       =\frac{A_aB_b+B_aA_b}{s^3}.
}\end{equation} The tensor \(dX^{aa'}\wedge dX^b{}_{a'}\) is symmetric
in \(a,b\), because exchanging its two exterior factors and its dotted
alternating factor gives two minus signs. Multiplication by
\((2\pi)^{-1}2\pi i\) therefore gives exactly \(2\omega_0\). This
calculation uses the printed prefactor and the stated contour
orientation. Thus the left side of (Adamo et al. 2026a, equation (2.32))
evaluates to twice its printed right side under the conventions just
specified. Replacing the contour prefactor by \(1/(4\pi)\) would instead
give \(\omega_0\). The two-form components below are bilinear components
in this convention, so they contain \(A_aB_b+B_aA_b\), rather than
\(A_{(a}B_{b)}\); the factor \(1/2\) belongs to the component expansion
of the two-form.

The two-form components and the undeformed endomorphism are
\begin{equation}\protect\phantomsection\label{eq:sh-r1-form-components}{\begin{gathered}
(\omega_0)_{aa',bb'}=
 \frac{i}{s^3}(A_aB_b+B_aA_b)L_{a'b'}, \\
J_0{}^{bb'}{}_{aa'}=-\frac i s(B_aA^b+A_aB^b)\delta_{a'}^{b'}, \\
J=J_0-2i\mathsf N.
\end{gathered}}\end{equation} Using \(A_aB_b-B_aA_b=sL_{ab}\), contract the second
slot of this two-form with \(J\). The undeformed contraction is
\(L_{ab}L_{a'b'}/s^2\). The deformation contraction uses \(A_cB^c=-s\)
and \(L_{a'c'}\varphi_{b'}{}^{c'}=\varphi_{a'b'}\). Consequently
\begin{equation}\protect\phantomsection\label{eq:sh-r1-unspecialized-metric}{
 G_0=\frac{\delta}{s^2}
       -\frac{2c}{s^4}B_aB_b\varphi_{a'b'}dX^{aa'}dX^{bb'}.
}\end{equation} All its terms are symmetric. The coefficient two here is
the evaluation coefficient of a two-form; it is additional to the
contour factor two already computed.

For the remaining tensor reduction put \(a_s=B_a\alpha_{a'}dX^{aa'}\).
The arrays give
\begin{equation}\protect\phantomsection\label{eq:sh-r1-null-form}{
 \begin{aligned}
 a_s=\frac{q^2}{\sqrt2}\bigl[&(z_--z_+)d\tau
   +i(1-z_-z_+)dx+(1+z_-z_+)dy
   +i(z_-+z_+)dz\bigr]\\
 =\frac1{\sqrt2}&(x\,dx+y\,dy+z\,dz-i\varrho\,d\tau)
 =\frac{\varrho}{\sqrt2}(d\varrho-i\,d\tau).
 \end{aligned}
}\end{equation} Here \(z_--z_+=2\varrho/u_s\), \(z_-+z_+=2z/u_s\),
\(z_-z_+=-(x-iy)/u_s\) and \(\varrho\,d\varrho=x\,dx+y\,dy+z\,dz\).
These relations are valid on both radial sheets. Substitution into
eq.~\ref{eq:sh-r1-unspecialized-metric} gives
\begin{equation}\protect\phantomsection\label{eq:sh-r1-metric-reduction}{
 G_0=-\frac\delta{\varrho^2}+\frac{8\mu_{\rm src}}{\varrho^5}a_s^2
 =-\frac1{\varrho^2}
   \left[\delta+\frac{4\mu_{\rm src}}{\varrho}
                          (d\tau+i\,d\varrho)^2\right].
}\end{equation} This proves the metric formula at general points,
without using finitely many test points.

For closure and integrability one can retain the equally explicit form
calculation. The symmetric matrix \(S_{ab}=A_aB_b+B_aA_b\) is
\begin{equation}\protect\phantomsection\label{eq:sh-r1-symmetric-spinor-matrix}{
 S=i\begin{pmatrix}x-iy&z\\z&-x-iy\end{pmatrix}.
}\end{equation} Substitute this matrix and \(dX\) into
eq.~\ref{eq:sh-r1-form-components}. It gives the Cartesian expression
\begin{equation}\protect\phantomsection\label{eq:sh-r1-cartesian-form}{
 \omega_0=\frac1{\varrho^3}
 \left[d\tau\wedge(x\,dx+y\,dy+z\,dz)
       -x\,dy\wedge dz-y\,dz\wedge dx-z\,dx\wedge dy\right].
}\end{equation} On a local polar chart
\((x,y,z)=\varrho(\sin\Theta\cos\psi,\sin\Theta\sin\psi,\cos\Theta)\) it
becomes
\begin{equation}\protect\phantomsection\label{eq:sh-r1-polar-form}{
 \omega_0=\varrho^{-2}d\tau\wedge d\varrho
                -\sin\Theta\,d\Theta\wedge d\psi.
}\end{equation} The one-form \(d\tau+i\,d\varrho\) is \(\delta\)-null,
so the rank-one determinant formula in Cartesian coordinates gives
\(\det G_0=\varrho^{-8}\). Thus \(\omega_0\) is nondegenerate throughout
the source chart. The polar formula proves closure on its regular
subcharts, and the holomorphic Cartesian formula extends closure across
the remaining points. Since \(G_0=\omega_0(\cdot,J\cdot)\) is symmetric
and \(J^2=-1\), the claimed compatibility follows as well. In the radial
and angular coordinate orders the same equality determines
\begin{equation}\protect\phantomsection\label{eq:sh-r1-source-complex-structure}{
 J_{\rm rad}=\begin{pmatrix}ia&1-a\\-(1+a)&-ia\end{pmatrix},\qquad
 J_{\rm ang}=\begin{pmatrix}0&-\sin\Theta\\\csc\Theta&0\end{pmatrix},
 \qquad a=4\mu_{\rm src}/\varrho.
}\end{equation} The \(-i\) eigenspace is spanned by
\((1-a)\partial_\tau-i(1+a)\partial_\varrho\) and
\(\partial_\Theta+i\csc\Theta\partial_\psi\); the \(+i\) eigenspace is
spanned by \(\partial_\tau+i\partial_\varrho\) and
\(\partial_\Theta-i\csc\Theta\partial_\psi\). The two generators in each
pair commute. Their radial generators never vanish and the two radial
generators have determinant \(2i\) in the displayed order. Hence both
distributions are involutive, including where \(1+a=0\). The Nijenhuis
tensor is holomorphic in the Cartesian primitives; its vanishing on the
dense regular polar subcharts therefore extends across their coordinate
degeneracies. This proves integrability without a static logarithm or an
external signature-dependent theorem. ∎

The chosen graph also has its stated primitive origin. In the source
convention the quadric is exactly
\begin{equation}\protect\phantomsection\label{eq:sh-r1-explicit-quadric}{\begin{aligned}
Q&=-\frac i2\left[u_s\lambda_0^2+2z\lambda_0\lambda_1
                         -(x-iy)\lambda_1^2\right] \\
&=\frac1{\sqrt2}\left[-X^{01'}\lambda_0^2
 +(X^{00'}-X^{11'})\lambda_0\lambda_1+X^{10'}\lambda_1^2\right].
\end{aligned}}\end{equation} The vertical \(\partial/\partial\hat\lambda\) term of
the ambient operator vanishes because \(Q\) is holomorphic in
\(\lambda\) and independent of \(\hat\lambda\). Base derivatives are
taken with \(\lambda\) fixed before graph substitution. They give
\(\lambda^a\partial_{aa'}Q=0\) and
\begin{equation}\protect\phantomsection\label{eq:sh-r1-coincidence-check}{
 \alpha^{b'}\partial_{bb'}Q
   =\frac q{\sqrt2}(\lambda_0+z_-\lambda_1)(\lambda_0,\lambda_1)_b.
}\end{equation} The latter vanishes at \(\lambda=A\). It kills the
deformation term of the ambient operator (2.21), while the former kills
its flat horizontal term. Thus \(Q=\overline\nabla Q=0\) on the selected
graph. On the projective chart \(A_0\ne0\), its coordinate is
\(A_1/A_0=-e^{i\psi}\tan(\Theta/2)\), whose differential annihilates the
displayed \(-i\) distribution. This is a local check of the selected
graph and restricted primitive. It introduces no assertion about a
global coincidence surface or completeness of a positive-definite SDTN
slice.

\pandocstatementheading{\protect\phantomsection\label{thm:sh-r1-calibrated-horizon-map}Theorem K.15 (Calibrated primitive-to-Schwarzschild map).}

Let \(U\) be a sufficiently small simply connected complex EF parent of
\(o=(0,2m,\pi/2,0)\) on which \(r\), \(\sin\theta\) and the following
source units are nonzero. Use \(\mu_{\rm src}=m/2\), take
\(\varrho=-r\), and define
\begin{equation}\protect\phantomsection\label{eq:sh-r1-horizon-map}{
 \Phi_{\rm pr}(v,r,\theta,\phi)=
 \left(\begin{aligned}\tau=i(r-v),\ x=r\sin\theta\cos\phi,\\
       \ y=r\sin\theta\sin\phi,\ z=r\cos\theta\end{aligned}\right).
}\end{equation} The source branches of \(q\) and the contour pole
\([A]\) are fixed as above. Then \(\Phi_{\rm pr}\) and its actual
inverse are regular across \(r=2m\), and
\begin{equation}\protect\phantomsection\label{eq:sh-r1-calibrated-pullback}{
 h_*:=\Phi_{\rm pr}^*G_{\rm int}=\frac2{r^2}g_m,\qquad
 \sigma_*=\frac{\sqrt2}{r},\qquad \sigma_*^{-2}h_*=g_m.
}\end{equation} With the density root and full spin lift fixed by their
marked central values, this is an actual calibrated comparison. In the
vector-middle-slot normal-tractor convention its auxiliary components
are
\begin{equation}\protect\phantomsection\label{eq:sh-r1-vector-tractor}{
 I_*=
 \left(\frac{\sqrt2}{r},\frac{\partial_v+f\partial_r}{\sqrt2},
                         \frac f{2\sqrt2r}\right),\qquad
 DI_*=0,\qquad H(I_*,I_*)=0.
}\end{equation} \emph{Proof.} The Jacobian in coordinate orders
\((\tau,x,y,z)\) and \((v,r,\theta,\phi)\) is \(-ir^2\sin\theta\), a
unit at the mark. The inverse is
\begin{equation}\protect\phantomsection\label{eq:sh-r1-horizon-inverse}{
 r=-\varrho,\qquad v=r+i\tau,\qquad
 \theta=\arccos(z/r),\qquad
 \phi=-i\log\left(\frac{x+iy}{r\sin\theta}\right),
}\end{equation} with the inverse branches initialized at \(o\). Shrink
\(U\) so these are mutual analytic inverses. No inverse power of \(f\)
occurs. Direct substitution gives
\begin{equation}\protect\phantomsection\label{eq:sh-r1-pullback-identities}{\begin{gathered}
d\tau+i\,d\varrho=-i\,dv, \\
\Phi_{\rm pr}^*\delta=-dv^2+2\,dv\,dr+r^2d\Omega^2, \\
4\mu_{\rm src}/\varrho=-2m/r.
\end{gathered}}\end{equation} The bracket in eq.~\ref{eq:sh-r1-general-source-metric}
therefore pulls back to \(-g_m\), proving the factor \(2/r^2\) in
eq.~\ref{eq:sh-r1-calibrated-pullback}. The source density coefficient
\(\sigma_{\rm src}=-\sqrt2/\varrho\) pulls back to the stated positive
\(\sigma_*\). Thus
\(\Phi_{\rm pr}^*(\sigma_{\rm src}^{-2}G_{\rm int})=g_m\) with no
constant homothety left over.

The sign of the radial sheet is an analytic continuation choice.
Starting from a real spatial point with positive radius and \(u_s\ne0\),
the finite path
\begin{equation}\protect\phantomsection\label{eq:sh-r1-sheet-path}{
 (x_t,y_t,z_t)=e^{i\pi t}(x_0,y_0,z_0),\qquad
 \varrho_t=e^{i\pi t}\varrho_0,\qquad
 q_t=e^{i\pi t/2}q_0,\qquad 0\le t\le1
}\end{equation} has no zero of \(u_s,\varrho,q,s\) or \([\beta\alpha]\).
It ends at the antipodal real spatial point on the negative sheet.
Choosing the initial antipode of the desired point gives the required
endpoint. The ratios \(z_\pm\) are constant along the path, so the
continued poles remain distinct, and the coordinate contour
\(|w|=\epsilon\) in \(\lambda=A+wB\) transports its positive orientation
throughout. A finite chain of local neighborhoods covers this compact
path. No global real SDTN interpretation is needed for this
continuation.

One can also check the continued tensors directly:
\begin{equation}\protect\phantomsection\label{eq:sh-r1-continued-structure}{\begin{gathered}
\Phi_{\rm pr}^*\omega_0=\frac i{r^2}dv\wedge dr
                          +\sin\theta\,d\theta\wedge d\phi,
 \quad
 J_*^{\rm rad}=\begin{pmatrix}-i&0\\-if&i\end{pmatrix}, \\
J_*^{\rm ang}=\begin{pmatrix}0&\sin\theta\\-\csc\theta&0\end{pmatrix}.
\end{gathered}}\end{equation} Their product is \(r^{-2}g_m\). The angular sign follows
from the antipodal source angular coordinates \(\Theta=\pi-\theta\),
\(\psi=\phi+\pi\). These are complex continued tensors on the Lorentz
slice; a real orthogonal complex structure for a real Lorentz metric is
not asserted.

To track the density itself, let \(e_g\) be the unit frame in the
physical scale. For \(h_0=r^{-2}g_m\) and \(h_*=2h_0\), choose
\begin{equation}\protect\phantomsection\label{eq:sh-r1-density-frame}{\begin{gathered}
\boldsymbol g=g_me_g^2=h_0e_c^2=h_*e_*^2, \\
e_c=re_g,\quad e_*=(r/\sqrt2)e_g, \\
\boldsymbol\sigma=e_g=\frac1r e_c=\frac{\sqrt2}{r}e_*.
\end{gathered}}\end{equation} These are components of one section. The preceding scale
calculation gives, in \(h_0\), the vector-slot tractor
\((1/r,\partial_v+f\partial_r,f/(2r))\). Under a constant auxiliary
rescaling \(\widehat h=\Omega^2h\) its components transform as
\((\Omega\sigma,\Omega^{-1}\mu,\Omega^{-1}\rho)\). Taking
\(\Omega=\sqrt2\) gives eq.~\ref{eq:sh-r1-vector-tractor}. In contrast,
the corresponding covector component is
\begin{equation}\protect\phantomsection\label{eq:sh-r1-vector-covector-dictionary}{
 d\sigma_*=-\frac{\sqrt2}{r^2}dr,
 \qquad \mu_*=(d\sigma_*)^{\sharp_{h_*}}
               =\frac{\partial_v+f\partial_r}{\sqrt2}.
}\end{equation} Thus the earlier covector \(-dr/r^2\) belongs to
\(h_0\), not to the untransformed middle vector slot of \(I_*\). The
normal equations and norm are preserved by the specified splitting
change. Explicitly \(2\sigma_*\rho_*=f/r^2\) and
\(h_*(\mu_*,\mu_*)=-f/r^2\). At \(r=2m\) the tractor is
\((\sqrt2/(2m),\partial_v/\sqrt2,0)\), regular and nonzero. Changing to
the physical splitting by \(\Omega=r/\sqrt2\) gives \((1,0,0)\).

If a coefficient equal to one at \(o\) is required, put \(r_0=r(o)\) and
simultaneously use
\begin{equation}\protect\phantomsection\label{eq:sh-r1-point-normalization}{
 h_o=\frac{r_0^2}{2}h_* =\frac{r_0^2}{r^2}g_m,
 \qquad \sigma_o=\frac{r_0}{r},\qquad
 \sigma_o(o)=1,\qquad \sigma_o^{-2}h_o=g_m.
}\end{equation} The physical affine calibration is retained by this
simultaneous change. In the final Einstein scale the scalar optical
section has Cauchy data \((1,0)\). Its root is the retained square
identification \(\beta:\mathcal E[1]^2\to\mathcal B\), satisfying
\(\beta(\boldsymbol\sigma^2)=\Theta^{-1}\) in the calibrated affine
torsor. The marked positive value fixes the density-root sign. Replacing
\(\sigma_{\rm src}\) by \(+\sqrt2/\varrho\) would give the same metric
but the opposite selected density and tractor, violating the fixed
central root value.

Finally take the regular physical tetrad \(l=\partial_v+f\partial_r/2\),
\(n=-\partial_r\), \(e_2=r^{-1}\partial_\theta\),
\(e_3=(r\sin\theta)^{-1}\partial_\phi\) and \(e_0=(l+n)/\sqrt2\),
\(e_1=(l-n)/\sqrt2\). Declare \(e_0\) future and choose its full spin
lift at \(o\). Push this tetrad through \(d\Phi_{\rm pr}\) to the
selected source metric, and extend the chosen spin lift on the
contractible parent. This specifies an actual marked isometry and
inverse. The central Spin sign is retained separately from the density
sign and the choice of \(q\); no principal-spinor array is silently
treated as the physical marked frame. Normal transport and the
local-twistor exterior-square identification therefore preserve the same
calibrated infinity pair of Proposition I.3. ∎

The scalar factors distinguish the two native tensors. The
evaluated-form tensor has
\(\operatorname{Scal}(G_0)=24\mu_{\rm src}/\varrho\) and
\(\Delta_{G_0}\varrho=4\mu_{\rm src}\). The literal-contour tensor has
half those quantities. Indeed eq.~\ref{eq:sh-r1-general-source-metric}
is \(-\varrho^{-2}\) times the complex Schwarzschild tensor of formal
mass \(-2\mu_{\rm src}\), and the radial scalar and Laplacian
calculation of Proposition K.2 applies with that parameter and the
displayed metric constant. Hence
\begin{equation}\protect\phantomsection\label{eq:sh-r1-calibrated-scalar}{
 \operatorname{Scal}(G_{\rm int})=\frac{12\mu_{\rm src}}{\varrho},
 \qquad \operatorname{Scal}(h_*)=-\frac{6m}{r},\qquad
 \sigma_*=-\frac{\sqrt2}{6m}\operatorname{Scal}(h_*).
}\end{equation} Theorem K.15 is restricted to \(U\) and its image with
\(u_s\ne0\) and the specified angular, root and inverse branches. It
does not extend the primitive comparison over the full spatial angular
sphere of the finite core in Section 10. Complete complex point skies
arise instead from the selected metric's null directions. The finite
standard-metric atlas and its whole skies in K.2 remain independent of
any extension of the primitive across \(u_s=0\).

\protect\phantomsection\label{prop:schw-bc-native-generated-composite}

\textbf{Enriched comparison on the same local source chart.} Use the
primitive recipe with source parameter \(\mu_{\rm src}=m/2>0\), the
continued sheet \(\varrho=-r\), the selected pole and contour
orientation, and the local angular and root branches of the primitive
tensor calculation. On that actual source chart put \[
 \sigma_{\rm src}=-\frac{\sqrt2}{\varrho},\qquad
 g_{\rm src}=\sigma_{\rm src}^{-2}G_{\rm int}.
\] Let \(\Phi_{\rm pr}\) be the physical-to-source map
\begin{equation}\protect\phantomsection\label{eq:schw-bc-native-chart-map}{
 \begin{aligned}
 \tau&=\mathrm i(r-v),\quad x=r\sin\theta\cos\phi,\\
 y&=r\sin\theta\sin\phi,\quad z=r\cos\theta,\qquad\varrho=-r.
 \end{aligned}
}\end{equation} On a neighborhood where its stated inverse branches,
\(u_{\rm src}=x+\mathrm iy\ne0\), \(q\), \(\varrho\) and the retained
roots are defined, the primitive tensor identity is
\(\Phi_{\rm pr}^*G_{\rm int}=2r^{-2}g_m\). It follows that
\(\Phi_{\rm pr}^*g_{\rm src}=g_m\) with the calibrated density, point
and full-spin lift. If \(F:g[A_g]\to(g_m,o,s)\) is the initialized
metric comparison of Theorem K.7, the completed native comparison and
inverse are
\begin{equation}\protect\phantomsection\label{eq:schw-bc-native-generated-composite}{
 F_{\rm src}=\Phi_{\rm pr}\circ F,\qquad
 F_{\rm src}^{-1}=F^{-1}\circ\Phi_{\rm pr}^{-1}.
}\end{equation} They induce the actual enriched comparisons of that
theorem on common smaller representatives.

The map has Jacobian \(-\mathrm ir^2\sin\theta\) and the explicit local
inverse \(r=-\varrho\), \(v=r+\mathrm i\tau\), with the retained inverse
angular branches. It is nonsingular at the chosen horizon point.
Pullback of \(\sigma_{\rm src}\) is \(\sqrt2/r\), so the selected
physical metric is exactly \(g_m\), including its constant. The positive
central density-root value and the specified full-spin lift fix these
comparisons; the source's principal spinor arrays are not directly
identified with a marked orthonormal frame.

The pullback of the Levi--Civita connection of \(g_{\rm src}\) is
torsion free and preserves \(g_m\), so it is its Levi--Civita
connection. Inductively, \[
 \Phi_{\rm pr}^*((\nabla^{g_{\rm src}})^jR^{g_{\rm src}})
       =(\nabla^{g_m})^jR^{g_m}\qquad(j\ge0).
\] The transported full-spin frame therefore intertwines the exact Weyl
contractions and averaged spin symmetrizations at every order. The
source records are the same complete records used to define \(A_g\).
Normal uniqueness gives \(F\) and its inverse. Composition now gives the
displayed two-sided inverse identities on the common restricted source
and target domains.

The induced null-cotangent map preserves the tautological form and null
characteristics, so it gives the ray contact maps and their inverse hit
maps. Its scalar optical solution map preserves value, slope, Wronskian,
density-square calibration and the selected affine evaluation. The
normal tractor and full-spin local-twistor maps preserve \(I,X,DX\),
optical flags, point planes and both infinity forms. H.4 and Theorem J.6
then supply the projective operator and whole compact point-module maps,
including the two-associator and framed-module equations
eq.~\ref{eq:schw-bc-compatible-enriched-comparisons}. G.5, Lemma F.8 and
framed-Hom uniqueness give their finite inverse and composition
coherence. These enrichments are canonically constructed from the
selected native metric; an operator presentation of the original
coincidence surface is not assumed as input. ∎

This native comparison is restricted to the primitive recipe's actual
angular/root chart and its common normal representatives. It is not a
native atlas over all spacetime angular directions in the \(S^2\) factor
of the standard core \(K\). A complete compact point sky
\(\mathbb P^1\times\mathbb P^1\) consists of complex null directions at
one base point; it is independent of that spacetime angular \(S^2\).
Complete-sky comparisons over this local base chart therefore do not
supply a full-angular native gluing theorem. The standard metric
finite-cover argument above still covers the original \(K\) with
finitely many regular angular charts. No identification of an
independently selected ordinary background with the fixed SDTN
background, direct \(Q\to\beta_Q\) map, global ray quotient or global
operator ring is used.

\protect\phantomsection\label{lem:schw-bc-compatible-local-amplitude}

\subsubsection{K.6. Compatible physical amplitude
families}\label{k.6.-compatible-physical-amplitude-families}

\pandocstatementheading{Lemma K.16 (Compatible local physical family).}

Fix a marked analytic Einstein germ with numerical constant zero and
complete records \(U\), a compact amplitude disk
\(C_\varepsilon=\{|\varepsilon|\le T_b\}\) with \(T_b>1\) lying strictly
inside a retained larger disk, and a finite list of physical comparison
operations and their required inverses and composites. There are
positive common point and passive-displacement neighborhoods on which
all these operations have actual representatives over \(C_\varepsilon\).
Every vertex is obtained by rebasing the same central metric family. The
source construction at a vertex uses a separate internal amplitude.
These representatives retain the complete compact skies, full-spin
labels, both root calibrations, actual operator associators and framed
point-module maps.

\emph{Proof.} Use the complete source family
\begin{equation}\protect\phantomsection\label{eq:schw-bc-central-amplitude-family}{
 U_\varepsilon=(\varepsilon U^+,\varepsilon U^-;0),\qquad
 g_\varepsilon=g[U_\varepsilon].
}\end{equation} Propositions 6.1--6.2 give its marked normal geometry,
finite derivatives, normal inverses and initialized spin transport on
common smaller parents, uniformly on the entire compact disk. At this
stage no final displaced marks or final affine interval have been
chosen.

Let finitely many passive parameters \((a_i,A_i)\) vary in a small
neighborhood of \((0,1)\) and put
\begin{equation}\protect\phantomsection\label{eq:schw-bc-compatible-rebased-records}{
 \begin{aligned}
 x_i(\varepsilon)&=\operatorname{Exp}^{g_\varepsilon}_o(a_i),\\
 e_i(\varepsilon)&=\nu_i(\varepsilon)A_i,\\
 V_i(\varepsilon)&=R_{a_i,A_i}(U_\varepsilon).
 \end{aligned}
}\end{equation} Here \(\nu_i\) is the initialized full-spin transport
along the radial geodesic. Theorem 8.3 identifies \(V_i(\varepsilon)\)
with the entire curvature record of this same metric at
\(x_i(\varepsilon)\) in frame \(e_i(\varepsilon)\). The record-flow
estimates and the null-geodesic/parallel-spin Cauchy argument of Lemma
K.6 give a uniform factorial bound on a smaller record radius. Thus the
exact encoders of Theorem 3.1 make \(V_i\) bounded holomorphic passive
source families.

The initialized metric maps are
\begin{equation}\protect\phantomsection\label{eq:schw-bc-compatible-metric-maps}{
 \Xi_i^\varepsilon:g[V_i(\varepsilon)]\longrightarrow g_\varepsilon,
 \qquad
 y\longmapsto\operatorname{Exp}^{g_\varepsilon}_{x_i(\varepsilon)}
                                      (e_i(\varepsilon)y).
}\end{equation} At \(a_i=0,A_i=1\) these maps are the identity in the
marked normal representatives for every
\(\varepsilon\in C_\varepsilon\). The finite endpoint equations, their
inverse equations and the specified finite composites are jointly
holomorphic. Their central determinants and the scalar coefficients
actually used as units are nonzero. Compactness of \(C_\varepsilon\)
therefore permits one sufficiently small passive neighborhood, then
smaller compact passive subcores, such that all maps, inverses and
stipulated composite images retain positive parent buffers and bounded
inverse matrices. Only scalar line-transition coefficients that are
inverted, or frame determinants, require positive modulus minima;
individual entries of an invertible frame matrix may vanish.

On these common representatives define the physical map from vertex
\(j\) to vertex \(i\) by \[
 f_{ij}^\varepsilon=(\Xi_i^\varepsilon)^{-1}\Xi_j^\varepsilon.
\] It has actual inverse \(f_{ji}^\varepsilon\) and satisfies
\(f_{ij}^\varepsilon f_{jk}^\varepsilon=f_{ik}^\varepsilon\). Its
full-spin lift uses the same initialized maps. Expressed as source
instructions, composition includes the radial-versus-broken-path
correction \(K_{U_\varepsilon}\) of Theorem 8.3. These chart identities
do not assert path independence of spin parallel transport.

The nonlinear source action must be retained:
\begin{equation}\protect\phantomsection\label{eq:schw-bc-amplitude-rebase-distinction}{
 V_i(\varepsilon)=R_{a_i,A_i}(\varepsilon U;0),\qquad
 V_i(\varepsilon)\ne\varepsilon R_{a_i,A_i}(U;0)
 \quad\hbox{in general}.
}\end{equation} For the prescribed construction at vertex \(i\), regard
\(V_i(\varepsilon)\), together with the passive displacement and frame
parameters, as the bounded input family. Apply §§6 and F--G using the
independent internal source amplitude \[
 (s_{\rm in}V_i^+(\varepsilon),s_{\rm in}V_i^-(\varepsilon);0),
\] and set \(s_{\rm in}=1\) after construction. The physical variable
\(\varepsilon\) is still retained. There is no assertion that these
vertex metrics are physically compared for arbitrary \(s_{\rm in}\) by
scaling the rebased records linearly. Theorem G.5 uses exactly this
independent-internal-amplitude evaluation and quantizes the inverse
pullback of the ray map of \(f_{ij}^\varepsilon\). Its reference at
\(\varepsilon=0\) is the flat translation/full-spin action, with its
half-density calibration.

The resulting chart algebra maps, normalized comparison units and framed
module maps obey
\begin{equation}\protect\phantomsection\label{eq:schw-bc-compatible-enriched-comparisons}{
 \begin{aligned}
 k_i\varphi^s_{ij}
    &=\operatorname{Ad}(T_{ij})\varphi^t_{ij}k_j,\\
 T_{ij}\varphi^t_{ij}(T_{jk})C^t_{ijk}
    &=k_i(C^s_{ijk})T_{ik},\\
 F_i\tau^s_{ij}&=T_{ij}\tau^t_{ij}F_j.
 \end{aligned}
}\end{equation} Both original associators, the optical density-square
identification, the module half-canonical root with any base factor, and
the entire parameter framing are retained. G.5 and H.4 give the
normalized comparisons between selected composites and composites of the
selected lifts. Lemma F.8 and framed-Hom uniqueness imply the identity,
inverse and all specified finite coherence equations, over the same
compact physical amplitude and passive parameter sets.

For clarity, the complete-sky restriction is uniform in these
parameters. Let \(\mathcal F\to B\times P\) be the proper sky incidence
over the common four-dimensional normal base \(B\) and a parameter
parent \(P\) containing the retained compact parameter set \(C\). Let
\(\mathcal U\subset\mathcal F\) be the open locus where all stipulated
direction charts, hit maps and finite composite conditions have their
required parents. The preceding finite restrictions make it contain
every central whole sky over \(\{0\}\times C\). Properness makes the
image of \(\mathcal F\setminus\mathcal U\) closed locally in
\(B\times P\). Its projection over the compact set \(C\) is closed in
\(B\) and avoids \(0\). A smaller positive base neighborhood therefore
retains all the required chart coverage for every parameter in \(C\).
Equivalently one may take a finite subcover of the compact
central-sky/parameter set. This supplies a common base neighborhood,
rather than intersecting infinitely many unrelated neighborhoods without
a compactness argument. H.4's finite closed-star refinement restricts
all root, operator and module data together over that same set. This
proves the lemma. ∎

There is a separate, stronger common-family assertion on a sufficiently
short subsystem about the horizon mark \(o\). Use the single
\(U_\varepsilon=\varepsilon U(o)\) and the positive displacement radius
\(a_{\rm pass}\) and base parents already obtained above. Also fix the
central metric-readout parent and its smaller \(C^3\) core if an
amplitude error estimate is to be used. Let \(\chi\) be the exact
amplitude-one normal-to-EF coordinate map, fixed independently of
amplitude. Choose a dimensionless EF polydisc of radius \(r_c>0\) whose
doubled polydisc is inside the geometric parents, the required base
comparison parents and the dimensionless image of that \(C^3\) core.
This intersects four-dimensional base parents; it does not map a
five-dimensional contact chart into spacetime.

For the nonradial affine geodesic initialized at \(o\), let
\(t=\lambda/m\) and let \(k_0=d\gamma/dt|_0\) in the fixed normal
tangent norm. Keep a finite coordinate-speed bound \(V_\infty\) on the
exact short geometric ray. After a further positive restriction of
\(r_c\), choose
\begin{equation}\protect\phantomsection\label{eq:schw-bc-common-subsystem-length}{
 0<d\le\min\left\{
 d_{\rm geom},\frac{r_c}{8\max(1,V_\infty)},
 \frac{a_{\rm pass}}{2\|k_0\|}
 \right\}.
}\end{equation} Here \(d_{\rm geom}>0\) is an independently established
finite geodesic control length and \(k_0\ne0\). Only now specialize
\(a_i=t_i k_0\) for \(t_i=-d,0,d\), and take \(A_i=1\) in the parallel
full-spin frames, or the fixed bounded initialized changes to the
desired final frames. Since \(\gamma\) is affine through \(o\),
\(\operatorname{Exp}_o(t_i k_0)=\gamma(t_i)\) on the retained
representatives. The choices satisfy both the passive displacement
margin and the fixed base/core inclusion. Transversality at the horizon
gives an exterior point, the horizon point and an interior point for
every sufficiently small positive \(d\).

For the ray of §K.7 and its speed bound \(V_\infty=2\), one may
equivalently follow its order of choices: first require
\(r_c\|k_0\|\le8a_{\rm pass}\), then put \(d=\min\{1/8,r_c/16\}\). The
inequalities \(d\|k_0\|\le a_{\rm pass}/2\) and \(V_\infty d\le r_c/8\)
follow. The length is positive but generally unevaluated. Its optical
formulas use \(h=2d\); numerical constants derived specifically for
\(d=1/8\) do not become constants for an unevaluated common source
radius.

Thus the quantifiers are: fix the bounded record family and the compact
amplitude/parameter sets; choose common parents, passive displacement
and inverse/root margins; choose the common base and readout core; then
choose the positive short length and final marks. The optical amplitude
truncation on this subsystem remains the single central Taylor
polynomial in \(\varepsilon\), read out before evaluation at one. It
neither requires linearity of the rebased records in \(\varepsilon\) nor
asserts an exact cocycle for independently truncated vertex metrics.

\protect\phantomsection\label{sec:sch-r1-rational-optics}

\subsubsection{K.7. A separately initialized rational optical
example}\label{k.7.-a-separately-initialized-rational-optical-example}

The family of §K.3 retains its entrance point, affine clock and
polar-first screen. Here we choose a different ray and a horizon
initialization. Let the physical energy be \(E=1\), angular momentum
\(L=m\), and physical affine parameter \(\lambda\). Put \(t=\lambda/m\),
\(\rho=r/m\), \(V=v/m\), and initialize
\((V,\rho,\theta,\phi)=(0,2,\pi/2,0)\) at \(t=0\). The physical point is
\(o=(0,2m,\pi/2,0)\), also the mark for the source and readout in this
example. All primes in this subsection denote \(d/dt\). Then
\begin{equation}\protect\phantomsection\label{eq:sch-r1-central-ray}{\begin{gathered}
\rho'=-w, \\
w=\sqrt{1-(\rho-2)/\rho^3}, \\
V'=\frac1{\rho^2(1+w)}, \\
\phi'=\rho^{-2}, \\
\theta'=0.
\end{gathered}}\end{equation} These are the regular null first integrals, with the
positive square root. Differentiating gives
\(\rho''=\rho^{-3}-3\rho^{-4}\) and the geodesic equation, so this is
the stated affine clock. On \(7/4\le\rho\le9/4\) one has \[
 \frac{713}{729}\le w^2\le\frac{359}{343},\qquad
 \frac{15}{16}<w<\frac{17}{16}.
\] A first exit from that radial interval takes at least \(4/17>1/8\)
units of \(t\). Hence the ray exists for \(|t|\le1/8\), with
\(|\rho-2|\le17/128\) and radial endpoint margins \(15/128\). The
inequalities \(0<V'\le256/1519\) and \(0<\phi'\le16/49\) give
\(|V|\le32/1519\) and \(|\phi|\le2/49\). Thus this is a finite nonradial
future-horizon crossing in a regular angular chart.

Write \(g_m=m^2\bar g\) in these dimensionless coordinates and
\(K=d\gamma/dt\). A parallel unit screen for \(\bar g\), in the order
in-plane then polar, is
\begin{equation}\protect\phantomsection\label{eq:sch-r1-central-screen}{\begin{gathered}
E_2=\rho^{-1}\partial_\theta, \\
E_1=-\rho(1+w)\partial_\rho+\rho^{-1}\partial_\phi+A_0K, \\
A_0(t)=\int_0^t(1+w(s))\,ds.
\end{gathered}}\end{equation} Indeed the first two terms of \(E_1\) have covariant
derivative \(-(1+w)K\), while \(\nabla_KE_2=0\); addition of \(A_0K\)
makes the representatives parallel. Their Gram matrix is \(-I_2\) and
they are orthogonal to \(K\). Physical unit vectors are
\(e_i=m^{-1}E_i\). Rotating the orbital plane gives the polar Jacobi
solutions \(\rho\cos\phi\), \(\rho\sin\phi\), whose Wronskian is
\(\rho^2\phi'=1\). They satisfy \(j_2''=-3\rho^{-5}j_2\). Reflection
makes the off-diagonal tidal entries zero and Ricci flatness makes their
trace zero. Thus
\begin{equation}\protect\phantomsection\label{eq:sch-r1-central-tide}{\begin{gathered}
j''+\mathcal Rj=0, \\
\mathcal R=\operatorname{diag}(-Q,Q), \\
Q=3\rho^{-5}, \\
\frac{3072}{59049}\le Q\le\frac{3072}{16807}<\frac15.
\end{gathered}}\end{equation} The screen order is reversed from §K.3, and \(L/m=1\)
lies outside that earlier family. Consequently none of the following
numerical values changes its initialization or endpoint bounds.

\pandocstatementheading{\protect\phantomsection\label{ex:sch-r1-rational-observable}Example K.17.}

For any \(0<d\le1/8\), put \(a=-d\), \(b=d\), \(h=2d\). Let \(F(t,s)\)
be the full transfer of \((j,j')\), initialized by \(F(s,s)=I_4\), and
let \(B\) be its upper-right block. The observable
\(\mathcal O=B_{11}(b,a)\) obeys
\begin{equation}\protect\phantomsection\label{eq:sch-r1-rational-general}{
 0<\mathcal O-\left(h+\frac{h^3}{64}\right)<\frac{79h^5}{9000}.
}\end{equation} In particular, on the full exactly controlled interval
\(d=1/8\),
\begin{equation}\protect\phantomsection\label{eq:sch-r1-rational-enclosure}{
 \frac{1025}{4096}<\mathcal O<\frac{1025}{4096}+\frac{79}{9216000},
 \qquad \left|\mathcal O-\frac{1025}{4096}\right|<10^{-5}.
}\end{equation} \emph{Proof.} The phase generator is
\(H=\left(\begin{smallmatrix}0&I_2\\-\mathcal R&0\end{smallmatrix}\right)\),
with \(\|H\|_2=1\) and \(H^TJ_0+J_0H=0\). Differentiation therefore
proves symplecticity, \(\|F(t,s)\|\le e^{|t-s|}\) and exact finite
composition. The defocusing component \(u=B_{11}(t,a)\) satisfies \[
 u(t)=t-a+\int_a^t(t-s)Q(s)u(s)\,ds.
\] Its positive Volterra series gives \(u(t)\ge t-a\). The first iterate
is \(\mathcal O_1=h+\int_{-d}^d(d^2-s^2)Q(s)\,ds\). Direct
differentiation of the ray equations gives \[
 Q'=15w\rho^{-6},\qquad
 Q''=15(3-\rho)\rho^{-10}+90w^2\rho^{-7}.
\] On the retained radial parent, \(\rho^{-7}<1/50\),
\(\rho^{-10}<1/250\), \(w^2<17/16\) and \(0<3-\rho\le5/4\). Hence
\(0<Q''<159/80<2\). Taylor's integral formula yields
\(0<Q(s)-Q(0)-Q'(0)s<s^2\) for \(s\ne0\). The odd term integrates to
zero against \(d^2-s^2\), and \(Q(0)=3/32\), so \[
 0<\mathcal O_1-\left(h+h^3/64\right)<h^5/120.
\] The full transfer bound gives \(u(s)\le(s-a)e^h\). Substitution in
its Volterra equation then gives \(u(s)-(s-a)\le Q_{\max}e^h(s-a)^3/6\).
Substituting this once more gives \[
 0<\mathcal O-\mathcal O_1
 \le Q_{\max}^2 e^h h^5/120<4h^5/(75\cdot120).
\] Addition gives \(79h^5/9000\), and \(h=1/4\) gives the rational
enclosure. This is an analytic interval bound, without numerical ODE
integration. ∎

For \(0<h\le1/4\), the same Volterra calculation gives
\(\|B-hI\|<h/360\), while \(\|A-I\|<1/120\). Thus
\(\sigma_{\min}(B)>(359/360)h\) and \(\sigma_{\min}(A)>119/120\). These
are projection margins for this example, in addition to the full
symplectic transfer. With physical input slope \(dj/d\lambda\), its
position-response block is \(mB\).

The exact interval \(|t|\le1/8\) is independent of any source-domain
estimate. To apply a single central amplitude approximation, first
choose the common source, incidence and inverse domains in §K.6 and the
inner normal \(C^3\) core of K.9, denoted here by \(P_2\). Let \(\chi\)
be the exact initialized map from that normal core to EF coordinates.
Choose a doubled dimensionless EF polydisc inside \(\chi(P_2)\) and
inside all retained base neighborhoods, with positive radius
\(r_{\rm common}\). Require also
\(r_{\rm common}\|k_0\|\le8a_{\rm pass}\), where
\(k_0=d\chi_o^{-1}K(0)\) is expressed in the fixed normal tangent norm
for the passive displacement radius in §K.6. Set
\(d=\min\{1/8,r_{\rm common}/16\}\). The exact coordinate speed is less
than two, so the ray remains inside the polydisc with a positive buffer
and its three marks are admitted passive displacements. This asserts
only a strictly positive interval; it supplies no numerical lower bound
on \(r_{\rm common}\).

Write \(\epsilon_N=D_gT_*^{-N-1}/(1-T_*^{-1})\), where \(T_*>1\) and
\(D_g\) includes the independent spatial Cauchy gap and the fixed
coordinate change through order four. For normalized approximations with
the same central null tangent and screen, K.9 and the state-stability
argument of K.12, applied on this example's own tube with its displayed
velocity and screen bounds, give a finite tidal constant \(C_K\) such
that, under their inverse, state and crossing margins,
\begin{equation}\protect\phantomsection\label{eq:sch-r1-rational-amplitude-error}{
 \|F_N(b,a)-F(b,a)\|\le
 \mathcal E_N:=hC_K\epsilon_N e^{(1+C_K\epsilon_N)h},\qquad
 |\mathcal O_N-\mathcal O|\le\mathcal E_N.
}\end{equation} To check the initialization, the normal metric satisfies
\(g_\varepsilon(0)=\eta\) for every amplitude. Thus each coefficient
polynomial has the same central metric, and the exact central null and
screen data are valid for every \(g_N\); parallel transport is
integrated in both directions. The \(C^3\) tail bounds the Christoffel
difference in \(C^2\) and curvature difference in \(C^1\). On a compact
buffered state tube the integral equation gives a state error
\(C_Z\epsilon_N\), by Gronwall. Replacing each factor in
\(-\bar g(E_i,R(E_j,K)K)\) bounds the tidal difference by
\(C_K\epsilon_N\). Requiring
\(C_Z\epsilon_N<\min\{\eta_Z,15d/32,15/32\}\) retains the tube and
transverse crossing of the fixed surface \(r=2m\). Duhamel's formula
then proves the displayed transfer bound; the approximate tidal matrix
is symmetric for its own real metric and its transfer remains symplectic
even if \(g_N\) is not Einstein.

Combining the two independently specified approximations gives \[
 |\mathcal O_N-(h+h^3/64)|<\mathcal E_N+79h^5/9000.
\] Inverse-position bounds additionally require
\(\mathcal E_N< (359/720)h\). All margins can be met for sufficiently
large \(N\) because \(d\) and the gaps are fixed and positive. The
constants \(D_g,C_K,r_{\rm common}\) and a numerical sufficient \(N\)
have not been evaluated. The full-interval rational numbers above cannot
be assigned to this common source interval unless \(d=1/8\) has
independently been justified.

\protect\phantomsection\label{sec:sh-finite-readout-concatenation}

\subsubsection{K.8. Finite readouts along the complete crossing
family}\label{k.8.-finite-readouts-along-the-complete-crossing-family}

We now prove the finite-readout alternative of Proposition 10.5 on the
entire original family. The local inputs are K.6--K.9 and the explicit
calculations in K.3; local flows use K.10--K.12. The physical
comparisons stay in the original exact atlas, while the approximate
Cauchy evolution is defined explicitly below.

\protect\phantomsection\label{sec:fro-a-domains-readouts}

\textbf{Finite readout domains.}

Fix \(m>0\) and the ingoing Schwarzschild family of §K.3.2, with \[
 P=[1,2]\times[1/4,1/2],\qquad
 \lambda=(E,\ell),\qquad L=m\ell,\qquad t=u/m.
\] Use the dimensionless EF coordinates
\(x=(V,\rho,\theta,\phi)=(v/m,r/m,\theta,\phi)\). Write \(g=g_0=\bar g\)
for \(m^{-2}g_m\) after expressing \(g_m\) in these coordinates with its
Jacobian factors. This explicit fixed unit conversion leaves the
Levi--Civita connection unchanged. Every metric readout and tail norm in
this section is taken after this conversion. The selected physical
metric and the physical enriched atlas remain \(g_m\); their
interpretation undoes the same fixed conversion. All coordinate and
tensor norms are auxiliary Euclidean norms, and all constants may depend
on this fixed \(m\). Let \(x_\lambda(t)\) be the exact trajectory and
let \(\tau(\lambda)\) be its first time at \(\rho=3/2\). Thus
\begin{equation}\protect\phantomsection\label{eq:fro-a-duration}{
 \tau(\lambda)=\int_{3/2}^{5/2}
   \frac{\mathrm d\rho}
    {\sqrt{E^2-\ell^2(\rho-2)/\rho^3}},
 \qquad
 \frac{10}{21}<\tau(\lambda)<\frac{100}{99}.
}\end{equation} The parameter \(s=t/\tau(\lambda)\) will be used only to
choose a finite schedule. The metric, geodesic, parallel-frame and
Jacobi equations retain the affine clock \(t\).

\pandocstatementheading{\protect\phantomsection\label{thm:fro-a-finite-domain-schedule}Theorem K.18 (finite buffered readout schedule).}

The whole family from \(\rho=5/2\) to \(\rho=3/2\), with every
\(\lambda\in P\), has the following data, chosen before the amplitude
truncation degree \(N\).

There are finitely many exact local normal-to-EF maps \(\chi_i\) and
nested complex EF domains
\begin{equation}\protect\phantomsection\label{eq:fro-a-nested-domains}{
 V_i^0\Subset V_i^1\Subset V_i^2\Subset V_i^3,
}\end{equation} where \(V_i^2\) lies in a \(C^3\) metric-control core of
a local readout of Lemma K.9 and \(V_i^3\) lies in its fixed-coordinate
readout parent. The metric and all inverse units needed for that readout
have their retained bounds. There are compact parameter cores \(P_a\)
covering \(P\), larger compact parameter neighborhoods \(P_a^+\), a
single integer \(J\ge1\), and a choice \(i=i(a,j)\) for \(1\le j\le J\),
such that the whole exact enlarged slab
\begin{equation}\protect\phantomsection\label{eq:fro-a-enlarged-slab}{
 \left\{
  x_\lambda\bigl(\tau(\lambda)s\bigr):
   \lambda\in P_a^+,\quad
   \frac{j-3}{J}\le s\le\frac{j+2}{J}
 \right\}
 \Subset V_{i(a,j)}^0 .
}\end{equation} The extensions in this display are short extensions of
the exact family. The cuts for \(\lambda\in P_a\) are
\begin{equation}\protect\phantomsection\label{eq:fro-a-affine-cuts}{
 t_j(\lambda)=\tau(\lambda)\frac jJ,\qquad
 \frac{10}{21J}<t_j-t_{j-1}<\frac{100}{99J}.
}\end{equation} Order the finite parameter cores once and, for each
\(\lambda\in P\), use the least index \(a\) with \(\lambda\in P_a\) and
its assigned readout list. Every estimate is uniform over all cores; no
analytic dependence across these selection seams is required. Adjacent
readouts contain a common open time neighborhood of their exact node,
with a positive common position buffer. There are also compact enlarged
state tubes, uniform vector bounds and positive state buffers for the
exact tangent, null companion and parallel screen frame on every
assigned slab. All these choices are independent of \(N\).

\emph{Proof.} The radicand in eq.~\ref{eq:fro-a-duration} is bounded
below by \(124/125\) on the original compact parameter/radial rectangle.
It stays nonzero on a slightly larger real rectangle and on a complex
neighborhood of that rectangle. Fix the positive square-root branch
there. The quadratures of §K.3.2, followed by the inverse with
\(\partial t/\partial\rho=-1/q\), prove real-analytic dependence of
\(x_\lambda(t)\) and \(\tau(\lambda)\) on their arguments near the full
compact family.

Choose a compact parameter rectangle \(\widehat P\) with
\(P\subset\operatorname{int}\widehat P\), sufficiently small that this
branch and the same initial data remain defined. The regular EF parent
of §K.3.2 has positive radial, advanced-time and angular buffers around
the original trajectories. Existence and continuous dependence for the
exact ODE, compactness and the transverse radial derivative therefore
give \(\alpha>0\) such that \[
 \Gamma(\lambda,s)=x_\lambda(\tau(\lambda)s)
\] is real analytic on an open neighborhood of
\(\widehat X=\widehat P\times[-\alpha,1+\alpha]\), and its image remains
in a compact subcore of that regular EF parent. The companion and
screen-frame ODEs have the same property: their coefficients and initial
data are analytic, all original initial denominators have positive lower
bounds, and their solutions are bounded on this compact set. A small
further restriction of \(\widehat P\) and \(\alpha\) makes these choices
simultaneous. The sharp rational estimates of §K.3 apply on the original
\(P\) and its original intervals; the added extensions need only the
finite analytic bounds just obtained. The assertion does not ask one ray
to remain in the radial strip until the larger terminal time of another
ray.

We next construct the readout domains; this step does not assume that
one old readout contains any long segment. The compact image
\(\Gamma(\widehat X)\) lies inside \(r>0\) and inside a finite physical
core of Theorem K.8 chosen to contain a small neighborhood of this
image. Such a core can be chosen with radial endpoints strictly outside
\([3m/2,5m/2]\), advanced-time endpoints strictly outside the displayed
trajectory range and the whole angular sphere. For the original family
choose \(m/2<\delta<2m\), \(v_-<0\) and \(v_+>25m/396\) in that core
definition; take the parameter and time extensions small enough to
retain strict buffers inside it. Its finite angular charts include the
regular strip used here.

At each point of this compact image, choose a marked full-spin frame in
one of the finite analytic frame parents. Lemma K.6 extracts a bounded
all-record family on a neighborhood of that mark. Propositions 6.1--6.2
and Theorem 6.3 supply the prescribed stored representative over an
amplitude disk of radius greater than one, on a sufficiently small
spatial neighborhood. Lemma K.9 supplies its holomorphic physical metric
readout, bounded inverse units, initialized normal development and
independent spatial gaps. This is applicable to these actual bounded
original inputs; the amplitude parent is furnished by §§6.1--6.3, rather
than added as an assumption about the long tube. The required
restrictions decrease space and retain the compact amplitude disk.

At amplitude one, Lemma K.6 and the initialized normal metric comparison
preceding Theorem K.7 give an exact normal-to-EF map \(\chi\) with
analytic inverse \(\kappa\). Choose sufficiently small concentric
complex EF polydiscs around the marked point. Four nested such polydiscs
can be placed inside the physical image of the normal readout parent, so
that the second larger polydisc \(V^2\) and its closure lie in the image
of the normal \(C^3\) control core and the largest \(V^3\) has a
positive gap to the readout boundary. Reserve any further spatial Cauchy
gap needed below before making this choice. The real locus of the
smallest \(V^0\) contains the marked point.

These local choices may also be made smaller than the local
comparison-star radii in Theorem K.8. Taking a finite subcover of the
compact image gives finitely many nested domains as in
eq.~\ref{eq:fro-a-nested-domains}, with all their exact maps,
inverse-unit margins and derivative bounds fixed. In particular the real
open sets \[
 O_i=\Gamma^{-1}(V_i^0)\cap\widehat X
\] cover \(\widehat X\). Give \(\widehat X\) the maximum product metric.
By the Lebesgue-number lemma there is \(d>0\) such that every subset of
\(\widehat X\) of diameter less than \(d\) lies in some \(O_i\). This
use of compactness is on the parameter-time product, after the readout
domains have been fixed.

Let \(\Delta_P=\operatorname{dist}_\infty(P,\partial\widehat P)>0\).
Choose
\begin{equation}\protect\phantomsection\label{eq:fro-a-product-mesh}{
 0<r<\min\{d/16,\Delta_P/4\},
 \qquad
 \frac1J<\min\{d/10,\alpha/3\}.
}\end{equation} Choose finitely many \(\lambda_a\in P\) whose open
maximum-norm balls of radius \(r/2\) cover \(P\), and put \[
 P_a=P\cap\overline B_\infty(\lambda_a,r),\qquad
 P_a^+=\overline B_\infty(\lambda_a,2r)
                              \subset\operatorname{int}\widehat P.
\] The core \(P_a\) has a relative open neighborhood contained in
\(P_a^+\), and the \(P_a\) cover \(P\). For each \(a,j\) consider the
compact product \[
 C_{aj}=P_a^+\times
          \left[\frac{j-3}{J},\frac{j+2}{J}\right].
\] It is contained in \(\widehat X\): its time endpoints lie between
\(-2/J\) and \(1+2/J\). Its product diameter is at most
\(\max\{4r,5/J\}<d\). The Lebesgue property therefore assigns a single
\(i(a,j)\) with \(C_{aj}\subset O_{i(a,j)}\). Its image is compact in
\(V_{i(a,j)}^0\), which proves eq.~\ref{eq:fro-a-enlarged-slab}. This
supplies an actual common schedule for the whole slab at every parameter
in a parameter core, rather than a different unspecified cover of each
individual ray. Equation eq.~\ref{eq:fro-a-affine-cuts} follows from
eq.~\ref{eq:fro-a-duration}.

For adjacent cells, both enlarged intervals contain \(|s-j/J|\le1/J\).
Thus their exact images over \(P_a^+\) lie in both assigned readouts on
this whole node neighborhood. In particular the images over the smaller
compact set \(P_a\times\{|s-j/J|\le1/(2J)\}\) are compact in the
intersection of the two smallest domains. Their distances to the
complements of the two \(V^1\) domains have a positive minimum. There
are only finitely many such nodes and cells. Taking the minimum of their
positive margins gives a common position buffer. Compactness likewise
gives a margin between every enlarged slab image and the boundary of its
assigned \(V^1\).

For clarity, state buffers are chosen separately from these
four-dimensional point domains. Use the fixed EF coordinates, affine
clock \(t\), and the \(\bar g\)-normalized state
\((K,\mathcal N,E_1,E_2)\), where \(K=\mathrm dx/\mathrm dt\),
\(\bar g(K,\mathcal N)=1\) and \(\bar g(E_A,E_B)=-\delta_{AB}\). The
physical vectors are \((k,n,e_A)=(K,\mathcal N,E_A)/m\), with the same
coordinate interpretation; in particular \(k=\mathrm dx/\mathrm du\).
All four vectors are transported in their own exact normalized frame.
This is the fixed mass unit conversion, not a rescaling by
\(\tau(\lambda)\). Denote the joint exact state by \[
 Z_\lambda(t)=(x_\lambda(t),K_\lambda(t),
                       \mathcal N_\lambda(t),E_{1,\lambda}(t),E_{2,\lambda}(t)).
\] Its values on the finitely many enlarged slabs form compact sets.
Choose a position radius smaller than one quarter of the common position
buffer, and vector radii at most one. Shrink these radii and fix a
positive metric-perturbation tolerance \(\beta_g\) so that the
frame-normalization denominators and Gram-matrix neighborhoods used in
Lemma K.11 retain their strict margins for every state in the resulting
tubes and every symmetric metric within \(\beta_g\) of the exact metric
there. This is possible by joint continuity at the exact normalized
compact frame family; the normalizing formulas are analytic on an open
neighborhood of that family. A common smaller radius \(\eta_Z>0\) and a
common finite vector bound \(V\) exist by taking the finite minima and
maxima. The point projections of all these state tubes lie inside the
corresponding \(V^1\), and hence inside the prescribed \(C^3\) cores.
The radii, \(\beta_g\), \(V\), all domain gaps and the schedule have now
been fixed without choosing \(N\). ∎

\pandocstatementheading{\protect\phantomsection\label{prop:fro-a-finite-readout-tail}Proposition K.19 (local finite readouts and common-coordinate tails).}

For every readout retained in Theorem K.18 there is a radius \(T_i>1\)
and a constant \(D_i<\infty\) such that its degree-\(N\) finite readout
defines a holomorphic symmetric tensor \(g_{i,N}\) on \(V_i^2\), real on
the real slice, with
\begin{equation}\protect\phantomsection\label{eq:fro-a-local-tail}{
 \|g_{i,N}-g\|_{C^3(V_i^2)}
 \le\delta_{i,N}
 :=D_i\frac{T_i^{-N-1}}{1-T_i^{-1}}.
}\end{equation} Here the norm is the maximum of the tensor norms and the
multilinear derivative norms through order three. The finite readout is
performed over its own finite amplitude algebra, followed by
coefficient-polynomial evaluation at one. There is a common
\(T_*=\min_iT_i>1\) and a finite \(D_*\) with
\begin{equation}\protect\phantomsection\label{eq:fro-a-uniform-tail}{
 \max_i\delta_{i,N}\le D_*T_*^{-N-1}.
}\end{equation} \emph{Proof.} For each fixed original input at the
selected mark, use its own amplitude variable \(\varepsilon_i\), stored
analytic tuple \(Y_i\) and finite tuple \(X_{i,N}=q_{i,N}Y_i\), where
\(q_{i,N}\) is reduction modulo \(\varepsilon_i^{N+1}\). Choose nested
amplitude radii \(1<T_i<T_{c,i}<T_{b,i}\) before taking any truncation.
Lemma K.9 proves, including its finite inverse and initialized ODE
recursions, the coefficient identity
\begin{equation}\protect\phantomsection\label{eq:fro-a-finite-readout-identity}{
 \mathcal R_{i,N}(X_{i,N})
       =q_{i,N}h_i(\varepsilon_i),
 \qquad
 h_{i,N}=\operatorname{ev}^{\mathrm{poly}}_1
                             (q_{i,N}h_i).
}\end{equation} Here \(h_i\) is \(m^{-2}\) times the physical metric in
the initialized normal coordinates, and \(\mathcal R_{i,N}\) includes
this final fixed scalar multiplication. It commutes with finite
reduction and coefficient-polynomial evaluation; the bound \(H_i\) below
includes the factor \(m^{-2}\). This identity is an equality of
coefficient germs on their common representatives. Inverses use finite
nilpotent sums; compositions use only the finitely many amplitude
coefficients and derivatives needed at the retained degree; initialized
ODE coefficients satisfy unique finite linear inhomogeneous recursions.
These are the operations in the readout, so the identity includes its
nonlinear metric reconstruction and normal development. Evaluating the
stored tuple first and then applying a nonlinear readout would not give
this identity.

Write \(\kappa_i=\chi_i^{-1}\) for the exact amplitude-one EF-to-normal
map. It is fixed for all \(\varepsilon_i\) and \(N\). The initialized
metric isometry gives
\begin{equation}\protect\phantomsection\label{eq:fro-a-fixed-coordinate-readout}{
 g=\kappa_i^*h_i(1),\qquad
 g_{i,N}
   =\kappa_i^*h_{i,N}
   =\operatorname{ev}^{\mathrm{poly}}_1
                      q_{i,N}\bigl(\kappa_i^*h_i(\varepsilon_i)\bigr).
}\end{equation} The equality on the right holds because this fixed
pullback is linear in the tensor and independent of amplitude. Its
finite coordinate factors have their exact values; this is not an
approximate coordinate map recomputed from \(h_{i,N}\).

Keep a positive normal spatial gap \(\rho_i\) around
\(\kappa_i(\overline V_i^2)\) inside the normal readout parent. Lemma
K.9 and amplitude Cauchy estimates give a finite \(H_i\) with \[
 \|D^q(h_i(1)-h_{i,N})\|
 \le 4^q q!\rho_i^{-q}
                  \frac{H_iT_i^{-N-1}}{1-T_i^{-1}},
 \qquad 0\le q\le3,
\] on this retained normal core. The first spatial loss has already made
the nonlinear readout and its inverses bounded; this \(\rho_i\) is the
independent loss for differentiating its amplitude tail.

We record explicitly the coordinate cost. Put \[
 B_i=\max\left\{1,\sup_{\overline V_i^2}
                           \|D^q\kappa_i\|:1\le q\le4\right\}<\infty.
\] Finiteness follows from analyticity on \(V_i^3\) and the retained
gaps. The coordinate map \(\chi_i\) also has finite derivatives through
order four on the corresponding compact normal cores, with bounded
inverse differential. Either direction of the coordinate comparison is
therefore controlled. For a covariant two-tensor \(r\), differentiating
\((\kappa_i^*r)(x)(v,w)
=r(\kappa_i(x))(D\kappa_i(x)v,D\kappa_i(x)w)\) three times uses no
derivative of \(\kappa_i\) above order four. If
\(R=\max_{q\le3}\|D^qr\|\), the chain-rule bounds on
\(D^a(r\circ\kappa_i)\), for \(a=0,1,2,3\), are respectively \(R\),
\(RB_i\), \(2RB_i^2\) and \(5RB_i^3\). The Leibniz rule for these three
factors has coefficient sum \(3^j\le27\) at order \(j\le3\).
Consequently the safe bound
\begin{equation}\protect\phantomsection\label{eq:fro-a-coordinate-c4-cost}{
 \|\kappa_i^*r\|_{C^3(V_i^2)}
       \le135B_i^5
                    \max_{q\le3}\|D^qr\|
}\end{equation} holds in the stated operator norms. One may therefore
take \[
 D_i=135B_i^5H_i
              \max_{0\le q\le3}(4^q q!\rho_i^{-q}).
\] This proves eq.~\ref{eq:fro-a-local-tail}. The Lorentz reality
assertion in Lemma K.9 and the exact real coordinate map give real
symmetric coefficient tensors and hence real \(g_{i,N}\).

There are only finitely many readouts. Thus \(T_*=\min_iT_i>1\) and
\(D_*=\max_i D_i/(1-T_i^{-1})<\infty\) give
eq.~\ref{eq:fro-a-uniform-tail}. The constants and domains were fixed
before \(N\). In particular all strict metric inverse,
frame-normalization and state first-exit smallness conditions that use
finitely many of these constants can subsequently be imposed by one
sufficiently large \(N\). This last observation does not itself perform
the finite state matching or its error propagation. ∎

\pandocstatementheading{\protect\phantomsection\label{prop:fro-a-overlaps-exact-atlas}Proposition K.20 (overlap defect and the retained exact atlas).}

On every common EF overlap of two domains used in the schedule,
\begin{equation}\protect\phantomsection\label{eq:fro-a-overlap-defect}{
 \|g_{i,N}-g_{k,N}\|_{C^3(V_i^2\cap V_k^2)}
     \le\delta_{i,N}+\delta_{k,N}.
}\end{equation} At amplitude one the exact local metric comparisons,
full-spin lifts, contact and affine comparisons, structured operator
comparisons and compact point-module comparisons are those of the finite
physical atlas in Theorem K.8, with its actual finite coherence. They
can be retained on the compact node neighborhoods of Theorem K.18.
Independently truncated metrics need satisfy neither an exact metric
cocycle nor a common physical amplitude family.

\emph{Proof.} Both exact amplitude-one tensors in the fixed EF
coordinates equal \(g\). The triangle inequality applied to the two
differences in eq.~\ref{eq:fro-a-local-tail}, and to their derivatives,
proves eq.~\ref{eq:fro-a-overlap-defect}. There is no
coordinate-transition factor in this formula because the two tensors and
norms are already in the same EF coordinates. In normal coordinates the
equivalent assertion must retain the pullback by the exact transition
\(\kappa_k\chi_i\), and its coordinate bound; it is not a literal
difference of arrays from different coordinate systems.

The exact charts satisfy \(h_i(1)=\chi_i^*g\). Undoing the fixed unit
conversion gives the original physical metrics and initialized
isometries, with their original root and full-spin calibration. Their
coordinate transitions are the maps \(\chi_k^{-1}\chi_i\), and their
specified full-spin lifts are the changes of the same fixed spin bundle.
By Theorems K.7--K.8 they induce the retained contact, affine, infinity,
operator and compact point-module comparisons, with actual comparison
units and module maps. On a triple their comparisons are the existing
finite coherence comparisons, not separately chosen principal-symbol
identifications.

Only finitely many chart pairs and node neighborhoods occur in Theorem
K.18. Their exact compact node images lie in the interiors of both
assigned chart parents and retain positive buffers. Keeping these
compact cores and their whole complex skies, apply Theorem K.8's
simultaneous finite closed-star refinement to this finite list of maps,
inverses and stipulated finite composites. Its parent restrictions
retain the required compact sets and all sky fibers. Existing atlas maps
are restricted, not altered; an additional local normal representative
is compared through its initialized exact metric isometry as in Theorem
K.7. Thus the required exact old comparisons and their
root/full-spin/module calibrations remain available at every switching
node.

No identity between \(\varepsilon_iU_i\) and the record family at a
different mark is used. The \(\varepsilon_i\) are internal construction
parameters of the separate readouts; only their amplitude-one exact
metrics have been identified. Polynomial evaluation at one preserves
their approximation estimates but does not turn the overlap defect into
zero or supply a cocycle for the truncated metrics. The finite schedule
and node buffers make them suitable inputs for a separate finite
matching argument; they do not claim that argument for free. ∎

\protect\phantomsection\label{sec:fro-b-switches}

\textbf{Metric and Cauchy switches.}

\pandocstatementheading{\protect\phantomsection\label{lem:fro-b-cauchy-switch}Lemma K.21 (Metric and Cauchy switches).}

Use the finite buffered schedule and local readouts of Theorem K.18 and
Propositions K.19--K.20. On a fixed neighborhood of each exact framed
node, the formulas below define a locally analytic null-frame reset, a
metric-to-metric tangent isometry and an oriented screen identification.
Their state and phase errors are bounded by the incoming state error and
the two readout metric errors, with constants and positive margins
independent of the truncation degree. The prescribed covariant
Cauchy-jet rule is well defined on the screen quotient, preserves the
same affine-clock Wronskian, and has the symplectic phase matrix
\(\operatorname{diag}(R,R)\) with \(R\in SO(2)\). At exact reference
data this phase matrix is the identity. All previously accumulated phase
data are transported by that matrix.

\emph{Proof.} We give the normalization, quotient map and phase
calculation explicitly.

\textbf{Fixed coordinates, clocks and reference data.}

Fix \(m>0\) and common dimensionless EF coordinates
\(x=(V,\rho,\theta,\phi)\), where \(V=v/m\) and \(\rho=r/m\). Put \[
 g_0=\bar g=m^{-2}g_m,\qquad t=u/m.
\] The components of \(g_m\) here are expressed in these dimensionless
coordinates. Its connection is the connection of \(g_0\), since their
factor \(m^2\) is constant. The reference state is \[
 Z_*(\lambda,t)=(x_*,K_*,\mathcal N_*,E_{*,1},E_{*,2}),\qquad
 \dot x_*=K_*,
\] with dot denoting \(d/dt\). The vectors are parallel along the
reference ray and have Gram matrix \[
 Q=
 \begin{pmatrix}
 0&1&0&0\\1&0&0&0\\0&0&-1&0\\0&0&0&-1
 \end{pmatrix}.
\] In particular, \(g_0(K_*,\mathcal N_*)=1\). Choose the same global
polar-first parallel Sachs screen and null companion as in §K.3,
rescaled to this convention. The physical tangent, companion and unit
screen vectors are \(K_*/m,\mathcal N_*/m,E_{*,A}/m\). The parameter
\(\lambda\) ranges over the fixed compact family of Proposition 10.2.
Its exact dimensionless terminal time is \(\tau(\lambda)\).

These reference data are supplied by the already proved first integrals,
quadratures and explicit parallel Sachs screen of §K.3, rather than by
an unknown Jacobi transfer. For example, with \(\ell=L/m\), \[
 \tau(E,\ell)=\int_{3/2}^{5/2}
 \frac{d\rho}{\sqrt{E^2-(1-2/\rho)\ell^2/\rho^2}}.
\] The reference companion can also be recovered algebraically from that
screen. Put \(\widetilde n=-\partial_\rho/\dot V_*\), whose denominator
has the positive lower bound of §K.3, and set \[
 \widetilde{\mathcal N}=\widetilde n+
       \sum_Ag_0(\widetilde n,E_{*,A})E_{*,A},\qquad
 \mathcal N_*=\widetilde{\mathcal N}-\tfrac12g_0(\widetilde{\mathcal N},\widetilde{\mathcal N})K_*.
\] It is the unique null vector orthogonal to the screen with
\(g_0(K_*,\mathcal N_*)=1\). Differentiating these conditions, using the
parallel \(K_*,E_{*,A}\), proves that \(\mathcal N_*\) is parallel. Thus
specifying the reference columns or the cuts requires no optical
observable or Jacobi transfer and is not circular.

Fix a finite integer \(J\) and exact cut times
\(t_j(\lambda)=j\tau(\lambda)/J\). At each cut, the incoming and
outgoing coordinate parents have a common overlap containing a fixed
positive position buffer around \(x_*(\lambda,t_j)\). Retain the
finitely many interfaces over compact parameter cores with these uniform
margins. A fixed compact state neighborhood and its positive buffer are
also retained. Only an interface whose incoming approximate position
remains in this overlap is used.

Let \(g=g_{i,N}\) and \(h=g_{k,N}\) be the incoming and outgoing real
symmetric dimensionless metric readouts. They are measured in these same
coordinates. On their respective buffered position domains assume \[
 \|g_{i,N}-g_0\|_{C^3}\le\delta_{i,N},\qquad
 \|g_{k,N}-g_0\|_{C^3}\le\delta_{k,N},\qquad
 d_{ik,N}:=\delta_{i,N}+\delta_{k,N}.
\] Thus \(\|h-g\|_{C^0}\) on the overlap is at most \(d_{ik,N}\). The
switch below uses only this \(C^0\) difference; the \(C^3\) bounds are
used for the actual flows and tidal estimates between cuts. All
constants below may depend on the fixed mass, coordinate and frame
bounds, finite interface system and compact reference family, but not on
\(N\).

\textbf{The null-frame reset.}

Suppose the incoming state at the cut is
\(Z^-=(x,K,\mathcal N,E_1,E_2)\) and the incoming frame
\(F=(K,\mathcal N,E_1,E_2)\) has Gram matrix \(Q\) for \(g(x)\). Write
\(\mathcal E^-=\|Z^--Z_*(\lambda,t_j)\|\) in the maximum of the five
fixed coordinate-block norms. All reference vector components in this
comparison are taken at the exact cut.

At the approximate position \(x\) apply Lemma K.11 with the incoming
metric and its actual normalized frame as the reference, and with
\(h(x)\) as the perturbed metric. Its formulas are as follows. Put \[
 \begin{gathered}
q=h(K,K), \\
b_0=h(K,\mathcal N), \\
d=h(\mathcal N,\mathcal N), \\
a=-\frac{q}{b_0+\sqrt{b_0^2-dq}}, \\
K^+=K+a\mathcal N.
\end{gathered}
\] The square root is the positive branch near \(1\). This expression is
valid also when \(d=0\) and is the root tending to zero of
\(h(K+a\mathcal N,K+a\mathcal N)=0\). Define \[
 \begin{aligned}
b&=h(K^+,\mathcal N), \\
\mathcal N^+&=b^{-1}\left(\mathcal N-\frac{h(\mathcal N,\mathcal N)}{2b}K^+\right), \\
P_h(w)&=w-h(w,\mathcal N^+)K^+-h(w,K^+)\mathcal N^+, \\
W&=(P_h(E_1),P_h(E_2)), \\
H_{\rm temp}&=-W^ThW, \\
E_{\rm temp}&=W H_{\rm temp}^{-1/2}.
\end{aligned}
\] The last inverse square root is the positive symmetric one. The
formulas give \(h(K^+,K^+)=h(\mathcal N^+,\mathcal N^+)=0\),
\(h(K^+,\mathcal N^+)=1\), and \[
 E_{\rm temp}^ThE_{\rm temp}=-I_2,\qquad
 h(E_{{\rm temp},A},K^+)=h(E_{{\rm temp},A},\mathcal N^+)=0.
\] There are constants \(\epsilon_*>0\) and \(C_{\rm n}<\infty\), common
to all the finitely many interfaces, such that if
\(\mathcal E^-+d_{ik,N}<\epsilon_*\) these formulas are defined and \[
 \begin{gathered}
b_0>\tfrac12, \\
b>\tfrac12, \\
\|H_{\rm temp}-I_2\|<\tfrac12, \\
\|F_{\rm temp}-F\|\le C_{\rm n}d_{ik,N},
 \quad F_{\rm temp}=(K^+,\mathcal N^+,E_{\rm temp}).
\end{gathered}
\] Here and below a four-column frame norm is the induced Euclidean
operator norm, with fixed finite conversions to the state-block norm. To
prove the assertions uniformly, at \(h=g\) one has \(q=d=a=0\),
\(b_0=b=1\), \(H_{\rm temp}=I_2\) and \(F_{\rm temp}=F\). The displayed
formulas are real analytic on a fixed neighborhood of the compact
reference data. A smaller closed neighborhood keeps the denominators and
eigenvalues in the stated margins. Its derivative bounds are finite.
Applying the mean-value formula with the source frame held fixed gives
the last estimate. For instance, \(|q|\le\|h-g\|\|K\|^2\) and the
displayed root formula gives \(|a|\le2|q|\) once \(b_0>1/2\); no inverse
power of a metric error is used. This is precisely the uniform version
of Lemma K.11, now with the actual incoming frame as reference.

There is no change of the cut time or of the affine clock: the next
segment starts at the same \((x,t_j)\) with tangent \(K^+\). In
particular, \(K^+\) is generally different from \(K\), but it is not
multiplied by an additional affine-scale factor.

Define the full-frame map at this point by \[
 L=F_{\rm temp}F^{-1}.
\] It is an actual isometry from \((T_xM,g(x))\) to \((T_xM,h(x))\): \[
 L^ThL=g,\qquad LK=K^+,\qquad L\mathcal N=\mathcal N^+,\qquad LE_A=E_{{\rm temp},A}.
\] Moreover \(\|L-I_4\|\le C_Ld_{ik,N}\). Indeed \(F^{-1}=Q^{-1}F^Tg\),
so its norm is uniformly bounded by the retained metric and frame
bounds, and \(L-I_4=(F_{\rm temp}-F)F^{-1}\). On the small connected
neighborhood this map preserves the retained orientation and time
orientation. The path obtained by replacing \(h\) with \(g+s(h-g)\)
starts at the identity and specifies its near-identity full spin lift.
An existing central spin label is not discarded.

\textbf{The target screen and the recorded phase.}

Let \(\mathsf E_*=(E_{*,1}(\lambda,t_j),E_{*,2}(\lambda,t_j))\) be the
coordinate-component columns of the known exact parallel screen at the
cut. For this pointwise operation extend those columns locally as
constant coordinate vector fields. Evaluate them at the approximate
position \(x\), and set \[
 W_*=(P_h(E_{*,1}),P_h(E_{*,2})),\qquad
 H_*=-W_*^ThW_*,\qquad
 E^+=W_*H_*^{-1/2}.
\] The constant-coordinate extension is a specified use of the common EF
chart, not an assertion that these extended fields are parallel. After
the switch, the initial columns \(E^+\) are parallel transported by
\(h\) along the next approximate segment.

At the exact state and exact metric these formulas give
\(E^+=E_{\rm temp}=\mathsf E_*\). Compactness and a smaller common
\(\epsilon_*\) give \(\|H_*-I_2\|<1/2\) throughout the retained
neighborhood. Thus \(E^+\) is an oriented orthonormal basis of the same
two-plane \[
 \mathcal S^+=\{w:h(w,K^+)=h(w,\mathcal N^+)=0\}
\] as \(E_{\rm temp}\). Define the computed transition matrix \[
 R_{AB}=-h(E_A^+,E_{{\rm temp},B}).
\] Then \[
 E_{\rm temp}=E^+R,\qquad R^TR=I_2,\qquad \det R=1.
\] The first identity follows by orthonormal expansion in the common
negative-definite plane, and the second by taking its Gram matrix. The
sign of the determinant is positive at the reference data and remains so
on the retained neighborhood; equivalently require \(\|R-I_2\|<1\).
Hence \(R\in SO(2)\). Its near-identity spin lift, when retained, is
specified by the same continuous choice. The target full-spin frame is
the lifted temporary frame followed by the lift of
\(\operatorname{diag}(I_2,R^{-1})\), since
\(F^+=F_{\rm temp}\operatorname{diag}(I_2,R^{-1})\). Thus its central
label is transported from the source lift, not recovered from the target
frame columns alone.

There are common constants \(A_{\rm sw},B_{\rm sw},C_{\rm ph}\) such
that the reset state \(Z^+=(x,K^+,\mathcal N^+,E_1^+,E_2^+)\) satisfies
\[
 \begin{aligned}
 \mathcal E^+
   &\le A_{\rm sw}\mathcal E^-+B_{\rm sw}d_{ik,N},\\
 \eta_{ik,N}:=\|\operatorname{diag}(R,R)-I_4\|
   &\le C_{\rm ph}(\mathcal E^-+d_{ik,N}).
 \end{aligned}
\] For the first bound, the temporary frame differs from the incoming
frame by \(C_{\rm n}d_{ik,N}\). The target-screen formula is a fixed
smooth function of \(h(x),K^+,\mathcal N^+\) and the exact reference
columns. At the reference data its value is \(\mathsf E_*\). Also \[
 \|h(x)-g_0(x_*)\|
 \le\delta_{k,N}+\|Dg_0\|\,|x-x_*|.
\] The retained compact derivative bounds therefore give
\(\|E^+-\mathsf E_*\|\le C(\mathcal E^-+d_{ik,N})\), proving the state
bound after including the unchanged position. For the second bound, use
\[
 R-I_2=-(E^+)^Th(E_{\rm temp}-E^+)
\] and the same bounds for both screen matrices. All these constants are
suprema on a fixed compact neighborhood, chosen before \(N\). The number
of interfaces is finite, so taking their maxima and the smallest of
their positive margins preserves this independence.

\textbf{Covariant Jacobi jets and the phase switch.}

At a normalized state use the screen quotient \[
 \mathscr S(K)=K^{\perp_g}/\mathbb RK
\] with positive metric induced by \(-g\). The incoming \(E_A\) identify
it with \(\mathbb R^2\); the companion \(\mathcal N\) merely chooses
representatives. Along an actual metric-null geodesic, the transverse
Jacobi phase consists of the value class \([J]\) and covariant slope
class \([V]\), with \(V=\nabla_KJ\). Its phase coordinates are
\(z=(q,p)\), where \[
 [J]=[Eq],\qquad [V]=[Ep],\qquad p=\dot q
\] in the parallel screen. Actual Jacobi lifts with these classes may
have longitudinal components. The transverse equation and its Cauchy
phase do not depend on those components.

Every pair of screen classes has a full Jacobi solution with
representatives in \(K^\perp\) initially: solve the ordinary Jacobi
Cauchy problem with the chosen representatives of its two slots. Metric
compatibility and the curvature symmetries keep \(g(J,K)=g(V,K)=0\).
Changing the initial representatives by multiples of \(K\) adds a
longitudinal Jacobi field \((a+bt)K\) and hence leaves the transverse
solution unchanged.

Prescribe the interface map on covariant Cauchy jets by \[
 (J,V)^+=(LJ,LV).
\] Since \(L\) is an isometry and \(LK=K^+\), it maps \(K^{\perp_g}\) to
\((K^+)^{\perp_h}\) and \(\mathbb RK\) to \(\mathbb RK^+\). Consequently
the rule induces a well-defined isometry on the screen quotient for each
of the two jet slots. Its coordinate rule in the target screen is
exactly \[
 z^+=\mathsf C_{ik,N}z^-,
 \qquad
 \mathsf C_{ik,N}=\operatorname{diag}(R,R).
\] Indeed \(LEq=E_{\rm temp}q=E^+Rq\), and the same computation holds
for \(V\). Longitudinal representatives map to longitudinal
representatives, so they introduce no correction to this formula.

With
\(\mathbb J=\left(\begin{smallmatrix}0&I_2\\-I_2&0\end{smallmatrix}\right)\),
\[
 \mathsf C_{ik,N}^T\mathbb J\mathsf C_{ik,N}=\mathbb J.
\] Thus the switch is symplectic and preserves the transverse Jacobi
Wronskian. This is also the direct consequence of \(L^ThL=g\) applied to
\(-g(J,\widetilde V)+g(V,\widetilde J)\). There is no
derivative-of-\(R\) term: the rule specifies Cauchy data at one cut, and
the outgoing screen is initialized there and then parallel transported.
It is not a time-dependent rotation of a single screen along one
uninterrupted segment.

The coordinate derivative corresponding to this covariant jet reset is
explicitly \[
 \dot J^{+\,\mu}=(LV)^\mu
       -\Gamma(h)^\mu_{\alpha\beta}(K^+)^\alpha(LJ)^\beta.
\] Therefore copying a raw coordinate derivative across a metric switch
would generally give a different rule. The formula prescribes a valid
Jacobi Cauchy jet for the outgoing metric. It does not differentiate the
nonlinear null-frame reset map, nor does it assert
\(\nabla^h_{K^+}(LJ)=LV\) for a previously chosen field of maps \(L\)
away from this cut.

The reference switch is the identity: at \(Z^-=Z_*\) and \(g=h=g_0\) one
has \(a=0\), \(L=I_4\), \(E^+=E_{\rm temp}=\mathsf E_*\) and \(R=I_2\).
This identifies the exact comparison phase without a new screen gauge.
The rule is nevertheless nontrivial. If \(g=h=g_0\), \(K,\mathcal N\)
are exact and the incoming screen is \(\mathsf E_*R_0\) for a small
nonidentity \(R_0\in SO(2)\), then \(L=I_4\), the target screen is
\(\mathsf E_*\), and the recorded phase switch is
\(\operatorname{diag}(R_0,R_0)\). The physical jet is unchanged while
its coordinates are correctly carried into the target frame.

\textbf{Piecewise propagation and scope.}

On each subinterval integrate the null geodesic and parallel frame of
its own readout metric, together with the transverse Jacobi equation of
that metric. Its tidal matrix is real symmetric, so the corresponding
four-dimensional transfer \(U_{j,N}\) is symplectic. At a cut use the
state reset above and record the computed matrix \(\mathsf C_{j,N}\).
The full optical algorithm uses, in the stipulated order, \[
 U_{J-1,N}\mathsf C_{J-1,N}U_{J-2,N}
       \cdots \mathsf C_{1,N}U_{0,N}.
\] Every factor and hence the product is symplectic. Earlier frame and
phase discrepancies remain in the incoming state, in \(R\), and in this
accumulated product. Target-frame reinitialization retains
\(\mathsf C_{j,N}\) and the accumulated transfer.

The local state-flow estimates of Theorem K.12 apply on each retained
interval with its initial error. The two switch inequalities propagate
this error and the interface contributions through the finite iteration,
using the same affine clock.

In the fixed scalar position units used above, \(p=dq/dt=m\,dq/du\). If
physical unit-screen coefficients are used instead of coefficients in
\(E_A\), both \(q\) and \(p\) acquire the same constant factor \(m\),
which leaves the four-by-four transfer unchanged. Once the position
units are fixed, the conversion between \((q,dq/du)\) and \((q,dq/dt)\)
is the single global conjugation \(\operatorname{diag}(I_2,mI_2)\) of
§K.4. No additional clock conversion occurs at an interface.

The construction is a piecewise optical approximation with continuous
position. Each open segment is a null geodesic of its own readout
metric; the prescribed small interface resets join tangents, frames and
covariant Jacobi data. Its optical transfer uses the covariant Cauchy
prescription above. The exact cuts and parallel Sachs frame from §K.3,
together with the coordinate comparisons fixed in Theorem K.18, are
benchmark data independent of the transfer being approximated. The
bounds are relative to these inputs and the retained analytic margins. ∎

\protect\phantomsection\label{sec:fro-c-accumulation}

\textbf{Finite accumulation and optical observables.}

Use the finite readouts and reset maps just constructed. The reference
is the fixed positive-mass Schwarzschild family of §K.3 with signature
\((+,-,-,-)\). All comparisons use its original individual affine
intervals and exact terminal times.

\protect\phantomsection\label{sec:fr-opt-c-setting}

\textbf{Fixed finite data and local bounds.}

Write \(\lambda=(E,\ell)\in[1,2]\times[1/4,1/2]\) for the parameters of
the original nonradial family. Keep its physical affine parameter \(u\),
dimensionless time \(t=u/m\), entrance \(r=5m/2\) and exact final time
\(\tau_\lambda\), so that
\begin{equation}\protect\phantomsection\label{eq:fr-opt-c-time-bounds}{
 \tau_-:=\frac{10}{21}<\tau_\lambda<\tau_+:=\frac{100}{99}.
}\end{equation} Keep the polar-first parallel Sachs frame of §K.3 and a
parallel null companion. In the common dimensionless EF coordinates
\(X=(v/m,r/m,\theta,\phi)\), let
\begin{equation}\protect\phantomsection\label{eq:fr-opt-c-dimensionless-convention}{\begin{gathered}
g_0=\bar g=m^{-2}g_m, \\
Z_\lambda=(X,K,\mathcal N,E_1,E_2), \\
K=dX/dt, \\
g_0(K,\mathcal N)=1, \\
g_0(E_A,E_B)=-\delta_{AB}.
\end{gathered}}\end{equation} The displayed \(g_m\) includes the coordinate Jacobian
factors. The physical tangent, companion and unit screen vectors are
\((K,\mathcal N,E_1,E_2)/m\), and \(u=mt\) throughout. These states and
their prescribed coordinate descriptions range over a compact set. The
screen phase is a different space:
\(y=(\xi,p)\in\mathbb R^2\oplus\mathbb R^2\), with
\(p=d\xi/dt=m\,d\xi/du\). Using physical unit-screen position
coefficients instead of the coefficients in \(E_A\) multiplies both
phase slots by the same constant \(m\) and leaves the four-by-four
transfer unchanged. Fix one such position convention throughout.

The finite schedule has \(J\) slabs and nodes
\begin{equation}\protect\phantomsection\label{eq:fr-opt-c-schedule}{
 t_j(\lambda)=\frac{j\tau_\lambda}{J},\qquad
 h_j=t_{j+1}-t_j,\qquad 0\le j<J.
}\end{equation} The parameter set may have finitely many compact cells,
with a specified finite list of readouts on each cell. Every bound below
is a maximum over this finite list and its compact parameter sets.
Neither this list nor \(J\), the nodes, parents, exact coordinate
comparisons or reset rules depend on \(N\).

Here slab \(j=0,\ldots,J-1\) uses Theorem K.18's label \(i_j=i(a,j+1)\);
its nodes and enlarged domains are unchanged.

The exact times \(\tau_\lambda\) are the quadratures of §K.3.2, and the
reference Sachs frame is the explicit parallel frame of §K.3.3. They,
the specified companion and the exact readout-to-physical coordinate
comparisons are retained benchmark data. They are not values of the
unknown Jacobi transfer obtained in advance: the new \(U_N\) below is
obtained from the local approximate Jacobi equations and the specified
switch maps. The conclusion is an approximation theorem with these exact
inputs, not numerical certification from finite source coefficients
alone.

For readout \(i\), use Proposition K.19 and its finite-algebra readout
followed by coefficient-polynomial evaluation at one. In the fixed
dimensionless EF coordinates let \(g_{i,N}\) be the real tensor obtained
after the same constant rescaling by \(m^{-2}\). Its exact value at
amplitude one is the reference metric \(g_0\). The local result
supplies, on an inflated inner position domain with a positive
additional buffer,
\begin{equation}\protect\phantomsection\label{eq:fr-opt-c-local-metric-tail}{
 \|D^s(g_{i,N}-g_0)\|\le D_{i,s}\beta_N,
 \qquad 0\le s\le3,\qquad
 \beta_N=\frac{T_{\min}^{-N-1}}{1-T_{\min}^{-1}},
 \quad T_{\min}>1.
}\end{equation} All \(D_{i,s}\) are finite and independent of \(N\).
Exact coordinate/frame comparisons and the fixed metric factor
\(m^{-2}\) are included in these bounds before taking finite maxima.
Different \(g_{i,N}\) need not agree on overlaps. There is no use of a
single approximate spacetime metric or of a physical gluing of different
source amplitude families away from amplitude one.

On the retained compact state tubes fix uniform metric inverse,
connection, derivative and vector bounds. Choose an initial tail
threshold \(\bar\beta>0\) so that the local real-signature, inverse-unit
and null/frame normalizations hold whenever \(\beta_N\le\bar\beta\).
With \(D_0=\max_iD_{i,0}\), also take \(D_0\bar\beta<\beta_g\); this
condition is automatic if \(D_0=0\). These are possible because there
are finitely many strictly positive margins. The reference state is
interior to each slab and switch domain, and the position projection of
the larger state tube is inside that readout's inner \(C^3\) domain.

Here are the local state bounds used in the proof. Set
\(E_j^+=\|Z_N(t_j^+)-Z(t_j)\|\). Up to a possible first exit on slab
\(j\), the joint geodesic/parallel-frame equation gives
\begin{equation}\protect\phantomsection\label{eq:fr-opt-c-local-state-flow}{
 \sup_{t_j\le t\le t_{j+1}}
       \|Z_N(t)-Z(t)\|
 \le a_j E_j^+ +b_j\beta_N,
 \qquad a_j\ge1,\quad b_j\ge0.
}\end{equation} One can take, for example, \(a_j=e^{L_j\tau_+/J}\) and
\(b_j=e^{L_j\tau_+/J}(\tau_+/J)F_j\), where \(L_j\) is a Lipschitz bound
for the dimensionless-\(t\) state equation on the larger tube and
\(F_j\beta_N\) bounds the difference of its vector fields there. This
follows from the integral inequality with initial error \(E_j^+\), not
only from a zero-error initialization. At time \(t_{j+1}\) write
\(E_{j+1}^-\) for the incoming error.

The actual reset normalizes this incoming state with the target metric
and the same specified affine-clock rule. It uses the small null root,
null-companion normalization, positive Gram square root and the
prescribed oriented screen alignment. Its compact-neighborhood estimate
is
\begin{equation}\protect\phantomsection\label{eq:fr-opt-c-local-reset}{
 E_{j+1}^+\le r_{j+1}E_{j+1}^-+c_{j+1}\beta_N,
 \qquad r_{j+1}\ge1,\quad c_{j+1}\ge0.
}\end{equation} These constants exist for the actual finite-matrix
reset, rather than for a replacement by the exact state. Indeed its null
equation has nonzero derivative at the reference root; its scalar
denominators and Gram matrices have the retained nonzero or positive
margins. The small root and positive square root are analytic on a
common smaller neighborhood. The target screen is obtained by projecting
the globally parallel exact Sachs frame into the target screen and
normalizing its positive Gram matrix. At the exact metric and exact
state all these operations return that same reference state and frame.
On a fixed smaller compact neighborhood, suprema of the first
derivatives of this finite collection of maps give \(r_{j+1}\) and
\(c_{j+1}\) by the mean-value formula. The metric perturbations at the
reset are bounded by the source and target versions of
eq.~\ref{eq:fr-opt-c-local-metric-tail}. The incoming error is retained
in eq.~\ref{eq:fr-opt-c-local-reset}.

For example Lemma K.21's constants permit
\(r_{j+1}=\max\{1,A_{{\rm sw},j+1}\}\) and
\(c_{j+1}=B_{{\rm sw},j+1}(D_{i_j,0}+D_{i_{j+1},0})\) at an actual
switch. The final endpoint convention below has no target readout and
uses no such reset.

At the first entrance keep the full initialization estimate
\begin{equation}\protect\phantomsection\label{eq:fr-opt-c-initial-error}{
 E_0^+\le c_0\beta_N.
}\end{equation} For example \(c_0=C_{\rm init}D_{i_0,0}\), with any
fixed coordinate/frame constants included, is permitted by Lemma K.11.
This term is not set to zero merely because a later slab has a different
central mark.

\pandocstatementheading{\protect\phantomsection\label{prop:fr-opt-c-state-accumulation}Proposition K.22 (Finite state accumulation).}

There are a finite \(C_Z\) and an integer \(N_0\), depending only on the
fixed finite schedule, its readouts and buffers, such that every
\(N\ge N_0\) admits every slab evolution and every prescribed reset on
its full scheduled domain. Throughout every slab, and on both sides of
every reset,
\begin{equation}\protect\phantomsection\label{eq:fr-opt-c-global-state}{
 \|Z_N-Z\|\le C_Z\beta_N.
}\end{equation} \emph{Proof.} For the last endpoint put \(r_J=1,c_J=0\),
meaning that there is no reset there. For \(0\le j<J\) set
\begin{equation}\protect\phantomsection\label{eq:fr-opt-c-state-majorant-recursion}{
 A_j=r_{j+1}a_j,\qquad
 B_j=r_{j+1}b_j+c_{j+1},\qquad
 H_0=c_0,\qquad H_{j+1}=A_jH_j+B_j.
}\end{equation} The finite recurrence has the explicit solution
\begin{equation}\protect\phantomsection\label{eq:fr-opt-c-state-product-sum}{\begin{gathered}
H_j=c_0\prod_{\ell=0}^{j-1}A_\ell
     +\sum_{k=0}^{j-1}B_k
                    \prod_{\ell=k+1}^{j-1}A_\ell, \\
C_Z=\max\left\{\max_{0\le j\le J}H_j,
                  \max_{0\le j<J}(a_jH_j+b_j)\right\}.
\end{gathered}}\end{equation} Empty products equal one and empty sums zero. Thus all
constants are fixed before \(N\) is chosen. In particular this bound
includes \(c_0\) and the error carried through every switch; it does not
restart the error estimate at zero.

Choose a single positive tail threshold small enough that \(C_Z\beta_N\)
is less than half of every retained slab state buffer. Also require it
to be less than the position buffers after applying their fixed
coordinate comparison bounds. There are finitely many reset-domain
tests: their errors are bounded by finite linear combinations of the
incoming state error and the source and target metric errors. Decrease
the same threshold so that each such combination is less than half its
positive margin. More explicitly, if a test requires
\(\alpha_qE+ d_q\beta_N<\mu_q\) with \(\mu_q>0\), impose
\begin{equation}\protect\phantomsection\label{eq:fr-opt-c-margin-test}{
 (\alpha_q C_Z+d_q)\beta_N<\mu_q/2.
}\end{equation} Tests with zero left-hand coefficient need no
restriction. Include \(\beta_N\le\bar\beta\). Their finite minimum is a
strictly positive threshold \(\beta_*\).

Now prove existence and the bounds together by finite induction. The
initial normalization exists at this threshold and
eq.~\ref{eq:fr-opt-c-initial-error} gives \(E_0^+\le H_0\beta_N\). If
the solution on slab \(j\) had a first exit before \(t_{j+1}\), the
integral estimate valid up to that exit would give error at most
\((a_jH_j+b_j)\beta_N\le C_Z\beta_N\), strictly inside its retained
state and position buffers. This contradicts a first exit. The ODE
therefore continues to the slab endpoint inside its compact parent, and
\(E_{j+1}^-\le(a_jH_j+b_j)\beta_N\).

The reset tests are satisfied by eq.~\ref{eq:fr-opt-c-margin-test}, so
the actual reset exists. Its bound gives
\(E_{j+1}^+\le[A_jH_j+B_j]\beta_N=H_{j+1}\beta_N\), again strictly
inside the next slab's initialization neighborhood. This proves the
induction and all switch assertions. Since \(T_{\min}>1\),
\(\beta_N\to0\); choose \(N_0\) so that \(\beta_N<\beta_*\) for all
\(N\ge N_0\). Equivalently it suffices to take
\(T_{\min}^{N+1}>[(1-T_{\min}^{-1})\beta_*]^{-1}\). No numerical value
of this sufficient order is asserted. ∎

\protect\phantomsection\label{sec:fr-opt-c-tidal}

\textbf{Tidal coefficients on the accepted slabs.}

All local estimates now hold on the whole scheduled slabs. Apply the
generic calculations of results K.10--K.12 directly to \(g_0\) and
\(g_{i,N}\) in the \(t\)-clock and unit-\(g_0\) frame convention above.
Let \(D_s=\max_iD_{i,s}\) and retain common \(G,K_s,V,R_0,R_1\) in these
coordinates. Thus \(\|g_0^{-1}\|\le G\), \(\|D^sg_0\|\le K_s\), the
reference and approximate vector components have norm at most \(V\), and
\(R_0,R_1\) bound the exact curvature operator and its first coordinate
derivative. Fixed mass factors are already included in the coordinate
readout constants.

The metric \(C^3\) estimate gives a curvature \(C^1\) estimate on these
tubes. To check this directly, express \(\Gamma\) by \(g_0^{-1}Dg_0\)
and \(R\) by \(D\Gamma+\Gamma\Gamma\), and use the same formulas for
each \(g_{i,N}\). Each component of \((R,DR)\) is a finite sum of
products of inverse metric entries and metric derivatives of order at
most three. On the retained inverse margins, replace one factor at a
time and use eq.~\ref{eq:fr-opt-c-local-metric-tail}. The result is a
finite constant \(C_{R,1}\), independent of \(N\), with
\(\|R_{i,N}-R\|_{C^1}\le C_{R,1}\beta_N\). In particular the exact
\(R_1\) used to move the evaluation point is bounded on the same tube;
this is not inferred from a curvature \(C^0\) estimate alone.

For an explicit admissible \(C^0\) curvature constant define
\begin{equation}\protect\phantomsection\label{eq:fr-opt-c-tidal-input-constants}{
 \begin{aligned}
 \gamma_0&=3G(K_1+1),\\
 d_0^*&=6G^2(K_1+1)D_0+3GD_1,\\
 d_1^*&=\{24G^3(K_1+1)^2+6G^2(K_2+1)\}D_0
          +12G^2(K_1+1)D_1+3GD_2,\\
 d_R^*&=2d_1^*+4\gamma_0d_0^*.
 \end{aligned}
}\end{equation} The smallness conditions in Lemma K.10 are included when
choosing \(N_0\). Its proof gives \(\|R_{i,N}-R\|\le d_R^*\beta_N\). In
the present convention the exact tidal matrix is
\(\mathcal R_{AB,t}=-g_0(R(E_A,K)K,E_B)\). Replacing its metric,
curvature and four vector factors and moving the position, exactly as in
Theorem K.12, gives on each slab a dimensionless coefficient bound
\begin{equation}\protect\phantomsection\label{eq:fr-opt-c-coefficient-bound}{
 \|\mathsf A_{j,N}(t)-\mathsf A_\lambda(t)\|
 \le\kappa_N:=K_{\rm tid}\beta_N,
}\end{equation} where one valid finite choice is
\begin{equation}\protect\phantomsection\label{eq:fr-opt-c-tidal-constant}{
 \begin{aligned}
 K_{\rm tid}=\bigl[&2V^4\{R_0D_0+(K_0+1)d_R^*\}\\
 &+2\{V^4(K_1R_0+K_0R_1)+4K_0R_0V^3\}C_Z\bigr].
 \end{aligned}
}\end{equation} Each state norm and coordinate comparison is the one
used in obtaining \(C_Z\). There is no additional \(m^2\) factor here:
\(-g_0(R(E_A,K)K,E_B)\) is already the dimensionless tidal coefficient,
equal to \(m^2\) times its physical-\(u\), unit-physical-frame
counterpart. More generally the same calculation before substituting
Proposition K.22 gives a bound \(C(\beta_N+\|Z_N-Z\|)\).

On a slab use its own real metric, null tangent and parallel orthonormal
screen. Its tidal matrix is symmetric, even though \(g_{i,N}\) need not
be Einstein. Therefore
\begin{equation}\protect\phantomsection\label{eq:fr-opt-c-local-hamiltonian}{
 \mathsf A_{j,N}(t)=
 \begin{pmatrix}0&I_2\\-\mathcal R_{j,N,t}(t)&0\end{pmatrix},
 \qquad
 \mathsf A_{j,N}^{\mathsf T}\mathsf J+
                  \mathsf J\mathsf A_{j,N}=0,
 \qquad
 \mathsf J=\begin{pmatrix}0&I_2\\-I_2&0\end{pmatrix}.
}\end{equation} The matrices in eq.~\ref{eq:fr-opt-c-coefficient-bound}
are expressed in the corresponding ordered screen components. The state
comparison includes the chosen alignment to the global exact Sachs
frame; arbitrary independent screen rotations are not silently
suppressed.

\pandocstatementheading{\protect\phantomsection\label{prop:fr-opt-c-phase-accumulation}Proposition K.23 (Piecewise Jacobi Cauchy transport).}

Let \(V_{j,N}\) be the exact fundamental matrix on slab \(j\),
initialized by the identity at \(t_j\). At switch \(j\), use the
oriented orthogonal screen identification supplied by the reset
construction. In Cauchy components its phase map is
\begin{equation}\protect\phantomsection\label{eq:fr-opt-c-screen-jump}{
 C_{j,N}=\operatorname{diag}(R_{j,N},R_{j,N}),
 \qquad R_{j,N}\in\operatorname{SO}(2),
 \qquad 1\le j<J.
}\end{equation} Its exact reference is \(C_j=I_4\), because both the
temporary and target screen reduce to the same globally parallel exact
Sachs frame when the metric errors and incoming state error vanish.
There are finite \(H_j^{\rm rot}\) such that
\begin{equation}\protect\phantomsection\label{eq:fr-opt-c-jump-error}{
 \eta_{j,N}:=\|C_{j,N}-I_4\|
 \le H_j^{\rm rot}\beta_N,\qquad
 H_{\rm rot}:=\sum_{j=1}^{J-1}H_j^{\rm rot}<\infty.
}\end{equation} Define the piecewise Jacobi Cauchy transport by solving
the slab Jacobi equation and imposing \(y(t_j^+)=C_{j,N}y(t_j^-)\). Its
complete endpoint matrix is
\begin{equation}\protect\phantomsection\label{eq:fr-opt-c-global-transport}{
 U_N=V_{J-1,N}C_{J-1,N}\cdots C_{1,N}V_{0,N}.
}\end{equation} Then \(U_N\) is symplectic and
\begin{equation}\protect\phantomsection\label{eq:fr-opt-c-product-error}{
 \|U_N-\Phi_\lambda(\tau_\lambda,0)\|
 \le e^{(1+\kappa_N)\tau_\lambda}
       \left(\tau_\lambda\kappa_N+
                    \sum_{j=1}^{J-1}\eta_{j,N}\right).
}\end{equation} Here \(\Phi_\lambda\) is the §K.3 transport in its
polar-first screen and dimensionless phase. The Cauchy transport is the
one defined in Lemma K.21.

\emph{Proof.} The oriented screen identification is analytic near the
exact framed state: projection, positive Gram normalization and the
oriented orthogonal factor are finite matrix operations on the retained
positive margins. At the exact state its rotation is the identity. Its
derivative is uniformly bounded on the finite compact reset
neighborhoods. Its mean-value estimate is
\begin{equation}\protect\phantomsection\label{eq:fr-opt-c-rotation-input-error}{
 \eta_{j,N}\le c_j^{\rm rot}
    \bigl(E_j^-+D_{i_{j-1},0}\beta_N+D_{i_j,0}\beta_N\bigr),
}\end{equation} with any finite matrix/coordinate factors included.
Proposition K.22 thus permits
\(H_j^{\rm rot}=c_j^{\rm rot}(C_Z+D_{i_{j-1},0}+D_{i_j,0})\). In
particular the jump estimate carries the incoming state error; it does
not erase it.

With the common \(g_0\) and \(t\)-clock convention, one may take
\(c_j^{\rm rot}=C_{{\rm ph},j}\) from Lemma K.21. The fixed
coordinate/state-block conversions are included in that constant.

Symmetry in eq.~\ref{eq:fr-opt-c-local-hamiltonian} proves
\(V_{j,N}^{\mathsf T}\mathsf JV_{j,N}=\mathsf J\) by differentiation.
Orthogonality gives \(C_{j,N}^{\mathsf T}\mathsf JC_{j,N}=\mathsf J\)
and \(\|C_{j,N}\|=1\). Thus the full product is symplectic. The
definition explicitly transports both displacement and slope components
with the same orthogonal identification, in the common physical clock
\(u=mt\). It is a well-defined linear phase model attached to the reset
ray/frame. It does not require a derivative of that reset on a family of
spacetime rays.

The exact Sachs generator has norm one. By
eq.~\ref{eq:fr-opt-c-coefficient-bound}, the approximate generator has
norm at most \(1+\kappa_N\). On a slab of length \(h_j\) the two
fundamental matrices satisfy
\begin{equation}\protect\phantomsection\label{eq:fr-opt-c-slab-matrix-error}{
 \|V_{j,N}\|\le e^{(1+\kappa_N)h_j},\qquad
 \|V_j\|\le e^{h_j},\qquad
 \|V_{j,N}-V_j\|\le h_j\kappa_N e^{(1+\kappa_N)h_j}.
}\end{equation} The last inequality is Duhamel's integral formula.
Compare the finite products factor by factor, inserting \(I_4\) at each
exact switch. A slab-difference term is bounded by
\(h_j\kappa_N e^{(1+\kappa_N)\tau_\lambda}\); a jump-difference term is
bounded by \(\eta_{j,N}e^{(1+\kappa_N)\tau_\lambda}\). The remaining
jump factors all have norm one, so no product of \(1+\eta_{j,N}\) is
introduced. Summing and using \(\sum_jh_j=\tau_\lambda\) gives
eq.~\ref{eq:fr-opt-c-product-error}. The same argument applies to every
partial product, with its elapsed time and switches. This also proves
exact composition of this piecewise Cauchy transport at any retained
cut, including a cut just before or just after a reset. ∎

\pandocstatementheading{\protect\phantomsection\label{cor:fr-opt-c-effective-error}Corollary K.24 (A uniform effective coefficient error).}

Put
\begin{equation}\protect\phantomsection\label{eq:fr-opt-c-effective-epsilon}{
 \varepsilon_N=\kappa_N+\frac{H_{\rm rot}\beta_N}{\tau_-}
  =\left(K_{\rm tid}+\frac{H_{\rm rot}}{\tau_-}\right)\beta_N.
}\end{equation} For \(N\ge N_0\) large enough that
\(\varepsilon_N\le1/10\), the full transfer satisfies
\begin{equation}\protect\phantomsection\label{eq:fr-opt-c-four-epsilon}{
 \|U_N-\Phi_\lambda(\tau_\lambda,0)\|
 \le4\varepsilon_N=:\delta_N,\qquad \|U_N\|<4,
}\end{equation} uniformly for the entire original compact ray family.
The transfer inequality is strict when \(\varepsilon_N>0\).

\emph{Proof.} Since \(\tau_\lambda>\tau_-\),
\(\tau_\lambda\kappa_N+\sum_j\eta_{j,N}\le\tau_\lambda\varepsilon_N\).
Moreover \(\kappa_N\le\varepsilon_N\). Consequently Proposition K.23
bounds the transfer difference by
\(\tau_\lambda\varepsilon_Ne^{(1+\varepsilon_N)\tau_\lambda}\). Here
\((1+\varepsilon_N)\tau_\lambda<10/9\). The elementary estimates
\(e<11/4\) and \(e^{1/9}\le9/8\) give \(e^{10/9}<99/32\), whence \[
 \tau_\lambda e^{(1+\varepsilon_N)\tau_\lambda}
 <\frac{100}{99}\frac{99}{32}=\frac{25}{8}<4.
\] The product norm is at most \(e^{(1+\kappa_N)\tau_\lambda}<99/32<4\)
because every jump has norm one. If \(\varepsilon_N=0\), the coefficient
and jump errors vanish and the transfer difference is zero. This proves
both assertions. All choices remain independent of \(N\), and increasing
\(N\) meets the final smallness condition. ∎

\pandocstatementheading{\protect\phantomsection\label{prop:fr-opt-c-observables}Proposition K.25 (Forward observations and conditional inverse
observations).}

Under the hypotheses of Corollary K.24, in particular \(N\ge N_0\) and
\(\varepsilon_N\le1/10\), write
\(\Phi_\lambda=\left(\begin{smallmatrix}A&B\\C&D\end{smallmatrix}\right)\)
and
\(U_N=\left(\begin{smallmatrix}A_N&B_N\\C_N&D_N\end{smallmatrix}\right)\)
in the endpoint phases just specified. Let \(\delta_N=4\varepsilon_N\).
Then the outgoing phase of an initial component vector of norm at most
\(J_0\) has error at most \(J_0\delta_N\). If the initial phase itself
is approximated with error \(\eta_0\), its total outgoing error is at
most \(4\eta_0+J_0\delta_N\).

For point-source Cauchy data in this piecewise phase model, \(\xi(0)=0\)
and \(p(0)=m\,d\xi/du(0)\), so its endpoint response block to the
physical slope is \(mB_N\). The oriented area observations
\(\mathcal O_N=m^2\det B_N\) and \(\mathcal O=m^2\det B\) satisfy
\begin{equation}\protect\phantomsection\label{eq:fr-opt-c-area-error}{
 |\mathcal O_N-\mathcal O|\le9m^2\delta_N.
}\end{equation} If additionally \(\delta_N\le3/14\), then
\begin{equation}\protect\phantomsection\label{eq:fr-opt-c-inverse-error}{
 \sigma_{\min}(B_N)\ge\frac3{14},\qquad
 \|B_N^{-1}\|\le\frac{14}{3},\qquad
 \|B_N^{-1}-B^{-1}\|\le\frac{98}{9}\delta_N.
}\end{equation} The corresponding inverse difference for \(mB_N\) is at
most \(98\delta_N/(9m)\). On the same margin the dimensionless expansion
trace also satisfies
\begin{equation}\protect\phantomsection\label{eq:fr-opt-c-expansion-error}{
 |\operatorname{tr}(D_NB_N^{-1})-
                  \operatorname{tr}(DB^{-1})|
 \le4\delta_N\left\{\frac73+\|D\|\frac{49}{9}\right\}.
}\end{equation} The physical-\(u\) trace has this bound divided by
\(m\).

\emph{Proof.} Block projections have norm one, so the global transfer
bound gives each block error at most \(\delta_N\). The outgoing phase
bounds follow from \(U_Ny_N-\Phi y=U_N(y_N-y)+(U_N-\Phi)y\). The phase
identification is oriented orthogonal at the switches; it introduces
neither a determinant sign nor an affine-clock dilation. Thus the same
physical normalization of the original point-source observation gives
\(mB_N\).

Since \(\varepsilon_N\le1/10\), \(\delta_N\le2/5<1\). The two-column
determinant identity gives \[
 |\det B_N-\det B|
 \le2\|B\|\delta_N+\delta_N^2.
\] For the original exact family, §K.3.7 gives \(\|B\|<100/33\). Hence
the right-hand side is at most \((200/33+1)\delta_N\le9\delta_N\),
proving the area bound; the final inequality is strict when
\(\delta_N>0\). This use of the constant nine is justified by the new
full product estimate, not by an assumption that different metric
readouts glue.

The same exact family has \(\sigma_{\min}(B)>3/7\). Therefore
\(\delta_N\le3/14\) implies \(\sigma_{\min}(B_N)\ge3/14\) and the
inverse norm bound. The inverse identity yields \[
 B_N^{-1}-B^{-1}=-B_N^{-1}(B_N-B)B^{-1},\qquad
 \|B_N^{-1}-B^{-1}\|
 \le\frac{14}{3}\frac73\delta_N=\frac{98}{9}\delta_N.
\] For the trace, subtract the products as
\((D_N-D)B_N^{-1}+D(B_N^{-1}-B^{-1})\) and use
\(|\operatorname{tr}M|\le2\|M\|\). This gives exactly
eq.~\ref{eq:fr-opt-c-expansion-error}. The physical clock conversion
follows from \(p=m\,d\xi/du\), with the same fixed positive mass on
every slab. ∎

\protect\phantomsection\label{sec:fr-opt-c-scope}

\textbf{Singular position projections.}

For every \(N\) under consideration, \(U_N\) is an invertible symplectic
map. Its graph is a four-dimensional Lagrangian relation in the product
phase space with forms \(-\Omega\) and \(\Omega\). It carries every
initial Lagrangian plane to a Lagrangian plane even when a position
projection is singular. The piecewise Cauchy definition and the forward
product and determinant estimates remain meaningful there. An inverse
endpoint block or expansion trace requires the explicit singular-value
margin in Proposition K.25; it does not follow from symplecticity alone.
For the original point-source family that margin persists for
sufficiently large \(N\) by eq.~\ref{eq:fr-opt-c-inverse-error}. No
persistence at a prescribed caustic time of a different initial
Lagrangian is asserted.

More explicitly, represent any fixed initial Lagrangian plane by a
rank-two matrix \(Y_0\in\mathbb R^{4\times2}\) with
\(Y_0^{\mathsf T}\mathsf JY_0=0\). Then \(U_NY_0\) has rank two and the
same isotropy identity. With
\(Q=\operatorname{pr}_{\xi}\Phi_\lambda Y_0\) and
\(Q_N=\operatorname{pr}_{\xi}U_NY_0\), one has
\begin{equation}\protect\phantomsection\label{eq:fr-opt-c-lagrangian-projection-error}{
 \|Q_N-Q\|\le\delta_N\|Y_0\|,\qquad
 |\det Q_N-\det Q|
 \le2\|Q\|\delta_N\|Y_0\|+\delta_N^2\|Y_0\|^2.
}\end{equation} These forward inequalities do not divide by \(\det Q\)
and therefore remain valid when that projection has a caustic. The full
transported plane, rather than an inverse of \(Q\), is retained.

The final times are the exact scheduled \(t_J=\tau_\lambda\),
corresponding to \(u=m\tau_\lambda\). The approximate endpoint need not
lie exactly on \(r=3m/2\). Comparing approximate hitting events instead
would require additional time-of-hit and matching estimates; none is
used here. The approximate piecewise state need not conserve the exact
Schwarzschild \(E\) and \(L\), although its reference ray is exactly the
original family with those parameters.

All finite readout constants, reset derivative bounds, buffers, \(C_Z\),
\(K_{\rm tid}\) and \(H_{\rm rot}\) are theoretical finite constants.
The finite product/sum and the sufficient tail inequalities specify them
and a sufficient \(N\) in terms of the chosen parents and local data.
They have not been numerically evaluated. The rational constants \(4\),
\(9\), \(14/3\) and \(98/9\) are explicit and apply by the proofs above.
No whole-long-family containment in a single local readout, global
approximate metric, off-one physical-amplitude gluing, or variational
interpretation of the reset derivative is required by these conclusions.

\protect\phantomsection\label{app:gg-notation}

\protect\phantomsection\label{sec:gg-notation-index}

\subsection{Appendix L. Notation
index}\label{appendix-l.-notation-index}

The following tables collect the notation used in the manuscript. It
introduces no additional assumptions or definitions. Chiral labels refer
to the two spin factors in the fixed convention, and the numerical
Einstein label is always the one in \(\operatorname{Ric}=\Lambda g\).
The Lorentz signature is \((+,-,-,-)\); the curvature contraction and
primary printed-label dictionary are collected in §2.

\protect\phantomsection\label{sec:gg-notation-source}

\subsubsection{L.1. Source and presentation
spaces}\label{l.1.-source-and-presentation-spaces}

{\def\LTcaptype{none} % do not increment counter
\begin{longtable}[]{@{}
  >{\raggedright\arraybackslash}p{(\linewidth - 4\tabcolsep) * \real{0.3333}}
  >{\raggedright\arraybackslash}p{(\linewidth - 4\tabcolsep) * \real{0.3333}}
  >{\raggedright\arraybackslash}p{(\linewidth - 4\tabcolsep) * \real{0.3333}}@{}}
\toprule\noalign{}
\begin{minipage}[b]{\linewidth}\raggedright
Symbol
\end{minipage} & \begin{minipage}[b]{\linewidth}\raggedright
Meaning
\end{minipage} & \begin{minipage}[b]{\linewidth}\raggedright
First definition
\end{minipage} \\
\midrule\noalign{}
\endhead
\bottomrule\noalign{}
\endlastfoot
\(T_0\) & The germ of \(\operatorname{Tot}(\mathcal O(1)^{\oplus2})\)
along its whole zero line. & §2 \\
\(\mathscr H_m(T_0)\) & Analytic two-chart Čech classes of weight \(m\),
on the whole line germ. & Definition 2.1 \\
\(\mathcal A_\Lambda\) & The original independent pair
\(\mathscr H_2(T_0)\oplus\mathscr H_{-6}(T_0)\) at numerical
\(\Lambda\). & Definition 2.1 \\
\(U_r^+,U_r^-\) & The separately symmetric free curvature records of
types \((r+4,r)\) and \((r,r+4)\), with the displayed component
factorials. & §3 and Theorem 3.1 \\
\(\mathfrak E_\Lambda\) & Marked local analytic Einstein germs, with
physical point and full-spin instructions retained. & Definition 2.1 \\
\(\mathfrak S_\Lambda\) & Independently recognized regular-normal affine
ray objects. & Definition 2.2 \\
\(\mathfrak W_\Lambda\) & Independent ordinary threefold data,
calibrated actual form and unrestricted opposite-chirality class. &
Definition 2.3 \\
\(\mathfrak C_\Lambda\) & The calibrated affine presentation category,
with independent projective-extension admission. & Definition 5.6 and
Theorem H.4 \\
\(\mathfrak O_\Lambda\) & Independently admitted structured analytic
operator objects and their physical comparison classes. & Definition
H.1 \\
\end{longtable}
}

\protect\phantomsection\label{sec:gg-notation-bundles}

\subsubsection{L.2. Bundles, roots and
operators}\label{l.2.-bundles-roots-and-operators}

{\def\LTcaptype{none} % do not increment counter
\begin{longtable}[]{@{}
  >{\raggedright\arraybackslash}p{(\linewidth - 4\tabcolsep) * \real{0.3333}}
  >{\raggedright\arraybackslash}p{(\linewidth - 4\tabcolsep) * \real{0.3333}}
  >{\raggedright\arraybackslash}p{(\linewidth - 4\tabcolsep) * \real{0.3333}}@{}}
\toprule\noalign{}
\begin{minipage}[b]{\linewidth}\raggedright
Symbol
\end{minipage} & \begin{minipage}[b]{\linewidth}\raggedright
Meaning
\end{minipage} & \begin{minipage}[b]{\linewidth}\raggedright
Location
\end{minipage} \\
\midrule\noalign{}
\endhead
\bottomrule\noalign{}
\endlastfoot
\(Q_x\) & The compact Legendrian sky \(\mathbb P^1\times\mathbb P^1\) of
a physical point. & §2 \\
\(p:F\to M\), \(q:F\to Z\) & The complete sky incidence; \(Z\) is the
contact fivefold. & §2 \\
\(\mathcal L\) & The contact line, restricting to \(\mathcal O(1,1)\) on
each sky. & §5 and Definition H.1 \\
\(\mathcal E_J\) & The intrinsic rank-two projective optical solution
bundle, with its specified Wronskian isomorphism
\(\bigwedge^2\mathcal E_J\simeq\mathcal L\). & Proposition 5.3 and
Appendix D.3 \\
\(\mathcal E[1]\) & The calibrated conformal density line; its optical
square identification is retained. & Appendix D \\
\(\mathcal V\), \(\mathcal T\) & Respectively the descended rank-six ray
bundle and its reconstructed normal tractor bundle. & Definition 2.2 and
§4 \\
\(W_{\mathrm{scr}}\) & The rank-four middle bundle
\((\mathbb CI+\mathbb CX)^\perp\), identified with \(TM\) by the solder;
its tangent Spin type is \(S\otimes S'\). & Definition 2.2 and
Proposition 4.3 \\
\(\mathcal T_{\mathrm{LT}}\) & The rank-four normal local-twistor
bundle, with volume and marked conformal splitting \(S\oplus S'\); the
tractor identification is
\(\mathcal T\simeq\Lambda^2\mathcal T_{\mathrm{LT}}^*\). & Proposition
I.3 and Appendix J.1 \\
\(L\), \(\alpha\), \(\kappa\) & The ordinary degree-one line, its actual
fourth-root identification and central degree-one calibration in the
retained full-spin mark. & Definition 2.3 \\
\(R\), \(R^2=K_{F/S}\) & The module half-canonical root. Here \(S\to B\)
is the total point family over the passive base, with \(S=M\) at fixed
parameter; retain \(K_{F/B}\simeq K_{F/S}\otimes p^*K_{S/B}\), its
square comparison and every calibrated base factor. & Appendix G.3 \\
\(\mathcal E_{T_0}^{1/2}\) & The analytic half-density microdifferential
operator sheaf on the derived projective cotangent carrier. & Appendix
F \\
\(\mathscr U_1\) & The truncated-unit sheaf
\(1+(F_{-1}\mathcal E/F_{-2}\mathcal E)\) inside
\((F_0\mathcal E/F_{-2}\mathcal E)^\times\), identified with the
additive sheaf \(\mathcal L^{-1}\). & Appendix F.3 \\
\(\varphi_{ij}\), \(C_{ijk}\) & Actual operator comparisons and their
normalized inner associators. & Appendices F--H \\
\(\tau_{ij}\) & Module descent comparisons satisfying
\(\tau_{ij}\tau_{jk}=C_{ijk}\tau_{ik}\). & Appendix G \\
\end{longtable}
}

The operator sheaf \(\mathcal E_{T_0}^{1/2}\), density line
\(\mathcal E[1]\) and optical bundle \(\mathcal E_J\) are different
objects. The module root \(R\) is also different from the optical
density root. A module comparison's principal scalar and its whole
local-system frame have separate roles.

\protect\phantomsection\label{sec:gg-notation-execution}

\subsubsection{L.3. Execution, inverses and scalar
conventions}\label{l.3.-execution-inverses-and-scalar-conventions}

{\def\LTcaptype{none} % do not increment counter
\begin{longtable}[]{@{}
  >{\raggedright\arraybackslash}p{(\linewidth - 4\tabcolsep) * \real{0.3333}}
  >{\raggedright\arraybackslash}p{(\linewidth - 4\tabcolsep) * \real{0.3333}}
  >{\raggedright\arraybackslash}p{(\linewidth - 4\tabcolsep) * \real{0.3333}}@{}}
\toprule\noalign{}
\begin{minipage}[b]{\linewidth}\raggedright
Symbol
\end{minipage} & \begin{minipage}[b]{\linewidth}\raggedright
Meaning
\end{minipage} & \begin{minipage}[b]{\linewidth}\raggedright
Location
\end{minipage} \\
\midrule\noalign{}
\endhead
\bottomrule\noalign{}
\endlastfoot
\(A_N\) & The amplitude algebra
\(\mathbb C[\varepsilon]/(\varepsilon^{N+1})\). & §5 \\
\(o_m\), \(d_m\), \(L_m\) & The optical obstruction derivative, its
compatibility map and the unique inverse on the zero-free-record
summand. & §5, Theorem 5.5 and Appendices C--D \\
\(\mathcal C_N\) & A finite source contact/optical stage before the
selected storage normalization. & §§5--6 \\
\(\mathcal N_P\), \(X_N\), \(Y_C\) & The fixed storage normalizer, its
finite output and the one joint analytic comparison tuple, with
\(X_N=j_\varepsilon^N Y_C\). & §6 \\
\(\mathcal G\), \(\mathcal E\) & The source-to-ray constructor and
full-record decoder; the latter is a map, not an operator sheaf. &
Theorem 2.4 \\
\(\Gamma\), \(d_{\mathrm{rec}}\) & Ordinary registration and its
independently defined record decoder. & §7 \\
\(P_Z\), \(K_Z\) & The curved negative contour and its one fixed bounded
representative inverse. & Appendix E \\
\(\mathcal S_{\mathrm{op}}\), \(\mathcal D_{\mathrm{op}}\) & The actual
source/operator constructor and decoder, with both independent-domain
inverse comparisons. & Theorem 9.11 \\
\(m_{r,\nu}\), \(m_r\), \(\gamma_r\) & Block atomic diagram masses, the
one weighted mass and its bounded-product rescaling; all refer to the
same recursively ordered coefficients. & Appendix B.2 and Lemma B.3 \\
\(F\), \(Q\) & Classical extraction and structured quantization, with
actual units and counits on the independent domains. & Theorem H.4 \\
\(I\) & The selected normal tractor, or its identified local-twistor
form; it is parallel on Einstein-admitted objects. The pre-admission
defect is given in J.1. & §§4--5 and Appendices I--J \\
\(H(I,I)\) & On Einstein-admitted objects, the selected normal-tractor
norm is \(-\Lambda/3\). & Definition 2.2 and Appendix J.1 \\
\(c\) & The scalar \(\Lambda/6\) in the rank-four spin realization; it
is distinct from a contact overlap factor \(c_{ij}\). & Proposition
I.3 \\
\(\Lambda_P\) & The primary printed label under the declared
curvature-component dictionary: \(\Lambda_P=-\Lambda\). The prescribed
native \(\Lambda\) is unchanged. & §2 and Proposition I.3 \\
\end{longtable}
}

Both original sources have their gravitational weights \(2\) and \(-6\).
The projective ket Euler sector has weight \(-2\); it is not a
relabeling of either gravitational source. The microlocal cone
\(\mathcal O(-1,-1)^\times\) and the Lorentz phase-scaled spinor bundle
are distinct typed objects.

\protect\phantomsection\label{sec:gg-notation-simultaneous-diagram}

\subsubsection{L.4. The simultaneous projective
diagram}\label{l.4.-the-simultaneous-projective-diagram}

{\def\LTcaptype{none} % do not increment counter
\begin{longtable}[]{@{}
  >{\raggedright\arraybackslash}p{(\linewidth - 4\tabcolsep) * \real{0.3333}}
  >{\raggedright\arraybackslash}p{(\linewidth - 4\tabcolsep) * \real{0.3333}}
  >{\raggedright\arraybackslash}p{(\linewidth - 4\tabcolsep) * \real{0.3333}}@{}}
\toprule\noalign{}
\begin{minipage}[b]{\linewidth}\raggedright
Symbol
\end{minipage} & \begin{minipage}[b]{\linewidth}\raggedright
Meaning
\end{minipage} & \begin{minipage}[b]{\linewidth}\raggedright
Location
\end{minipage} \\
\midrule\noalign{}
\endhead
\bottomrule\noalign{}
\endlastfoot
\(L\sigma\) & The canonical tractor prolongation of the density selected
by the same affine torsor and optical root. This \(L\) is a differential
operator, distinct from either line bundle denoted \(L\) or
\(\mathcal L\). & Appendix J.1 \\
\(\mathcal T_{\mathrm{LT,ray}}\) & The rank-four local-twistor bundle on
the ray space, obtained by descending \(\mathcal T_{\mathrm{LT}}\) along
the retained short null rays, with its volume. It is distinct from the
middle bundle \(W_{\mathrm{scr}}\) and the rank-six tractor
\(\mathcal T\). & Theorem J.2 \\
\(\ell\subset h\) & The descended rank-one/rank-three flag;
\(\ell|_{Q_x}=\mathcal O(0,-1)\) and
\(\operatorname{Ann}(h)|_{Q_x}=\mathcal O(-1,0)\). & Theorem J.2 \\
\(\mathcal F_{\mathrm{spin}}\) & The exterior-square identification
\(\mathcal T\simeq\Lambda^2\mathcal T_{\mathrm{LT}}^*\) of the normal
tractor with alternating forms on the local-twistor bundle. &
Proposition I.3 and Appendix J.1 \\
\(V_\gamma\) & The parallel totally null optical two-plane
\(\operatorname{span}(X,D_kX)\), with
\(V_\gamma^*\simeq\mathcal E_{J,\gamma}\). & Proposition J.4 \\
\(K_{x,\gamma}\) & The rank-two point plane
\(\ker\mathcal F_{\mathrm{spin}}(X(x))\) for a ray through \(x\). &
Proposition J.4 \\
\(\mathcal R(O)\), \(g_A\) & The independently decoded marked Einstein
germ and the canonical germ for the original input records. & Appendix
J.5 \\
\(\mathcal P(O)\), \(\operatorname{Can}(g)\) & The full derived
projective diagram of a supplied operator object, and the canonical
diagram of an Einstein germ enriched by the specified quantization. &
Appendix J.5 \\
\(\beta_O\), \(\alpha_A\) & The actual independent-domain and source
diagram comparisons. The subscript distinguishes \(\beta_O\) from the
optical root-square map \(\beta\). & Theorem J.6 \\
\(J_{\mathrm{LT}}\) & The antilinear Lorentz duality from
\((\mathcal T_{\mathrm{LT,ray}})_y\) to
\((\mathcal T_{\mathrm{LT,ray}}^*)_{\iota(y)}\). The target is the
paired dual fiber; the complex-linear operator formal adjoint has a
different type. & Proposition J.7 \\
\(H_{\mathrm{LT}}\) & The Hermitian pairing matrix of the local-twistor
Lorentz comparison in compatible marked spin frames. & Proposition
J.7 \\
\end{longtable}
}

Parameter families are based on finite-dimensional complex manifold
germs. Their finite jets mean actual Taylor restrictions, as specified
in Definition 2.1.

\protect\phantomsection\label{sec:gg-notation-local-conventions}

\subsubsection{L.5. Component and locally reused
notation}\label{l.5.-component-and-locally-reused-notation}

{\def\LTcaptype{none} % do not increment counter
\begin{longtable}[]{@{}
  >{\raggedright\arraybackslash}p{(\linewidth - 4\tabcolsep) * \real{0.3333}}
  >{\raggedright\arraybackslash}p{(\linewidth - 4\tabcolsep) * \real{0.3333}}
  >{\raggedright\arraybackslash}p{(\linewidth - 4\tabcolsep) * \real{0.3333}}@{}}
\toprule\noalign{}
\begin{minipage}[b]{\linewidth}\raggedright
Symbol
\end{minipage} & \begin{minipage}[b]{\linewidth}\raggedright
Meaning and scope
\end{minipage} & \begin{minipage}[b]{\linewidth}\raggedright
Location
\end{minipage} \\
\midrule\noalign{}
\endhead
\bottomrule\noalign{}
\endlastfoot
\(\epsilon_{AB},\epsilon^{AB}\) & The mutually inverse alternating
matrices; raising and lowering are left contractions. The primed
matrices have the same convention. & §2 \\
\(C_{abcd},\Psi^\pm,\mathfrak p_\pm\) & Last-slot covariant Weyl tensor,
its signed chiral readouts, and their explicitly inverse component
contractions. & §2 \\
\(\mathcal D_z\) & The ray contact plane, identified with orthogonal
Jacobi solutions modulo \((a+bu)k\). It is distinct from the source
decoder \(\mathcal E\). & Proposition 5.3 \\
\(S(u)\) & A matrix graph of the tangent-sky curve; its symbol is local
to the optical calculation. & Proposition 5.3 and Appendix D.2 \\
\(S\to B\), \(M\) & Total physical point family over passive parameters;
\(M\) is a fiber at a fixed passive parameter. These bases are distinct
from the spin spaces \(S,S'\) of §2. & Appendices G.3 and H.1 \\
\(X\), \(X_N\), \(Y_C\) & Position tractor in the tractor constructions,
with locally declared reuse for leaf graphs, solder matrices and EF
coordinates; \(Y_C\) is the joint analytic tuple and \(X_N\) its
amplitude Taylor polynomial. & §§4 and 6, Appendices E.6 and K.8, and
Proposition J.7 \\
\(P\), \(\mathsf P\), \(P_Z\) & The Schouten tensor, its flat linearized
metric operator, and the curved contour transform, respectively. & §2,
Appendix C.3, Appendix E \\
\(P_m=\mathcal P_m\), \(\mathcal P(O)\) & The space
\(\operatorname{Sym}^mV^*\), zero for \(m<0\), and the derived
projective diagram of an operator object, respectively. & Appendix D.4
and Appendix J.5 \\
\(J\), \(\mathfrak j\) & \(J=\operatorname{Scal}/6\) is the Schouten
trace; a separately declared highest square-zero ideal is a
coefficient-algebra ideal. & §§5 and J.1 \\
\(\iota\), \(\mathcal K\), \(B\) & Vacuum inclusion, ray/density or
kernel objects, and block/parameter symbols are used only with their
displayed local domains. The passive base \(B\) in G--H is distinct from
the optical matrices \(B_{ij}\). & §§5--6 and Appendices C--H \\
\(g_0=\bar g\), \(K,\mathcal N,E_A\) & Dimensionless EF metric
\(m^{-2}g_m\) and its normalized full null frame; the physical vectors
are \((K,\mathcal N,E_A)/m\). & §K.8 \\
\(\beta_N\), \(\varepsilon_N\) & The common local-readout geometric tail
and the effective tidal-plus-switch error. & Proposition 10.5, §K.8 \\
\(L\), \(R\), \(C_{j,N}\) & The pointwise metric-to-metric tangent
isometry, recorded oriented screen rotation and symplectic Cauchy
switch; these symbols are local to the finite-readout construction. &
Lemma K.21 \\
\(U_N\) & The complete piecewise metric Jacobi Cauchy transport at the
original exact affine endpoint. & §K.8 \\
\end{longtable}
}

\protect\phantomsection\label{sec:gg-ai-use-disclosure}

\protect\phantomsection\label{ai-use-disclosure}

\textbf{AI-use disclosure.} AI systems were used extensively throughout
this research and the preparation of this manuscript. The models used
were GPT-6 Pro, gpt-6.1-sol (xhigh), and Claude Opus 5.5 (xhigh).
Selected symbolic identities and explicit metric examples were checked
using SageMath and SymPy; the code, inputs, computational environment
and recorded results are provided in the computational supplement
(Isthmus 2026).

\section*{References}

\protect\phantomsection\label{refs}
\begin{CSLReferences}{1}{1}
\sloppy
\bibitem[\citeproctext]{ref-Adamo2018SpinorConventions}
Adamo, Tim. 2018. {``Lectures on Twistor Theory.''} \emph{Proceedings of
Science} Modave2017: 003. \url{https://doi.org/10.22323/1.323.0003}.

\bibitem[\citeproctext]{ref-AdamoEtAl2026v2}
Adamo, Tim, Bernardo Araneda, Sean Seet, and Atul Sharma. 2026a.
\emph{Schwarzschild Black Holes from Twistor Space}.
\url{https://arxiv.org/abs/2607.06236v2}.

\bibitem[\citeproctext]{ref-AdamoEtAl2026}
Adamo, Tim, Bernardo Araneda, Sean Seet, and Atul Sharma. 2026b.
\emph{Schwarzschild Black Holes from Twistor Space}. Version 1.
\url{https://arxiv.org/abs/2607.06236v1}.

\bibitem[\citeproctext]{ref-BastonMason1987}
Baston, Robert J., and Lionel J. Mason. 1987. {``Conformal Gravity, the
{Einstein} Equations and Spaces of Complex Null Geodesics.''}
\emph{Classical and Quantum Gravity} 4: 815--26.
\url{https://doi.org/10.1088/0264-9381/4/4/018}.

\bibitem[\citeproctext]{ref-DAgnoloSchapira2007}
D'Agnolo, Andrea, and Pierre Schapira. 2007. {``Quantization of Complex
Lagrangian Submanifolds.''} \emph{Advances in Mathematics} 213 (1):
358--79. \url{https://doi.org/10.1016/j.aim.2006.12.009}.

\bibitem[\citeproctext]{ref-Douady1966}
Douady, Adrien. 1966. {``{Le problème des modules pour les sous-espaces
analytiques compacts d'un espace analytique donné}.''} \emph{Annales de
l'Institut Fourier} 16 (1): 1--95.
\url{https://doi.org/10.5802/aif.226}.

\bibitem[\citeproctext]{ref-Friedrich2013}
Friedrich, Helmut. 2013. {``The {Taylor} Expansion at Past Time-Like
Infinity.''} \emph{Communications in Mathematical Physics} 324:
263--300. \url{https://doi.org/10.1007/s00220-013-1803-1}.

\bibitem[\citeproctext]{ref-Gover2004}
Gover, A. Rod. 2004. \emph{Almost Conformally {Einstein} Manifolds and
Obstructions}. \url{https://arxiv.org/abs/math/0412393v1}.

\bibitem[\citeproctext]{ref-Grauert1960}
Grauert, Hans. 1960. {``{Ein Theorem der analytischen Garbentheorie und
die Modulräume komplexer Strukturen}.''} \emph{Publications
Mathématiques de l'IHÉS} 5: 5--64.
\url{https://doi.org/10.1007/BF02684746}.

\bibitem[\citeproctext]{ref-Grauert1963}
Grauert, Hans. 1963. {``{Berichtigung: {Ein Theorem der analytischen
Garbentheorie und die Modulräume komplexer Strukturen}}.''}
\emph{Publications Mathématiques de l'IHÉS} 16: 35--36.
\url{https://doi.org/10.1007/BF02684291}.

\bibitem[\citeproctext]{ref-Isthmus2026LocalComputations}
Isthmus, Lior. 2026. \emph{Computational Supplement for a Projective
Twistor Correspondence for Local Analytic Einstein Germs}. Versions
0.2.0, revision 3.
\url{https://github.com/lioristhmus/projective-twistor-einstein-germs/tree/paper-v1.1/computations}.

\bibitem[\citeproctext]{ref-Kashiwara1986}
Kashiwara, Masaki. 1986. \emph{Introduction to Microlocal Analysis}.
L'Enseignement Mathématique, Université de Genève.
\url{https://www.kurims.kyoto-u.ac.jp/~kenkyubu/kashiwara/micro.pdf}.

\bibitem[\citeproctext]{ref-Kashiwara1996}
Kashiwara, Masaki. 1996. {``Quantization of Contact Manifolds.''}
\emph{Publications of the Research Institute for Mathematical Sciences}
32 (1): 1--7. \url{https://doi.org/10.2977/PRIMS/1195163179}.

\bibitem[\citeproctext]{ref-Kodaira1962}
Kodaira, Kunihiko. 1962. {``A Theorem of Completeness of Characteristic
Systems for Analytic Families of Compact Submanifolds of Complex
Manifolds.''} \emph{Annals of Mathematics}, 2nd series, vol. 75 (1):
146--62. \url{https://doi.org/10.2307/1970424}.

\bibitem[\citeproctext]{ref-Kowalevsky1875}
Kowalevsky, Sophie von. 1875. {``{Zur Theorie der partiellen
Differentialgleichungen}.''} \emph{{Journal für die reine und angewandte
Mathematik}} 80: 1--32.

\bibitem[\citeproctext]{ref-LeBrun1983}
LeBrun, Claude. 1983. {``Spaces of Complex Null Geodesics in
Complex-{Riemannian} Geometry.''} \emph{Transactions of the American
Mathematical Society} 278 (1): 209--31.
\url{https://doi.org/10.1090/S0002-9947-1983-0697071-9}.

\bibitem[\citeproctext]{ref-LeBrun1985}
LeBrun, Claude. 1985. {``Ambi-Twistors and {Einstein}'s Equations.''}
\emph{Classical and Quantum Gravity} 2 (4): 555--63.
\url{https://doi.org/10.1088/0264-9381/2/4/020}.

\bibitem[\citeproctext]{ref-LeBrun1991}
LeBrun, Claude. 1991. {``Thickenings and Conformal Gravity.''}
\emph{Communications in Mathematical Physics} 139 (1): 1--43.
\url{https://doi.org/10.1007/BF02102727}.

\bibitem[\citeproctext]{ref-Penrose1976}
Penrose, Roger. 1976. {``Non-Linear Gravitons and Curved Twistor
Theory.''} \emph{General Relativity and Gravitation} 7: 31--52.
\url{https://doi.org/10.1007/BF00762011}.

\bibitem[\citeproctext]{ref-Penrose1980}
Penrose, Roger. 1980. {``Golden Oldie: Null Hypersurface Initial Data
for Classical Fields of Arbitrary Spin and for General Relativity.''}
\emph{General Relativity and Gravitation} 12 (3): 225--64.
\url{https://doi.org/10.1007/BF00756234}.

\bibitem[\citeproctext]{ref-Penrose2015}
Penrose, Roger. 2015. {``Palatial Twistor Theory and the Twistor Googly
Problem.''} \emph{Philosophical Transactions of the Royal Society A} 373
(2047): 20140237. \url{https://doi.org/10.1098/rsta.2014.0237}.

\bibitem[\citeproctext]{ref-Penrose2026BiTwistors}
Penrose, Roger. 2026. {``Bi-Twistors: An Extension of Twistor Theory for
General Spacetimes.''} \emph{International Journal of Modern Physics D}
35 (6): 2540015. \url{https://doi.org/10.1142/S0218271825400152}.

\bibitem[\citeproctext]{ref-Polesello2005}
Polesello, Pietro. 2007. {``Uniqueness of {Kashiwara}'s Quantization of
Complex Contact Manifolds.''} In \emph{Geometric and Topological Methods
for Quantum Field Theory}, vol. 434. Contemporary Mathematics. American
Mathematical Society. \url{https://arxiv.org/abs/math/0512501v1}.

\bibitem[\citeproctext]{ref-Ward1980}
Ward, Richard S. 1980. {``Self-Dual Space-Times with Cosmological
Constant.''} \emph{Communications in Mathematical Physics} 78 (1):
1--17. \url{https://doi.org/10.1007/BF01941967}.

\bibitem[\citeproctext]{ref-Wolf2007}
Wolf, Martin. 2007. {``Self-Dual Supergravity and Twistor Theory.''}
\emph{Classical and Quantum Gravity} 24 (24): 6287--328.
\url{https://doi.org/10.1088/0264-9381/24/24/010}.

\end{CSLReferences}

\end{document}
